% Generated mechanically from frozen input inputs/p08/ja/spaces-morphisms.tex.
% Do not edit this generated file; edit the bound locale source instead.

% OK, start here.
%

\setcounter{chapter}{66}
\chapter{代数空間の射}



\phantomsection
\label{spaces-morphisms-section-phantom}


\section{はじめに}
\label{spaces-morphisms-section-introduction}

\noindent
本章では、代数空間の射のいくつかの型を導入する。
文献として \cite{Kn} がある。

\medskip\noindent
目標は、スキームおよびスキームの射に関する各章で定義された
射の各型の定義を、代数空間の圏へ拡張することである。
各場合には少しずつ違いがあるため、すべてを個別に扱うのが最善と思われる。





\section{規約}
\label{spaces-morphisms-section-conventions}

\noindent
以下では、すべてのスキームが大 fppf サイト $\Sch_{fppf}$ に
含まれるものとする。また、考えるすべての環 $A$ は、
$\Spec(A)$ がこの大サイトの対象と（同型で）あるという性質をもつものとする。

\medskip\noindent
$S$ をスキームとし、$X$ を $S$ 上の代数空間とする。
本章および以下の各章では、$X$ とそれ自身との積（$S$ 上の代数空間の
圏における積）を、$X \times X$ ではなく $X \times_S X$ と書く。





\section{表現可能な射の性質}
\label{spaces-morphisms-section-representable}

\noindent
$S$ をスキームとする。
$f : X \to Y$ を代数空間の表現可能な射とする。Spaces, Section
\ref{spaces-section-representable-properties} では、$\mathcal{P}$ が
次を満たすスキームの射の性質である場合に、$f$ が性質
$\mathcal{P}$ をもつとは何かを定義した。
\begin{enumerate}
\item 任意の基底変換で保たれる。Schemes, Definition
\ref{schemes-definition-preserved-by-base-change} を参照。
\item 基底上 fppf 局所的である。Descent, Definition
\ref{descent-definition-property-morphisms-local} を参照。
\end{enumerate}
すなわちこの場合、任意のスキーム $U$ と任意の射 $U \to Y$ に対して、
スキームの射 $X \times_Y U \to U$ が性質 $\mathcal{P}$ をもつとき、
かつそのときに限り、$f$ は性質 $\mathcal{P}$ をもつという。

\medskip\noindent
Spaces, Section \ref{spaces-section-lists} の一覧によれば、
これは次の性質に適用される：
(1)(a) 閉埋め込み、
(1)(b) 開埋め込み、
(1)(c) 準コンパクトな埋め込み、
(2) 準コンパクト、
(3) 普遍閉、
(4) （準）分離的、
(5) 単射、
(6) 全射、
(7) 普遍的単射、
(8) アフィン、
(9) 準アフィン、
(10) （局所）有限型、
(11) （局所）準有限、
(12) （局所）有限表示、
(13) 相対次元 $d$ の局所有限型、
(14) 普遍開、
(15) 平坦、
(16) syntomic、
(17) smooth、
(18) 非分岐（それぞれ G-非分岐）、
(19) エタール、
(20) 固有、
(21) 有限または整、
(22) 有限局所自由、
(23) universally submersive、
(24) 普遍同相写像、および
(25) 埋め込み。

\medskip\noindent
本章では、必ずしも表現可能でない代数空間の射に対して、
これらの概念を再定義する。その際には常に、曖昧さを避けるため、
新しい定義が旧来の定義と一致することを確かめる。

\medskip\noindent
上の定義は、$X$ がスキームである場合には適用できることに注意する。
実際、スキームから代数空間への射は表現可能である。
特に、$X$ と $Y$ がともにスキームである場合にも適用できる。
Spaces, Lemma
\ref{spaces-lemma-morphism-schemes-gives-representable-transformation-property}
では、この場合に各定義が一致し、曖昧さが生じないことを見た。

\medskip\noindent
さらに、Spaces, Lemma
\ref{spaces-lemma-base-change-representable-transformations-property}
では、このように定義された代数空間の表現可能な射の性質が、
代数空間の射による任意の基底変換で安定であることを見た。
最後に、Spaces, Lemmas
\ref{spaces-lemma-composition-representable-transformations-property} および
\ref{spaces-lemma-product-representable-transformations-property}
では、$\mathcal{P}$ が合成で安定であれば、すなわち上の性質
(1)(a), (1)(b), (1)(c), (2) -- (25) のうち (13) を除くものについては、
表現可能な射の積が性質 $\mathcal{P}$ を保ち、表現可能な射の合成も
性質 $\mathcal{P}$ を保つことを見た。

\medskip\noindent
以下ではこれらの事実を用い、そのたびに単に本節を参照する。














\section{分離公理}
\label{spaces-morphisms-section-separation-axioms}

\noindent
代数空間の射の対角射に関して、先験的ないくつかの性質を
列挙しておくことには意味がある。

\begin{lemma}
\label{spaces-morphisms-lemma-properties-diagonal}
$S$ を $\Sch_{fppf}$ に含まれるスキームとする。
$f : X \to Y$ を $S$ 上の代数空間の射とする。
$\Delta_{X/Y} : X \to X \times_Y X$ を対角射とする。このとき
\begin{enumerate}
\item $\Delta_{X/Y}$ は表現可能である。
\item $\Delta_{X/Y}$ は局所有限型である。
\item $\Delta_{X/Y}$ は単射である。
\item $\Delta_{X/Y}$ は分離的である。
\item $\Delta_{X/Y}$ は局所準有限である。
\end{enumerate}
\end{lemma}

\begin{proof}
$\Delta_{X/S}$ が（代数空間の定義により）表現可能であり、
さらに性質 (2) -- (5) を満たすという事実を用いる。
Spaces, Lemma \ref{spaces-lemma-properties-diagonal} を参照。
対角射 $\Delta_{X/S} : X \to X \times_S X$ は
$$
X
\longrightarrow
X \times_Y X
\longrightarrow
X \times_S X
$$
と分解することに注意する。$X \times_Y X \to X \times_S X$ は
単射であり、$\Delta_{X/S}$ は表現可能なので、形式的に
$\Delta_{X/Y}$ が表現可能であることが従う。特に残りの主張にも
意味がある。Section \ref{spaces-morphisms-section-representable} を参照。

\medskip\noindent
$U$ がスキームであるような全射エタール射 $U \to X$ を選ぶ。
図式
$$
\xymatrix{
R = U \times_X U \ar[r] \ar[d] &
U \times_Y U \ar[d] \ar[r] &
U \times_S U \ar[d] \\
X \ar[r] & X \times_Y X \ar[r] & X \times_S X
}
$$
を考える。二つの正方形はともに Cartesian であり、したがって外側の
長方形もそうである。上段はスキームからなり、垂直矢印は全射エタール射である。
Spaces, Lemma \ref{spaces-lemma-representable-morphisms-spaces-property}
により、$\Delta_{X/Y}$ の性質 (2) -- (5) は
$R \to U \times_Y U$ の同じ性質と同値である。
Spaces, Lemma \ref{spaces-lemma-properties-diagonal} の証明で、
$R \to U \times_S U$ が性質 (2) -- (5) をもつことを見た。
スキームの射 $U \times_Y U \to U \times_S U$ は単射である。
これらの事実から、$R \to U \times_Y U$ が性質 (2) -- (5) をもつことが従う。

\medskip\noindent
具体的には次のとおりである。(3) については、合成
$R \to U \times_S U$ が単射であるため、$R \to U \times_Y U$
も単射である。(2) については、合成 $R \to U \times_S U$ が
局所有限型なので（Morphisms, Lemma
\ref{morphisms-lemma-permanence-finite-type}）、
$R \to U \times_Y U$ も局所有限型である。
局所有限型の単射は有限なファイバーをもつため局所準有限であり
（Morphisms, Lemma \ref{morphisms-lemma-finite-fibre}）、したがって (5) が従う。
単射は分離的なので（Schemes, Lemma
\ref{schemes-lemma-monomorphism-separated}）、(4) が従う。
\end{proof}

\begin{definition}
\label{spaces-morphisms-definition-separated}
$S$ をスキームとする。
$f : X \to Y$ を $S$ 上の代数空間の射とする。
$\Delta_{X/Y} : X \to X \times_Y X$ を対角射とする。
\begin{enumerate}
\item $\Delta_{X/Y}$ が閉埋め込みであるとき、$f$ は {\it 分離的}であるという。
\item $\Delta_{X/Y}$ が埋め込みであるとき、$f$ は
{\it 局所分離的}\footnote{文献では、この用語はしばしば準分離的かつ
局所分離的な射を指す。}であるという。
\item $\Delta_{X/Y}$ が準コンパクトであるとき、$f$ は
{\it 準分離的}であるという。
\end{enumerate}
\end{definition}

\noindent
この定義には意味がある。実際、$\Delta_{X/Y}$ は表現可能なので、
定義中の各性質をもつとは何かが分かっている。
以下で（Lemma \ref{spaces-morphisms-lemma-match-separated}）、この定義がスキームの射および
表現可能な射に対して既にある定義と一致することを見る。

\begin{lemma}
\label{spaces-morphisms-lemma-trivial-implications}
$S$ をスキームとする。$f : X \to Y$ を $S$ 上の代数空間の射とする。
$f$ が分離的ならば、$f$ は局所分離的であり、$f$ は準分離的である。
\end{lemma}

\begin{proof}
これは一般原理 Spaces, Lemma
\ref{spaces-lemma-representable-transformations-property-implication}
から従う。実際、スキームの閉埋め込みは埋め込みであり、準コンパクトである。
\end{proof}

\begin{lemma}
\label{spaces-morphisms-lemma-base-change-separated}
Definition \ref{spaces-morphisms-definition-separated} に挙げたすべての分離公理は、
基底変換で安定である。
\end{lemma}

\begin{proof}
$f : X \to Y$ と $Y' \to Y$ を代数空間の射とする。
$f' : X' \to Y'$ を $Y' \to Y$ による $f$ の基底変換とする。このとき
$\Delta_{X'/Y'}$ は、射 $X' \times_{Y'} X' \to X \times_Y X$ による
$\Delta_{X/Y}$ の基底変換である。Section \ref{spaces-morphisms-section-representable}
の結果により、Definition \ref{spaces-morphisms-definition-separated} で用いた対角射の
各性質は基底変換で安定である。よって補題が成り立つ。
\end{proof}

\begin{lemma}
\label{spaces-morphisms-lemma-fibre-product-after-map}
\begin{slogan}
代数空間の「magic diagram」の上辺は、埋め込みに似た良い性質をもち、
分離性の仮定のもとではそれらの性質はさらに強くなる。
\end{slogan}
$S$ をスキームとする。$f : X \to Z$, $g : Y \to Z$ および $Z \to T$
を $S$ 上の代数空間の射とする。誘導される射
$i : X \times_Z Y \to X \times_T Y$ を考える。このとき
\begin{enumerate}
\item $i$ は表現可能、局所有限型、局所準有限、分離的、かつ単射である。
\item $Z \to T$ が局所分離的ならば、$i$ は埋め込みである。
\item $Z \to T$ が分離的ならば、$i$ は閉埋め込みである。
\item $Z \to T$ が準分離的ならば、$i$ は準コンパクトである。
\end{enumerate}
\end{lemma}

\begin{proof}
一般の圏論により、図式
$$
\xymatrix{
X \times_Z Y \ar[r]_i \ar[d] & X \times_T Y \ar[d] \\
Z \ar[r]^-{\Delta_{Z/T}} \ar[r] & Z \times_T Z
}
$$
はファイバー積図式である。したがって $i$ は対角射
$\Delta_{Z/T}$ の基底変換である。よって補題は
Lemma \ref{spaces-morphisms-lemma-properties-diagonal} と
Section \ref{spaces-morphisms-section-representable} の内容から従う。
\end{proof}

\begin{lemma}
\label{spaces-morphisms-lemma-semi-diagonal}
\begin{slogan}
代数空間の射のグラフがもつ性質は、終域の分離性質から従う。
\end{slogan}
$S$ をスキームとする。$T$ を $S$ 上の代数空間とする。
$g : X \to Y$ を $T$ 上の代数空間の射とする。
$g$ のグラフ $i : X \to X \times_T Y$ を考える。このとき
\begin{enumerate}
\item $i$ は表現可能、局所有限型、局所準有限、分離的、かつ単射である。
\item $Y \to T$ が局所分離的ならば、$i$ は埋め込みである。
\item $Y \to T$ が分離的ならば、$i$ は閉埋め込みである。
\item $Y \to T$ が準分離的ならば、$i$ は準コンパクトである。
\end{enumerate}
\end{lemma}

\begin{proof}
これは Lemma \ref{spaces-morphisms-lemma-fibre-product-after-map} を
射 $X = X \times_Y Y \to X \times_T Y$ に適用した特別な場合である。
\end{proof}

\begin{lemma}
\label{spaces-morphisms-lemma-section-immersion}
$S$ をスキームとする。
$f : X \to T$ を $S$ 上の代数空間の射とする。
$s : T \to X$ を $f$ の切断とする（式では $f \circ s = \text{id}_T$）。このとき
\begin{enumerate}
\item $s$ は表現可能、局所有限型、局所準有限、分離的、かつ単射である。
\item $f$ が局所分離的ならば、$s$ は埋め込みである。
\item $f$ が分離的ならば、$s$ は閉埋め込みである。
\item $f$ が準分離的ならば、$s$ は準コンパクトである。
\end{enumerate}
\end{lemma}

\begin{proof}
これは Lemma \ref{spaces-morphisms-lemma-semi-diagonal} を $g = s$、したがって
$i = s : T \to T \times_T X$ に適用した特別な場合である。
\end{proof}

\begin{lemma}
\label{spaces-morphisms-lemma-composition-separated}
Definition \ref{spaces-morphisms-definition-separated} に挙げたすべての分離公理は、
射の合成で安定である。
\end{lemma}

\begin{proof}
$f : X \to Y$ と $g : Y \to Z$ を、考えている公理が適用される
代数空間の射とする。対角射 $\Delta_{X/Z}$ は合成
$$
X \longrightarrow X \times_Y X \longrightarrow X \times_Z X.
$$
である。考えている分離公理は、対角射がある性質 $\mathcal{P}$ を
もつことによって定義される。上の Lemma
\ref{spaces-morphisms-lemma-fibre-product-after-map} により、第二の矢印もこの性質をもつ。
したがって、性質 $\mathcal{P}$ をもつ（表現可能な）射の合成も
性質 $\mathcal{P}$ をもつため、補題が従う。
Section \ref{spaces-morphisms-section-representable} を参照。
\end{proof}

\begin{lemma}
\label{spaces-morphisms-lemma-separated-over-separated}
$S$ をスキームとする。
$f : X \to Y$ を $S$ 上の代数空間の射とする。
\begin{enumerate}
\item $Y$ が分離的で $f$ が分離的ならば、$X$ は分離的である。
\item $Y$ が準分離的で $f$ が準分離的ならば、$X$ は準分離的である。
\item $Y$ が局所分離的で $f$ が局所分離的ならば、$X$ は局所分離的である。
\item $Y$ が $S$ 上分離的で $f$ が分離的ならば、$X$ は $S$ 上分離的である。
\item $Y$ が $S$ 上準分離的で $f$ が準分離的ならば、$X$ は $S$ 上準分離的である。
\item $Y$ が $S$ 上局所分離的で $f$ が局所分離的ならば、
$X$ は $S$ 上局所分離的である。
\end{enumerate}
\end{lemma}

\begin{proof}
(4), (5), (6) は Lemma \ref{spaces-morphisms-lemma-composition-separated} と
Spaces, Definition \ref{spaces-definition-separated} から直ちに従う。
(1), (2), (3) は、$X$ と $Y$ を $\Spec(\mathbf{Z})$ 上の代数空間と
みなすことにより (4), (5), (6) に帰着する。
Properties of Spaces, Definition
\ref{spaces-properties-definition-separated} を参照。
\end{proof}

\begin{lemma}
\label{spaces-morphisms-lemma-compose-after-separated}
$S$ をスキームとする。
$f : X \to Y$ と $g : Y \to Z$ を $S$ 上の代数空間の射とする。
\begin{enumerate}
\item $g \circ f$ が分離的ならば、$f$ も分離的である。
\item $g \circ f$ が局所分離的ならば、$f$ も局所分離的である。
\item $g \circ f$ が準分離的ならば、$f$ も準分離的である。
\end{enumerate}
\end{lemma}

\begin{proof}
$g \circ f$ の対角射の分解
$$
X \to X \times_Y X \to X \times_Z X
$$
を考える。いずれの場合にも最後の射は単射である。したがって、
任意のスキーム $T$ と射 $T \to X \times_Y X$ に対して等式
$$
X \times_{(X \times_Y X)} T = X \times_{(X \times_Z X)} T.
$$
を得る。よって結果は明らかである。
\end{proof}

\begin{lemma}
\label{spaces-morphisms-lemma-separated-implies-morphism-separated}
$S$ をスキームとし、$X$ を $S$ 上の代数空間とする。
\begin{enumerate}
\item $X$ が分離的ならば、$X$ は $S$ 上分離的である。
\item $X$ が局所分離的ならば、$X$ は $S$ 上局所分離的である。
\item $X$ が準分離的ならば、$X$ は $S$ 上準分離的である。
\end{enumerate}
$f : X \to Y$ を $S$ 上の代数空間の射とする。
\begin{enumerate}
\item[(4)] $X$ が $S$ 上分離的ならば、$f$ は分離的である。
\item[(5)] $X$ が $S$ 上局所分離的ならば、$f$ は局所分離的である。
\item[(6)] $X$ が $S$ 上準分離的ならば、$f$ は準分離的である。
\end{enumerate}
\end{lemma}

\begin{proof}
(4), (5), (6) は Lemma \ref{spaces-morphisms-lemma-compose-after-separated} と
Spaces, Definition \ref{spaces-definition-separated} から直ちに従う。
(1), (2), (3) は、$X$ と $Y$ を $\Spec(\mathbf{Z})$ 上の代数空間と
みなすことにより (4), (5), (6) から従う。
Properties of Spaces, Definition
\ref{spaces-properties-definition-separated} を参照。
\end{proof}

\begin{lemma}
\label{spaces-morphisms-lemma-separated-local}
$S$ をスキームとする。
$f : X \to Y$ を $S$ 上の代数空間の射とする。
$\mathcal{P}$ を Definition \ref{spaces-morphisms-definition-separated} の分離公理の
いずれかとする。次は同値である。
\begin{enumerate}
\item $f$ は $\mathcal{P}$ である。
\item 任意のスキーム $Z$ と射 $Z \to Y$ に対して、$f$ の基底変換
$Z \times_Y X \to Z$ は $\mathcal{P}$ である。
\item 任意のアフィンスキーム $Z$ と任意の射 $Z \to Y$ に対して、
$f$ の基底変換 $Z \times_Y X \to Z$ は $\mathcal{P}$ である。
\item 任意のアフィンスキーム $Z$ と任意の射 $Z \to Y$ に対して、
代数空間 $Z \times_Y X$ は $\mathcal{P}$ である（Properties of Spaces,
Definition \ref{spaces-properties-definition-separated} を参照）。
\item あるスキーム $V$ と全射エタール射 $V \to Y$ が存在して、
基底変換 $V \times_Y X \to V$ が $\mathcal{P}$ をもつ。
\item Zariski 被覆 $Y = \bigcup Y_i$ が存在して、各射
$f^{-1}(Y_i) \to Y_i$ が $\mathcal{P}$ をもつ。
\end{enumerate}
\end{lemma}

\begin{proof}
以下では Lemma \ref{spaces-morphisms-lemma-base-change-separated} を繰り返し用いるが、
その都度は断らない。特に、(1) が (2) を含み、(2) が (3) を含むことは明らかである。

\medskip\noindent
(3) と (4) が同値であることを示す。$Z$ がアフィンスキームならば、
射 $Z \to \Spec(\mathbf{Z})$ は $\Spec(\mathbf{Z})$ 上の代数空間の射として
分離的であることに注意する。$Z \times_Y X \to Z$ が $\mathcal{P}$ ならば、
合成として $Z \times_Y X \to \Spec(\mathbf{Z})$ も $\mathcal{P}$ である
（Lemma \ref{spaces-morphisms-lemma-composition-separated} を参照）。
したがって、代数空間 $Z \times_Y X$ は $\mathcal{P}$ である。
逆に、代数空間 $Z \times_Y X$ が $\mathcal{P}$ ならば、
$Z \times_Y X \to \Spec(\mathbf{Z})$ は $\mathcal{P}$ であり、
Lemma \ref{spaces-morphisms-lemma-compose-after-separated} により
$Z \times_Y X \to Z$ は $\mathcal{P}$ である。

\medskip\noindent
(3) が (5) を含むことを示す。(3) を仮定する。$V$ をスキームとし、
$V \to Y$ を全射エタールとする。$V \times_Y X \to V$ が性質
$\mathcal{P}$ をもつことを示さなければならない。言い換えると、射
$$
V \times_Y X \longrightarrow
(V \times_Y X) \times_V (V \times_Y X) = V \times_Y X \times_Y X
$$
が対応する性質（すなわち、閉埋め込み、埋め込み、または準コンパクト）を
もつことを示せばよい。$V = \bigcup V_j$ を $V$ のアフィン開被覆とする。
仮定により、各射
$$
V_j \times_Y X \longrightarrow V_j \times_Y X \times_Y X
$$
は対応する性質をもつ。閉埋め込み、埋め込み、準コンパクトな埋め込み、
または準コンパクトであることは終域上 Zariski 局所的であり、
$V_j$ は $V$ を被覆するので、所望の結論を得る。

\medskip\noindent
(5) が (1) を含むことを示す。$V \to Y$ を (5) のとおりとする。
このときファイバー積図式
$$
\xymatrix{
V \times_Y X \ar[r] \ar[d] &
X \ar[d] \\
V \times_Y X \times_Y X \ar[r] &
X \times_Y X
}
$$
を得る。仮定により左の垂直矢印は、閉埋め込み、埋め込み、
準コンパクトな埋め込み、または準コンパクトである。
Spaces, Lemma \ref{spaces-lemma-descent-representable-transformations-property}
から、右の垂直矢印も、閉埋め込み、埋め込み、準コンパクトな埋め込み、
または準コンパクトであることが従う。

\medskip\noindent
$Y = Y$ という被覆を取れば、(1) が (6) を含むことは明らかである。
$Y = \bigcup Y_i$ を (6) のとおりと仮定する。スキーム $V_i$ と
全射エタール射 $V_i \to Y_i$ を選ぶ。射 $V_i \times_Y X \to V_i$ は、
射 $f^{-1}(Y_i) \to Y_i$ の基底変換なので $\mathcal{P}$ をもつことに注意する。
$V = \coprod V_i$ とおく。このとき $V \to Y$ は (5) のような射である
（細部は省略する）。したがって (6) は (5) を含み、証明は完了する。
\end{proof}

\begin{lemma}
\label{spaces-morphisms-lemma-match-separated}
$S$ をスキームとする。
$f : X \to Y$ を $S$ 上の代数空間の表現可能な射とする。
\begin{enumerate}
\item 射 $f$ は局所分離的である。
\item 射 $f$ が上の Definition \ref{spaces-morphisms-definition-separated} の意味で
（準）分離的であることと、Section \ref{spaces-morphisms-section-representable} の意味で
$f$ が（準）分離的であることとは同値である。
\end{enumerate}
特に、$f : X \to Y$ が $S$ 上のスキームの射ならば、
$f$ が Definition \ref{spaces-morphisms-definition-separated} の意味で（準）分離的であることと、
$f$ がスキームの射として（準）分離的であることとは同値である。
\end{lemma}

\begin{proof}
これは Lemma \ref{spaces-morphisms-lemma-separated-local} の (1) と (2) の同値性を、
任意のスキームの射が局所分離的であるという事実と組み合わせたものである。
Schemes, Lemma \ref{schemes-lemma-diagonal-immersion} を参照。
\end{proof}











\section{全射な射}
\label{spaces-morphisms-section-surjective}

\noindent
代数空間の表現可能な射が全射であるとは何かは、既に
Section \ref{spaces-morphisms-section-representable} で定義した。

\begin{lemma}
\label{spaces-morphisms-lemma-surjective-representable}
$S$ をスキームとし、$f : X \to Y$ を $S$ 上の代数空間の表現可能な射とする。
このとき、$f$ が（Section \ref{spaces-morphisms-section-representable} の意味で）全射であることと、
$|f| : |X| \to |Y|$ が全射であることとは同値である。
\end{lemma}

\begin{proof}
実際、$f : X \to Y$ が表現可能ならば、それが全射であることと、
任意のスキーム $T$ と任意の射 $T \to Y$ に対して、$f$ の基底変換
$f_T : T \times_Y X \to T$ がスキームの全射射であること、言い換えると
$|f_T|$ が全射であることとは同値である。Properties of Spaces, Lemma
\ref{spaces-properties-lemma-points-cartesian} により、写像
$|T \times_Y X| \to |T| \times_{|Y|} |X|$ は常に全射である。
したがって、$|f| : |X| \to |Y|$ が全射ならば
$|f_T| : |T \times_Y X| \to |T|$ は全射である。
逆に、上の任意の $T \to Y$ に対して $|f_T|$ が全射ならば、
$T$ を体のスペクトルとして取ることにより、$|X| \to |Y|$ が全射であると結論する。
\end{proof}

\noindent
これにより、次の定義を与える準備が整った。

\begin{definition}
\label{spaces-morphisms-definition-surjective}
$S$ をスキームとし、$f : X \to Y$ を $S$ 上の代数空間の射とする。
付随する位相空間の写像 $|f| : |X| \to |Y|$ が全射であるとき、
$f$ は {\it 全射}であるという。
\end{definition}

\begin{lemma}
\label{spaces-morphisms-lemma-surjective-local}
$S$ をスキームとする。
$f : X \to Y$ を $S$ 上の代数空間の射とする。次は同値である：
\begin{enumerate}
\item $f$ は全射である。
\item 任意のスキーム $Z$ と任意の射 $Z \to Y$ に対して、射
$Z \times_Y X \to Z$ は全射である。
\item 任意のアフィンスキーム $Z$ と任意の射 $Z \to Y$ に対して、射
$Z \times_Y X \to Z$ は全射である。
\item あるスキーム $V$ と全射エタール射 $V \to Y$ が存在して、
$V \times_Y X \to V$ が全射である。
\item あるスキーム $U$ と全射エタール射 $\varphi : U \to X$ が存在して、
合成 $f \circ \varphi$ が全射である。
\item 可換図式
$$
\xymatrix{
U \ar[d] \ar[r] & V \ar[d] \\
X \ar[r] & Y
}
$$
が存在し、$U$, $V$ はスキーム、垂直矢印は全射エタールであり、
上の水平矢印は全射である。
\item Zariski 被覆 $Y = \bigcup Y_i$ が存在して、各射
$f^{-1}(Y_i) \to Y_i$ は全射である。
\end{enumerate}
\end{lemma}

\begin{proof}
省略する。
\end{proof}












\begin{lemma}
\label{spaces-morphisms-lemma-composition-surjective}
全射の合成は全射である。
\end{lemma}

\begin{proof}
定義から直ちに従う。
\end{proof}

\begin{lemma}
\label{spaces-morphisms-lemma-base-change-surjective}
全射の基底変換は全射である。
\end{lemma}

\begin{proof}
Properties of Spaces, Lemma
\ref{spaces-properties-lemma-points-cartesian} から直ちに従う。
\end{proof}










\section{開射}
\label{spaces-morphisms-section-open}

\noindent
代数空間の表現可能な射が普遍開であるとは何かは、既に
Section \ref{spaces-morphisms-section-representable} で定義した。そこで自然な定義を
与える前に、表現可能な場合にこの定義と一致することを確かめる。

\begin{lemma}
\label{spaces-morphisms-lemma-characterize-representable-universally-open}
$S$ をスキームとし、$f : X \to Y$ を $S$ 上の代数空間の表現可能な射とする。
次は同値である。
\begin{enumerate}
\item $f$ は（Section \ref{spaces-morphisms-section-representable} の意味で）普遍開である。
\item 代数空間の任意の射 $Z \to Y$ に対して、位相空間の写像
$|Z \times_Y X| \to |Z|$ は開である。
\end{enumerate}
\end{lemma}

\begin{proof}
(1) を仮定し、$Z \to Y$ を (2) のとおりとする。スキーム $V$ と
全射エタール射 $V \to Y$ を選ぶ。仮定により、スキームの射
$V \times_Y X \to V$ は普遍開である。Properties of Spaces,
Section \ref{spaces-properties-section-points} により、可換図式
$$
\xymatrix{
|V \times_Y X| \ar[r] \ar[d] & |Z \times_Y X| \ar[d] \\
|V| \ar[r] & |Z|
}
$$
の水平矢印は開かつ全射であり、さらに
$$
|V \times_Y X| \longrightarrow |V| \times_{|Z|} |Z \times_Y X|
$$
は全射である。したがって、左の垂直矢印が開であることから、
右の垂直矢印も開である。これで (2) が示された。
(2) $\Rightarrow$ (1) は定義から直ちに従う。
\end{proof}

\noindent
よって次の自然な定義を用いることができる。

\begin{definition}
\label{spaces-morphisms-definition-open}
$S$ をスキームとし、$f : X \to Y$ を $S$ 上の代数空間の射とする。
\begin{enumerate}
\item 位相空間の写像 $|f| : |X| \to |Y|$ が開であるとき、
$f$ は {\it 開}であるという。
\item 代数空間の任意の射 $Z \to Y$ に対して、位相空間の写像
$$
|Z \times_Y X| \to |Z|
$$
が開であるとき、すなわち基底変換 $Z \times_Y X \to Z$ が開であるとき、
$f$ は {\it 普遍開}であるという。
\end{enumerate}
\end{definition}

\noindent
代数空間のエタール射は普遍開であることに注意する。
Properties of Spaces, Definition
\ref{spaces-properties-definition-etale} および Lemmas
\ref{spaces-properties-lemma-etale-open} と
\ref{spaces-properties-lemma-base-change-etale} を参照。

\begin{lemma}
\label{spaces-morphisms-lemma-base-change-universally-open}
代数空間の普遍開射を代数空間の任意の射で基底変換した射は、普遍開である。
\end{lemma}

\begin{proof}
定義から直ちに従う。
\end{proof}

\begin{lemma}
\label{spaces-morphisms-lemma-composition-open}
代数空間の二つの（普遍）開射の合成は（普遍）開である。
\end{lemma}

\begin{proof}
省略する。
\end{proof}

\begin{lemma}
\label{spaces-morphisms-lemma-universally-open-local}
$S$ をスキームとし、$f : X \to Y$ を $S$ 上の代数空間の射とする。
次は同値である。
\begin{enumerate}
\item $f$ は普遍開である。
\item 任意のスキーム $Z$ と任意の射 $Z \to Y$ に対して、射影
$|Z \times_Y X| \to |Z|$ は開である。
\item 任意のアフィンスキーム $Z$ と任意の射 $Z \to Y$ に対して、射影
$|Z \times_Y X| \to |Z|$ は開である。
\item あるスキーム $V$ と全射エタール射 $V \to Y$ が存在して、
$V \times_Y X \to V$ は代数空間の普遍開射である。
\item Zariski 被覆 $Y = \bigcup Y_i$ が存在して、各射
$f^{-1}(Y_i) \to Y_i$ は普遍開である。
\end{enumerate}
\end{lemma}

\begin{proof}
(1) が (2) を含み、(2) が (3) を含むことの証明は省略する。

\medskip\noindent
(3) を仮定する。全射エタール射 $V \to Y$ を選ぶ。
$V \times_Y X \to V$ が代数空間の普遍開射であることを示す。
$Z \to V$ を代数空間から $V$ への射とする。$W = \coprod W_i$ が
アフィンスキームの非交和であるような全射エタール射 $W \to Z$ を取る。
Properties of Spaces, Lemma
\ref{spaces-properties-lemma-cover-by-union-affines} を参照。
このとき可換図式
$$
\xymatrix{
\coprod_i |W_i \times_Y X| \ar@{=}[r] \ar[d] &
|W \times_Y X| \ar[r] \ar[d] &
|Z \times_Y X| \ar[d] \ar@{=}[r] &
|Z \times_V (V \times_Y X)| \ar[ld] \\
\coprod |W_i| \ar@{=}[r] &
|W| \ar[r] &
|Z|
}
$$
を得る。南東向きの矢印が開であることを示せばよい。
中央の水平矢印は全射かつ開である（Properties of Spaces, Lemma
\ref{spaces-properties-lemma-etale-open}）。仮定 (3) と $W_i$ が
アフィンであることから、左の垂直矢印は開である。
したがって右の垂直矢印も開である。

\medskip\noindent
$V \to Y$ を (4) のとおりと仮定する。$f$ が普遍開であることを示す。
$Z \to Y$ を代数空間の射とし、図式
$$
\xymatrix{
|(V \times_Y Z) \times_V (V \times_Y X)| \ar@{=}[r] \ar[rd] &
|V \times_Y X| \ar[r] \ar[d] &
|Z \times_Y X| \ar[d] \\
 &
|V \times_Y Z| \ar[r] &
|Z|
}
$$
を考える。南西向きの矢印は仮定により開である。
対応する代数空間の射がエタールなので、水平矢印は全射かつ開である
（Properties of Spaces, Lemma
\ref{spaces-properties-lemma-etale-open} を参照）。
したがって右の垂直矢印は開である。

\medskip\noindent
もちろん、被覆 $Y = Y$ を取れば (1) は (5) を含む。
$Y = \bigcup Y_i$ を (5) のとおりと仮定する。任意の $Z \to Y$ に対して、
対応する Zariski 被覆 $Z = \bigcup Z_i$ を得て、$f$ の $Z_i$ への
基底変換は開である。簡単な位相的議論により
$Z \times_Y X \to Z$ は開である。よって (1) が成り立つ。
\end{proof}

\begin{lemma}
\label{spaces-morphisms-lemma-space-over-field-universally-open}
$S$ をスキームとする。$p : X \to \Spec(k)$ を $S$ 上の代数空間の射とし、
$k$ は体であるとする。このとき $p : X \to \Spec(k)$ は普遍開である。
\end{lemma}

\begin{proof}
スキーム $U$ と全射エタール射 $U \to X$ を選ぶ。
合成 $U \to \Spec(k)$ は、スキームの射として普遍開である。
Morphisms, Lemma
\ref{morphisms-lemma-scheme-over-field-universally-open} を参照。
$Z \to \Spec(k)$ をスキームの射とする。このとき
$U \times_{\Spec(k)} Z \to X \times_{\Spec(k)} Z$ は全射である。
Lemma \ref{spaces-morphisms-lemma-base-change-surjective} を参照。
したがって写像
$$
|U \times_{\Spec(k)} Z| \to |X \times_{\Spec(k)} Z| \to |Z|
$$
の第一の写像は全射である。上により合成は開なので、
第二の写像も開であると結論する。したがって Lemma
\ref{spaces-morphisms-lemma-universally-open-local} により $p$ は普遍開である。
\end{proof}









\section{Submersive な射}
\label{spaces-morphisms-section-submersive}

\noindent
代数空間の表現可能な射については、Section
\ref{spaces-morphisms-section-representable} において、universally submersive であることの
意味をすでに定義した。そこで自然な定義を与える前に、表現可能な場合には
それがこの定義と一致することを確かめる。

\begin{lemma}
\label{spaces-morphisms-lemma-characterize-representable-universally-submersive}
$S$ をスキームとし、$f : X \to Y$ を $S$ 上の代数空間の表現可能な射とする。
次は同値である。
\begin{enumerate}
\item $f$ は universally submersive である
（Section \ref{spaces-morphisms-section-representable} の意味で）。
\item 代数空間の任意の射 $Z \to Y$ に対して、位相空間の写像
$|Z \times_Y X| \to |Z|$ は submersive である。
\end{enumerate}
\end{lemma}

\begin{proof}
(1) を仮定し、$Z \to Y$ を (2) のとおりとする。スキーム $V$ と
全射エタール射 $V \to Y$ を選ぶ。仮定により、スキームの射
$V \times_Y X \to V$ は universally submersive である。
Properties of Spaces, Section
\ref{spaces-properties-section-points} により、可換図式
$$
\xymatrix{
|V \times_Y X| \ar[r] \ar[d] & |Z \times_Y X| \ar[d] \\
|V| \ar[r] & |Z|
}
$$
において水平矢印は開かつ全射であり、さらに
$$
|V \times_Y X| \longrightarrow |V| \times_{|Z|} |Z \times_Y X|
$$
は全射である。したがって左の垂直矢印が submersive であることから、
右の垂直矢印も submersive である。(2) が示された。
(2) $\Rightarrow$ (1) は定義から直ちに従う。
\end{proof}

\noindent
したがって、次の自然な定義を用いることができる。

\begin{definition}
\label{spaces-morphisms-definition-submersive}
$S$ をスキームとし、$f : X \to Y$ を $S$ 上の代数空間の射とする。
\begin{enumerate}
\item 連続写像 $|X| \to |Y|$ が submersive であるとき、$f$ は
{\it submersive} であるという\footnote{これは微分多様体における
submersion の概念とは大きく異なる。}。Topology, Definition
\ref{topology-definition-submersive} を参照。
\item 代数空間の任意の射 $Y' \to Y$ に対して基底変換
$Y' \times_Y X \to Y'$ が submersive であるとき、$f$ は
{\it universally submersive} であるという。
\end{enumerate}
\end{definition}

\noindent
submersive な射は、とくに全射であることに注意する。

\begin{lemma}
\label{spaces-morphisms-lemma-base-change-universally-submersive}
代数空間の universally submersive な射を代数空間の任意の射で
基底変換したものは universally submersive である。
\end{lemma}

\begin{proof}
これは定義から直ちに従う。
\end{proof}

\begin{lemma}
\label{spaces-morphisms-lemma-composition-universally-submersive}
代数空間の (universally) submersive な二つの射の合成は
(universally) submersive である。
\end{lemma}

\begin{proof}
省略する。
\end{proof}












\section{準コンパクト射}
\label{spaces-morphisms-section-quasi-compact}

\noindent
Section \ref{spaces-morphisms-section-representable} により、代数空間の表現可能な射が
準コンパクトであることの意味はわかっている。代数空間の一般の射について
定義を述べるため、次のことを確認する。

\begin{lemma}
\label{spaces-morphisms-lemma-characterize-representable-quasi-compact}
$S$ をスキームとし、$f : X \to Y$ を $S$ 上の代数空間の表現可能な射とする。
次は同値である。
\begin{enumerate}
\item $f$ は準コンパクトである
（Section \ref{spaces-morphisms-section-representable} の意味で）。
\item 任意の準コンパクト代数空間 $Z$ と任意の射 $Z \to Y$ に対して、
代数空間 $Z \times_Y X$ は準コンパクトである。
\end{enumerate}
\end{lemma}

\begin{proof}
(1) を仮定し、$Z$ が準コンパクトであるような代数空間の射
$Z \to Y$ を取る。Properties of Spaces, Definition
\ref{spaces-properties-definition-quasi-compact} により、準コンパクトな
スキーム $U$ と全射エタール射 $U \to Z$ が存在する。$f$ は表現可能かつ
準コンパクトなので、定義により $U \times_Y X$ はスキームであり、
$U \times_Y X \to U$ は準コンパクトである。したがって
$U \times_Y X$ は準コンパクトなスキームである。射
$U \times_Y X \to Z \times_Y X$ は、表現可能、エタールかつ全射な射
$U \to Z$ の基底変換なので、エタールかつ全射である
（Section \ref{spaces-morphisms-section-representable} を参照）。
ゆえに定義から $Z \times_Y X$ は準コンパクトである。

\medskip\noindent
(2) を仮定する。$Z$ をスキームとして、射 $Z \to Y$ を取る。
$p : Z \times_Y X \to Z$ が準コンパクトであることを示さなければならない。
$U \subset Z$ をアフィン開集合とする。このとき
$p^{-1}(U) = U \times_Y Z$ であり、スキーム $U \times_Y Z$ は
仮定 (2) により準コンパクトである。したがって $p$ は準コンパクトである。
Schemes, Section \ref{schemes-section-quasi-compact} を参照。
\end{proof}

\noindent
これを動機として次のように定義する。

\begin{definition}
\label{spaces-morphisms-definition-quasi-compact}
$S$ をスキームとし、$f : X \to Y$ を $S$ 上の代数空間の射とする。
任意の準コンパクト代数空間 $Z$ と射 $Z \to Y$ に対してファイバー積
$Z \times_Y X$ が準コンパクトであるとき、$f$ は {\it 準コンパクト}
であるという。
\end{definition}

\noindent
上の Lemma \ref{spaces-morphisms-lemma-characterize-representable-quasi-compact} により、
表現可能な代数空間の射について、これは既存の概念と一致する。

\begin{lemma}
\label{spaces-morphisms-lemma-quasi-compact-is-quasi-compact}
$S$ をスキームとする。$f : X \to Y$ が $S$ 上の代数空間の
準コンパクト射ならば、位相空間の基礎写像
$|f| : |X| \to |Y|$ は準コンパクトである。
\end{lemma}

\begin{proof}
$V \subset |Y|$ を準コンパクト開集合とする。Properties of Spaces,
Lemma \ref{spaces-properties-lemma-open-subspaces} により、
$V = |Y'|$ となる開部分空間 $Y' \subset Y$ が存在する。このとき
Properties of Spaces, Lemma
\ref{spaces-properties-lemma-quasi-compact-space} により $Y'$ は
準コンパクトな代数空間であり、したがって Definition
\ref{spaces-morphisms-definition-quasi-compact} により
$X' = Y' \times_Y X$ は準コンパクトな代数空間である。
一方、$X' \subset X$ は開部分空間であり
（Spaces, Lemma \ref{spaces-lemma-base-change-immersions}）、
Properties of Spaces, Lemma
\ref{spaces-properties-lemma-points-cartesian} により
$|X'| = |f|^{-1}(|X'|) = |f|^{-1}(V)$ である。
再び Properties of Spaces, Lemma
\ref{spaces-properties-lemma-quasi-compact-space} を用いると、
$|X'|$ は望みどおり $|X|$ の準コンパクト開集合である。
\end{proof}

\begin{lemma}
\label{spaces-morphisms-lemma-base-change-quasi-compact}
代数空間の準コンパクト射を代数空間の任意の射で基底変換したものは
準コンパクトである。
\end{lemma}

\begin{proof}
省略する。ヒント：ファイバー積の推移性。
\end{proof}

\begin{lemma}
\label{spaces-morphisms-lemma-composition-quasi-compact}
代数空間の二つの準コンパクト射の合成は準コンパクトである。
\end{lemma}

\begin{proof}
省略する。ヒント：ファイバー積の推移性。
\end{proof}

\begin{lemma}
\label{spaces-morphisms-lemma-surjection-from-quasi-compact}
\begin{slogan}
準コンパクト代数空間の全射による像は準コンパクトである。
\end{slogan}
$S$ をスキームとする。
\begin{enumerate}
\item $X \to Y$ が $S$ 上の代数空間の全射であり、
$X$ が準コンパクトならば、$Y$ は準コンパクトである。
\item
$$
\xymatrix{
X \ar[rr]_f \ar[rd]_p & &
Y \ar[dl]^q \\
& Z
}
$$
が $S$ 上の代数空間の射の可換図式であり、$f$ が全射かつ $p$ が
準コンパクトならば、$q$ は準コンパクトである。
\end{enumerate}
\end{lemma}

\begin{proof}
$X$ が準コンパクトで $X \to Y$ が全射であると仮定する。
Definition \ref{spaces-morphisms-definition-surjective} により写像 $|X| \to |Y|$ は
全射なので、Properties of Spaces, Lemma
\ref{spaces-properties-lemma-quasi-compact-space} と、連続写像による
準コンパクト空間の像が準コンパクトであるという位相的事実から、
$Y$ は準コンパクトである。Topology, Lemma
\ref{topology-lemma-image-quasi-compact} を参照。
$f, p, q$ を (2) のとおりとする。始域が準コンパクトな代数空間である
射 $T \to Z$ を取る。仮定により $T \times_Z X$ は準コンパクトである。
Lemma \ref{spaces-morphisms-lemma-base-change-surjective} により、射
$T \times_Z X \to T \times_Z Y$ は全射である。したがって (1) により
$T \times_Z Y$ も準コンパクトである。ゆえに $q$ は準コンパクトである。
\end{proof}

\begin{lemma}
\label{spaces-morphisms-lemma-descent-quasi-compact}
$S$ をスキームとする。$f : X \to Y$ を $S$ 上の代数空間の射とする。
$g : Y' \to Y$ を universally open かつ全射な代数空間の射とし、
基底変換 $f' : X' \to Y'$ が準コンパクトであるとする。
このとき $f$ は準コンパクトである。
\end{lemma}

\begin{proof}
$Z$ が準コンパクトであるような代数空間の射 $Z \to Y$ を取る。
$g$ は universally open かつ全射なので、
$Y' \times_Y Z \to Z$ は開かつ全射である。
$|Y' \times_Y Z|$ の各点は準コンパクト開近傍の基本系をもつので
（Properties of Spaces, Lemma
\ref{spaces-properties-lemma-space-locally-quasi-compact} を参照）、
$Z$ 上へ全射となる準コンパクト開集合
$W \subset |Y' \times_Y Z|$ を取ることができる。
$f'' : W \times_Y X \to W$ を、$W \to Y'$ による $f'$ の
基底変換と書く。仮定により $W \times_Y X$ は準コンパクトである。
$W \to Z$ は全射なので、
$W \times_Y X \to Z \times_Y X$ は全射である。
したがって Lemma \ref{spaces-morphisms-lemma-surjection-from-quasi-compact} により
$Z \times_Y X$ は準コンパクトである。ゆえに $f$ は準コンパクトである。
\end{proof}

\begin{lemma}
\label{spaces-morphisms-lemma-quasi-compact-local}
\begin{slogan}
代数空間の準コンパクト射は引き戻しで保たれ、標的上局所的である。
\end{slogan}
$S$ をスキームとし、$f : X \to Y$ を $S$ 上の代数空間の射とする。
次は同値である。
\begin{enumerate}
\item $f$ は準コンパクトである。
\item 任意のスキーム $Z$ と任意の射 $Z \to Y$ に対して、
代数空間の射 $Z \times_Y X \to Z$ は準コンパクトである。
\item 任意のアフィンスキーム $Z$ と任意の射 $Z \to Y$ に対して、
代数空間 $Z \times_Y X$ は準コンパクトである。
\item あるスキーム $V$ と全射エタール射 $V \to Y$ が存在し、
$V \times_Y X \to V$ は代数空間の準コンパクト射である。
\item 代数空間の全射エタール射 $Y' \to Y$ が存在し、
$Y' \times_Y X \to Y'$ は代数空間の準コンパクト射である。
\item Zariski 被覆 $Y = \bigcup Y_i$ が存在し、各射
$f^{-1}(Y_i) \to Y_i$ は準コンパクトである。
\end{enumerate}
\end{lemma}

\begin{proof}
以下では Lemma \ref{spaces-morphisms-lemma-base-change-quasi-compact} を断りなく用いる。
(1) が (2) を含み、(2) が (3) を含むことは明らかである。
(3) を仮定する。$Z$ を $S$ 上の準コンパクトな代数空間とし、
射 $Z \to Y$ を取る。Properties of Spaces, Lemma
\ref{spaces-properties-lemma-quasi-compact-affine-cover} により、
アフィンスキーム $U$ と全射エタール射 $U \to Z$ が存在する。
このとき $U \times_Y X \to Z \times_Y X$ は代数空間の全射である。
Lemma \ref{spaces-morphisms-lemma-base-change-surjective} を参照。仮定により
$|U \times_Y X|$ は準コンパクトである。これは
$|Z \times_Y X|$ 上へ全射なので、Topology, Lemma
\ref{topology-lemma-image-quasi-compact} により
$|Z \times_Y X|$ は準コンパクトであると結論する。
これで (3) が (1) を含むことが示された。

\medskip\noindent
(1) $\Rightarrow$ (4) および (4) $\Rightarrow$ (5) は明らかである。
(5) $\Rightarrow$ (1) は Lemma \ref{spaces-morphisms-lemma-descent-quasi-compact} と、
代数空間のエタール射が universally open であることから従う
（Definition \ref{spaces-morphisms-definition-open} に続く議論を参照）。

\medskip\noindent
もちろん、被覆 $Y = Y$ を取れば (1) は (6) を含む。
$Y = \bigcup Y_i$ を (6) のとおりと仮定する。$Z$ をアフィンとし、
射 $Z \to Y$ を取る。このとき、各 $Z_j \to Y$ がある $i_j$ に対して
$Y_{i_j}$ を経由するような有限標準アフィン被覆
$Z = Z_1 \cup \ldots \cup Z_n$ が存在する。したがって代数空間
$$
Z_j \times_Y X = Z_j \times_{Y_{i_j}} f^{-1}(Y_{i_j})
$$
は準コンパクトである。
$Z \times_Y X = \bigcup_{j = 1, \ldots, n} Z_j \times_Y X$ は
Zariski 被覆なので、
$|Z \times_Y X| = \bigcup_{j = 1, \ldots, n} |Z_j \times_Y X|$
は準コンパクト空間の有限和、したがって準コンパクトである
（Properties of Spaces, Lemma
\ref{spaces-properties-lemma-open-subspaces} を参照）。
ゆえに (6) は (3) を含む。
\end{proof}

\noindent
次の補題（およびその次の補題）は、とくに $X$ が準コンパクトな
代数空間であれば射 $X \to \Spec(A)$ が準コンパクトであることを保証する。

\begin{lemma}
\label{spaces-morphisms-lemma-quasi-compact-permanence}
$S$ をスキームとし、$f : X \to Y$ および $g : Y \to Z$ を
$S$ 上の代数空間の射とする。$g \circ f$ が準コンパクトで
$g$ が準分離ならば、$f$ は準コンパクトである。
\end{lemma}

\begin{proof}
$f$ は
$(1, f) : X \to X \times_Z Y \to Y$ という合成に等しい。
第一の写像は準分離射 $X \times_Z Y \to X$ の切断なので、
Lemma \ref{spaces-morphisms-lemma-section-immersion} により準コンパクトである
（この射は $g$ の基底変換である。Lemma
\ref{spaces-morphisms-lemma-base-change-separated} を参照）。
第二の写像は $g \circ f$ の基底変換なので、
Lemma \ref{spaces-morphisms-lemma-base-change-quasi-compact} により準コンパクトである。
また、Lemma \ref{spaces-morphisms-lemma-composition-quasi-compact} により
準コンパクト射の合成は準コンパクトである。
\end{proof}

\begin{lemma}
\label{spaces-morphisms-lemma-quasi-compact-quasi-separated-permanence}
$f : X \to Y$ をスキーム $S$ 上の代数空間の射とする。
\begin{enumerate}
\item $X$ が準コンパクトで $Y$ が準分離ならば、$f$ は準コンパクトである。
\item $X$ が準コンパクトかつ準分離で $Y$ が準分離ならば、
$f$ は準コンパクトかつ準分離である。
\item 準コンパクトかつ準分離な代数空間のファイバー積は、
準コンパクトかつ準分離である。
\end{enumerate}
\end{lemma}

\begin{proof}
(1) は $Z = S = \Spec(\mathbf{Z})$ として
Lemma \ref{spaces-morphisms-lemma-quasi-compact-permanence} を用いれば従う。
(2) は (1) と Lemma \ref{spaces-morphisms-lemma-compose-after-separated} から従う。
(3) について、$X \to Y$ および $Z \to Y$ を準コンパクトかつ準分離な
代数空間の射とする。このとき $X \times_Y Z \to Z$ は、
(2) と Lemmas \ref{spaces-morphisms-lemma-base-change-quasi-compact} および
\ref{spaces-morphisms-lemma-base-change-separated} を用いれば、$X \to Y$ の
基底変換として準コンパクトかつ準分離である。
したがって $X \times_Y Z$ は、$Z$ 上準コンパクトかつ準分離な
代数空間として、準コンパクトかつ準分離である。Lemmas
\ref{spaces-morphisms-lemma-separated-over-separated} および
\ref{spaces-morphisms-lemma-composition-quasi-compact} を参照。
\end{proof}


\section{普遍閉射}
\label{spaces-morphisms-section-universally-closed}

\noindent
代数空間の表現可能な射については、Section
\ref{spaces-morphisms-section-representable} において、普遍閉であることの意味を
すでに定義した。そこで自然な定義を与える前に、表現可能な場合には
それがこの定義と一致することを確かめる。

\begin{lemma}
\label{spaces-morphisms-lemma-characterize-representable-universally-closed}
$S$ をスキームとし、$f : X \to Y$ を $S$ 上の代数空間の表現可能な射とする。
次は同値である。
\begin{enumerate}
\item $f$ は普遍閉である
（Section \ref{spaces-morphisms-section-representable} の意味で）。
\item 代数空間の任意の射 $Z \to Y$ に対して、位相空間の写像
$|Z \times_Y X| \to |Z|$ は閉である。
\end{enumerate}
\end{lemma}

\begin{proof}
(1) を仮定し、$Z \to Y$ を (2) のとおりとする。スキーム $V$ と
全射エタール射 $V \to Y$ を選ぶ。仮定により、スキームの射
$V \times_Y X \to V$ は普遍閉である。Properties of Spaces,
Section \ref{spaces-properties-section-points} により、可換図式
$$
\xymatrix{
|V \times_Y X| \ar[r] \ar[d] & |Z \times_Y X| \ar[d] \\
|V| \ar[r] & |Z|
}
$$
において水平矢印は開かつ全射であり、さらに
$$
|V \times_Y X| \longrightarrow |V| \times_{|Z|} |Z \times_Y X|
$$
は全射である。左の垂直矢印は閉なので、右の垂直矢印も閉である。
(2) が示された。(2) $\Rightarrow$ (1) は定義から直ちに従う。
\end{proof}

\noindent
したがって、次の自然な定義を用いることができる。

\begin{definition}
\label{spaces-morphisms-definition-closed}
$S$ をスキームとし、$f : X \to Y$ を $S$ 上の代数空間の射とする。
\begin{enumerate}
\item 位相空間の写像 $|X| \to |Y|$ が閉であるとき、
$f$ は {\it 閉}であるという。
\item 代数空間の任意の射 $Z \to Y$ に対して、位相空間の写像
$$
|Z \times_Y X| \to |Z|
$$
が閉であるとき、すなわち基底変換 $Z \times_Y X \to Z$ が
閉であるとき、$f$ は {\it 普遍閉}であるという。
\end{enumerate}
\end{definition}

\begin{lemma}
\label{spaces-morphisms-lemma-base-change-universally-closed}
代数空間の普遍閉射を代数空間の任意の射で基底変換したものは普遍閉である。
\end{lemma}

\begin{proof}
これは定義から直ちに従う。
\end{proof}

\begin{lemma}
\label{spaces-morphisms-lemma-composition-universally-closed}
代数空間の二つの（普遍）閉射の合成は（普遍）閉である。
\end{lemma}

\begin{proof}
省略する。
\end{proof}

\begin{lemma}
\label{spaces-morphisms-lemma-universally-closed-local}
$S$ をスキームとし、$f : X \to Y$ を $S$ 上の代数空間の射とする。
次は同値である。
\begin{enumerate}
\item $f$ は普遍閉である。
\item 任意のスキーム $Z$ と任意の射 $Z \to Y$ に対して、射影
$|Z \times_Y X| \to |Z|$ は閉である。
\item 任意のアフィンスキーム $Z$ と任意の射 $Z \to Y$ に対して、射影
$|Z \times_Y X| \to |Z|$ は閉である。
\item あるスキーム $V$ と全射エタール射 $V \to Y$ が存在し、
$V \times_Y X \to V$ は代数空間の普遍閉射である。
\item Zariski 被覆 $Y = \bigcup Y_i$ が存在し、各射
$f^{-1}(Y_i) \to Y_i$ は普遍閉である。
\end{enumerate}
\end{lemma}

\begin{proof}
(1) が (2) を含み、(2) が (3) を含むことの証明は省略する。

\medskip\noindent
(3) を仮定する。全射エタール射 $V \to Y$ を選ぶ。
$V \times_Y X \to V$ が代数空間の普遍閉射であることを示す。
$Z \to V$ を代数空間から $V$ への射とする。
$W = \coprod W_i$ がアフィンスキームの非交和であるような
全射エタール射 $W \to Z$ を取る。Properties of Spaces,
Lemma \ref{spaces-properties-lemma-cover-by-union-affines} を参照。
このとき可換図式
$$
\xymatrix{
\coprod_i |W_i \times_Y X| \ar@{=}[r] \ar[d] &
|W \times_Y X| \ar[r] \ar[d] &
|Z \times_Y X| \ar[d] \ar@{=}[r] &
|Z \times_V (V \times_Y X)| \ar[ld] \\
\coprod |W_i| \ar@{=}[r] &
|W| \ar[r] &
|Z|
}
$$
を得る。南東向きの矢印が閉であることを示さなければならない。
中央の水平矢印は全射かつ開である
（Properties of Spaces, Lemma
\ref{spaces-properties-lemma-etale-open}）。
仮定 (3) と $W_i$ がアフィンであることから、左の垂直矢印は閉である。
したがって右の垂直矢印も閉である。

\medskip\noindent
(4) を仮定する。$f$ が普遍閉であることを示す。
$Z \to Y$ を代数空間の射とし、図式
$$
\xymatrix{
|(V \times_Y Z) \times_V (V \times_Y X)| \ar@{=}[r] \ar[rd] &
|V \times_Y X| \ar[r] \ar[d] &
|Z \times_Y X| \ar[d] \\
 &
|V \times_Y Z| \ar[r] &
|Z|
}
$$
を考える。南西向きの矢印は仮定により閉である。
対応する代数空間の射がエタールなので、水平矢印は全射かつ開である
（Properties of Spaces, Lemma
\ref{spaces-properties-lemma-etale-open} を参照）。
したがって右の垂直矢印は閉である。

\medskip\noindent
もちろん、被覆 $Y = Y$ を取れば (1) は (5) を含む。
$Y = \bigcup Y_i$ を (5) のとおりと仮定する。任意の $Z \to Y$ に対して、
対応する Zariski 被覆 $Z = \bigcup Z_i$ を得て、$f$ の $Z_i$ への
基底変換は閉である。簡単な位相的議論により
$Z \times_Y X \to Z$ は閉である。よって (1) が成り立つ。
\end{proof}

\begin{example}
\label{spaces-morphisms-example-strange-universally-closed}
普遍閉射の奇妙な例を挙げる。
$\mathbf{Q} \subset k$ を標数零の体とする。
Spaces, Example \ref{spaces-example-affine-line-translation} と同様に
$X = \mathbf{A}^1_k/\mathbf{Z}$ とする。構造射
$p : X \to \Spec(k)$ は普遍閉であると主張する。
実際、$Z/k$ がスキームで $T \subset |X \times_k Z|$ が閉ならば、
$T$ は $T' \subset |\mathbf{A}^1 \times Z|$ という
$\mathbf{Z}$-不変閉部分集合に対応する。これから $T'$ が
$Z$ のある部分集合 $T''$ の逆像であることは容易に分かる。
Morphisms, Lemma \ref{morphisms-lemma-fpqc-quotient-topology} により、
$T'' \subset Z$ は閉である。もちろん $T''$ は $T$ の像である。
したがって Lemma \ref{spaces-morphisms-lemma-universally-closed-local} により
$p$ は普遍閉である。
\end{example}

\begin{lemma}
\label{spaces-morphisms-lemma-universally-closed-quasi-compact}
$S$ をスキームとする。$S$ 上の代数空間の普遍閉射は準コンパクトである。
\end{lemma}

\begin{proof}
この証明はスキームの場合の証明の繰り返しである。Morphisms, Lemma
\ref{morphisms-lemma-universally-closed-quasi-compact} を参照。
$f : X \to Y$ を $S$ 上の代数空間の射とし、$f$ が準コンパクトでないと
仮定する。目標は $f$ が普遍閉でないことを示すことである。
Lemma \ref{spaces-morphisms-lemma-quasi-compact-local} により、あるアフィンスキーム $Z$ と
射 $Z \to Y$ が存在して $Z \times_Y X \to Z$ は準コンパクトでない。
目標を達成するには $Z \times_Y X \to Z$ が普遍閉でないことを
示せば十分なので、ある環 $B$ に対して $Y = \Spec(B)$ と仮定してよい。

\medskip\noindent
$X_i$ を $X$ の準コンパクト開部分空間として
$X = \bigcup_{i \in I} X_i$ と書く。例えば、$U$ がスキームであるような
全射エタール射 $U \to X$ を選び、アフィン開被覆
$U = \bigcup U_i$ を選び、$X_i \subset X$ を $U_i$ の像とすればよい。
後で、射 $X_i \to Y$ が準コンパクトであることを用いる。
Lemma \ref{spaces-morphisms-lemma-quasi-compact-permanence} を参照。
$T = \Spec(B[a_i ; i \in I])$ とし、$T_i = D(a_i) \subset T$ とする。
$Z \subset T \times_Y X$ を、基礎となる点の閉集合が
$|T \times_Y Z| \setminus \bigcup_{i \in I} |T_i \times_Y X_i|$
である被約閉部分空間とする。Properties of Spaces, Lemma
\ref{spaces-properties-lemma-reduced-closed-subspace} を参照。
（$T_i \times_Y X_i$ は $T \times_Y X$ の開部分空間である。これは
$T_i \to T$ と $X_i \to X$ が開埋め込みだからである。Spaces, Lemmas
\ref{spaces-lemma-base-change-immersions} および
\ref{spaces-lemma-composition-immersions} を参照。）図式は
$$
\xymatrix{
Z \ar[r] \ar[rd] &
T \times_Y X \ar[d]^{f_T} \ar[r]_q &
X \ar[d]^f \\
& T \ar[r]^p & Y
}
$$
である。像 $f_T(|Z|)$ が $|T|$ で閉でないことを示せば十分である。

\medskip\noindent
ある点 $y \in Y$ が存在して、$Y$ における $y$ のアフィン開近傍
$V$ で $X_V$ が準コンパクトとなるものは存在しないと主張する。
そうでなければ有限個のそのような $V$ で $Y$ を被覆でき、各 $V$ に対して
射 $Y_V \to V$ は Lemma \ref{spaces-morphisms-lemma-quasi-compact-permanence} により
準コンパクトであり、次いで Lemma \ref{spaces-morphisms-lemma-quasi-compact-local} により
$f$ は準コンパクトとなって矛盾する。主張のとおりの $y \in Y$ を固定する。

\medskip\noindent
$t \in T$ を $y$ 上の点で、$\kappa(t) = \kappa(y)$ かつすべての $i$ に
対して $\kappa(t)$ で $a_i = 1$ となるものとする。
$f_T(z) = t$ となる $z \in |Z|$ が存在すると仮定する。
このとき、ある $i$ に対して $q(t) \in X_i$ である。
したがって $Z$ の構成により $f_T(z) \not \in T_i$ となるが、
構成により $t \in T_i$ であることと矛盾する。
ゆえに $t \in |T| \setminus f_T(|Z|)$ である。

\medskip\noindent
$f_T(|Z|)$ が $|T|$ で閉であると仮定する。このとき、
$f_T(|Z|) \subset V(g)$ かつ $t \not \in V(g)$ となる元
$g \in B[a_i; i \in I]$ が存在する。したがって $\kappa(t)$ における
$g$ の像は零でない。とくに $g$ のある係数は $\kappa(y)$ で非零の像をもつ。
ゆえに、この係数は $y$ のあるアフィン開近傍 $V$ 上で可逆である。
$J$ を、変数 $a_j$ が $g$ に現れるような $j \in I$ の有限集合とする。
$X_V$ は準コンパクトでなく、各 $X_{i, V}$ は準コンパクトなので、
点 $x \in |X_V| \setminus \bigcup_{j \in J} |X_{j, V}|$ を選べる。
換言すれば、$x \in |X| \setminus \bigcup_{j \in J} |X_j|$ であり、
$x$ はある $v \in V$ 上にある。$g$ は $V$ 上可逆な係数をもつので、
$v$ 上にある点 $t' \in T$ で、$t' \not \in V(g)$ かつ
すべての $i \notin J$ に対して $t' \in V(a_i)$ となるものを取れる。
これは
$V(a_i; i \in I \setminus J) = \Spec(B[a_j; j\in J])$
であり、このスキームの $v$ 上の点の集合は
$\Spec(\kappa(v)[a_j; j \in J])$ と全単射で、構成により
$g$ はこの多項式環の非零元に制限されるからである。
換言すれば、各 $i \not \in J$ に対して $t' \not \in T_i$ である。
Properties of Spaces, Lemma
\ref{spaces-properties-lemma-points-cartesian} により、
$x \in X$ と $t' \in T$ に写る $X \times_Y T$ の点 $z$ を取れる。
$j \in J$ に対して $x \not \in |X_j|$ であり、
$i \in I \setminus J$ に対して $t' \not \in T_i$ なので、
$z \in |Z|$ である。一方、
$f_T(z) = t' \not \in V(g)$ であり、$f_T(Z) \subset V(g)$ に矛盾する。
したがって「$f_T(|Z|)$ は閉である」という仮定は誤りであり、
望みどおり $f_T$ は閉でない。
\end{proof}

\noindent
分離代数空間の全射かつ普遍閉な射による標的は分離的である。

\begin{lemma}
\label{spaces-morphisms-lemma-image-universally-closed-separated}
$S$ をスキームとし、$B$ を $S$ 上の代数空間とする。
$f : X \to Y$ を $B$ 上の代数空間の全射かつ普遍閉な射とする。
\begin{enumerate}
\item $X$ が準分離ならば、$Y$ は準分離である。
\item $X$ が分離的ならば、$Y$ は分離的である。
\item $X$ が $B$ 上準分離ならば、$Y$ は $B$ 上準分離である。
\item $X$ が $B$ 上分離的ならば、$Y$ は $B$ 上分離的である。
\end{enumerate}
\end{lemma}

\begin{proof}
(1) と (2) は $S = B = \Spec(\mathbf{Z})$ とした (3) と (4) の帰結である
（Properties of Spaces, Definition
\ref{spaces-properties-definition-separated} を参照）。可換図式
$$
\xymatrix{
X \ar[d] \ar[rr]_{\Delta_{X/B}} & & X \times_B X \ar[d] \\
Y \ar[rr]^{\Delta_{Y/B}} & & Y \times_B Y
}
$$
を考える。左の垂直矢印は全射（すなわち普遍的に全射）である。
右の垂直矢印は普遍閉射
$X \times_B X \to X \times_B Y \to Y \times_B Y$ の合成として
普遍閉である。したがって Lemma
\ref{spaces-morphisms-lemma-universally-closed-quasi-compact} により準コンパクトでもある。

\medskip\noindent
$X$ が $B$ 上準分離、すなわち $\Delta_{X/B}$ が準コンパクトであると仮定する。
$Z$ が準コンパクトで $Z \to Y \times_B Y$ が射ならば、
$Z \times_{Y \times_B Y} X \to Z \times_{Y \times_B Y} Y$ は全射であり、
上の注意により $Z \times_{Y \times_B Y} X$ は準コンパクトである。
したがって $\Delta_{Y/B}$ は準コンパクト、すなわち
$Y$ は $B$ 上準分離である。

\medskip\noindent
$X$ が $B$ 上分離的、すなわち $\Delta_{X/B}$ が閉埋め込みであると仮定する。
$Z$ がアフィンで $Z \to Y \times_B Y$ が射ならば、
$Z \times_{Y \times_B Y} X \to Z \times_{Y \times_B Y} Y$ は全射であり、
上の注意により $Z \times_{Y \times_B Y} X \to Z$ は普遍閉である。
したがって $\Delta_{Y/B}$ は普遍閉である。
Lemma \ref{spaces-morphisms-lemma-properties-diagonal} により $\Delta_{Y/B}$ は
表現可能、局所有限型、かつ単射であり、さらに普遍閉なので閉埋め込みである。
\'Etale Morphisms, Lemma
\ref{etale-lemma-characterize-closed-immersion} を参照
（また抽象的原理 Spaces, Lemma
\ref{spaces-lemma-representable-transformations-property-implication} も参照）。
ゆえに $Y$ は $B$ 上分離的である。
\end{proof}













\section{単射}
\label{spaces-morphisms-section-monomorphisms}

\noindent
代数空間の表現可能な射 $X \to Y$ が Section
\ref{spaces-morphisms-section-representable} の意味で単射であるとは、任意のスキーム
$Z$ と射 $Z \to Y$ に対して、射 $Z \times_Y X \to Z$ が
スキームの単射によって表現されることをいう。これはちょうど
$Z \times_Y X \to Z$ が $(\Sch/S)_{fppf}$ 上の層の単射であることを意味する。
これがすべての $Z$ とすべての写像 $Z \to Y$ について成り立つべきなので、
さらにこれは、写像 $X \to Y$ が $(\Sch/S)_{fppf}$ 上の層の単射であることと
同値である。したがって、代数空間の（必ずしも表現可能でない
\footnote{代数空間の任意の単射が表現可能かどうかは分かっていない。
議論については More on Morphisms of Spaces, Section
\ref{spaces-more-morphisms-section-monomorphisms} を参照。}）射の単射を
次のように定義できる。

\begin{definition}
\label{spaces-morphisms-definition-monomorphism}
$S$ をスキームとする。$S$ 上の代数空間の射が層の単射、すなわち
$(\Sch/S)_{fppf}$ 上の層の圏における単射であるとき、その射を
{\it 単射}という。
\end{definition}

\noindent
次の補題は、これが $S$ 上の代数空間の圏における単射でもあることを示す。

\begin{lemma}
\label{spaces-morphisms-lemma-monomorphism}
$S$ をスキームとし、$j : X \to Y$ を $S$ 上の代数空間の射とする。
次は同値である。
\begin{enumerate}
\item $j$ は単射である（Definition \ref{spaces-morphisms-definition-monomorphism} の意味で）。
\item $j$ は $S$ 上の代数空間の圏における単射である。
\item 対角射 $\Delta_{X/Y} : X \to X \times_Y X$ は同型である。
\end{enumerate}
\end{lemma}

\begin{proof}
$X \times_Y X$ は $(\Sch/S)_{fppf}$ 上の層の圏におけるファイバー積であると
同時に、$S$ 上の代数空間の圏におけるファイバー積でもある。
Spaces, Lemma \ref{spaces-lemma-fibre-product-spaces} を参照。
(1) と (3) の同値性は、任意のサイト上の層の単射の一般的特徴づけである。
(2) と (3) の同値性は、ファイバー積をもつ任意の圏における単射の
特徴づけである。
\end{proof}

\begin{lemma}
\label{spaces-morphisms-lemma-monomorphism-separated}
代数空間の単射は分離的である。
\end{lemma}

\begin{proof}
これは同型が閉埋め込みであることと、上の Lemma
\ref{spaces-morphisms-lemma-monomorphism} から従う。
\end{proof}

\begin{lemma}
\label{spaces-morphisms-lemma-composition-monomorphism}
単射の合成は単射である。
\end{lemma}

\begin{proof}
層の単射の合成が単射であることから従う。
\end{proof}

\begin{lemma}
\label{spaces-morphisms-lemma-base-change-monomorphism}
単射の基底変換は単射である。
\end{lemma}

\begin{proof}
これは層の圏におけるファイバー積についての一般的事実である。
\end{proof}

\begin{lemma}
\label{spaces-morphisms-lemma-monomorphism-local}
$S$ をスキームとし、$f : X \to Y$ を $S$ 上の代数空間の射とする。
次は同値である。
\begin{enumerate}
\item $f$ は単射である。
\item 任意のスキーム $Z$ と射 $Z \to Y$ に対して、$f$ の基底変換
$Z \times_Y X \to Z$ は単射である。
\item 任意のアフィンスキーム $Z$ と任意の射 $Z \to Y$ に対して、
$f$ の基底変換 $Z \times_Y X \to Z$ は単射である。
\item あるスキーム $V$ と全射エタール射 $V \to Y$ が存在し、
基底変換 $V \times_Y X \to V$ は単射である。
\item Zariski 被覆 $Y = \bigcup Y_i$ が存在し、各射
$f^{-1}(Y_i) \to Y_i$ は単射である。
\end{enumerate}
\end{lemma}

\begin{proof}
単射の基底変換が単射であることを、以下では断りなく用いる。
Lemma \ref{spaces-morphisms-lemma-base-change-monomorphism} を参照。とくに
(1) $\Rightarrow$ (2) $\Rightarrow$ (3) $\Rightarrow$ (4)
は明らかである（$V$ として $Y$ 上エタールなアフィンスキームの
非交和を取る。Properties of Spaces, Lemma
\ref{spaces-properties-lemma-cover-by-union-affines} を参照）。
$V$ をスキームとし、$V \to Y$ を全射エタール射とする。
$V \times_Y X \to V$ が単射ならば、$X \to Y$ も単射である。
実際、層の任意の Cartesian 図式
$$
\vcenter{
\xymatrix{
\mathcal{F} \ar[r]_a \ar[d]_b & \mathcal{G} \ar[d]^c \\
\mathcal{H} \ar[r]^d & \mathcal{I}
}
}
\quad
\quad
\mathcal{F} = \mathcal{H} \times_\mathcal{I} \mathcal{G}
$$
において、$c$ が層の全射で $a$ が単射ならば、$d$ も単射である。
したがって (4) は (1) を含む。(5) と (1) の同値性の証明は省略する。
\end{proof}

\begin{lemma}
\label{spaces-morphisms-lemma-immersions-monomorphisms}
代数空間の埋め込みは単射である。とくに、任意の埋め込みは分離的である。
\end{lemma}

\begin{proof}
$f : X \to Y$ を代数空間の埋め込みとする。$Z$ が表現可能であるような
任意の射 $Z \to Y$ に対して、基底変換 $Z \times_Y X \to Z$ は
スキームの埋め込み、したがって単射である。Schemes, Lemma
\ref{schemes-lemma-immersions-monomorphisms} を参照。
ゆえに $f$ は表現可能かつ単射である。
\end{proof}

\noindent
次の補題は Decent Spaces, Lemma
\ref{decent-spaces-lemma-monomorphism-toward-disjoint-union-dim-0-rings}
で強化される。

\begin{lemma}
\label{spaces-morphisms-lemma-monomorphism-toward-field}
$S$ をスキームとする。$k$ を体とし、$Z \to \Spec(k)$ を
$S$ 上の代数空間の単射とする。このとき
$Z = \emptyset$ または $Z = \Spec(k)$ である。
\end{lemma}

\begin{proof}
Lemmas \ref{spaces-morphisms-lemma-monomorphism-separated} および
\ref{spaces-morphisms-lemma-separated-over-separated} により、$Z$ は分離代数空間である。
したがってスキームである稠密開部分空間 $Z' \subset Z$ が存在する。
Properties of Spaces, Proposition
\ref{spaces-properties-proposition-locally-quasi-separated-open-dense-scheme}
を参照。Schemes, Lemma
\ref{schemes-lemma-mono-towards-spec-field} により、
$Z' = \emptyset$ または $Z' \cong \Spec(k)$ である。
第一の場合には $Z = \emptyset$、第二の場合には
$Z \to \Spec(k)$ が $Z'$ 上で同型となる単射であるため
$Z' = Z = \Spec(k)$ と結論する。
\end{proof}

\begin{lemma}
\label{spaces-morphisms-lemma-monomorphism-injective-points}
$S$ をスキームとする。$X \to Y$ が $S$ 上の代数空間の単射ならば、
$|X| \to |Y|$ は単射である。
\end{lemma}

\begin{proof}
定義から直ちに従う。
\end{proof}















\section{準連接層の順像}
\label{spaces-morphisms-section-pushforward}

\noindent
まず、代数空間の射に対する加群の層の順像と、スキームの射に対する
加群の層の順像を関係づける簡単な補題を証明する。

\begin{lemma}
\label{spaces-morphisms-lemma-compute-pushforward}
$S$ をスキームとし、$f : X \to Y$ を $S$ 上の代数空間の射とする。
$U \to X$ をスキームから $X$ への全射エタール射とする。
$R = U \times_X U$ とおき、通常どおり射影を $t, s : R \to U$ と書く。
誘導される射を $a : U \to Y$ および $b : R \to Y$ と書く。
$\textit{Mod}(\mathcal{O}_X)$ の任意の対象 $\mathcal{F}$ に対して完全列
$$
0 \to f_*\mathcal{F} \to a_*(\mathcal{F}|_U) \to b_*(\mathcal{F}|_R)
$$
が存在し、第二の矢印は差 $t^* - s^*$ である。
\end{lemma}

\begin{proof}
$\mathcal{F}$ という記号を、その $X_{spaces, \etale}$ 上の
加群の層への拡張にも用いる。Properties of Spaces, Remark
\ref{spaces-properties-remark-explain-equivalence} を参照。
$V \to Y$ を $Y_\etale$ の対象とする。このとき $V \times_Y X$ は
$X_{spaces, \etale}$ の対象であり、定義により
$f_*\mathcal{F}(V) = \mathcal{F}(V \times_Y X)$ である。
$U \to X$ は全射エタールなので、
$\{V \times_Y U \to V \times_Y X\}$ は被覆である。また
$(V \times_Y U) \times_X (V \times_Y U) = V \times_Y R$ である。
したがって、$X_{spaces, \etale}$ 上の $\mathcal{F}$ の層条件により
短完全列
$$
0 \to \mathcal{F}(V \times_Y X)
\to \mathcal{F}(V \times_Y U) \to \mathcal{F}(V \times_Y R)
$$
を得る。ここで第二の矢印は $t$ または $s$ による制限の差である。
この完全列は $V$ に関して関手的なので、補題を得る。
\end{proof}

\noindent
$S$ をスキームとする。$f : X \to Y$ を、$S$ 上の表現可能な代数空間
$X$ および $Y$ の準コンパクトかつ準分離な射とする。
Descent, Proposition
\ref{descent-proposition-equivalence-quasi-coherent-functorial} により、関手
$f_* : \QCoh(\mathcal{O}_X) \to \QCoh(\mathcal{O}_Y)$ は、
$X$ と $Y$ をスキームとみなしたときの通常の関手と一致する。

\medskip\noindent
より一般に、$f : X \to Y$ を $S$ 上の代数空間の表現可能、準コンパクト、
かつ準分離な射とする。$V$ をスキームとし、
$V \to Y$ を全射エタール射とする。$U = V \times_Y X$ とおき、
$f' : U \to V$ を $f$ の基底変換とする。このとき任意の準連接
$\mathcal{O}_X$-加群 $\mathcal{F}$ に対して
\begin{equation}
\label{spaces-morphisms-equation-representable-pushforward}
f'_*(\mathcal{F}|_U) = (f_*\mathcal{F})|_V,
\end{equation}
である。Properties of Spaces, Lemma
\ref{spaces-properties-lemma-pushforward-etale-base-change-modules} を参照。
$f' : U \to V$ はスキームの準コンパクトかつ準分離な射なので、
前段落の注意により、$\mathcal{F}|_U$ をスキーム $U$ 上の準連接層、
$f'$ をスキームの射とみなすことで $f'_*(\mathcal{F}|_U)$ を計算できる。
以下ではこれをしばしば断りなく用いる。

\medskip\noindent
次の一般性の段階では、代数空間の任意の準コンパクトかつ準分離な射を考える。

\begin{lemma}
\label{spaces-morphisms-lemma-pushforward}
$S$ をスキームとし、$f : X \to Y$ を $S$ 上の代数空間の射とする。
$f$ が準コンパクトかつ準分離ならば、$f_*$ は準連接
$\mathcal{O}_X$-加群を準連接 $\mathcal{O}_Y$-加群へ写す。
\end{lemma}

\begin{proof}
$\mathcal{F}$ を $X$ 上の準連接層とする。$f_*\mathcal{F}$ が $Y$ 上の
準連接層であることを示さなければならない。そのためには、任意の
アフィンスキーム $V$ とエタール射 $V \to Y$ に対して、
$f_*\mathcal{F}$ の $V$ への制限が準連接であることを示せば十分である。
Properties of Spaces, Lemma
\ref{spaces-properties-lemma-characterize-quasi-coherent} を参照。
$f' : V \times_Y X \to V$ を $V \to Y$ による $f$ の基底変換とする。
Lemmas \ref{spaces-morphisms-lemma-base-change-quasi-compact} および
\ref{spaces-morphisms-lemma-base-change-separated} により、$f'$ も準コンパクトかつ
準分離であることに注意する。
(\ref{spaces-morphisms-equation-representable-pushforward}) により、
$f_*\mathcal{F}$ の $V$ への制限は、$\mathcal{F}$ の
$V \times_Y X$ への制限の $f'_*$ である。したがって $f$ を $f'$ で
置き換え、$Y$ がアフィンスキームであると仮定してよい。

\medskip\noindent
$Y$ がアフィンスキームであると仮定する。$f$ は準コンパクトなので
$X$ は準コンパクトである。したがって、アフィンスキーム $U$ と
全射エタール射 $U \to X$ を選べる。Properties of Spaces, Lemma
\ref{spaces-properties-lemma-quasi-compact-affine-cover} を参照。
Lemma \ref{spaces-morphisms-lemma-compute-pushforward} により完全列
$$
0 \to f_*\mathcal{F} \to a_*(\mathcal{F}|_U) \to b_*(\mathcal{F}|_R).
$$
を得る。ここで $R = U \times_X U$ である。$X \to Y$ は準分離なので、
$R \to U \times_Y U$ は準コンパクトな単射である。このことから
$R$ は準コンパクトな分離スキームである
（この時点で $U$ と $Y$ はアフィンである）。
したがって $a : U \to Y$ および $b : R \to Y$ は
スキームの準コンパクトかつ準分離な射である。ゆえに Descent,
Proposition \ref{descent-proposition-equivalence-quasi-coherent-functorial}
により、層 $a_*(\mathcal{F}|_U)$ および $b_*(\mathcal{F}|_R)$ は
準連接である（この補題の前の議論も参照）。
よって $f_*\mathcal{F}$ は準連接加群間の射の核であり、それ自身準連接である。
Properties of Spaces, Lemma
\ref{spaces-properties-lemma-properties-quasi-coherent} を参照。
\end{proof}

\noindent
高次順像については Cohomology of Spaces, Section
\ref{spaces-cohomology-section-higher-direct-image} で論じる。









\section{埋め込み}
\label{spaces-morphisms-section-immersions}

\noindent
代数空間の開埋め込み、閉埋め込み、および局所閉埋め込みは
Spaces, Section \ref{spaces-section-Zariski} で定義した。
すなわち、代数空間の射が {\it 閉埋め込み}
（それぞれ {\it 開埋め込み}、{\it 埋め込み}）であるとは、
それが表現可能であり、Section \ref{spaces-morphisms-section-representable} の意味で
閉埋め込み（それぞれ開埋め込み、埋め込み）であることをいう。

\medskip\noindent
とくに、これらの型の射は基底変換で安定であり、$S$ 上の代数空間の圏で
射の合成について安定である。Spaces, Lemmas
\ref{spaces-lemma-composition-immersions} および
\ref{spaces-lemma-base-change-immersions} を参照。

\begin{lemma}
\label{spaces-morphisms-lemma-closed-immersion-local}
$S$ をスキームとし、$f : X \to Y$ を $S$ 上の代数空間の射とする。
次は同値である。
\begin{enumerate}
\item $f$ は閉埋め込み（それぞれ開埋め込み、埋め込み）である。
\item 任意のスキーム $Z$ と任意の射 $Z \to Y$ に対して、射
$Z \times_Y X \to Z$ は閉埋め込み（それぞれ開埋め込み、埋め込み）である。
\item 任意のアフィンスキーム $Z$ と任意の射 $Z \to Y$ に対して、射
$Z \times_Y X \to Z$ は閉埋め込み（それぞれ開埋め込み、埋め込み）である。
\item あるスキーム $V$ と全射エタール射 $V \to Y$ が存在し、
$V \times_Y X \to V$ は閉埋め込み（それぞれ開埋め込み、埋め込み）である。
\item Zariski 被覆 $Y = \bigcup Y_i$ が存在し、各射
$f^{-1}(Y_i) \to Y_i$ は閉埋め込み（それぞれ開埋め込み、埋め込み）である。
\end{enumerate}
\end{lemma}

\begin{proof}
閉埋め込み（それぞれ開埋め込み、埋め込み）の基底変換が同じ型であることを
用いれば、(1) が (2) を含み、(2) が (3) を含むことは明らかである。
また、Properties of Spaces, Lemma
\ref{spaces-properties-lemma-cover-by-union-affines} により、
$V$ をアフィンの非交和として取れるので、(3) は (4) を含む。

\medskip\noindent
$V \to Y$ を (4) のとおりと仮定する。$\mathcal{P}$ を
スキームの射の性質「閉埋め込み（それぞれ開埋め込み、埋め込み）」とする。
性質 $\mathcal{P}$ は任意の基底変換で保たれ、基底上 fppf 局所的である
（Section \ref{spaces-morphisms-section-representable} を参照）。
さらに、型 $\mathcal{P}$ の射は分離的かつ局所準有限である
（三つの各場合について Schemes, Lemma
\ref{schemes-lemma-immersions-monomorphisms} および Morphisms, Lemma
\ref{morphisms-lemma-immersion-locally-quasi-finite} を参照）。
したがって More on Morphisms, Lemma
\ref{more-morphisms-lemma-separated-locally-quasi-finite-morphisms-fppf-descend}
により、型 $\mathcal{P}$ の射は fppf 被覆に関する降下を満たす。
ゆえに Spaces, Lemma
\ref{spaces-lemma-morphism-sheaves-with-P-effective-descent-etale} を適用でき、
$X \to Y$ は表現可能で性質 $\mathcal{P}$ をもつ、すなわち (1) が成り立つ。

\medskip\noindent
(1) と (5) の同値性は、$\mathcal{P}$ が標的上 Zariski 局所的であることから
従う（実際、上で $\mathcal{P}$ は標的上 fppf 局所的であることを見た）。
\end{proof}

\begin{lemma}
\label{spaces-morphisms-lemma-immersion-permanence}
$S$ をスキームとし、$Z \to Y \to X$ を $S$ 上の代数空間の射とする。
\begin{enumerate}
\item $Z \to X$ が表現可能、局所有限型、局所準有限、分離的、かつ単射ならば、
$Z \to Y$ は表現可能、局所有限型、局所準有限、分離的、かつ単射である。
\item $Z \to X$ が埋め込みで $Y \to X$ が局所分離的ならば、
$Z \to Y$ は埋め込みである。
\item $Z \to X$ が閉埋め込みで $Y \to X$ が分離的ならば、
$Z \to Y$ は閉埋め込みである。
\end{enumerate}
\end{lemma}

\begin{proof}
各場合に、可換図式
$$
\xymatrix{
Z \ar[r] \ar[rd] & Y \times_X Z \ar[r] \ar[d] & Z \ar[d] \\
& Y \ar[r] & X
}
$$
を考えればよい。上の水平矢印の合成は恒等射である。
(1) を証明しよう。第一の水平矢印は $Y \times_X Z \to Z$ の切断なので、
Lemma \ref{spaces-morphisms-lemma-section-immersion} により、表現可能、局所有限型、
局所準有限、分離的、かつ単射である。矢印 $Y \times_X Z \to Y$ は
$Z \to X$ の基底変換なので、表現可能、局所有限型、局所準有限、
分離的、かつ単射である
（スキームの射のこれらの性質はいずれも基底変換で安定である。
Spaces, Remark \ref{spaces-remark-list-properties-stable-base-change} を参照）。
したがって合成についても同じことが成り立つ
（スキームの射のこれらの性質はいずれも合成で安定である。
Spaces, Remark \ref{spaces-remark-list-properties-stable-composition} を参照）。
(1) が示された。他の結果もまったく同様に証明される。
\end{proof}

\begin{lemma}
\label{spaces-morphisms-lemma-immersion-when-closed}
$S$ をスキームとし、$i : Z \to X$ を $S$ 上の代数空間の埋め込みとする。
このとき $|i| : |Z| \to |X|$ は局所閉部分集合の上への同相写像であり、
$i$ が閉埋め込みであることと、像
$|i|(|Z|) \subset |X|$ が閉部分集合であることは同値である。
\end{lemma}

\begin{proof}
第一の主張は Properties of Spaces, Lemma
\ref{spaces-properties-lemma-subspace-induced-topology} である。
$U$ をスキームとし、$U \to X$ を全射エタール射とする。
仮定により $T = U \times_X Z$ はスキームであり、射 $j : T \to U$ は
スキームの埋め込みである。Lemma \ref{spaces-morphisms-lemma-closed-immersion-local} により、
$i$ が閉埋め込みであることと $j$ が閉埋め込みであることは同値である。
Schemes, Lemma \ref{schemes-lemma-immersion-when-closed} により、
これは $j(T)$ が $U$ で閉であることと同値である。
しかし、部分集合 $j(T) \subset U$ は
$|i|(|Z|) \subset |X|$ の逆像である。Properties of Spaces, Lemma
\ref{spaces-properties-lemma-points-cartesian} を参照。これで証明が完了する。
\end{proof}

\begin{remark}
\label{spaces-morphisms-remark-immersion}
$S$ をスキームとし、$i : Z \to X$ を $S$ 上の代数空間の埋め込みとする。
$i$ は単射なので、$|Z|$ を $|X|$ の部分集合とみなしてよい。
この注意の残りではそうする。$\partial |Z|$ を、位相空間 $|X|$ における
$|Z|$ の境界とする。式では
$$
\partial |Z| = \overline{|Z|} \setminus |Z|.
$$
である。$\partial Z$ を、被約誘導閉部分空間構造を取って得られる
$X$ の被約閉部分空間で $|\partial Z| = \partial |Z|$ となるものとする。
Properties of Spaces, Definition
\ref{spaces-properties-definition-reduced-induced-space} を参照。
構成により、$|Z|$ は
$|X| \setminus |\partial Z| = |X \setminus \partial Z|$ で閉である。
したがって、代数空間の任意の埋め込みは、閉埋め込みとそれに続く開埋め込みに
分解できる（一般には逆順にはできない。Morphisms, Example
\ref{morphisms-example-thibaut} を参照）。
\end{remark}

\begin{remark}
\label{spaces-morphisms-remark-space-structure-locally-closed-subset}
$S$ をスキームとし、$X$ を $S$ 上の代数空間とする。
$T \subset |X|$ を局所閉部分集合とする。$\partial T$ を
位相空間 $|X|$ における $T$ の境界とする。式では
$$
\partial T = \overline{T} \setminus T.
$$
である。$U \subset X$ を、$|U| = |X| \setminus \partial T$ となる
$X$ の開部分空間とする。Properties of Spaces, Lemma
\ref{spaces-properties-lemma-open-subspaces} を参照。
$Z$ を、被約誘導閉部分空間構造を取って得られる $U$ の被約閉部分空間で
$|Z| = T$ となるものとする。Properties of Spaces, Definition
\ref{spaces-properties-definition-reduced-induced-space} を参照。
構成により $Z \to U$ は代数空間の閉埋め込みであり、$U \to X$ は
開埋め込みなので、$Z \to X$ は $S$ 上の代数空間の埋め込みである
（Spaces, Lemma \ref{spaces-lemma-composition-immersions} を参照）。
$Z$ は被約代数空間であり、$|X|$ の部分集合として $|Z| = T$ であることに
注意する。$Z$ を $T$ 上の {\it 被約誘導部分空間構造}と呼ぶことがある。
\end{remark}

\begin{lemma}
\label{spaces-morphisms-lemma-factor-the-other-way}
$S$ をスキームとし、$Z \to X$ を $S$ 上の代数空間の埋め込みとする。
$Z \to X$ が準コンパクトであると仮定する。このとき
$Z \to \overline{Z} \to X$ という分解が存在し、
$Z \to \overline{Z}$ は開埋め込み、$\overline{Z} \to X$ は閉埋め込みである。
\end{lemma}

\begin{proof}
$U$ をスキームとし、$U \to X$ を全射エタール射とする。通常どおり
$R = U \times_X U$ と書き、射影を $s, t : R \to U$ とする。
$T = Z \times_U X$ とおく。$\overline{T} \subset U$ を
$T \to U$ のスキーム論的像とする。準コンパクト射のスキーム論的像を取ることは
平坦基底変換と可換なので
$s^{-1}\overline{T} = t^{-1}\overline{T}$ であることに注意する。
Morphisms, Lemma
\ref{morphisms-lemma-flat-base-change-scheme-theoretic-image} を参照。
したがって、$U$ への引き戻しが $\overline{T}$ となる閉部分空間
$\overline{Z} \subset X$ を得る。Properties of Spaces, Lemma
\ref{spaces-properties-lemma-subspaces-presentation} を参照。
Morphisms, Lemma \ref{morphisms-lemma-quasi-compact-immersion} により、
射 $T \to \overline{T}$ は開埋め込みである。
したがって $Z \to \overline{Z}$ は開埋め込みであり、証明が完了する。
\end{proof}


















\section{閉埋め込み}
\label{spaces-morphisms-section-closed-immersions}

\noindent
本節では、代数空間の埋め込みについて先に得た結果のいくつかを詳しく述べる。
Spaces, Section \ref{spaces-section-Zariski} および本章の Section
\ref{spaces-morphisms-section-immersions} を参照。本節は代数空間に対する
Morphisms, Section \ref{morphisms-section-closed-immersions} の類似である。

\begin{lemma}
\label{spaces-morphisms-lemma-closed-immersion-ideals}
$S$ をスキームとし、$X$ を $S$ 上の代数空間とする。
任意の閉埋め込み $i : Z \to X$ に対して、層 $i_*\mathcal{O}_Z$ は
準連接 $\mathcal{O}_X$-加群であり、写像
$i^\sharp : \mathcal{O}_X \to i_*\mathcal{O}_Z$ は全射で、その核は
準連接なイデアル層である。規則
$Z \mapsto \Ker(\mathcal{O}_X \to i_*\mathcal{O}_Z)$ は
包含関係を逆転する全単射
$$
\begin{matrix}
\text{閉部分空間}\\
Z \subset X
\end{matrix}
\longrightarrow
\begin{matrix}
\text{準連接な層}\\
\text{イデアル }\mathcal{I} \subset \mathcal{O}_X
\end{matrix}
$$
を定める。さらに、準連接なイデアル層
$\mathcal{I} \subset \mathcal{O}_X$ に対応する閉部分スキーム $Z$ が
与えられたとき、代数空間の射 $h : Y \to X$ が $Z$ を経由するための
必要十分条件は、写像
$h^*\mathcal{I} \to h^*\mathcal{O}_X = \mathcal{O}_Y$ が零であることである。
\end{lemma}

\begin{proof}
$U \to X$ を、始域がスキームである全射エタール射とする。図式
$$
\xymatrix{
U \times_X Z \ar[r] \ar[d]_{i'} & Z \ar[d]^i \\
U \ar[r] & X
}
$$
を考える。Lemma \ref{spaces-morphisms-lemma-closed-immersion-local} により、
$i$ が閉埋め込みであることと $i'$ が閉埋め込みであることは同値である。
Properties of Spaces, Lemma
\ref{spaces-properties-lemma-pushforward-etale-base-change-modules} により、
$i'_*\mathcal{O}_{U \times_X Z}$ は $i_*\mathcal{O}_Z$ の $U$ への制限である。
したがって $\mathcal{O}_X \to i_*\mathcal{O}_Z$ に関する主張は、
$\mathcal{O}_U \to i'_*\mathcal{O}_{U \times_X Z}$ に関する対応する主張と
同値である。$i'$ はスキームの閉埋め込みなので、これらの結果は
Morphisms, Lemma \ref{morphisms-lemma-closed-immersion} から従う。

\medskip\noindent
準連接なイデアル層 $\mathcal{I} \subset \mathcal{O}_X$ が与えられたとき、
式
$$
Z(T) = \{h : T \to X \mid h^*\mathcal{I} \to \mathcal{O}_T
\text{ は零である}\}
$$
が $X$ の閉部分空間を定めることを証明しよう。
これは明らかに $X$ の部分関手である。$Z \to X$ が閉埋め込みにより
表現されることを示すため、スキームから $X$ への射
$\varphi : U \to X$ を取る。このとき $Z \times_X U$ は、
イデアル層 $\Im(\varphi^*\mathcal{I} \to \mathcal{O}_U)$ に対応する
$U$ の同様な部分関手によって表現される。Properties of Spaces, Lemma
\ref{spaces-properties-lemma-pullback-quasi-coherent} により、
$\mathcal{O}_U$-加群 $\varphi^*\mathcal{I}$ は $U$ 上準連接であり、
したがって $\Im(\varphi^*\mathcal{I} \to \mathcal{O}_U)$ は
$U$ 上の準連接なイデアル層である。Schemes, Lemma
\ref{schemes-lemma-characterize-closed-subspace} により、
$Z \times_X U$ は $\Im(\varphi^*\mathcal{I} \to \mathcal{O}_U)$ に
付随する $U$ の閉部分スキームによって表現されると結論する。
ゆえに $Z$ は $X$ の閉部分空間である。

\medskip\noindent
上の $Z$ の式では、代数空間が $(\Sch/S)_{fppf}$ 上の層なので、
入力 $T$ はスキームである。同じ式が $T$ が代数空間の場合にも
成り立つことの確認は省略する。
\end{proof}

\begin{definition}
\label{spaces-morphisms-definition-inverse-image-closed-subspace}
$S$ をスキームとし、$f : Y \to X$ を $S$ 上の代数空間の射とする。
$Z \subset X$ を閉部分空間とする。{\it 閉部分空間 $Z$ の逆像
$f^{-1}(Z)$} とは、$Y$ の閉部分空間 $Z \times_X Y$ のことである。
\end{definition}

\noindent
この定義は Lemma \ref{spaces-morphisms-lemma-closed-immersion-local} により意味をもつ。
$\mathcal{I} \subset \mathcal{O}_X$ が Lemma
\ref{spaces-morphisms-lemma-closed-immersion-ideals} によって $Z$ に対応する準連接な
イデアル層ならば、
$f^{-1}\mathcal{I}\mathcal{O}_Y = \Im(f^*\mathcal{I} \to \mathcal{O}_Y)$
は $f^{-1}(Z)$ に対応するイデアル層である。

\begin{lemma}
\label{spaces-morphisms-lemma-closed-immersion-quasi-compact}
代数空間の閉埋め込みは準コンパクトである。
\end{lemma}

\begin{proof}
これは Schemes, Lemma
\ref{schemes-lemma-closed-immersion-quasi-compact} から一般原理によって従う。
Spaces, Lemma
\ref{spaces-lemma-representable-transformations-property-implication} を参照。
\end{proof}

\begin{lemma}
\label{spaces-morphisms-lemma-closed-immersion-separated}
代数空間の閉埋め込みは分離的である。
\end{lemma}

\begin{proof}
これは Schemes, Lemma \ref{schemes-lemma-immersions-monomorphisms} から
一般原理によって従う。Spaces, Lemma
\ref{spaces-lemma-representable-transformations-property-implication} を参照。
\end{proof}

\begin{lemma}
\label{spaces-morphisms-lemma-closed-immersion-push-pull}
$S$ をスキームとし、$i : Z \to X$ を $S$ 上の代数空間の閉埋め込みとする。
\begin{enumerate}
\item 関手
$$
i_{small, *} :
\Sh(Z_\etale)
\longrightarrow
\Sh(X_\etale)
$$
は充満忠実であり、その本質像は $X_\etale$ 上の集合の層
$\mathcal{F}$ で、その $X \setminus Z$ への制限が $*$ と同型なものからなる。
\item 関手
$$
i_{small, *} :
\textit{Ab}(Z_\etale)
\longrightarrow
\textit{Ab}(X_\etale)
$$
は充満忠実であり、その本質像は台が $|Z|$ に含まれる
$X_\etale$ 上のアーベル層からなる。
\end{enumerate}
いずれの場合も、$i_{small}^{-1}$ は関手 $i_{small, *}$ の左逆である。
\end{lemma}

\begin{proof}
$U$ をスキームとし、$U \to X$ を全射エタール射とする。
$V = Z \times_X U$ とおく。このとき $V$ はスキームであり、
$i' : V \to U$ はスキームの閉埋め込みである。Properties of Spaces,
Lemma \ref{spaces-properties-lemma-pushforward-etale-base-change} により、
$Z$ 上の任意の層 $\mathcal{G}$ に対して
$$
(i_{small}^{-1}i_{small, *}\mathcal{G})|_V =
(i')_{small}^{-1}i'_{small, *}(\mathcal{G}|_V)
$$
である。\'Etale Cohomology, Proposition
\ref{etale-cohomology-proposition-closed-immersion-pushforward} により、写像
$(i')_{small}^{-1}i'_{small, *}(\mathcal{G}|_V) \to \mathcal{G}|_V$
は同型である。$V \to Z$ は全射かつエタールなので、これは
$i_{small}^{-1}i_{small, *}\mathcal{G} \to \mathcal{G}$ が同型であることを
意味する。このことから $i_{small, *}$ が充満忠実であることは明らかである。
Sites, Lemma \ref{sites-lemma-exactness-properties} を参照。
本質像に関する主張を証明するため、$X_\etale$ 上の集合の層
$\mathcal{F}$ で、その $X \setminus Z$ への制限が $*$ と同型なものを考える。
\'Etale Cohomology, Proposition
\ref{etale-cohomology-proposition-closed-immersion-pushforward} の証明と同様に、
随伴写像
$$
\mathcal{F} \longrightarrow i_{small, *}i_{small}^{-1}\mathcal{F}.
$$
を考える。第一の部分と同様に、スキームの場合の対応する結果により、
この写像の $U$ への制限は同型である。$U$ は $X$ のエタール被覆なので、
これは同型であると結論する。
\end{proof}

\begin{lemma}
\label{spaces-morphisms-lemma-stalk-push-closed}
$S$ をスキームとし、$i : Z \to X$ を $S$ 上の代数空間の閉埋め込みとする。
$\overline{z}$ を $Z$ の幾何点とし、その $X$ における像を
$\overline{x}$ とする。このとき $Z_\etale$ 上の任意の層 $\mathcal{F}$ に対して
$(i_{small, *}\mathcal{F})_{\overline{z}} = \mathcal{F}_{\overline{x}}$
である。
\end{lemma}

\begin{proof}
$\overline{x}$ のエタール近傍 $(U, \overline{u})$ を選ぶ。このとき茎
$(i_{small, *}\mathcal{F})_{\overline{z}}$ は
$i_{small, *}\mathcal{F}|_U$ の $\overline{u}$ における茎である。
Properties of Spaces, Lemma
\ref{spaces-properties-lemma-pushforward-etale-base-change} により、
$X$ を $U$ で、$Z$ を $Z \times_X U$ で置き換えてよい。
すると $Z \to X$ はスキームの閉埋め込みであり、結果は
\'Etale Cohomology, Lemma
\ref{etale-cohomology-lemma-stalk-pushforward-closed-immersion} である。
\end{proof}

\noindent
次の補題は閉トポスの埋め込みという設定でより一般に成り立つ
（ここに将来の参照を挿入）。

\begin{lemma}
\label{spaces-morphisms-lemma-closed-immersion-rings}
$S$ をスキームとし、$i : Z \to X$ を $S$ 上の代数空間の閉埋め込みとする。
$\mathcal{A}$ を $X_\etale$ 上の環の層とし、$\mathcal{B}$ を
$Z_\etale$ 上の環の層とする。$\varphi : \mathcal{A} \to
i_{small, *}\mathcal{B}$ を環の層の準同型とし、これによって環付きトポスの射
$$
f : (\Sh(Z_\etale), \mathcal{B}) \longrightarrow (\Sh(X_\etale), \mathcal{A}).
$$
を得るとする。$\mathcal{A}$-加群の層 $\mathcal{F}$ と
$\mathcal{B}$-加群の層 $\mathcal{G}$ に対して、標準写像
$$
\mathcal{F} \otimes_\mathcal{A} f_*\mathcal{G}
\longrightarrow
f_*(f^*\mathcal{F} \otimes_\mathcal{B} \mathcal{G}).
$$
は同型である。
\end{lemma}

\begin{proof}
この写像は、$\text{id} : f^*\mathcal{F} \to f^*\mathcal{F}$ と随伴写像
$f^* f_*\mathcal{G} \to \mathcal{G}$ から得られる写像
$$
f^*\mathcal{F} \otimes_\mathcal{B}
f^* f_*\mathcal{G} =
f^*(\mathcal{F} \otimes_\mathcal{A} f_*\mathcal{G})
\longrightarrow
f^*\mathcal{F} \otimes_\mathcal{B} \mathcal{G}
$$
に随伴する写像である。これが同型であることは茎で確かめてよい
（Properties of Spaces, Theorem
\ref{spaces-properties-theorem-exactness-stalks}）。
$\overline{z} : \Spec(k) \to Z$ を幾何点とし、その像を
$\overline{x} = i \circ \overline{z} : \Spec(k) \to X$ とする。
写像が茎に及ぼす作用を書き下すと、示すべきことは
$$
\mathcal{F}_{\overline{x}}
\otimes_{\mathcal{A}_{\overline{x}}}
\mathcal{G}_{\overline{z}} =
(\mathcal{F}_{\overline{x}}
\otimes_{\mathcal{A}_{\overline{x}}}
\mathcal{B}_{\overline{z}}) \otimes_{\mathcal{B}_{\overline{z}}}
\mathcal{G}_{\overline{z}}
$$
となり、これは成り立つ。ここでは、テンソル積を取ることが茎を取ることと
可換であること、引き戻しのもとでの茎の振る舞い
Properties of Spaces, Lemma
\ref{spaces-properties-lemma-stalk-pullback}、および閉埋め込みに沿う
順像のもとでの茎の振る舞い Lemma \ref{spaces-morphisms-lemma-stalk-push-closed} を用いた。
\end{proof}






\section{閉埋め込みと準連接層}
\label{spaces-morphisms-section-closed-immersions-quasi-coherent}

\noindent
本節は Morphisms, Section
\ref{morphisms-section-closed-immersions-quasi-coherent} の類似である。

\begin{lemma}
\label{spaces-morphisms-lemma-i-star-equivalence}
$S$ をスキームとし、$i : Z \to X$ を $S$ 上の代数空間の閉埋め込みとする。
$\mathcal{I} \subset \mathcal{O}_X$ を $Z$ を切り出す準連接なイデアル層とする。
\begin{enumerate}
\item 任意の $\mathcal{O}_X$-加群 $\mathcal{F}$ に対して、随伴写像
$\mathcal{F} \to i_*i^*\mathcal{F}$ は同型
$\mathcal{F}/\mathcal{I}\mathcal{F} \cong i_*i^*\mathcal{F}$ を誘導する。
\item 関手 $i^*$ は $i_*$ の左逆である。すなわち任意の
$\mathcal{O}_Z$-加群 $\mathcal{G}$ に対して、随伴写像
$i^*i_*\mathcal{G} \to \mathcal{G}$ は同型である。
\item 関手
$$
i_* :
\QCoh(\mathcal{O}_Z)
\longrightarrow
\QCoh(\mathcal{O}_X)
$$
は完全かつ充満忠実であり、その本質像は
$\mathcal{I}\mathcal{F} = 0$ を満たす準連接
$\mathcal{O}_X$-加群 $\mathcal{F}$ からなる。
\end{enumerate}
\end{lemma}

\begin{proof}
この証明では小エタールサイト上の層だけを扱い、アーベル群の層の
順像および逆像を $i_{small, *}, i_{small}^{-1}$ ではなく
$i_*, i^{-1}, \ldots$ と書く。

\medskip\noindent
$\mathcal{F}$ を $\mathcal{O}_X$-加群とする。
Lemma \ref{spaces-morphisms-lemma-closed-immersion-rings} を
$\mathcal{A} = \mathcal{O}_X$ および
$\mathcal{G} = \mathcal{B} = \mathcal{O}_Z$ として適用すると、
$i_*i^*\mathcal{F} = \mathcal{F} \otimes_{\mathcal{O}_X} \mathcal{O}_Z$
を得る。Lemma \ref{spaces-morphisms-lemma-closed-immersion-ideals} により短完全列
$$
0 \to \mathcal{I} \to \mathcal{O}_X \to i_*\mathcal{O}_Z \to 0
$$
を得る。テンソル積の性質から
$\mathcal{F} \otimes_{\mathcal{O}_X} i_*\mathcal{O}_Z
= \mathcal{F}/\mathcal{I}\mathcal{F}$ が従う。これで (1) が示された
（ただし写像が随伴写像から誘導されることの確認は省略する）。

\medskip\noindent
$\mathcal{G}$ を任意の $\mathcal{O}_Z$-加群とする。Lemma
\ref{spaces-morphisms-lemma-closed-immersion-push-pull} により
$i^{-1}i_*\mathcal{G} = \mathcal{G}$ である。したがって (2) を示すには、
標準写像
$\mathcal{G} \otimes_{i^{-1}\mathcal{O}_X} \mathcal{O}_Z \to \mathcal{G}$
が同型であることを示さなければならない。これは
$i^{-1}\mathcal{O}_X \to \mathcal{O}_Z$ が全射であることを示せば、
テンソル積の一般的性質から従う。Lemma
\ref{spaces-morphisms-lemma-closed-immersion-push-pull} により、写像
$i_*i^{-1}\mathcal{O}_X \to i_*\mathcal{O}_Z$ が全射であることを
証明すれば十分である。全射 $\mathcal{O}_X \to i_*\mathcal{O}_Z$ は
この写像を経由するので、(2) が成り立つ。

\medskip\noindent
最後に、補題の最も興味深い部分である (3) を証明する。
閉埋め込みは準コンパクトかつ分離的である。Lemmas
\ref{spaces-morphisms-lemma-closed-immersion-quasi-compact} および
\ref{spaces-morphisms-lemma-closed-immersion-separated} を参照。
したがって Lemma \ref{spaces-morphisms-lemma-pushforward} を適用でき、
$Z$ 上の準連接層の順像は実際に $X$ 上の準連接層である。
これにより関手
$i^{QCoh}_* : \QCoh(\mathcal{O}_Z)
\to \QCoh(\mathcal{O}_X)$ を得る。
(2) により $i^{QCoh}_*$ は左逆 $i^*$ をもつので、充満忠実である。

\medskip\noindent
次に関手 $i_*$ の本質像を記述する。$\mathcal{I}$ は写像
$\mathcal{O}_X \to i_*\mathcal{O}_Z$ の核であり、この写像を使って
$i_*\mathcal{G}$ に $\mathcal{O}_X$-加群構造を入れるので、任意の
$\mathcal{O}_Z$-加群について $\mathcal{I}(i_*\mathcal{G}) = 0$ は明らかである。
次に、$\mathcal{I}\mathcal{F} = 0$ を満たす任意の準連接
$\mathcal{O}_X$-加群 $\mathcal{F}$ を考える。
$i_*\mathcal{O}_Z = \mathcal{O}_X/\mathcal{I}$ なので、
$\mathcal{F}$ は $i_*\mathcal{O}_Z$-加群である。とくに、その台は
$|Z|$ に含まれる。Lemma \ref{spaces-morphisms-lemma-closed-immersion-push-pull} を適用すると、
ある $\mathcal{O}_Z$-加群 $\mathcal{G}$ に対して
$\mathcal{F} \cong i_*\mathcal{G}$ である。残る小さな点は
$\mathcal{G}$ が準連接である理由である。これは (2) と
Properties of Spaces, Lemma
\ref{spaces-properties-lemma-pullback-quasi-coherent} により
$\mathcal{G} \cong i^*\mathcal{F}$ だからである。
\end{proof}

\noindent
$i : Z \to X$ を代数空間の閉埋め込みとする。
上の補題により、記号の濫用として、$X$ 上の層 $i_*\mathcal{F}$ を
$\mathcal{F}$ と書くことが多い。

\begin{lemma}
\label{spaces-morphisms-lemma-largest-quasi-coherent-subsheaf}
$S$ をスキームとし、$X$ を $S$ 上の代数空間とする。
$\mathcal{F}$ を準連接 $\mathcal{O}_X$-加群とし、
$\mathcal{G} \subset \mathcal{F}$ を $\mathcal{O}_X$-部分加群とする。
次の性質をもつ一意な準連接 $\mathcal{O}_X$-部分加群
$\mathcal{G}' \subset \mathcal{G}$ が存在する。任意の準連接
$\mathcal{O}_X$-加群 $\mathcal{H}$ に対して写像
$$
\Hom_{\mathcal{O}_X}(\mathcal{H}, \mathcal{G}')
\longrightarrow
\Hom_{\mathcal{O}_X}(\mathcal{H}, \mathcal{G})
$$
は全単射である。とくに $\mathcal{G}'$ は、$\mathcal{G}$ に含まれる
$\mathcal{F}$ の最大の準連接 $\mathcal{O}_X$-部分加群である。
\end{lemma}

\begin{proof}
$\mathcal{G}_a$, $a \in A$ を、$\mathcal{G}$ に含まれる準連接
$\mathcal{O}_X$-部分加群の集合とする。このとき
$$
\bigoplus\nolimits_{a \in A} \mathcal{G}_a \longrightarrow \mathcal{F}
$$
の像 $\mathcal{G}'$ は準連接である。実際、$X$ 上の準連接層の射の像は
準連接であり、準連接層の直和も準連接である。Properties of Spaces,
Lemma \ref{spaces-properties-lemma-properties-quasi-coherent} を参照。
加群 $\mathcal{G}'$ は $\mathcal{G}$ に含まれる。したがってこれは
$\mathcal{G}$ に含まれる最大の準連接 $\mathcal{O}_X$-加群である。

\medskip\noindent
式を証明するため、$\mathcal{H}$ を準連接 $\mathcal{O}_X$-加群とし、
$\alpha : \mathcal{H} \to \mathcal{G}$ を $\mathcal{O}_X$-加群写像とする。
合成 $\mathcal{H} \to \mathcal{G} \to \mathcal{F}$ の像は、
準連接層の射の像として準連接である。したがってそれは
$\mathcal{G}'$ に含まれる。ゆえに望みどおり $\alpha$ は
$\mathcal{G}'$ を経由する。
\end{proof}

\begin{lemma}
\label{spaces-morphisms-lemma-i-upper-shriek}
$S$ をスキームとし、$i : Z \to X$ を $S$ 上の代数空間の閉埋め込みとする。
$i_*$ の右随伴となる関手\footnote{これはおそらく標準的でない記法である。}
$i^! : \QCoh(\mathcal{O}_X) \to \QCoh(\mathcal{O}_Z)$
が存在する。（Modules, Lemma
\ref{modules-lemma-i-star-right-adjoint} と比較せよ。）
\end{lemma}

\begin{proof}
準連接 $\mathcal{O}_X$-加群 $\mathcal{G}$ が与えられたとき、
$\mathcal{I}$ によって零化される局所切断からなる $\mathcal{G}$ の部分層
$\mathcal{H}_Z(\mathcal{G})$ を考える。Lemma
\ref{spaces-morphisms-lemma-largest-quasi-coherent-subsheaf} により、標準的な最大の
準連接 $\mathcal{O}_X$-部分加群 $\mathcal{H}_Z(\mathcal{G})'$ がある。
構成により、任意の準連接 $\mathcal{O}_Z$-加群 $\mathcal{F}$ に対して
$$
\Hom_{\mathcal{O}_X}(i_*\mathcal{F}, \mathcal{H}_Z(\mathcal{G})')
=
\Hom_{\mathcal{O}_X}(i_*\mathcal{F}, \mathcal{G})
$$
である。したがって
$i^!\mathcal{G} = i^*(\mathcal{H}_Z(\mathcal{G})')$ とおける。
詳細は省略する。
\end{proof}

\noindent
準連接なイデアル層と閉部分空間の $1$ 対 $1$ 対応
（Lemma \ref{spaces-morphisms-lemma-closed-immersion-ideals} を参照）を用いると、
閉部分スキームのスキーム論的交叉と和を定義できる。

\begin{definition}
\label{spaces-morphisms-definition-scheme-theoretic-intersection-union}
$S$ をスキームとし、$X$ を $S$ 上の代数空間とする。
$Z, Y \subset X$ を、準連接なイデアル層
$\mathcal{I}, \mathcal{J} \subset \mathcal{O}_X$ に対応する閉部分空間とする。
$Z$ と $Y$ の {\it スキーム論的交叉}とは、
$\mathcal{I} + \mathcal{J}$ によって切り出される $X$ の閉部分空間である。
$Z$ と $Y$ の {\it スキーム論的和}とは、
$\mathcal{I} \cap \mathcal{J}$ によって切り出される $X$ の閉部分空間である。
\end{definition}

\noindent
スキーム論的交叉を作ることがエタール局所化と可換であることは明らかであり、
スキーム論的和についても同様である。

\begin{lemma}
\label{spaces-morphisms-lemma-scheme-theoretic-intersection}
$S$ をスキームとし、$X$ を $S$ 上の代数空間とする。
$Z, Y \subset X$ を閉部分空間とし、$Z \cap Y$ を $Z$ と $Y$ の
スキーム論的交叉とする。このとき $Z \cap Y \to Z$ および
$Z \cap Y \to Y$ は閉埋め込みであり、
$$
\xymatrix{
Z \cap Y \ar[r] \ar[d] & Z \ar[d] \\
Y \ar[r] & X
}
$$
は $S$ 上の代数空間の Cartesian 図式である。すなわち
$Z \cap Y = Z \times_X Y$ である。
\end{lemma}

\begin{proof}
射 $Z \cap Y \to Z$ および $Z \cap Y \to Y$ は Lemma
\ref{spaces-morphisms-lemma-closed-immersion-ideals} により閉埋め込みである。
スキーム論的交叉を作ることはエタール局所化と可換なので、
スキームの場合から図式が Cartesian であると結論する。
Morphisms, Lemma
\ref{morphisms-lemma-scheme-theoretic-intersection} を参照。
\end{proof}

\begin{lemma}
\label{spaces-morphisms-lemma-scheme-theoretic-union}
$S$ をスキームとし、$X$ を $S$ 上の代数空間とする。
$Y, Z \subset X$ を閉部分空間とする。
$Y \cup Z$ を $Y$ と $Z$ のスキーム論的和とし、
$Y \cap Z$ を $Y$ と $Z$ のスキーム論的交叉とする。
このとき $Y \to Y \cup Z$ および $Z \to Y \cup Z$ は閉埋め込みであり、
$\mathcal{O}_Z$-加群の短完全列
$$
0 \to \mathcal{O}_{Y \cup Z} \to \mathcal{O}_Y \times \mathcal{O}_Z
\to \mathcal{O}_{Y \cap Z} \to 0
$$
が存在し、図式
$$
\xymatrix{
Y \cap Z \ar[r] \ar[d] & Y \ar[d] \\
Z \ar[r] & Y \cup Z
}
$$
は $S$ 上の代数空間の圏で余 Cartesian である。すなわち
$Y \cup Z = Y \amalg_{Y \cap Z} Z$ である。
\end{lemma}

\begin{proof}
射 $Y \to Y \cup Z$ および $Z \to Y \cup Z$ は Lemma
\ref{spaces-morphisms-lemma-closed-immersion-ideals} により閉埋め込みである。
短完全列では Lemma \ref{spaces-morphisms-lemma-i-star-equivalence} の同値を用い、
$X$ の閉部分空間上の準連接加群を $X$ 上の準連接加群とみなす。
列の第一の写像には標準写像
$\mathcal{O}_{Y \cup Z} \to \mathcal{O}_Y$ および
$\mathcal{O}_{Y \cup Z} \to \mathcal{O}_Z$ を用い、
第二の写像には標準写像
$\mathcal{O}_Y \to \mathcal{O}_{Y \cap Z}$ と、標準写像
$\mathcal{O}_Z \to \mathcal{O}_{Y \cap Z}$ の負を用いる。
完全性はエタール局所的に確認でき、スキームの場合
（Morphisms, Lemma
\ref{morphisms-lemma-scheme-theoretic-union}）から従う。

\medskip\noindent
図式が余 Cartesian であることを示すため、$S$ 上の代数空間 $T$ と、
$Y \cap Z \to T$ として一致する射 $f : Y \to T$ および
$g : Z \to T$ が与えられたと仮定する。目標は、$f$ と $g$ に一致する
一意な射 $h : Y \cup Z \to T$ が存在することを示すことである。
$h$ の構成は $Y \cup Z$ 上エタール局所的に行ってよい
（$Y \cup Z$ は代数空間なのでエタール層である）。
したがって $X$ はスキームであると仮定してよい。この場合、
Morphisms, Lemma
\ref{morphisms-lemma-scheme-theoretic-union} により、
$Y \cup Z$ はスキームの圏で $Y \cap Z$ に沿う $Y$ と $Z$ の押し出しである。
スキーム $T'$ と全射エタール射 $T' \to T$ を選ぶ。
$Y' = T' \times_{T, f} Y$ および
$Z' = T' \times_{T, g} Z$ とおく。このとき $Y'$ と $Z'$ はスキームであり、
スキームの標準同型
$\varphi : Y' \times_Y (Y \cap Z) \to Z' \times_Z (Y \cap Z)$
がある。More on Morphisms, Lemma
\ref{more-morphisms-lemma-pushout-along-closed-immersions} により、押し出し
$W' = Y' \amalg_{Y' \times_Y (Y \cap Z), \varphi} Z'$ は
スキームの圏に存在する。More on Morphisms, Lemma
\ref{more-morphisms-lemma-pushout-along-closed-immersions-properties-above}
により射 $W' \to Y \cup Z$ はエタールである。
$Y' \to Y$ と $Z' \to Z$ が全射なので、これは全射である。
射 $f' : Y' \to T'$ および $g' : Z' \to T'$ は、
一意なスキームの射 $h' : W' \to T'$ に貼り合わされる。
一意性により、合成 $W' \to T' \to T$ は望む射
$h : Y \cup Z \to T$ に降下する。いくつかの詳細は省略する。
\end{proof}






\section{加群の台}
\label{spaces-morphisms-section-support}

\noindent
本節では、代数空間上の準連接加群の台に関するいくつかの初等的結果を集める。
$X$ を代数空間とする。$X_\etale$ 上のアーベル層の台は
Properties of Spaces, Section
\ref{spaces-properties-section-support} で定義した。
加群の台にも同じ定義を用いる。次の補題は、これがスキーム上の
準連接加群について定義された概念と一致することを述べる。

\begin{lemma}
\label{spaces-morphisms-lemma-support-covering}
$S$ をスキームとし、$X$ を $S$ 上の代数空間とする。
$\mathcal{F}$ を準連接 $\mathcal{O}_X$-加群とする。
$U$ をスキームとし、$\varphi : U \to X$ をエタール射とする。このとき
$$
\text{Supp}(\varphi^*\mathcal{F}) = |\varphi|^{-1}(\text{Supp}(\mathcal{F}))
$$
である。ここで左辺はスキーム $U$ 上の準連接加群としての
$\varphi^*\mathcal{F}$ の台である。
\end{lemma}

\begin{proof}
$u\in U$ を（通常の）点とし、$\overline{x}$ を $u$ 上の幾何点とする。
Properties of Spaces, Lemma
\ref{spaces-properties-lemma-stalk-quasi-coherent} により
$(\varphi^*\mathcal{F})_u \otimes_{\mathcal{O}_{U, u}}
\mathcal{O}_{X, \overline{x}} = \mathcal{F}_{\overline{x}}$ である。
Properties of Spaces, Lemma
\ref{spaces-properties-lemma-describe-etale-local-ring} により
$\mathcal{O}_{U, u} \to \mathcal{O}_{X, \overline{x}}$ は強 Hensel 化なので、
忠実平坦である（More on Algebra, Lemma
\ref{more-algebra-lemma-dumb-properties-henselization} を参照）。
したがって $(\varphi^*\mathcal{F})_u = 0$ であることと
$\mathcal{F}_{\overline{x}} = 0$ であることは同値である。
これで補題が示された。
\end{proof}

\noindent
有限型準連接加群の台は閉であり、ファイバー上で確認でき、
基底変換と可換である。

\begin{lemma}
\label{spaces-morphisms-lemma-support-finite-type}
$S$ をスキームとし、$X$ を $S$ 上の代数空間とする。
$\mathcal{F}$ を有限型準連接 $\mathcal{O}_X$-加群とする。このとき
\begin{enumerate}
\item $\mathcal{F}$ の台は閉である。
\item $x \in |X|$ 上の幾何点 $\overline{x}$ に対して
$$
x \in \text{Supp}(\mathcal{F})
\Leftrightarrow
\mathcal{F}_{\overline{x}} \not = 0
\Leftrightarrow
\mathcal{F}_{\overline{x}} \otimes_{\mathcal{O}_{X, \overline{x}}}
\kappa(\overline{x}) \not = 0.
$$
\item 代数空間の任意の射 $f : Y \to X$ に対して、引き戻し
$f^*\mathcal{F}$ も有限型であり、
$\text{Supp}(f^*\mathcal{F}) = f^{-1}(\text{Supp}(\mathcal{F}))$ である。
\end{enumerate}
\end{lemma}

\begin{proof}
スキーム $U$ と全射エタール射 $\varphi : U \to X$ を選ぶ。
Lemma \ref{spaces-morphisms-lemma-support-covering} により、$\mathcal{F}$ の台の逆像は
$\varphi^*\mathcal{F}$ の台であり、Morphisms, Lemma
\ref{morphisms-lemma-support-finite-type} により閉である。
したがって (1) は $|X|$ の位相の定義から従う。

\medskip\noindent
(2) の第一の同値性は台の定義である。第二の同値性は Nakayama の補題から
従う。Algebra, Lemma \ref{algebra-lemma-NAK} を参照。

\medskip\noindent
$f : Y \to X$ を (3) のとおりとする。Properties of Spaces,
Section \ref{spaces-properties-section-properties-modules} により、
$f^*\mathcal{F}$ は有限型である。最後の主張について、
$Y$ の幾何点 $\overline{y}$ が $X$ の幾何点 $\overline{x}$ に写るとする。
$$
(f^*\mathcal{F})_{\overline{y}} =
\mathcal{F}_{\overline{x}} \otimes_{\mathcal{O}_{X, \overline{x}}}
\mathcal{O}_{Y, \overline{y}},
$$
であることを思い出す。Properties of Spaces, Lemma
\ref{spaces-properties-lemma-stalk-pullback-quasi-coherent} を参照。
したがって
$(f^*\mathcal{F})_{\overline{y}} \otimes \kappa(\overline{y})$ が非零であることと
$\mathcal{F}_{\overline{x}} \otimes \kappa(\overline{x})$ が非零であることは
同値である。(2) により、これは
$x \in \text{Supp}(\mathcal{F})$ と
$y \in \text{Supp}(f^*\mathcal{F})$ の同値性、すなわち主張 (3) を与える。
\end{proof}

\noindent
次の課題は、有限型準連接加群のスキーム論的台
（Morphisms, Definition
\ref{morphisms-definition-scheme-theoretic-support} を参照）が
代数空間上の有限型準連接加群についても意味をもつことを示すことである。

\begin{lemma}
\label{spaces-morphisms-lemma-scheme-theoretic-support}
$S$ をスキームとし、$X$ を $S$ 上の代数空間とする。
$\mathcal{F}$ を有限型準連接 $\mathcal{O}_X$-加群とする。
準連接 $\mathcal{O}_Z$-加群 $\mathcal{G}$ で
$i_*\mathcal{G} \cong \mathcal{F}$ となるものが存在するような
最小の閉部分空間 $i : Z \to X$ が存在する。さらに次が成り立つ。
\begin{enumerate}
\item $U$ がスキームで $\varphi : U \to X$ がエタール射ならば、
$Z \times_X U$ は $\varphi^*\mathcal{F}$ のスキーム論的台である。
\item 準連接層 $\mathcal{G}$ は一意な同型を除いて一意である。
\item 準連接層 $\mathcal{G}$ は有限型である。
\item $\mathcal{G}$ と $\mathcal{F}$ の台は $|Z|$ である。
\end{enumerate}
\end{lemma}

\begin{proof}
スキーム $U$ と全射エタール射 $\varphi : U \to X$ を選ぶ。
$R = U \times_X U$ とし、射影を $s, t : R \to U$ とする。
$i' : Z' \to U$ を $\varphi^*\mathcal{F}$ のスキーム論的台とし、
$\mathcal{G}'$ を、一意な同型を除いて一意な有限型準連接
$\mathcal{O}_{Z'}$-加群で
$i'_*\mathcal{G}' = \varphi^*\mathcal{F}$ となるものとする。
Morphisms, Lemma \ref{morphisms-lemma-scheme-theoretic-support} を参照。
$s^*\varphi^*\mathcal{F} = t^*\varphi^*\mathcal{F}$ なので、
Morphisms, Lemma
\ref{morphisms-lemma-flat-pullback-support} により、
$R$ の閉部分スキームとして
$R' = s^{-1}Z' = t^{-1}Z'$ である。
したがって Properties of Spaces, Lemma
\ref{spaces-properties-lemma-subspaces-presentation} を適用し、
$U$ への引き戻しが $Z'$ となる閉部分空間 $i : Z \to X$ を得る。
射影を $s', t' : R' \to Z'$ と書き、与えられた閉埋め込みを
$j' : R' \to R$ と書くと、
$$
j'_* (s')^*\mathcal{G}' = s^* i'_*\mathcal{G}' =
s^*\varphi^*\mathcal{F} = t^*\varphi^*\mathcal{F} =
t^*i'_*\mathcal{G}' = j'_*(t')^*\mathcal{G}'
$$
を得る（最初と最後の等式には Cohomology of Schemes, Lemma
\ref{coherent-lemma-flat-base-change-cohomology} を用いた）。
そこで $R' \to R$ に Morphisms, Lemma
\ref{morphisms-lemma-flat-pullback-support} の一意性を適用すると、
$j'_*$ を介した標準同型
$t^*\varphi^*\mathcal{F} = s^*\varphi^*\mathcal{F}$ と整合する同型
$\alpha : (t')^*\mathcal{G}' \to (s')^*\mathcal{G}'$ を得る。
明らかに $\alpha$ は余サイクル条件を満たすので、Properties of Spaces,
Proposition \ref{spaces-properties-proposition-quasi-coherent} を適用し、
$Z'$ への制限が $\alpha$ と整合する $\mathcal{G}'$ である
$Z$ 上の準連接加群 $\mathcal{G}$ を得る。
再び上記命題の同値（今度は $X$ について）を用いると、
$i_*\mathcal{G} \cong \mathcal{F}$ と結論する。

\medskip\noindent
これで存在が示された。補題の他の性質は、Lemma
\ref{spaces-morphisms-lemma-support-covering} を用いてスキームの場合の結果と比較すれば従う。
詳細な証明は省略する。
\end{proof}

\begin{definition}
\label{spaces-morphisms-definition-scheme-theoretic-support}
$S$ をスキームとし、$X$ を $S$ 上の代数空間とする。
$\mathcal{F}$ を有限型準連接 $\mathcal{O}_X$-加群とする。
{\it $\mathcal{F}$ のスキーム論的台}とは、Lemma
\ref{spaces-morphisms-lemma-scheme-theoretic-support} で構成した閉部分空間
$Z \subset X$ のことである。
\end{definition}

\noindent
この状況では、Lemma \ref{spaces-morphisms-lemma-i-star-equivalence} の圏同値を介して、
$\mathcal{F}$ を $Z$ 上の有限型準連接層とみなすことが多い。







\section{スキーム論的像}
\label{spaces-morphisms-section-scheme-theoretic-image}

\noindent
注意：この節の内容の一部は極めて一般的なものであり、
予想とは異なる振る舞いをすることがある。

\begin{lemma}
\label{spaces-morphisms-lemma-scheme-theoretic-image}
\begin{slogan}
代数空間の射のスキーム論的像は存在する。
\end{slogan}
$S$ をスキームとする。$f : X \to Y$ を $S$ 上の代数空間の射とする。
閉部分空間 $Z \subset Y$ であって $f$ が $Z$ を経由し、かつ、任意の
他の閉部分空間 $Z' \subset Y$ に対して、$f$ が $Z'$ を経由するならば
$Z \subset Z'$ となるものが存在する。
\end{lemma}

\begin{proof}
$\mathcal{I} = \Ker(\mathcal{O}_Y \to f_*\mathcal{O}_X)$ とおく。
$\mathcal{I}$ が準連接ならば、Lemma \ref{spaces-morphisms-lemma-closed-immersion-ideals}
を参照して、$\mathcal{I}$ が定める閉部分スキームを $Z$ とすればよい。
一般の場合には、$\mathcal{I}$ に含まれる最大の準連接イデアル層
$\mathcal{I}'$ が存在することを示す必要がある。
これは Lemma \ref{spaces-morphisms-lemma-largest-quasi-coherent-subsheaf} から従う。
\end{proof}

\noindent
上の Lemma \ref{spaces-morphisms-lemma-scheme-theoretic-image} の状況で $X$ と $Y$ が
表現可能であるとする。このとき、補題で得られた閉部分空間
$Z \subset Y$ は、Morphisms, Lemma
\ref{morphisms-lemma-scheme-theoretic-image} で得られる閉部分スキーム
$Z \subset Y$ と一致する。その理由は、閉部分空間（または閉部分スキーム）と
$X_\etale$ および $X$ 上の準連接イデアル層との間に、包含を逆転する対応が
あるからである。$X_\etale$ 上と $X$ 上の準連接加群の圏は同じなので
（Properties of Spaces, Section
\ref{spaces-properties-section-quasi-coherent}）、結論を得る。
したがって次の定義は、スキームの射に対する先の定義と一致する。

\begin{definition}
\label{spaces-morphisms-definition-scheme-theoretic-image}
$S$ をスキームとする。$f : X \to Y$ を $S$ 上の代数空間の射とする。
$f$ の {\it スキーム論的像}とは、$f$ が経由する最小の閉部分空間
$Z \subset Y$ のことである。上の Lemma
\ref{spaces-morphisms-lemma-scheme-theoretic-image} を参照せよ。
\end{definition}

\noindent
$f$ のこの分解を単に $f : X \to Z$ と表すことが多い。
射 $f$ が準コンパクトでなければ、一般には、スキーム論的像の構成は
$Y$ の開部分空間への制限と可換しない。

\begin{lemma}
\label{spaces-morphisms-lemma-quasi-compact-scheme-theoretic-image}
$S$ をスキームとする。
$f : X \to Y$ を $S$ 上の代数空間の射とする。
$Z \subset Y$ を $f$ のスキーム論的像とする。
$f$ が準コンパクトならば、
\begin{enumerate}
\item イデアル層
$\mathcal{I} = \Ker(\mathcal{O}_Y \to f_*\mathcal{O}_X)$
は準連接であり、
\item スキーム論的像 $Z$ は $\mathcal{I}$ に対応する閉部分空間であり、
\item 任意のエタール射 $V \to Y$ に対して、
$X \times_Y V \to V$ のスキーム論的像は $Z \times_Y V$ に等しく、
\item 像 $|f|(|X|) \subset |Z|$ は $|Z|$ の稠密部分集合である。
\end{enumerate}
\end{lemma}

\begin{proof}
$\mathcal{I}$ の形成はエタール局所化と可換するので、(3) を示すには
(1) と (2) を示せば十分である。(1) が成り立てば、Lemma
\ref{spaces-morphisms-lemma-scheme-theoretic-image} の証明で (2) を示している。
$\mathcal{I}$ が準連接であることを示そう。準連接であるという性質は
エタール局所的なので、$Y$ はアフィンスキームであると仮定してよい。
$f$ は準コンパクトなので、アフィンスキーム $U$ と全射エタール射
$U \to X$ を取ることができる。合成 $U \to X \to Y$ を $f'$ と表す。
このとき $f_*\mathcal{O}_X$ は $f'_*\mathcal{O}_U$ の部分層であり、したがって
$\mathcal{I} = \Ker(\mathcal{O}_Y \to \mathcal{O}_{X'})$ である。
Lemma \ref{spaces-morphisms-lemma-pushforward} により、層 $f'_*\mathcal{O}_U$ は $Y$ 上準連接である。
ゆえに $\mathcal{I}$ は連接加群間の写像の核として準連接である。
最後に、$|f|(|X|)$ の閉包に含まれない $|Y|$ の任意の点では
イデアル $\mathcal{I}$ が単位イデアルとなるので、(4) は (1)、(2)、(3) から従う。
\end{proof}

\begin{lemma}
\label{spaces-morphisms-lemma-scheme-theoretic-image-reduced}
$S$ をスキームとする。$f : X \to Y$ を $S$ 上の代数空間の射とする。
$X$ が被約であると仮定する。このとき
\begin{enumerate}
\item $f$ のスキーム論的像 $Z$ は、$\overline{|f|(|X|)}$ 上に誘導される
被約代数空間構造であり、
\item 任意のエタール射 $V \to Y$ に対して、
$X \times_Y V \to V$ のスキーム論的像は $Z \times_Y V$ に等しい。
\end{enumerate}
\end{lemma}

\begin{proof}
(1) が成り立つのは、$\overline{|f|(|X|)}$ 上に誘導される被約代数空間構造が、
$f$ が経由する $Y$ の最小の閉部分空間だからである。Properties of Spaces,
Lemma \ref{spaces-properties-lemma-map-into-reduction} を参照せよ。
(2) は (1)、$|V| \to |Y|$ が開写像であること、および被約性が
エタール局所化で保たれることから従う。
\end{proof}

\begin{lemma}
\label{spaces-morphisms-lemma-reach-points-scheme-theoretic-image}
$S$ をスキームとする。
$f : X \to Y$ を $S$ 上の代数空間の準コンパクト射とする。
$Z$ を $f$ のスキーム論的像とする。$z \in |Z|$ とする。
分数体 $K$ をもつ付値環 $A$ と可換図式
$$
\xymatrix{
\Spec(K) \ar[rr] \ar[d] & & X \ar[d] \ar[ld] \\
\Spec(A) \ar[r] & Z \ar[r] & Y
}
$$
であって、$\Spec(A)$ の閉点が $z$ に写るものが存在する。
\end{lemma}

\begin{proof}
点 $z' \in V$ をもつアフィンスキーム $V$ と、$z'$ を $z$ に写す
エタール射 $V \to Y$ を取る。
$Z' \subset V$ を $X \times_Y V \to V$ のスキーム論的像とする。
Lemma \ref{spaces-morphisms-lemma-quasi-compact-scheme-theoretic-image} により
$Z' = Z \times_Y V$ である。したがって $z' \in Z'$ である。
$f$ は準コンパクトで $V$ はアフィンなので、$X \times_Y V$ は準コンパクトである。
したがって、アフィンスキーム $W$ と全射エタール射
$W \to X \times_Y V$ が存在する。このとき $Z' \subset V$ は
$W \to V$ のスキーム論的像でもある。
Morphisms, Lemma \ref{morphisms-lemma-reach-points-scheme-theoretic-image}
により、図式
$$
\xymatrix{
\Spec(K) \ar[r] \ar[d] &
W \ar[r] \ar[d] &
X \times_Y V \ar[d] \ar[r] &
X \ar[d] \\
\Spec(A) \ar[r] &
Z' \ar[r] &
V \ar[r] &
Y
}
$$
であって、$\Spec(A)$ の閉点が $z'$ に写るものを取ることができる。
$Z' \to Z$ および $W \to X \times_Y V \to X$ と合成すれば所望の図式を得る。
\end{proof}

\begin{lemma}
\label{spaces-morphisms-lemma-factor-factor}
$S$ をスキームとし、
$$
\xymatrix{
X_1 \ar[d] \ar[r]_{f_1} & Y_1 \ar[d] \\
X_2 \ar[r]^{f_2} & Y_2
}
$$
を $S$ 上の代数空間の可換図式とする。
$Z_i \subset Y_i$、$i = 1, 2$ を $f_i$ のスキーム論的像とする。
このとき射 $Y_1 \to Y_2$ は射 $Z_1 \to Z_2$ および可換図式
$$
\xymatrix{
X_1 \ar[r] \ar[d] & Z_1 \ar[d] \ar[r] & Y_1 \ar[d] \\
X_2 \ar[r] & Z_2 \ar[r] & Y_2
}
$$
を誘導する。
\end{lemma}

\begin{proof}
$Y_1$ における $Z_2$ のスキーム論的逆像は、$f_1$ が経由する
$Y_1$ の閉部分空間である。したがって $Z_1$ はこの逆像に含まれる。
これで補題が示された。
\end{proof}

\begin{lemma}
\label{spaces-morphisms-lemma-scheme-theoretic-image-of-partial-section}
$S$ をスキームとする。
$f : X \to Y$ を $S$ 上の代数空間の分離射とする。
$V \subset Y$ を、$V \to Y$ が準コンパクトであるような開部分空間とする。
$s : V \to X$ を $f \circ s = \text{id}_V$ を満たす射とする。
$Y'$ を $s$ のスキーム論的像とする。このとき $Y' \to Y$ は
$V$ 上で同型である。
\end{lemma}

\begin{proof}
Lemma \ref{spaces-morphisms-lemma-quasi-compact-permanence} により、射
$s : V \to X$ は準コンパクトである。したがって Lemma
\ref{spaces-morphisms-lemma-quasi-compact-scheme-theoretic-image} により、$s$ の
スキーム論的像 $Y'$ の構成は開集合への制限と可換する。
特に、$Y' \cap f^{-1}(V)$ は分離射 $f^{-1}(V) \to V$ の
切断のスキーム論的像であることが分かる。
分離射の切断は閉埋入なので（Lemma \ref{spaces-morphisms-lemma-section-immersion}）、
$Y' \cap f^{-1}(V) \to V$ は所望のように同型であると結論する。
\end{proof}







\section{スキーム論的閉包と稠密性}
\label{spaces-morphisms-section-scheme-theoretic-closure}

\noindent
この節は Morphisms, Section
\ref{morphisms-section-scheme-theoretic-closure} に対応するものである。

\begin{lemma}
\label{spaces-morphisms-lemma-scheme-theoretically-dense-representable}
$S$ をスキームとする。$W \subset S$ をスキーム論的に稠密な開部分スキーム
（Morphisms, Definition
\ref{morphisms-definition-scheme-theoretically-dense}）とする。
$f : X \to S$ を、平坦、局所有限表示かつ局所準有限であるスキームの射とする。
このとき $f^{-1}(W)$ は $X$ でスキーム論的に稠密である。
\end{lemma}

\begin{proof}
Morphisms, Lemma
\ref{morphisms-lemma-characterize-scheme-theoretically-dense} の特徴づけを用いる。
$V \subset X$ を開集合とし、$g \in \Gamma(V, \mathcal{O}_V)$ を
$f^{-1}(W) \cap V$ 上でゼロに制限される函数とする。
$g = 0$ を示さなければならない。矛盾を導くため $g \not = 0$ と仮定する。
More on Morphisms, Lemma
\ref{more-morphisms-lemma-go-down-with-annihilators} により、$V$ を縮小し、
可換図式
$$
\xymatrix{
V \ar[r] \ar[d]_\pi & X \ar[d]^f \\
U \ar[r] & S,
}
$$
に収まる開集合 $U \subset S$、準連接部分層
$\mathcal{F} \subset \mathcal{O}_U$、整数 $r > 0$、および像が
$g|_V$ を含む単射な $\mathcal{O}_U$-加群写像
$\mathcal{F}^{\oplus r} \to \pi_*\mathcal{O}_V$ を見いだせる。
$(g_1, \ldots, g_r) \in \Gamma(U, \mathcal{F}^{\oplus r})$ が $g$ に写るとする。
$g|_{f^{-1}W \cap V} = 0$ なので、$g_i|_{W \cap U} = 0$ であることが分かる。
$\mathcal{F} \subset \mathcal{O}_U$ であり、$W$ は $S$ で
スキーム論的に稠密なので、$g_i = 0$ である。
これは $g = 0$ を意味し、所望の矛盾である。
\end{proof}

\begin{lemma}
\label{spaces-morphisms-lemma-scheme-theoretically-dense}
$S$ をスキームとする。
$X$ を $S$ 上の代数空間とする。
$U \subset X$ を開部分空間とする。
次の条件は同値である。
\begin{enumerate}
\item 代数空間の任意のエタール射 $\varphi : V \to X$ に対して、
$V$ における $\varphi^{-1}(U)$ のスキーム論的閉包は $V$ に等しい。
\item スキーム $V$ と全射エタール射 $\varphi : V \to X$ であって、
$V$ における $\varphi^{-1}(U)$ のスキーム論的閉包が $V$ に等しいものが存在する。
\end{enumerate}
\end{lemma}

\begin{proof}
$V \to V'$ を $X$ 上エタールな代数空間の射とし、$Z \subset V$、それぞれ
$Z' \subset V'$ を、$V$、それぞれ $V'$ における
$U \times_X V$、それぞれ $U \times_X V'$ のスキーム論的閉包とする。
このとき $Z$ は $Z'$ の中へ写る。したがって $V \to V'$ が全射かつ
エタールならば、$Z = V$ は $Z' = V'$ を含意する。
次に、エタール射は平坦、局所有限表示かつ局所準有限であることに注意する
（Morphisms, Section \ref{morphisms-section-etale} 参照）。
したがって Lemma \ref{spaces-morphisms-lemma-scheme-theoretically-dense-representable}
により、$V$ と $V'$ がスキームならば、$Z' = V'$ は $Z = V$ を含意する。
すべての代数空間がスキームによるエタール被覆をもつことを用いる形式的議論から、
(1) と (2) は同値である。
\end{proof}

\noindent
Lemma \ref{spaces-morphisms-lemma-scheme-theoretically-dense} から、次の定義が
スキームの場合の定義と両立することが従う。

\begin{definition}
\label{spaces-morphisms-definition-scheme-theoretically-dense}
$S$ をスキームとする。
$X$ を $S$ 上の代数空間とする。
$U \subset X$ を開部分空間とする。
\begin{enumerate}
\item 射 $U \to X$ のスキーム論的像を、
{\it $X$ における $U$ のスキーム論的閉包}という。
\item Lemma \ref{spaces-morphisms-lemma-scheme-theoretically-dense} の同値な条件が満たされるとき、
$U$ は {\it $X$ でスキーム論的に稠密である}という。
\end{enumerate}
\end{definition}

\noindent
この定義では、$U$ が $X$ でスキーム論的に稠密であることと、
$U$ のスキーム論的閉包が $X$ であることとは、同値で\textbf{ない}。
これはいささか洗練されていない。しかし適切な有限性条件のもとでは、
両者が同値になることを見る。

\begin{lemma}
\label{spaces-morphisms-lemma-scheme-theoretically-dense-quasi-compact}
$S$ をスキームとする。$X$ を $S$ 上の代数空間とする。
$U \subset X$ を開部分空間とする。
$U \to X$ が準コンパクトならば、$U$ が $X$ でスキーム論的に稠密であることと、
$X$ における $U$ のスキーム論的閉包が $X$ であることとは同値である。
\end{lemma}

\begin{proof}
Lemma \ref{spaces-morphisms-lemma-quasi-compact-scheme-theoretic-image} の (3) から従う。
\end{proof}

\begin{lemma}
\label{spaces-morphisms-lemma-characterize-scheme-theoretically-dense}
$S$ をスキームとする。
$j : U \to X$ を $S$ 上の代数空間の開埋入とする。
$U$ が $X$ でスキーム論的に稠密であるための必要十分条件は、
$\mathcal{O}_X \to j_*\mathcal{O}_U$ が単射であることである。
\end{lemma}

\begin{proof}
$\mathcal{O}_X \to j_*\mathcal{O}_U$ が単射ならば、$X$ 上エタールな
任意の代数空間 $V$ への制限についても同じことが成り立つ。
したがって $V$ における $U \times_X V$ のスキーム論的閉包は $V$ に等しい。
Lemma \ref{spaces-morphisms-lemma-scheme-theoretic-image} の証明を参照せよ。
逆に、$X$ 上エタールなすべての $V$ に対して、$U \times_X V$ の
スキーム論的閉包が $V$ に等しいと仮定する。
$\mathcal{O}_X \to j_*\mathcal{O}_U$ が単射でないと仮定する。
すると $X$ 上エタールなアフィン、たとえば $V = \Spec(A)$ と、
非零元 $f \in A$ であって、$f$ が
$\Gamma(V \times_X U, \mathcal{O})$ でゼロに写るものを取れる。
この場合、$V$ における $V \times_X U$ のスキーム論的閉包は明らかに
$\Spec(A/(f))$ に含まれ、矛盾である。
\end{proof}

\begin{lemma}
\label{spaces-morphisms-lemma-intersection-scheme-theoretically-dense}
$S$ をスキームとする。$X$ を $S$ 上の代数空間とする。
$U$, $V$ が $X$ のスキーム論的に稠密な開部分空間ならば、
$U \cap V$ もそうである。
\end{lemma}

\begin{proof}
$W \to X$ を任意のエタール射とする。写像
$\mathcal{O}(W) \to \mathcal{O}(W \times_X V)
\to \mathcal{O}(W \times_X (V \cap U))$ を考える。
Lemma \ref{spaces-morphisms-lemma-characterize-scheme-theoretically-dense} により、
両方の写像は単射である。したがって合成も単射である。
ゆえに再び Lemma \ref{spaces-morphisms-lemma-characterize-scheme-theoretically-dense}
により、$U \cap V$ は $X$ でスキーム論的に稠密である。
\end{proof}

\begin{lemma}
\label{spaces-morphisms-lemma-quasi-compact-immersion}
$S$ をスキームとする。$h : Z \to X$ を $S$ 上の代数空間の埋入とする。
$Z \to X$ が準コンパクトであるか、または $Z$ が被約であると仮定する。
$\overline{Z} \subset X$ を $h$ のスキーム論的像とする。
このとき射 $Z \to \overline{Z}$ は開埋入であり、$Z$ を
$\overline{Z}$ のスキーム論的に稠密な開部分空間と同一視する。
さらに、$Z$ は $\overline{Z}$ で位相的に稠密である。
\end{lemma}

\begin{proof}
どちらの場合にも $\overline{Z}$ の形成はエタール局所化と可換する。
Lemmas \ref{spaces-morphisms-lemma-quasi-compact-scheme-theoretic-image} および
\ref{spaces-morphisms-lemma-scheme-theoretic-image-reduced} を参照せよ。
したがってこの補題はスキームの場合から従う。Morphisms, Lemma
\ref{morphisms-lemma-quasi-compact-immersion} を参照せよ。
\end{proof}

\begin{lemma}
\label{spaces-morphisms-lemma-equality-of-morphisms}
$S$ をスキームとする。$B$ を $S$ 上の代数空間とする。
$f, g : X \to Y$ を $B$ 上の代数空間の射とする。
$U \subset X$ を $f|_U = g|_U$ となる開部分空間とする。
$X$ における $U$ のスキーム論的閉包が $X$ であり、$Y \to B$ が分離ならば、
$f = g$ である。
\end{lemma}

\begin{proof}
$Y \to B$ は分離なので、ファイバー積
$Y \times_{\Delta, Y \times_B Y, (f, g)} X$ は閉部分空間 $Z \subset X$ である。
$f|_U = g|_U$ なので $U \subset Z$ である。$U$ は $X$ で
スキーム論的に稠密と仮定されているから、$Z = X$ である。
\end{proof}






\section{支配的射}
\label{spaces-morphisms-section-dominant}

\noindent
スキームの支配的射の定義をそのまま移して、代数空間の支配的射の概念を得る。
射が準コンパクトで、代数空間がいくつかの分離公理を満たさないかぎり、
この定義は良い振る舞いをしないことに注意されたい。

\begin{definition}
\label{spaces-morphisms-definition-dominant}
$S$ をスキームとする。$f : X \to Y$ を $S$ 上の代数空間の射とする。
$|f| : |X| \to |Y|$ の像が $|Y|$ で稠密ならば、この射を
{\it 支配的}という。
\end{definition}






\section{普遍的単射な射}
\label{spaces-morphisms-section-universally-injective}

\noindent
代数空間の表現可能射が普遍的単射であるとは何かは、すでに Section
\ref{spaces-morphisms-section-representable} で定義した。$S$ 上の体 $K$（これは構造射
$\Spec(K) \to S$ が与えられているという意味であることを想起せよ）と、
$S$ 上の代数空間 $X$ に対し、
$X(K) = \Mor_S(\Spec(K), X)$ と書く。まず、表現可能射に対する条件を
点の関手に対する条件へ翻訳する。

\begin{lemma}
\label{spaces-morphisms-lemma-universally-injective-representable}
$S$ をスキームとする。$f : X \to Y$ を $S$ 上の代数空間の表現可能射とする。
$f$ が Section \ref{spaces-morphisms-section-representable} の意味で普遍的単射であるための
必要十分条件は、すべての体 $K$ に対して写像 $X(K) \to Y(K)$ が単射であることである。
\end{lemma}

\begin{proof}
以下では Morphisms, Lemma
\ref{morphisms-lemma-universally-injective} を断りなく用いる。
$f$ が普遍的単射であると仮定する。このとき任意の体 $K$ と任意の射
$\Spec(K) \to Y$ に対して、スキームの射
$\Spec(K) \times_Y X \to \Spec(K)$ は普遍的単射である。
したがって射 $\Spec(K) \times_Y X \to \Spec(K)$ の切断は高々一つである。
ゆえに写像 $X(K) \to Y(K)$ は単射である。逆に、すべての体 $K$ に対して
写像 $X(K) \to Y(K)$ が単射であると仮定する。
$T \to Y$ をスキームから $Y$ への射とし、基底変換
$f_T : T \times_Y X \to T$ を考える。任意の体 $K$ に対して、
ファイバー積の定義により
$$
(T \times_Y X)(K) = T(K) \times_{Y(K)} X(K)
$$
である。したがって $X(K) \to Y(K)$ の単射性は
$(T \times_Y X)(K) \to T(K)$ の単射性を保証し、これは所望のように
$f_T$ が普遍的単射であることを意味する。
\end{proof}

\noindent
次に、体に値をもつ点の間の変換が単射であるという性質を、
より幾何学的な条件へ翻訳する。

\begin{lemma}
\label{spaces-morphisms-lemma-universally-injective}
$S$ をスキームとする。
$f : X \to Y$ を $S$ 上の代数空間の射とする。
次の条件は同値である。
\begin{enumerate}
\item $S$ 上のすべての体 $K$ に対して写像 $X(K) \to Y(K)$ が単射である。
\item $S$ 上の代数空間の任意の射 $Y' \to Y$ に対して、誘導写像
$|Y' \times_Y X| \to |Y'|$ が単射である。
\item 対角射 $X \to X \times_Y X$ が全射である。
\end{enumerate}
\end{lemma}

\begin{proof}
(1) を仮定する。$g : Y' \to Y$ を代数空間の射とし、$f$ の基底変換を
$f' : Y' \times_Y X \to Y'$ と表す。
$K_i$、$i = 1, 2$ を体とし、$\varphi_i : \Spec(K_i) \to Y' \times_Y X$ を、
$f' \circ \varphi_1$ と $f' \circ \varphi_2$ が $|Y'|$ の同じ元を
定めるような射とする。定義により、これは体 $\Omega$ と埋め込み
$\alpha_i : K_i \subset \Omega$ であって、二つの射
$f' \circ \varphi_i \circ \alpha_i : \Spec(\Omega) \to Y'$ が
等しくなるものが存在することを意味する。対応する可換図式は
$$
\xymatrix{
\Spec(\Omega) \ar@/_5ex/[ddrr] \ar[rd]^{\alpha_1} \ar[r]_{\alpha_2} &
\Spec(K_2) \ar[rd]^{\varphi_2} \\
& \Spec(K_1) \ar[r]^{\varphi_1} &
Y' \times_Y X  \ar[d]^{f'} \ar[r]^{g'} &
 X \ar[d]^f \\
& & Y' \ar[r]^g & Y.
}
$$
である。特に合成 $g \circ f' \circ \varphi_i \circ \alpha_i$ は等しい。
仮定 (1) により、射 $g' \circ \varphi_i \circ \alpha_i$ は等しい。
ここで $g' : Y' \times_Y X \to X$ は射影である。
ファイバー積の普遍性により、射
$\varphi_i \circ \alpha_i : \Spec(\Omega) \to Y' \times_Y X$ は等しい。
言い換えると、$\varphi_1$ と $\varphi_2$ は $Y' \times_Y X$ の同じ点を定める。
したがって (2) が成り立つ。

\medskip\noindent
(2) を仮定する。$K$ を $S$ 上の体とし、$a, b : \Spec(K) \to X$ を
$f \circ a = f \circ b$ を満たす二つの射とする。その共通値を
$c : \Spec(K) \to Y$ と表す。仮定により
$|\Spec(K) \times_{c, Y} X| \to |\Spec(K)|$ は単射である。
これは、体 $\Omega$ と埋め込み $\alpha_i : K \to \Omega$ であって、図式
$$
\xymatrix{
\Spec(\Omega) \ar[r]_{\alpha_1} \ar[d]_{\alpha_2} &
\Spec(K) \ar[d]^a \\
\Spec(K) \ar[r]^-b &
\Spec(K) \times_{c, Y} X
}
$$
が可換となるものが存在することを意味する。$\Spec(K)$ への射影と合成すると、
$\alpha_1 = \alpha_2$ であることが分かる。その共通値を $\alpha$ と表す。
すると $\{\alpha : \Spec(\Omega) \to \Spec(K)\}$ は $\Spec(K)$ の
fpqc 被覆であり、その被覆の各要素上で二つの射 $a, b$ は等しくなる。
Properties of Spaces, Proposition
\ref{spaces-properties-proposition-sheaf-fpqc} により $a = b$ と結論する。
したがって (1) が成り立つ。

\medskip\noindent
(3) を仮定する。$x, x' \in |X|$ を $|Y|$ で $f(x) = f(x')$ となる
二点とする。Properties of Spaces, Lemma
\ref{spaces-properties-lemma-points-cartesian} により、射影が $x$ と $x'$ である
$x'' \in |X \times_Y X|$ が存在する。仮定および Properties of Spaces,
Lemma \ref{spaces-properties-lemma-characterize-surjective} により、
$\Delta_{X/Y}(x''') = x''$ を満たす $x''' \in |X|$ が存在する。
したがって $x = x'$ である。言い換えると $f$ は単射である。
条件 (3) は基底変換で安定なので、$f$ は (2) を満たす。

\medskip\noindent
(2) を仮定する。このとき特に $|X \times_Y X| \to |X|$ は単射であり、
これは直ちに $|\Delta_{X/Y}| : |X| \to |X \times_Y X|$ が全射であることを含意する。
ゆえに Properties of Spaces, Lemma
\ref{spaces-properties-lemma-characterize-surjective} により、
$\Delta_{X/Y}$ は全射である。
\end{proof}

\noindent
以上二つの補題により、次の定義は、すでに定義した代数空間の
普遍的単射な表現可能射の概念と矛盾しない。

\begin{definition}
\label{spaces-morphisms-definition-universally-injective}
$S$ をスキームとする。$f : X \to Y$ を $S$ 上の代数空間の射とする。
任意の射 $Y' \to Y$ に対して誘導写像
$|Y' \times_Y X| \to |Y'|$ が単射であるとき、$f$ は
{\it 普遍的単射}であるという。
\end{definition}

\noindent
正確には、これは Lemma \ref{spaces-morphisms-lemma-universally-injective} の同値な条件の
いずれか、したがってすべてが成り立つことを意味する。

\begin{remark}
\label{spaces-morphisms-remark-universally-injective-not-separated}
スキームの普遍的単射は分離である。Morphisms, Lemma
\ref{morphisms-lemma-universally-injective-separated} を参照せよ。
代数空間の射についてはそうではない。実際、Spaces, Example
\ref{spaces-example-affine-line-involution} で構成された代数空間
$X = \mathbf{A}^1_k/\{x \sim -x \mid x \not = 0\}$ には射
$X \to \mathbf{A}^1_k$ が備わっており、座標 $x$ の点を座標 $x^2$ の点へ写す。
これは $0$ の外で同型であり、$0$ の上にある $X$ の点は一意である。
$X$ は分離でないので、これは分離でない代数空間の普遍的単射である。
\end{remark}

\begin{lemma}
\label{spaces-morphisms-lemma-base-change-universally-injective}
普遍的単射の基底変換は普遍的単射である。
\end{lemma}

\begin{proof}
省略する。ヒント：これは形式的である。
\end{proof}

\begin{lemma}
\label{spaces-morphisms-lemma-universally-injective-local}
$S$ をスキームとする。
$f : X \to Y$ を $S$ 上の代数空間の射とする。
次の条件は同値である。
\begin{enumerate}
\item $f$ は普遍的単射である。
\item 任意のスキーム $Z$ と任意の射 $Z \to Y$ に対して、射
$Z \times_Y X \to Z$ は普遍的単射である。
\item 任意のアフィンスキーム $Z$ と任意の射 $Z \to Y$ に対して、射
$Z \times_Y X \to Z$ は普遍的単射である。
\item スキーム $Z$ と全射 $Z \to Y$ であって、
$Z \times_Y X \to Z$ が普遍的単射となるものが存在する。
\item ザリスキ被覆 $Y = \bigcup Y_i$ であって、各射
$f^{-1}(Y_i) \to Y_i$ が普遍的単射となるものが存在する。
\end{enumerate}
\end{lemma}

\begin{proof}
この証明では、普遍的単射であることが基底変換で保たれること
（Lemma \ref{spaces-morphisms-lemma-base-change-universally-injective}）を断りなく用いる。
(1) $\Rightarrow$ (2) $\Rightarrow$ (3) $\Rightarrow$ (4) は明らかである。

\medskip\noindent
(4) のような $g : Z \to Y$ を仮定する。$y : \Spec(K) \to Y$ を
体のスペクトルから $Y$ への射とする。仮定により、拡大体
$\alpha : K \subset K'$ と射 $z : \Spec(K') \to Z$ であって、
$y \circ \alpha = g \circ z$ となるものを取れる（記法を明らかに濫用している）。
仮定により射 $Z \times_Y X \to Z$ は普遍的単射なので、
$g \circ z : \Spec(K') \to Y$ の $X$ への持ち上げは高々一つである。
$\{\alpha : \Spec(K') \to \Spec(K)\}$ は fpqc 被覆だから、これは
$y : \Spec(K) \to Y$ の $X$ への持ち上げが高々一つであることを意味する。
Properties of Spaces, Proposition
\ref{spaces-properties-proposition-sheaf-fpqc} を参照せよ。
したがって (1) が成り立つ。

\medskip\noindent
(5) と (1) が同値であることの確認は省略する。
\end{proof}

\begin{lemma}
\label{spaces-morphisms-lemma-composition-universally-injective}
普遍的単射の合成は普遍的単射である。
\end{lemma}

\begin{proof}
省略する。
\end{proof}


\section{アフィン射}
\label{spaces-morphisms-section-affine}

\noindent
代数空間の表現可能射がアフィンであるとは何かは、すでに Section
\ref{spaces-morphisms-section-representable} で定義した。

\begin{lemma}
\label{spaces-morphisms-lemma-affine-representable}
$S$ をスキームとする。$f : X \to Y$ を $S$ 上の代数空間の表現可能射とする。
$f$ が Section \ref{spaces-morphisms-section-representable} の意味でアフィンであるための
必要十分条件は、すべてのアフィンスキーム $Z$ と射 $Z \to Y$ に対して、
スキーム $X \times_Y Z$ がアフィンであることである。
\end{lemma}

\begin{proof}
これはスキームのアフィン射の定義から直ちに従う
（Morphisms, Definition \ref{morphisms-definition-affine}）。
\end{proof}

\noindent
これにより次の定義が可能になる。

\begin{definition}
\label{spaces-morphisms-definition-affine}
$S$ をスキームとする。
$f : X \to Y$ を $S$ 上の代数空間の射とする。
任意のアフィンスキーム $Z$ と射 $Z \to Y$ に対して、代数空間
$X \times_Y Z$ がアフィンスキームで表現可能であるとき、$f$ は
{\it アフィン}であるという。
\end{definition}

\begin{lemma}
\label{spaces-morphisms-lemma-affine-local}
$S$ をスキームとする。
$f : X \to Y$ を $S$ 上の代数空間の射とする。
次の条件は同値である。
\begin{enumerate}
\item $f$ は表現可能かつアフィンである。
\item $f$ はアフィンである。
\item 任意のアフィンスキーム $V$ とエタール射 $V \to Y$ に対して、
スキーム $X \times_Y V$ はアフィンである。
\item スキーム $V$ と全射エタール射 $V \to Y$ であって、
$V \times_Y X \to V$ がアフィンとなるものが存在する。
\item ザリスキ被覆 $Y = \bigcup Y_i$ であって、各射
$f^{-1}(Y_i) \to Y_i$ がアフィンとなるものが存在する。
\end{enumerate}
\end{lemma}

\begin{proof}
(1) が (2) を含意し、(2) が (3) を含意することは明らかである。
$V$ を $Y$ 上エタールなアフィンの非交和として取れば、(3) は (4) を含意する。
Properties of Spaces, Lemma
\ref{spaces-properties-lemma-cover-by-union-affines} を参照せよ。
$V \to Y$ が (4) のような射であると仮定する。このとき $V$ の任意の
アフィン開集合 $W$ に対して、$W \times_Y X$ は $V \times_Y X$ の
アフィン開集合である。したがって Properties of Spaces, Lemma
\ref{spaces-properties-lemma-subscheme} により、$V \times_Y X$ はスキームである。
さらに射 $V \times_Y X \to V$ はアフィンである。よって、アフィン射の類は
必要な性質をすべて満たすので、Spaces, Lemma
\ref{spaces-lemma-morphism-sheaves-with-P-effective-descent-etale}
を適用できる（Morphisms, Lemmas \ref{morphisms-lemma-base-change-affine}
および Descent, Lemmas \ref{descent-lemma-descending-property-affine}
と \ref{descent-lemma-affine} を参照せよ）。この補題を適用した結論は、
$f$ が表現可能かつアフィン、すなわち (1) が成り立つということである。

\medskip\noindent
(1) と (5) の同値性は、アフィンであることが標的上ザリスキ局所的であることから従う
（上の参照は、実際にはアフィン性が標的上 fpqc 局所的であることを示す）。
\end{proof}

\begin{lemma}
\label{spaces-morphisms-lemma-composition-affine}
アフィン射の合成はアフィンである。
\end{lemma}

\begin{proof}
省略する。ヒント：ファイバー積の推移性を用いよ。
\end{proof}

\begin{lemma}
\label{spaces-morphisms-lemma-base-change-affine}
アフィン射の基底変換はアフィンである。
\end{lemma}

\begin{proof}
省略する。ヒント：ファイバー積の推移性を用いよ。
\end{proof}

\begin{lemma}
\label{spaces-morphisms-lemma-closed-immersion-affine}
閉埋入はアフィンである。
\end{lemma}

\begin{proof}
スキームの射に対する対応する主張から直ちに従う。Morphisms, Lemma
\ref{morphisms-lemma-closed-immersion-affine} を参照せよ。
\end{proof}

\begin{lemma}
\label{spaces-morphisms-lemma-affine-equivalence-algebras}
$S$ をスキームとする。$X$ を $S$ 上の代数空間とする。
圏の反同値
$$
\begin{matrix}
\text{代数空間} \\
\text{上でアフィン：}X
\end{matrix}
\longleftrightarrow
\begin{matrix}
\text{準連接層} \\
\text{代数構造：}\mathcal{O}_X\text{-代数}
\end{matrix}
$$
があり、$f : Y \to X$ に層 $f_*\mathcal{O}_Y$ を対応させる。
さらに、この同値は任意の基底変換と両立する。
\end{lemma}

\begin{proof}
この補題は Morphisms, Lemma
\ref{morphisms-lemma-affine-equivalence-algebras} に対応するものである。
$\mathcal{A}$ を準連接 $\mathcal{O}_X$-代数層とする。
$\pi_*\mathcal{O}_Y \cong \mathcal{A}$ を満たす代数空間のアフィン射
$\pi : Y = \underline{\Spec}_X(\mathcal{A}) \to X$ を構成する。
そのため、スキーム $U$ と全射エタール射 $\varphi : U \to X$ を取る。
通常どおり $R = U \times_X U$ と書き、射影を $s, t : R \to U$ とする。
合成 $\psi = \varphi \circ s = \varphi \circ t$ を $\psi : R \to X$ と表す。
前述の補題により、スキームのアフィン射 $\pi_0 : V \to U$ と
$\pi_1 : W \to R$ であって、
$\pi_{0, *}\mathcal{O}_V \cong \varphi^*\mathcal{A}$ および
$\pi_{1, *}\mathcal{O}_W \cong \psi^*\mathcal{A}$ を満たすものが存在する。
この構成は基底変換と両立するので、図式
$$
\vcenter{
\xymatrix{
W \ar[r]_{s'} \ar[d] & V \ar[d] \\
R \ar[r]^s & U
}
}
\quad\text{および}\quad
\vcenter{
\xymatrix{
W \ar[r]_{t'} \ar[d] & V \ar[d] \\
R \ar[r]^t & U
}
}
$$
をカルテジアンにする射 $s', t' : W \to V$ が存在する。
したがって $s', t'$ はエタールである。以上から形式的に、
$(t', s') : W \to V \times_S V$ はモノ射である。
$W \to V \times_S V$ が同値関係であることの確認は省略する
（ヒント：$U \times_X U \times_X U = R \times_{s, U, t} R$ への
$\mathcal{A}$ の引き戻しを考えよ）。商層 $Y = V/W$ は代数空間である。
Spaces, Theorem \ref{spaces-theorem-presentation} を参照せよ。
Groupoids, Lemma \ref{groupoids-lemma-criterion-fibre-product} により、
$Y \times_X U \cong V$ である。したがって Lemma \ref{spaces-morphisms-lemma-affine-local}
により $Y \to X$ はアフィンである。最後に、同型
$$
(Y \times_X U \to U)_*\mathcal{O}_{Y \times_X U} =
\pi_{0, *}\mathcal{O}_V \cong \varphi^*\mathcal{A}
$$
は貼り合わせ同型と両立するので、Properties of Spaces, Proposition
\ref{spaces-properties-proposition-quasi-coherent} により
$(Y \to X)_*\mathcal{O}_Y \cong \mathcal{A}$ である。
この構成が基底変換と両立することの確認は省略する。
\end{proof}

\begin{definition}
\label{spaces-morphisms-definition-relative-spec}
$S$ をスキームとする。$X$ を $S$ 上の代数空間とする。
$\mathcal{A}$ を準連接 $\mathcal{O}_X$-代数層とする。
圏同値 Lemma \ref{spaces-morphisms-lemma-affine-equivalence-algebras} のもとで
$\mathcal{A}$ に対応するアフィン射 $\underline{\Spec}(\mathcal{A}) \to X$ を、
{\it $X$ 上の $\mathcal{A}$ の相対スペクトル}、または単に
{\it $X$ 上の $\mathcal{A}$ のスペクトル}という。
\end{definition}

\noindent
相対スペクトルの形成は任意の基底変換と可換する。

\begin{remark}
\label{spaces-morphisms-remark-factorization-quasi-compact-quasi-separated}
$S$ をスキームとする。$f : Y \to X$ を $S$ 上の代数空間の
準コンパクトかつ準分離な射とする。このとき $f$ は標準的な分解
$$
Y \longrightarrow \underline{\Spec}_X(f_*\mathcal{O}_Y) \longrightarrow X
$$
をもつ。Lemma \ref{spaces-morphisms-lemma-pushforward} により $f_*\mathcal{O}_Y$ は
準連接なので、これは意味をもつ。射
$Y \to \underline{\Spec}_X(f_*\mathcal{O}_Y)$ は、標準的な
$\mathcal{O}_Y$-代数写像 $f^*f_*\mathcal{O}_Y \to \mathcal{O}_Y$ から得られる。
この写像は $Y$ 上の標準射
$Y \to Y \times_X \underline{\Spec}_X(f_*\mathcal{O}_Y)$ に対応するので
（Lemma \ref{spaces-morphisms-lemma-affine-equivalence-algebras}）、上のような $f$ の分解を得る。
\end{remark}

\begin{lemma}
\label{spaces-morphisms-lemma-affine-equivalence-modules}
$S$ をスキームとする。
$f : Y \to X$ を $S$ 上の代数空間のアフィン射とする。
$\mathcal{A} = f_*\mathcal{O}_Y$ とおく。関手
$\mathcal{F} \mapsto f_*\mathcal{F}$ は圏同値
$$
\left\{
\begin{matrix}
\text{準連接加群の圏}\\
\mathcal{O}_Y\text{-加群}
\end{matrix}
\right\}
\longrightarrow
\left\{
\begin{matrix}
\text{準連接加群の圏}\\
\mathcal{A}\text{-加群}
\end{matrix}
\right\}
$$
を誘導する。さらに、$\mathcal{A}$-加群が $\mathcal{O}_X$-加群として
準連接であるための必要十分条件は、$\mathcal{A}$-加群として
準連接であることである。
\end{lemma}

\begin{proof}
省略する。
\end{proof}

\begin{lemma}
\label{spaces-morphisms-lemma-affine-permanence}
$S$ をスキームとする。$B$ を $S$ 上の代数空間とする。
$g : X \to Y$ を $B$ 上の代数空間の射とする。
\begin{enumerate}
\item $X$ が $B$ 上アフィンで、$\Delta : Y \to Y \times_B Y$ がアフィンならば、
$g$ はアフィンである。
\item $X$ が $B$ 上アフィンで、$Y$ が $B$ 上分離ならば、$g$ はアフィンである。
\item アフィンスキームから、$\mathbf{Z}$ 上アフィンな対角をもつ代数空間
（Properties of Spaces, Definition
\ref{spaces-properties-definition-separated} の意味）への射はアフィンである。
\item アフィンスキームから分離代数空間への射はアフィンである。
\end{enumerate}
\end{lemma}

\begin{proof}
(1) の証明。Lemma \ref{spaces-morphisms-lemma-base-change-affine} により、基底変換
$X \times_B Y \to Y$ はアフィンである。射
$(1, g) : X \to X \times_B Y$ は、射 $X \times_B Y \to Y \times_B Y$ による
$Y \to Y \times_B Y$ の基底変換である。したがって Lemma
\ref{spaces-morphisms-lemma-base-change-affine} によりこれはアフィンである。
アフィン射の合成はアフィンなので
（Lemma \ref{spaces-morphisms-lemma-composition-affine} 参照）、(1) が従う。
閉埋入はアフィンであり（Lemma \ref{spaces-morphisms-lemma-closed-immersion-affine} 参照）、
$Y/B$ が分離であるとは $\Delta$ が閉埋入であることを意味するので、
(2) は (1) から従う。(3) と (4) は (1) と (2) の特殊な場合である。
\end{proof}

\begin{lemma}
\label{spaces-morphisms-lemma-Artinian-affine}
$S$ をスキームとする。$X$ を $S$ 上の準分離代数空間とする。
$A$ をアルティン環とする。任意の射 $\Spec(A) \to X$ はアフィンである。
\end{lemma}

\begin{proof}
$U$ がアフィンであるようなエタール射 $U \to X$ を取る。
補題を示すには、Lemma \ref{spaces-morphisms-lemma-affine-local} により
$\Spec(A) \times_X U$ がアフィンであることを示さなければならない。
$X$ は準分離なので、スキーム $\Spec(A) \times_X U$ は準コンパクトである。
さらに射影 $\Spec(A) \times_X U \to \Spec(A)$ はエタールである。
したがってこの射のファイバーは有限離散であり、さらに
$\Spec(A)$ の位相は離散である。ゆえに $\Spec(A) \times_X U$ は
台位相空間が有限離散集合であるスキームである。
Schemes, Lemma \ref{schemes-lemma-scheme-finite-discrete-affine} により証明が完了する。
\end{proof}












\section{準アフィン射}
\label{spaces-morphisms-section-quasi-affine}

\noindent
代数空間の表現可能射が準アフィンであるとは何かは、すでに Section
\ref{spaces-morphisms-section-representable} で定義した。

\begin{lemma}
\label{spaces-morphisms-lemma-quasi-affine-representable}
$S$ をスキームとする。$f : X \to Y$ を $S$ 上の代数空間の表現可能射とする。
$f$ が Section \ref{spaces-morphisms-section-representable} の意味で準アフィンであるための
必要十分条件は、すべてのアフィンスキーム $Z$ と射 $Z \to Y$ に対して、
スキーム $X \times_Y Z$ が準アフィンであることである。
\end{lemma}

\begin{proof}
これはスキームの準アフィン射の定義から直ちに従う
（Morphisms, Definition \ref{morphisms-definition-quasi-affine}）。
\end{proof}

\noindent
これにより次の定義が可能になる。

\begin{definition}
\label{spaces-morphisms-definition-quasi-affine}
$S$ をスキームとする。
$f : X \to Y$ を $S$ 上の代数空間の射とする。
任意のアフィンスキーム $Z$ と射 $Z \to Y$ に対して、代数空間
$X \times_Y Z$ が準アフィンスキームで表現可能であるとき、$f$ は
{\it 準アフィン}であるという。
\end{definition}

\begin{lemma}
\label{spaces-morphisms-lemma-quasi-affine-local}
$S$ をスキームとする。
$f : X \to Y$ を $S$ 上の代数空間の射とする。
次の条件は同値である。
\begin{enumerate}
\item $f$ は表現可能かつ準アフィンである。
\item $f$ は準アフィンである。
\item スキーム $V$ と全射エタール射 $V \to Y$ であって、
$V \times_Y X \to V$ が準アフィンとなるものが存在する。
\item ザリスキ被覆 $Y = \bigcup Y_i$ であって、各射
$f^{-1}(Y_i) \to Y_i$ が準アフィンとなるものが存在する。
\end{enumerate}
\end{lemma}

\begin{proof}
(1) が (2) を含意することは明らかである。また、$V$ を $Y$ 上エタールな
アフィンの非交和として取れば、(2) は (3) を含意する。Properties of Spaces,
Lemma \ref{spaces-properties-lemma-cover-by-union-affines} を参照せよ。
$V \to Y$ が (3) のような射であると仮定する。このとき $V$ の任意の
アフィン開集合 $W$ に対して、$W \times_Y X$ は $V \times_Y X$ の
準アフィン開集合である。したがって Properties of Spaces, Lemma
\ref{spaces-properties-lemma-subscheme} により、$V \times_Y X$ はスキームである。
さらに射 $V \times_Y X \to V$ は準アフィンである。よって、準アフィン射の類は
必要な性質をすべて満たすので、Spaces, Lemma
\ref{spaces-lemma-morphism-sheaves-with-P-effective-descent-etale}
を適用できる（Morphisms, Lemmas
\ref{morphisms-lemma-base-change-quasi-affine} および Descent, Lemmas
\ref{descent-lemma-descending-property-quasi-affine} と
\ref{descent-lemma-quasi-affine} を参照せよ）。この補題を適用した結論は、
$f$ が表現可能かつ準アフィン、すなわち (1) が成り立つということである。

\medskip\noindent
(1) と (4) の同値性は、準アフィンであることが標的上ザリスキ局所的であることから従う
（上の参照は、実際には準アフィン性が標的上 fpqc 局所的であることを示す）。
\end{proof}

\begin{lemma}
\label{spaces-morphisms-lemma-composition-quasi-affine}
準アフィン射の合成は準アフィンである。
\end{lemma}

\begin{proof}
省略する。
\end{proof}

\begin{lemma}
\label{spaces-morphisms-lemma-base-change-quasi-affine}
準アフィン射の基底変換は準アフィンである。
\end{lemma}

\begin{proof}
省略する。
\end{proof}

\begin{lemma}
\label{spaces-morphisms-lemma-characterize-quasi-affine}
$S$ をスキームとする。代数空間の準コンパクトかつ準分離な射
$f : Y \to X$ が準アフィンであるための必要十分条件は、標準的な分解
$Y \to \underline{\Spec}_X(f_*\mathcal{O}_Y)$
（Remark \ref{spaces-morphisms-remark-factorization-quasi-compact-quasi-separated}）が
開埋入であることである。
\end{lemma}

\begin{proof}
$U$ がスキームであるような全射 $U \to X$ を取る。
$f$ が準アフィンかどうかは $U$ への基底変換後に調べられ
（Lemma \ref{spaces-morphisms-lemma-quasi-affine-local}）、$f_*\mathcal{O}_Y|_V$ は
$(Y \times_X U \to U)_*\mathcal{O}_{Y \times_X U}$ に等しく
（Properties of Spaces, Lemma
\ref{spaces-properties-lemma-pushforward-etale-base-change-modules}）、
相対スペクトルの形成は基底変換と可換するので
（Lemma \ref{spaces-morphisms-lemma-affine-equivalence-algebras}）、主張は $X$ が
スキームである場合に帰着する。$X$ がスキームで、$f$ が準アフィンであるか、
または $Y \to \underline{\Spec}_X(f_*\mathcal{O}_Y)$ が開埋入ならば、
$Y$ もスキームである。したがって Morphisms, Lemma
\ref{morphisms-lemma-characterize-quasi-affine} に帰着する。
\end{proof}







\section{始域・終域についてエタール局所的な射の型}
\label{spaces-morphisms-section-local-source-target}

\noindent
スキームの射の性質であって {\it 始域・終域についてエタール局所的}なもの
（Descent, Definition \ref{descent-definition-local-source-target} 参照）が
与えられたとする。下の補題の同値な条件のいずれかを課すことで、
対応する代数空間の射の性質を定義できる。

\begin{lemma}
\label{spaces-morphisms-lemma-local-source-target}
$\mathcal{P}$ を、始域・終域についてエタール局所的なスキームの射の性質とする。
$S$ をスキームとする。
$f : X \to Y$ を $S$ 上の代数空間の射とする。
可換図式
$$
\xymatrix{
U \ar[d]_a \ar[r]_h & V \ar[d]^b \\
X \ar[r]^f & Y
}
$$
であって、$U$ と $V$ はスキーム、垂直矢印はエタールであるものを考える。
次の条件は同値である。
\begin{enumerate}
\item 上のような任意の図式に対して、射 $h$ は性質 $\mathcal{P}$ をもつ。
\item 上のような図式で $a : U \to X$ が全射であり、射 $h$ が
性質 $\mathcal{P}$ をもつものが存在する。
\end{enumerate}
$X$ と $Y$ が表現可能ならば、これは $f$ がスキームの射として
性質 $\mathcal{P}$ をもつこととも同値である。
さらに $\mathcal{P}$ が任意の基底変換で保たれ、基底上 fppf 局所的でもあるならば、
表現可能射 $f$ については、Section \ref{spaces-morphisms-section-representable} の意味で
$f$ が性質 $\mathcal{P}$ をもつこととも同値である。
\end{lemma}

\begin{proof}
(1) と (2) の同値性を示そう。(1) $\Rightarrow$ (2) は直ちに従う
（Spaces, Lemma \ref{spaces-lemma-lift-morphism-presentations} を考慮せよ）。
補題のような二つの図式
$$
\xymatrix{
U \ar[d] \ar[r]_h & V \ar[d] \\
X \ar[r]^f & Y
}
\quad\quad
\xymatrix{
U' \ar[d] \ar[r]_{h'} & V' \ar[d] \\
X \ar[r]^f & Y
}
$$
が与えられたとする。$U \to X$ は全射で、$h$ は性質 $\mathcal{P}$ を
もつと仮定する。(2) が (1) を含意することを示すには、$h'$ が
$\mathcal{P}$ をもつことを示さなければならない。そのため図式
$$
\xymatrix{
U \ar[d]_h &
U \times_X U' \ar[l] \ar[d]^{(h, h')} \ar[r] &
U' \ar[d]^{h'} \\
V &
V \times_Y V' \ar[l] \ar[r] &
V'
}
$$
を考える。Descent, Lemma
\ref{descent-lemma-local-source-target-characterize} により、$h$ が
$\mathcal{P}$ をもつことは $(h, h')$ が $\mathcal{P}$ をもつことを含意する。
$U \times_X U' \to U'$ は全射なので、同じ補題によりこれは $h'$ が
$\mathcal{P}$ をもつことを含意する。

\medskip\noindent
$X$ と $Y$ が表現可能ならば、Descent, Lemma
\ref{descent-lemma-local-source-target-characterize} を適用でき、(1) と (2) が
$f$ が $\mathcal{P}$ をもつことと同値であることが分かる。

\medskip\noindent
最後に、$f$ が表現可能であり、$U, V, a, b, h$ が補題の (2) のようなもので、
$\mathcal{P}$ が任意の基底変換で保たれると仮定する。
任意のスキーム $Z$ と射 $Z \to X$ に対して、基底変換
$Z \times_Y X \to Z$ が性質 $\mathcal{P}$ をもつことを示さなければならない。
図式
$$
\xymatrix{
Z \times_Y U \ar[d] \ar[r] &
Z \times_Y V \ar[d] \\
Z \times_Y X \ar[r] &
Z
}
$$
を考える。上の水平矢印は $h$ の基底変換なので性質 $\mathcal{P}$ をもつ。
左の垂直矢印はエタールかつ全射であり、右の垂直矢印はエタールである。
したがって Descent, Lemma
\ref{descent-lemma-local-source-target-characterize} を再び適用すると、
$Z \times_Y X \to Z$ が性質 $\mathcal{P}$ をもつことが分かる。
\end{proof}

\begin{definition}
\label{spaces-morphisms-definition-P}
$S$ をスキームとする。$\mathcal{P}$ を、始域・終域について
エタール局所的なスキームの射の性質とする。
$S$ 上の代数空間の射 $f : X \to Y$ が Lemma
\ref{spaces-morphisms-lemma-local-source-target} の同値な条件を満たすとき、この射は
{\it 性質 $\mathcal{P}$ をもつ}という。
\end{definition}

\noindent
明らかな注意を二つ述べる。

\begin{remark}
\label{spaces-morphisms-remark-composition-P}
$S$ をスキームとする。$\mathcal{P}$ を、始域・終域について
エタール局所的なスキームの射の性質とする。さらに $\mathcal{P}$ が
合成で安定であると仮定する。このとき性質 $\mathcal{P}$ をもつ
代数空間の射の類は合成で安定である。
\end{remark}

\begin{remark}
\label{spaces-morphisms-remark-base-change-P}
$S$ をスキームとする。$\mathcal{P}$ を、始域・終域について
エタール局所的なスキームの射の性質とする。さらに $\mathcal{P}$ が
基底変換で安定であると仮定する。このとき性質 $\mathcal{P}$ をもつ
代数空間の射の類は基底変換で安定である。
\end{remark}

\noindent
スキームの芽の射の性質であって {\it 始域・終域についてエタール局所的}なもの
（Descent, Definition
\ref{descent-definition-local-source-target-at-point} 参照）が与えられたとする。
下の補題の同値な条件のいずれかを課すことで、点における対応する
代数空間の射の性質を定義できる。

\begin{lemma}
\label{spaces-morphisms-lemma-local-source-target-at-point}
$\mathcal{Q}$ を、始域・終域についてエタール局所的な芽の射の性質とする。
$S$ をスキームとする。
$f : X \to Y$ を $S$ 上の代数空間の射とする。
$x \in |X|$ を $X$ の点とする。
図式
$$
\xymatrix{
U \ar[d]_a \ar[r]_h & V \ar[d]^b \\
X \ar[r]^f & Y
}
\quad\quad
\xymatrix{
u \ar[d] \ar[r] & v \ar[d] \\
x \ar[r] & y
}
$$
であって、$U$ と $V$ はスキーム、$a, b$ はエタールで、$u, v, x, y$ は
対応する空間の点であるものを考える。次の条件は同値である。
\begin{enumerate}
\item 上のような任意の図式に対して
$\mathcal{Q}((U, u) \to (V, v))$ が成り立つ。
\item 上のようなある図式に対して
$\mathcal{Q}((U, u) \to (V, v))$ が成り立つ。
\end{enumerate}
$X$ と $Y$ が表現可能ならば、これは
$\mathcal{Q}((X, x) \to (Y, y))$ とも同値である。
\end{lemma}

\begin{proof}
省略する。ヒント：Lemma \ref{spaces-morphisms-lemma-local-source-target} の証明と非常によく似ている。
\end{proof}

\begin{definition}
\label{spaces-morphisms-definition-P-at-point}
$\mathcal{Q}$ を、始域・終域についてエタール局所的なスキームの芽の射の性質とする。
$S$ をスキームとする。$S$ 上の代数空間の射 $f : X \to Y$ と
点 $x \in |X|$ が与えられたとする。Lemma
\ref{spaces-morphisms-lemma-local-source-target-at-point} の同値な条件を満たすとき、$f$ は
{\it $x$ において性質 $\mathcal{Q}$ をもつ}という。
\end{definition}

\noindent
次の補題を、射の性質から点における射の性質へ無条件に移るために
用いてはならない。例えば $\mathcal{P}$ が平坦であるという性質ならば、
次の補題の性質 $Q$ は「$f$ が $x$ のある開近傍で平坦である」ことを意味し、
「$f$ が $x$ で平坦である」ことと同じではない。

\begin{lemma}
\label{spaces-morphisms-lemma-local-source-target-global-implies-local}
$\mathcal{P}$ を、始域・終域についてエタール局所的なスキームの射の性質とする。
Descent, Lemma
\ref{descent-lemma-local-source-target-global-implies-local} で
$\mathcal{P}$ に対応づけられた芽の射の性質 $\mathcal{Q}$ を考える。このとき
\begin{enumerate}
\item $\mathcal{Q}$ は始域・終域についてエタール局所的である。
\item 代数空間の射 $f : X \to Y$ と $x \in |X|$ が与えられると、
次の条件は同値である。
\begin{enumerate}
\item $f$ は $x$ において $\mathcal{Q}$ をもつ。
\item $x$ の開近傍 $X' \subset X$ であって、$X' \to Y$ が
$\mathcal{P}$ をもつものが存在する。
\end{enumerate}
\item 代数空間の射 $f : X \to Y$ が与えられると、次の条件は同値である。
\begin{enumerate}
\item $f$ は $\mathcal{P}$ をもつ。
\item すべての $x \in |X|$ に対して、射 $f$ は $x$ において
$\mathcal{Q}$ をもつ。
\end{enumerate}
\end{enumerate}
\end{lemma}

\begin{proof}
(1) については Descent, Lemma
\ref{descent-lemma-local-source-target-global-implies-local} を参照せよ。
Lemma \ref{spaces-morphisms-lemma-local-source-target-at-point} のような図式が与えられたとき、
$X' = a(U) \subset X$ とすれば、(1)(a) $\Rightarrow$ (2)(b) が従う。
(2)(b) $\Rightarrow$ (1)(a) は明らかである。
(3)(a) と (3)(b) の同値性は、スキームの射に対する対応する結果から従う。
Descent, Lemma
\ref{descent-lemma-local-source-target-local-implies-global} を参照せよ。
\end{proof}

\begin{remark}
\label{spaces-morphisms-remark-when-apply}
上の Lemma \ref{spaces-morphisms-lemma-local-source-target-global-implies-local} を、
Descent, Remark \ref{descent-remark-list-local-source-target} に列挙された
「平坦」を除くすべての場合に適用する。各場合に、$f$ が $x$ の近傍で
$\mathcal{P}$ をもつとき、かつそのときに限り $f$ が $x$ で
性質 $\mathcal{P}$ をもつと定義することで、これを行う。
\end{remark}




\section{有限型の射}
\label{spaces-morphisms-section-finite-type}

\noindent
スキームの射が「局所有限型」であるという性質は、始域・終域について
エタール局所的である。Descent, Remark
\ref{descent-remark-list-local-source-target} を参照せよ。
また、これは基底変換で安定で、標的上 fpqc 局所的でもある。
Morphisms, Lemma \ref{morphisms-lemma-base-change-finite-type} および
Descent, Lemmas
\ref{descent-lemma-descending-property-locally-finite-type} を参照せよ。
したがって上の Lemma \ref{spaces-morphisms-lemma-local-source-target} により、
代数空間の射が局所有限型であるとは何かを次のように定義できる。
射が表現可能ならば、これは Section \ref{spaces-morphisms-section-representable} で
すでに定義した概念と一致する。

\begin{definition}
\label{spaces-morphisms-definition-locally-finite-type}
$S$ をスキームとする。
$f : X \to Y$ を $S$ 上の代数空間の射とする。
\begin{enumerate}
\item Lemma \ref{spaces-morphisms-lemma-local-source-target} の同値な条件が
$\mathcal{P} = \text{locally of finite type}$ に対して成り立つとき、
$f$ は {\it 局所有限型}であるという。
\item $x \in |X|$ とする。$x$ の開近傍 $X' \subset X$ であって、
$f|_{X'} : X' \to Y$ が局所有限型となるものが存在するとき、$f$ は
{\it $x$ において有限型}であるという。
\item $f$ は、局所有限型かつ準コンパクトであるとき
{\it 有限型}であるという。
\end{enumerate}
\end{definition}

\noindent
Spaces, Example \ref{spaces-example-affine-line-translation} の代数空間
$\mathbf{A}^1_k/\mathbf{Z}$ を考える。射
$\mathbf{A}^1_k/\mathbf{Z} \to \Spec(k)$ は有限型である。

\begin{lemma}
\label{spaces-morphisms-lemma-composition-finite-type}
有限型射の合成は有限型である。局所有限型についても同様である。
\end{lemma}

\begin{proof}
Remark \ref{spaces-morphisms-remark-composition-P} および Morphisms, Lemma
\ref{morphisms-lemma-composition-finite-type} を参照せよ。
\end{proof}

\begin{lemma}
\label{spaces-morphisms-lemma-base-change-finite-type}
有限型射の基底変換は有限型である。局所有限型についても同様である。
\end{lemma}

\begin{proof}
Remark \ref{spaces-morphisms-remark-base-change-P} および Morphisms, Lemma
\ref{morphisms-lemma-base-change-finite-type} を参照せよ。
\end{proof}

\begin{lemma}
\label{spaces-morphisms-lemma-finite-type-local}
$S$ をスキームとする。
$f : X \to Y$ を $S$ 上の代数空間の射とする。
次の条件は同値である。
\begin{enumerate}
\item $f$ は局所有限型である。
\item すべての $x \in |X|$ に対して、射 $f$ は $x$ において有限型である。
\item 任意のスキーム $Z$ と任意の射 $Z \to Y$ に対して、射
$Z \times_Y X \to Z$ は局所有限型である。
\item 任意のアフィンスキーム $Z$ と任意の射 $Z \to Y$ に対して、射
$Z \times_Y X \to Z$ は局所有限型である。
\item スキーム $V$ と全射エタール射 $V \to Y$ であって、
$V \times_Y X \to V$ が局所有限型となるものが存在する。
\item スキーム $U$ と全射エタール射 $\varphi : U \to X$ であって、
合成 $f \circ \varphi$ が局所有限型となるものが存在する。
\item 任意の可換図式
$$
\xymatrix{
U \ar[d] \ar[r] & V \ar[d] \\
X \ar[r] & Y
}
$$
であって、$U$, $V$ はスキーム、垂直矢印はエタールであるものに対して、
上の水平矢印は局所有限型である。
\item 可換図式
$$
\xymatrix{
U \ar[d] \ar[r] & V \ar[d] \\
X \ar[r] & Y
}
$$
であって、$U$, $V$ はスキーム、垂直矢印はエタール、$U \to X$ は全射で、
上の水平矢印は局所有限型であるものが存在する。
\item ザリスキ被覆 $Y = \bigcup_{i \in I} Y_i$ および
$f^{-1}(Y_i) = \bigcup X_{ij}$ であって、各射 $X_{ij} \to Y_i$ が
局所有限型となるものが存在する。
\end{enumerate}
\end{lemma}

\begin{proof}
条件 (2)、(3)、(4)、(5)、(6)、(7)、(9) のそれぞれは、
直接的に条件 (8) を含意する。例えば (5) が成り立つならば、
アフィンの非交和であるスキーム $V$ と全射 $V \to Y$ を取れる
（Properties of Spaces, Lemma
\ref{spaces-properties-lemma-cover-by-union-affines} 参照）。
このとき (5) により $V \times_Y X \to V$ は局所有限型である。
スキーム $U$ と全射エタール射 $U \to V \times_Y X$ を取る。
すると Lemma \ref{spaces-morphisms-lemma-composition-finite-type} により $U \to V$ は
局所有限型である。したがって (8) が成り立つ。

\medskip\noindent
条件 (1)、(7)、(8) は定義により同値である。

\medskip\noindent
証明を完了するため、(1) が条件 (2)、(3)、(4)、(5)、(6)、(9) のすべてを
含意することを示す。(2) については直ちに従う。
(3)、(4)、(5)、(9) については、局所有限型であることが基底変換で保たれることから従う。
Lemma \ref{spaces-morphisms-lemma-base-change-finite-type} を参照せよ。
(6) については $U = X$ と取ればよく、証明が完了する。
\end{proof}

\begin{lemma}
\label{spaces-morphisms-lemma-locally-finite-type-locally-noetherian}
$S$ をスキームとする。
$f : X \to Y$ を $S$ 上の代数空間の射とする。
$f$ が局所有限型で、$Y$ が局所ネーターならば、$X$ は局所ネーターである。
\end{lemma}

\begin{proof}
$U$, $V$ がスキーム、垂直矢印が全射エタールである可換図式
$$
\xymatrix{
U \ar[d] \ar[r] & V \ar[d] \\
X \ar[r] & Y
}
$$
を取る。$f$ が局所有限型ならば、$U \to V$ は局所有限型である。
$Y$ が局所ネーターならば、$V$ は局所ネーターである。
Morphisms, Lemma \ref{morphisms-lemma-finite-type-noetherian} により
$U$ は局所ネーターであり、これは $X$ が局所ネーターであることを意味する。
\end{proof}

\begin{lemma}
\label{spaces-morphisms-lemma-permanence-finite-type}
$S$ をスキームとする。
$f : X \to Y$, $g : Y \to Z$ を $S$ 上の代数空間の射とする。
$g \circ f : X \to Z$ が局所有限型ならば、$f : X \to Y$ は局所有限型である。
\end{lemma}

\begin{proof}
$U$, $V$, $W$ がスキーム、垂直矢印がエタールかつ全射である図式
$$
\xymatrix{
U \ar[r] \ar[d] & V \ar[r] \ar[d] & W \ar[d] \\
X \ar[r] & Y \ar[r] & Z
}
$$
を取ることができる。Spaces, Lemma
\ref{spaces-lemma-lift-morphism-presentations} を参照せよ。
ここで Lemma \ref{spaces-morphisms-lemma-finite-type-local} および Morphisms, Lemma
\ref{morphisms-lemma-permanence-finite-type} を用いれば結論を得る。
\end{proof}

\begin{lemma}
\label{spaces-morphisms-lemma-immersion-locally-finite-type}
埋入は局所有限型である。
\end{lemma}

\begin{proof}
一般原理 Spaces, Lemma
\ref{spaces-lemma-representable-transformations-property-implication} および
Morphisms, Lemmas
\ref{morphisms-lemma-immersion-locally-finite-type} から従う。
\end{proof}







\section{点と幾何点}
\label{spaces-morphisms-section-points-fields}

\noindent
この節では、点と幾何点についていくつか注意する
（Properties of Spaces, Definition
\ref{spaces-properties-definition-geometric-point} 参照）。
$X$ の幾何点を考える一つの方法は、$S$ の幾何点
$\overline{s} : \Spec(k) \to S$ と、$\overline{s}$ の $X$ への射
$\overline{x}$ への持ち上げを考えることである。図式は
$$
\xymatrix{
\Spec(k) \ar[r]_-{\overline{x}} \ar[rd]_{\overline{s}} & X \ar[d] \\
& S.
}
$$
である。「$k$ を $S$ 上の代数閉体とする」という言い方を、
$\Spec(k)$ に射 $\Spec(k) \to S$ が備わっていることを表すためによく用いる。
この状況では、$X$ の $k$-値点の集合を
$$
X(k) = \Mor_S(\Spec(k), X) =
\{\overline{x} \in X\text{ は次の点の上にある：}\overline{s}\}
$$
と書く。この場合、写像 $X(k) \to |X|$ の像は部分集合
$|X_s| \subset |X|$ に含まれる。ここで
$X_s = \Spec(\kappa(s)) \times_S X$ であり、$s \in S$ は
$\overline{s}$ に対応する点である。$\Spec(\kappa(s)) \to S$ はモノ射なので、
基底変換 $X_s \to X$ もモノ射であり、実際 $|X_s|$ は $|X|$ の部分集合である。

\begin{lemma}
\label{spaces-morphisms-lemma-locally-finite-type-surjective-geometric-points}
$S$ をスキームとする。$f : X \to Y$ を $S$ 上の代数空間の射とする。
$f$ が局所有限型であると仮定する。次の条件は同値である。
\begin{enumerate}
\item $f$ は全射である。
\item $S$ 上の任意の代数閉体 $k$ に対して、誘導写像
$X(k) \to Y(k)$ は全射である。
\end{enumerate}
\end{lemma}

\begin{proof}
$U$, $V$ は $S$ 上のスキーム、垂直矢印は全射かつエタールである図式
$$
\xymatrix{
U \ar[d] \ar[r] & V \ar[d] \\
X \ar[r] & Y
}
$$
を取る。Spaces, Lemma
\ref{spaces-lemma-lift-morphism-presentations} を参照せよ。
$f$ は局所有限型なので、$U \to V$ は局所有限型である。

\medskip\noindent
(1) を仮定し、$\overline{y} \in Y(k)$ とする。このとき Lemmas
\ref{spaces-morphisms-lemma-composition-surjective} および
\ref{spaces-morphisms-lemma-composition-finite-type} により、$U \to Y$ は全射かつ局所有限型である。
$Z = U \times_{Y, \overline{y}} \Spec(k)$ とおく。これはスキームである。
Lemmas \ref{spaces-morphisms-lemma-base-change-surjective} および
\ref{spaces-morphisms-lemma-base-change-finite-type} により、射影 $Z \to \Spec(k)$ は
全射かつ局所有限型である。Varieties, Lemma
\ref{varieties-lemma-locally-finite-type-Jacobson} により、$Z$ は
$k$-値点 $\overline{z}$ をもつ。$\overline{z}$ の像
$\overline{x} \in X(k)$ は所望のように $\overline{y}$ に写る。

\medskip\noindent
(2) を仮定する。Properties of Spaces, Lemma
\ref{spaces-properties-lemma-characterize-surjective} により、
$|X| \to |Y|$ が全射であることを示せば十分である。
$y \in |Y|$ とする。$y$ に写る $u \in U$ を取る。
$k \supset \kappa(u)$ を代数閉包とする。対応する点を
$\overline{u} \in U(k)$、その像を $\overline{y} \in Y(k)$ と表す。
仮定により、$\overline{y}$ に写る $\overline{x} \in X(k)$ が存在する。
このとき $\overline{x}$ の像 $x \in |X|$ が $y$ に写ることは明らかである。
\end{proof}

\noindent
次の補題を述べるため、以下の記法を導入する。スキーム $T$ に対し、
$$
\lambda(T) = \sup\{\aleph_0, |\kappa(t)| ; t \in T\}.
$$
と書く。言い換えると、$\lambda(T)$ は $T$ のすべての剰余体の濃度を
上から抑える最小の無限基数である。$R$ が環ならば、$R$ の任意の剰余体
$\kappa(\mathfrak p)$ の濃度は $R$ の濃度で抑えられることに注意する
（詳細は省略する）。これは $\lambda(T) \leq \text{size}(T)$ を含意する。
ここで $\text{size}(T)$ は Sets, Section
\ref{sets-section-categories-schemes} で導入したスキーム $T$ の大きさである。
$L/K$ が有限生成体拡大ならば、
$|K| \leq |L| \leq \max\{\aleph_0, |K|\}$ である。
したがって、$T' \to T$ が局所有限型なスキームの射ならば
$\lambda(T') \leq \lambda(T)$ であり、$T' \to T$ がさらに全射ならば等号が成り立つ。
次に $S$ をスキーム、$X$ を $S$ 上の代数空間とする。この場合、
$$
\lambda(X) := \lambda(U)
$$
と定義する。ここで $U$ は、$X$ への全射エタール射をもつ任意の
$S$ 上のスキームである。これが $U$ の選択に依存しない理由は次のとおりである。
そのような一対のスキーム $U$ と $U'$ が与えられると、ファイバー積
$U \times_X U'$ はスキームであり、$U$ と $U'$ の両方への全射エタール射をもつ。
したがって上の議論により
$\lambda(U) = \lambda(U \times_X U') = \lambda(U')$ である。

\begin{lemma}
\label{spaces-morphisms-lemma-large-enough}
$S$ をスキームとする。$X$, $Y$ を $S$ 上の代数空間とする。
\begin{enumerate}
\item $k$ が $S$ 上のすべての代数閉体を走るとき、幾何点
$\overline{y} \in Y(k)$ の集まりは $|Y|$ 全体を被覆する。
\item $k$ が、$|k| \geq \lambda(Y)$ かつ $|k| > \lambda(X)$ を満たす
$S$ 上のすべての代数閉体を走るとき、幾何点 $\overline{y} \in Y(k)$ は
$|Y|$ 全体を被覆する。
\item 幾何点 $\overline{s} : \Spec(k) \to S$ であって、$k$ の濃度が
$> \lambda(X)$ である任意のものに対して、写像
$$
X(k) \longrightarrow |X_s|
$$
は全射である。
\item $X \to Y$ を $S$ 上の代数空間の射とする。
幾何点 $\overline{s} : \Spec(k) \to S$ であって、$k$ の濃度が
$> \lambda(X)$ である任意のものに対して、写像
$$
X(k) \longrightarrow |X| \times_{|Y|} Y(k)
$$
は全射である。
\item $X \to Y$ を $S$ 上の代数空間の射とする。次の条件は同値である。
\begin{enumerate}
\item 写像 $X \to Y$ は全射である。
\item $|k| > \lambda(X)$ かつ $|k| \geq \lambda(Y)$ を満たす
$S$ 上のすべての代数閉体 $k$ に対して、写像 $X(k) \to Y(k)$ は全射である。
\end{enumerate}
\end{enumerate}
\end{lemma}

\begin{proof}
(1) を示すため、$V$ がスキームであるような全射エタール射 $V \to Y$ を取る。
各 $v \in V$ に対して代数閉包 $\kappa(v) \subset k_v$ を取る。
射 $\overline{x} : \Spec(k_v) \to V \to Y$ を考える。
$|Y|$ の構成により、これらは $|Y|$ を被覆する。

\medskip\noindent
(2) を示すため、証明を省略する次の二つの事実を用いる。
(\romannumeral1) $K$ が体で $\overline{K}$ が代数閉包ならば、
$|\overline{K}| \leq \max\{\aleph_0, |K|\}$ である。
(\romannumeral2) 任意の代数閉体 $k$ と任意の基数 $\aleph$、
$\aleph \geq |k|$ に対して、$|k'| = \aleph$ となる代数閉体の拡大
$k'/k$ が存在する。ここで
$\aleph = \max\{\lambda(X), \lambda(Y)\}^+$ とおく。
$\lambda^+ > \lambda$ は次に大きい基数を表す。Sets, Section
\ref{sets-section-cardinals} を参照せよ。
すると (\romannumeral1) により、証明の第一段落で構成した体 $k_u$ の濃度は
すべて $\lambda(X)$ で抑えられる。したがって (\romannumeral2) により、
$|k'_u| = \aleph$ となる拡大 $k_u \subset k'_u$ を取れる。
射 $\overline{x}' : \Spec(k'_u) \to X$ は所望のように $|X|$ を被覆する。
(2) の証明を厳密に完了するには、スキーム $\Spec(k'_u)$ が
$\Sch_{fppf}$ の対象と同型であることを示す必要がある。
これは、すべてのスキームを $\Sch_{fppf}$ の対象とするという規約による。
集合論的考察に関心のない読者はこの段落の残りを飛ばしてよい。
構成により、$\lambda(X)$ と $\lambda(Y)$ が $\text{size}(T)$ で抑えられるような
$\Sch_{fppf}$ の対象 $T$ が存在する。Topologies, Definitions
\ref{topologies-definition-big-fppf-site} における圏 $\Sch_{fppf}$ の構成、
すなわち Sets, Lemma \ref{sets-lemma-construct-category} で構成した
圏 $\Sch_\alpha$ としての構成により、大きさが
$\leq \text{size}(T)^+$ である任意のスキームは $\Sch_{fppf}$ の対象と同型である。
Sets, Equation (\ref{sets-equation-bound}) の函数 $Bound$ の式を参照せよ。
$\aleph \leq \text{size}(T)^+$ なので結論を得る。

\medskip\noindent
(3) における記法 $X_s$ はファイバー積
$\Spec(\kappa(s)) \times_S X$ を意味する。ここで $s \in S$ は
$\overline{s}$ に対応する点である。したがって (2) は
$Y = \Spec(\kappa(s))$ とした (4) から従う。

\medskip\noindent
(4) を示そう。$X \to Y$ を $S$ 上の代数空間の射とする。
$k$ を、濃度が $> \lambda(X)$ である $S$ 上の代数閉体とする。
$\overline{y} \in Y(k)$ と $x \in |X|$ が $|Y|$ の同じ元 $y$ に写るとする。
$x$ と $\overline{y}$ に写る $\overline{x} \in X(k)$ を見いださなければならない。
$U$, $V$ は $S$ 上のスキーム、垂直矢印は全射かつエタールである可換図式
$$
\xymatrix{
U \ar[d] \ar[r] & V \ar[d] \\
X \ar[r] & Y
}
$$
を取る。Spaces, Lemma
\ref{spaces-lemma-lift-morphism-presentations} を参照せよ。
$x$ に写る $u \in |U|$ を取り、その像を $v \in |V|$ と表す。
$u = \Spec(\kappa(u))$ と $v = \Spec(\kappa(v))$ をスキームとみなす。
$V \times_Y \Spec(k)$ は $k$ 上エタールなスキームであることに注意する。
したがって Morphisms, Lemma \ref{morphisms-lemma-etale-over-field} により、
これは $k$ の有限分離拡大のスペクトルの非交和である。
$v$ は $y$ に写るので、$v \times_Y \Spec(k)$ は空でないスキームである。
$v \to V$ はモノ射なので、
$v \times_Y \Spec(k) \to V \times_Y \Spec(k)$ はモノ射である。
したがって Schemes, Lemma
\ref{schemes-lemma-mono-towards-spec-field} により、
$v \times_Y \Spec(k)$ は $k$ の有限分離拡大のスペクトルの非交和である。
ゆえに射 $v \times_Y \Spec(k) \to \Spec(k)$ は切断をもつ。
すなわち、$v$ と $\overline{y}$ の上にある射
$\overline{v} : \Spec(k) \to V$ を取れる。最後にスキーム
$$
u \times_{V, \overline{v}} \Spec(k)
=
\Spec(\kappa(u) \otimes_{\kappa(v)} k)
$$
を考える。ここで $\kappa(v) \to k$ は射 $\overline{v}$ を定める体準同型である。
仮定により $k$ の濃度は $\kappa(u)$ の濃度より大きいので、
Algebra, Lemma \ref{algebra-lemma-base-change-Jacobson} を適用すると、
任意の極大イデアル $\mathfrak m \subset \kappa(u) \otimes_{\kappa(v)} k$ の
剰余体は $k$ 上代数的であり、したがって $k$ に等しいことが分かる。
よってそのような極大イデアルは、$u$ の上にあり $\overline{v}$ に写る射
$\overline{u} : \Spec(k) \to U$ を与える。
合成 $\Spec(k) \to U \to X$ が所望の幾何点
$\overline{x} \in X(k)$ である。これで (4) の証明が完了する。

\medskip\noindent
(5) は (2)、(4) および Properties of Spaces, Lemma
\ref{spaces-properties-lemma-characterize-surjective} から形式的に従う。
\end{proof}








\section{有限型点}
\label{spaces-morphisms-section-points-finite-type}

\noindent
$S$ をスキームとする。$X$ を $S$ 上の代数空間とする。
有限型点 $x \in |X|$ とは、局所有限型な射 $\Spec(k) \to X$ で
表現できる点のことである。有限型点は、代数空間および代数スタックにおける
閉点の適切な代替である。例えば、常に「十分多く」存在する。

\begin{lemma}
\label{spaces-morphisms-lemma-point-finite-type}
$S$ をスキームとする。$X$ を $S$ 上の代数空間とする。
$x \in |X|$ とする。次の条件は同値である。
\begin{enumerate}
\item 局所有限型で $x$ を表す射 $\Spec(k) \to X$ が存在する。
\item スキーム $U$、閉点 $u \in U$、およびエタール射
$\varphi : U \to X$ であって、$\varphi(u) = x$ を満たすものが存在する。
\end{enumerate}
\end{lemma}

\begin{proof}
$u \in U$ と $U \to X$ が (2) のようなものとする。このとき
$\Spec(\kappa(u)) \to U$ は有限型であり、$U \to X$ は表現可能かつ局所有限型である
（一般原理 Spaces, Lemma
\ref{spaces-lemma-representable-transformations-property-implication} および
Morphisms, Lemmas
\ref{morphisms-lemma-etale-locally-finite-presentation} と
\ref{morphisms-lemma-finite-presentation-finite-type} による）。
したがって Lemma \ref{spaces-morphisms-lemma-composition-finite-type} により (1) が成り立つ。

\medskip\noindent
逆に、$\Spec(k) \to X$ が局所有限型で $x$ を表すと仮定する。
$U$ がスキームであるような全射エタール射 $U \to X$ を取る。
仮定により $U \times_X \Spec(k) \to U$ は局所有限型である。
$U \times_X \Spec(k)$ の有限型点 $v$ を取る
（少なくとも一つ存在する。Morphisms, Lemma
\ref{morphisms-lemma-identify-finite-type-points} を参照せよ）。
Morphisms, Lemma \ref{morphisms-lemma-finite-type-points-morphism} により、
$v$ の像 $u \in U$ は $U$ の有限型点である。したがって Morphisms, Lemma
\ref{morphisms-lemma-identify-finite-type-points} により、$U$ を縮小した後、
$u$ は $U$ の閉点であると仮定できる。すなわち (2) が成り立つ。
\end{proof}

\begin{definition}
\label{spaces-morphisms-definition-finite-type-point}
$S$ をスキームとする。$X$ を $S$ 上の代数空間とする。
Lemma \ref{spaces-morphisms-lemma-point-finite-type} の同値な条件を満たす点
$x \in |X|$ を {\it 有限型点}\footnote{「局所有限型点」と呼ぶ方が
おそらくより正確なので、これはわずかな言葉の濫用である。}という。
$X$ の有限型点の集合を $X_{\text{ft-pts}}$ と表す。
\end{definition}

\noindent
有限型点の集合は次のように記述できる。

\begin{lemma}
\label{spaces-morphisms-lemma-identify-finite-type-points}
$S$ をスキームとする。$X$ を $S$ 上の代数空間とする。このとき
$$
X_{\text{ft-pts}} =
\bigcup\nolimits_{\varphi : U \to X\text{ エタール }} |\varphi|(U_0)
$$
である。ここで $U_0$ は $U$ の閉点の集合である。
$U$ は $X$ 上エタールなすべてのスキームを走らせても、$X$ 上エタールな
すべてのアフィンスキームを走らせてもよい。
\end{lemma}

\begin{proof}
Lemma \ref{spaces-morphisms-lemma-point-finite-type} から直ちに従う。
\end{proof}

\begin{lemma}
\label{spaces-morphisms-lemma-finite-type-points-morphism}
$S$ をスキームとする。
$f : X \to Y$ を $S$ 上の代数空間の射とする。
$f$ が局所有限型ならば、$f(X_{\text{ft-pts}}) \subset Y_{\text{ft-pts}}$ である。
\end{lemma}

\begin{proof}
$x \in X_{\text{ft-pts}}$ とする。$x$ を局所有限型な射
$x : \Spec(k) \to X$ で表す。Lemma \ref{spaces-morphisms-lemma-composition-finite-type} により
$f \circ x$ は局所有限型である。したがって $f(x) \in Y_{\text{ft-pts}}$ である。
\end{proof}

\begin{lemma}
\label{spaces-morphisms-lemma-finite-type-points-surjective-morphism}
$S$ をスキームとする。
$f : X \to Y$ を $S$ 上の代数空間の射とする。
$f$ が局所有限型かつ全射ならば、
$f(X_{\text{ft-pts}}) = Y_{\text{ft-pts}}$ である。
\end{lemma}

\begin{proof}
Lemma \ref{spaces-morphisms-lemma-finite-type-points-morphism} により
$f(X_{\text{ft-pts}}) \subset Y_{\text{ft-pts}}$ である。
$y \in |Y|$ を有限型点とする。$y$ を局所有限型な射
$\Spec(k) \to Y$ で表す。$f$ は全射なので、代数空間
$X_k = \Spec(k) \times_Y X$ は空でない。したがって Lemma
\ref{spaces-morphisms-lemma-identify-finite-type-points} により、有限型点
$x \in |X_k|$ をもつ。射 $X_k \to X$ は $\Spec(k) \to Y$ の
基底変換として局所有限型である
（Lemma \ref{spaces-morphisms-lemma-base-change-finite-type}）。ゆえに Lemma
\ref{spaces-morphisms-lemma-finite-type-points-morphism} により $x$ の $X$ における像は
有限型点であり、構成により $y$ に写る。
\end{proof}

\begin{lemma}
\label{spaces-morphisms-lemma-enough-finite-type-points}
$S$ をスキームとする。$X$ を $S$ 上の代数空間とする。
任意の局所閉部分集合 $T \subset |X|$ に対して
$$
T \not = \emptyset
\Rightarrow
T \cap X_{\text{ft-pts}} \not = \emptyset.
$$
特に、任意の閉部分集合 $T \subset |X|$ に対して、
$T \cap X_{\text{ft-pts}}$ は $T$ で稠密である。
\end{lemma}

\begin{proof}
$i : Z \to X$ を $T$ 上の誘導被約部分空間構造とする。Remark
\ref{spaces-morphisms-remark-space-structure-locally-closed-subset} を参照せよ。
任意の埋入は局所有限型である。Lemma
\ref{spaces-morphisms-lemma-immersion-locally-finite-type} を参照せよ。
したがって Lemma \ref{spaces-morphisms-lemma-finite-type-points-morphism} により
$Z_{\text{ft-pts}} \subset X_{\text{ft-pts}} \cap T$ である。
最後に、$Z$ へのエタール射をもつ任意の空でないアフィンスキーム $U$ は
少なくとも一つの閉点をもつ。したがって Lemma
\ref{spaces-morphisms-lemma-identify-finite-type-points} により $Z$ は少なくとも一つの
有限型点をもつ。補題が従う。
\end{proof}

\noindent
代数空間上の有限型点について、別の、より技術的な特徴づけを述べる。

\begin{lemma}
\label{spaces-morphisms-lemma-point-finite-type-monomorphism}
$S$ をスキームとする。$X$ を $S$ 上の代数空間とする。
$x \in |X|$ とする。次の条件は同値である。
\begin{enumerate}
\item $x$ は有限型点である。
\item 台位相空間 $|Z|$ が一点集合である代数空間 $Z$ と、
局所有限型な射 $f : Z \to X$ であって、$\{x\} = |f|(|Z|)$ を満たすものが存在する。
\item 代数空間 $Z$ と射 $f : Z \to X$ であって、次の性質をもつものが存在する。
\begin{enumerate}
\item $k$ が体であるような全射エタール射 $z : \Spec(k) \to Z$ が存在する。
\item $f$ は局所有限型である。
\item $f$ はモノ射である。
\item $x = f(z)$ である。
\end{enumerate}
\end{enumerate}
\end{lemma}

\begin{proof}
$x$ が有限型点であると仮定する。アフィンスキーム $U$、閉点
$u \in U$、および $\varphi(u) = x$ を満たすエタール射
$\varphi : U \to X$ を取る。Lemma
\ref{spaces-morphisms-lemma-identify-finite-type-points} を参照せよ。
通常どおり $u = \Spec(\kappa(u))$ とおく。射影
$u \times_X u \to u$ は合成
$$
u \times_X u \to u \times_X U \to u \times_X X = u
$$
である。ここで第一の矢印は閉埋入（$u \to U$ の基底変換）で、
第二の矢印はエタール（エタール射 $U \to X$ の基底変換）である。
したがって $u \times_X U$ は $k$ の有限分離拡大のスペクトルの非交和である
（Morphisms, Lemma \ref{morphisms-lemma-etale-over-field} 参照）。
ゆえに閉部分スキーム $u \times_X u$ は $k$ の有限分離拡大の非交和である。
すなわち $u \times_X u \to u$ はエタールである。
Spaces, Theorem \ref{spaces-theorem-presentation} により、
$Z = u/u \times_X u$ は代数空間である。構成により図式
$$
\xymatrix{
u \ar[d] \ar[r] & U \ar[d] \\
Z \ar[r] & X
}
$$
は可換で、垂直矢印はエタールである。したがって Lemma
\ref{spaces-morphisms-lemma-finite-type-local} により、$Z \to X$ は局所有限型である。
構成により射 $Z \to X$ はモノ射であり、$z$ の像は $x$ である。
したがって (3) が成り立つ。

\medskip\noindent
(3) が (2) を含意することは明らかである。
(2) が成り立つならば、Lemma \ref{spaces-morphisms-lemma-finite-type-points-morphism} により
$x$ は $X$ の有限型点である（$Z_{\text{ft-pts}}$ が空でないこと、すなわち
$Z$ の一意な点が $Z$ の有限型点であることを見るには、Lemma
\ref{spaces-morphisms-lemma-enough-finite-type-points} を用いる）。
\end{proof}





\section{長田空間}
\label{spaces-morphisms-section-nagata}

\noindent
長田代数空間の定義については Properties of Spaces, Section
\ref{spaces-properties-section-types-properties} を参照せよ。

\begin{lemma}
\label{spaces-morphisms-lemma-finite-type-nagata}
$S$ をスキームとする。
$f : X \to Y$ を $S$ 上の代数空間の射とする。
$Y$ が長田で $f$ が局所有限型ならば、$X$ は長田である。
\end{lemma}

\begin{proof}
$V$ をスキームとし、$V \to Y$ を全射エタール射とする。
$U$ をスキームとし、$U \to X \times_Y V$ を全射エタール射とする。
$Y$ が長田ならば、$V$ は長田スキームである。
$X \to Y$ が局所有限型ならば、$U \to V$ は局所有限型である。
したがって Morphisms, Lemma
\ref{morphisms-lemma-finite-type-nagata} により、$V$ は長田スキームである。
すると定義により $X$ は長田である。
\end{proof}

\begin{lemma}
\label{spaces-morphisms-lemma-ubiquity-nagata}
次の型の代数空間は長田である。
\begin{enumerate}
\item 長田スキーム上局所有限型な任意の代数空間。
\item 体上局所有限型な任意の代数空間。
\item ネーター完備局所環上局所有限型な任意の代数空間。
\item $\mathbf{Z}$ 上局所有限型な任意の代数空間。
\item 標数零のデデキント環上局所有限型な任意の代数空間。
\item など。
\end{enumerate}
\end{lemma}

\begin{proof}
第一の性質は Lemma \ref{spaces-morphisms-lemma-finite-type-nagata} により成り立つ。
したがって残りも成り立つ。Morphisms, Lemma
\ref{morphisms-lemma-ubiquity-nagata} を参照せよ。
\end{proof}









\section{準有限射}
\label{spaces-morphisms-section-quasi-finite}

\noindent
スキームの射が「局所準有限」であるという性質は、始域・終域について
エタール局所的である。Descent, Remark
\ref{descent-remark-list-local-source-target} を参照せよ。
また、これは基底変換で安定で、標的上 fpqc 局所的でもある。
Morphisms, Lemma \ref{morphisms-lemma-base-change-quasi-finite} および
Descent, Lemma \ref{descent-lemma-descending-property-quasi-finite} を参照せよ。
したがって上の Lemma \ref{spaces-morphisms-lemma-local-source-target} により、
代数空間の射が局所準有限であるとは何かを次のように定義できる。
射が表現可能ならば、これは Section \ref{spaces-morphisms-section-representable} で
すでに定義した概念と一致する。

\begin{definition}
\label{spaces-morphisms-definition-locally-quasi-finite}
$S$ をスキームとする。
$f : X \to Y$ を $S$ 上の代数空間の射とする。
\begin{enumerate}
\item Lemma \ref{spaces-morphisms-lemma-local-source-target} の同値な条件が
$\mathcal{P} = \text{locally quasi-finite}$ に対して成り立つとき、
$f$ は {\it 局所準有限}であるという。
\item $x \in |X|$ とする。$x$ の開近傍 $X' \subset X$ であって、
$f|_{X'} : X' \to Y$ が局所準有限となるものが存在するとき、$f$ は
{\it $x$ において準有限}であるという。
\item 代数空間の射 $f : X \to Y$ は、局所準有限かつ準コンパクトであるとき
{\it 準有限}であるという。
\end{enumerate}
\end{definition}

\noindent
最後の部分は Morphisms, Lemma
\ref{morphisms-lemma-quasi-finite-locally-quasi-compact} により、
スキームの射の準有限性の概念と両立する。

\begin{lemma}
\label{spaces-morphisms-lemma-base-change-quasi-finite-locus}
$S$ をスキームとする。$f : X \to Y$ と $g : Y' \to Y$ を
$S$ 上の代数空間の射とする。$g$ による $f$ の基底変換を
$f' : X' \to Y'$ と表す。射影を $g' : X' \to X$ と表す。
$f$ は局所有限型であると仮定する。$W \subset |X|$、それぞれ
$W' \subset |X'|$ を、$f$、それぞれ $f'$ が準有限である点の集合とする。
\begin{enumerate}
\item $W \subset |X|$ と $W' \subset |X'|$ は開である。
\item $W' = (g')^{-1}(W)$ である。すなわち、$f$ が準有限である軌跡の形成は
基底変換と可換する。
\item 局所準有限射の基底変換は局所準有限である。
\item 準有限射の基底変換は準有限である。
\end{enumerate}
\end{lemma}

\begin{proof}
スキーム $V$ と全射エタール射 $V \to Y$ を取る。
スキーム $U$ と全射エタール射 $U \to V \times_Y X$ を取る。
スキーム $V'$ と全射エタール射 $V' \to Y' \times_Y V$ を取る。
$U' = V' \times_V U$ とおくと、$U' \to X'$ も全射エタール射である。
図式は
$$
\vcenter{
\xymatrix{
U' \ar[d] \ar[r] & U \ar[d] \\
V' \ar[r] & V
}
}
\quad\text{は次の図式の上にある}\quad
\vcenter{
\xymatrix{
X' \ar[d] \ar[r] & X \ar[d] \\
Y' \ar[r] & Y
}
}
$$
である。像が $x \in |X|$ である $u \in |U|$ を取る。
「局所準有限」であるという性質は始域・終域についてエタール局所的である。
Descent, Remark \ref{descent-remark-list-local-source-target} を参照せよ。
したがって Lemmas \ref{spaces-morphisms-lemma-local-source-target-at-point} と
\ref{spaces-morphisms-lemma-local-source-target-global-implies-local} を適用でき、
$f : X \to Y$ が $x$ において準有限であることと、$U \to V$ が
$u$ において準有限であることとは同値である。同様に
$f' : X' \to Y'$ と射 $U' \to V'$ についても同じことが成り立つ。
したがって (1)、(2)、(3) は Morphisms, Lemmas
\ref{morphisms-lemma-base-change-quasi-finite} と
\ref{morphisms-lemma-quasi-finite-points-open} に帰着する。
(4) は (3) と Lemma \ref{spaces-morphisms-lemma-base-change-quasi-compact} から従う。
\end{proof}

\begin{lemma}
\label{spaces-morphisms-lemma-composition-quasi-finite}
準有限射の合成は準有限である。局所準有限についても同様である。
\end{lemma}

\begin{proof}
Remark \ref{spaces-morphisms-remark-composition-P} および Morphisms, Lemma
\ref{morphisms-lemma-composition-quasi-finite} を参照せよ。
\end{proof}

\begin{lemma}
\label{spaces-morphisms-lemma-base-change-quasi-finite}
準有限射の基底変換は準有限である。局所準有限についても同様である。
\end{lemma}

\begin{proof}
Lemma \ref{spaces-morphisms-lemma-base-change-quasi-finite-locus} から直ちに従う。
\end{proof}

\noindent
次の補題は、局所準有限射を、局所有限型で「離散ファイバー」をもつ射として
特徴づける。しかし、これは $|X| \to |Y|$ に離散ファイバーを要求することと
同じではない。Examples, Section
\ref{examples-section-specializations-fibre-etale} の議論を参照せよ。

\begin{lemma}
\label{spaces-morphisms-lemma-locally-quasi-finite}
$S$ をスキームとする。$f : X \to Y$ を代数空間の射とする。
$f$ が局所有限型であると仮定する。次の条件は同値である。
\begin{enumerate}
\item $f$ は局所準有限である。
\item $k$ が体であるような任意の射 $\Spec(k) \to Y$ に対して、
空間 $|X_k|$ は離散である。ここで $X_k = \Spec(k) \times_Y X$ である。
\end{enumerate}
\end{lemma}

\begin{proof}
$f$ が局所準有限であると仮定する。(2) のような $\Spec(k) \to Y$ を取る。
$U$ がスキームであるような全射エタール射 $U \to X$ を取る。
すると Properties of Spaces, Lemma
\ref{spaces-properties-lemma-base-change-etale} により
$U_k = \Spec(k) \times_Y U \to X_k$ は代数空間のエタール射である。
Lemma \ref{spaces-morphisms-lemma-base-change-quasi-finite} により、
$X_k \to \Spec(k)$ は局所準有限である。定義により、これは
$U_k \to \Spec(k)$ が局所準有限であることを意味する。
したがって Morphisms, Lemma
\ref{morphisms-lemma-locally-quasi-finite-fibres} により $|U_k|$ は離散である。
$|U_k| \to |X_k|$ は全射かつ開なので、$|X_k|$ は離散である。

\medskip\noindent
逆に、(2) を仮定する。$V$ がスキームであるような全射エタール射
$V \to Y$ を取る。$U$ がスキームであるような全射エタール射
$U \to V \times_Y X$ を取る。$f$ は局所有限型なので、$U \to V$ は
局所有限型であることに注意する。図式は
$$
\xymatrix{
U \ar[r] \ar[rd] & X \times_Y V \ar[d] \ar[r] & V \ar[d] \\
& X \ar[r] & Y
}
$$
である。$f$ が局所準有限でないならば、$U \to V$ は局所準有限でない。
したがって Morphisms, Lemma
\ref{morphisms-lemma-quasi-finite-at-point-characterize} により、同じ点
$v \in V$ の上にある $u, u' \in U$ に対して特殊化
$u \leadsto u'$ が存在する。$u, u'$ は
$X_v = \Spec(\kappa(v)) \times_Y X$ に同じ像をもたないと主張する。
これは $|X_v|$ が離散であるという仮定に所望の矛盾を与える。
$d = \text{trdeg}_{\kappa(v)}(\kappa(u))$ および
$d' = \text{trdeg}_{\kappa(v)}(\kappa(u'))$ とおく。
Morphisms, Lemma
\ref{morphisms-lemma-dimension-fibre-specialization} により $d > d'$ である。
$U_v$（$v$ 上の $U \to V$ のファイバー）は $X \times_Y V$ 上の
$U$ と $X_v$ のファイバー積であることに注意する。したがって
$U_v \to X_v$ はエタールである（エタール射 $U \to X \times_Y V$ の基底変換）。
$u, u' \in U_v$ が $|X_v|$ の同じ元に写るならば、Properties of Spaces,
Lemma \ref{spaces-properties-lemma-points-cartesian} により、
$t(r) = u$ かつ $s(r) = u'$ を満たす点
$r \in R_v = U_v \times_{X_v} U_v$ が存在する。
$s, t : R_v \to U_v$ は $\kappa(v)$ 上のスキームのエタール射なので、
$\kappa(u) \subset \kappa(r) \supset \kappa(u')$ は $\kappa(v)$ 上の
体の有限分離拡大であることに注意する
（Morphisms, Lemma \ref{morphisms-lemma-etale-over-field} 参照）。
したがって超越次数は等しい。この矛盾で証明が完了する。
\end{proof}

\begin{lemma}
\label{spaces-morphisms-lemma-quasi-finite-local}
$S$ をスキームとする。
$f : X \to Y$ を $S$ 上の代数空間の射とする。
次の条件は同値である。
\begin{enumerate}
\item $f$ は局所準有限である。
\item すべての $x \in |X|$ に対して、射 $f$ は $x$ において準有限である。
\item 任意のスキーム $Z$ と任意の射 $Z \to Y$ に対して、射
$Z \times_Y X \to Z$ は局所準有限である。
\item 任意のアフィンスキーム $Z$ と任意の射 $Z \to Y$ に対して、射
$Z \times_Y X \to Z$ は局所準有限である。
\item スキーム $V$ と全射エタール射 $V \to Y$ であって、
$V \times_Y X \to V$ が局所準有限となるものが存在する。
\item スキーム $U$ と全射エタール射 $\varphi : U \to X$ であって、
合成 $f \circ \varphi$ が局所準有限となるものが存在する。
\item 任意の可換図式
$$
\xymatrix{
U \ar[d] \ar[r] & V \ar[d] \\
X \ar[r] & Y
}
$$
であって、$U$, $V$ はスキーム、垂直矢印はエタールであるものに対して、
上の水平矢印は局所準有限である。
\item 可換図式
$$
\xymatrix{
U \ar[d] \ar[r] & V \ar[d] \\
X \ar[r] & Y
}
$$
であって、$U$, $V$ はスキーム、垂直矢印はエタール、$U \to X$ は全射で、
上の水平矢印は局所準有限であるものが存在する。
\item ザリスキ被覆 $Y = \bigcup_{i \in I} Y_i$ および
$f^{-1}(Y_i) = \bigcup X_{ij}$ であって、各射 $X_{ij} \to Y_i$ が
局所準有限となるものが存在する。
\end{enumerate}
\end{lemma}

\begin{proof}
省略する。
\end{proof}

\begin{lemma}
\label{spaces-morphisms-lemma-immersion-quasi-finite}
埋入は局所準有限である。
\end{lemma}

\begin{proof}
省略する。
\end{proof}

\begin{lemma}
\label{spaces-morphisms-lemma-permanence-quasi-finite}
$S$ をスキームとする。$X \to Y \to Z$ を $S$ 上の代数空間の射とする。
$X \to Z$ が局所準有限ならば、$X \to Y$ は局所準有限である。
\end{lemma}

\begin{proof}
垂直矢印がエタールかつ全射である可換図式
$$
\xymatrix{
U \ar[d] \ar[r] & V \ar[d] \ar[r] & W \ar[d] \\
X \ar[r] & Y \ar[r] & Z
}
$$
を取る（Spaces, Lemma
\ref{spaces-lemma-lift-morphism-presentations} 参照）。
上の行に Morphisms, Lemma
\ref{morphisms-lemma-permanence-quasi-finite} を適用する。
\end{proof}

\begin{lemma}
\label{spaces-morphisms-lemma-quasi-finite-at-a-finite-number-of-points}
$S$ をスキームとする。$f : X \to Y$ を $S$ 上の代数空間の有限型射とする。
$y \in |Y|$ とする。$y$ の上にある $|X|$ の点であって、$f$ が
準有限であるものは高々有限個である。
\end{lemma}

\begin{proof}
$y$ が定める同値類に属する体 $k$ と射 $\Spec(k) \to Y$ を取る。
ファイバー $X_k = \Spec(k) \times_Y X$ は体上有限型の代数空間であり、
特に準コンパクトである。Properties of Spaces, Lemma
\ref{spaces-properties-lemma-points-cartesian} により、写像
$|X_k| \to |X|$ は、$y$ 上の $|X| \to |Y|$ のファイバー上へ全射である。
さらに Lemma \ref{spaces-morphisms-lemma-base-change-quasi-finite-locus} により、
$X_k \to \Spec(k)$ が準有限である点の集合は、$f$ が準有限である
$y$ 上の点の集合上へ写る。アフィンスキーム $U$ と全射エタール射
$U \to X_k$ を取る（Properties of Spaces, Lemma
\ref{spaces-properties-lemma-quasi-compact-affine-cover}）。
すると $U \to \Spec(k)$ は有限型射であり、この射が準有限である点は
高々有限個である。Morphisms, Lemma
\ref{morphisms-lemma-quasi-finite-at-a-finite-number-of-points} を参照せよ。
$X_k \to \Spec(k)$ が点 $x'$ で準有限であるための必要十分条件は、
$U \to \Spec(k)$ が準有限である $U$ の点の像であることなので、結論を得る。
\end{proof}

\begin{lemma}
\label{spaces-morphisms-lemma-monomorphism-loc-finite-type-loc-quasi-finite}
$S$ をスキームとする。
$f : X \to Y$ を $S$ 上の代数空間の射とする。
$f$ が局所有限型かつモノ射ならば、$f$ は分離かつ局所準有限である。
\end{lemma}

\begin{proof}
モノ射は分離である。Lemma \ref{spaces-morphisms-lemma-monomorphism-separated} を参照せよ。
Lemma \ref{spaces-morphisms-lemma-quasi-finite-local} により、$Z$ がアフィンであるような
$Z \to Y$ による基底変換を行った後で補題を示せば十分である。
したがって $Y$ はアフィンスキームであると仮定してよい。
アフィンスキーム $U$ とエタール射 $U \to X$ を取る。
$X \to Y$ は局所有限型なので、アフィンスキームの射 $U \to Y$ は有限型である。
$X \to Y$ はモノ射なので $U \times_X U = U \times_Y U$ である。
特に写像 $U \times_Y U \to U$ はエタールである。
$y \in Y$ とする。このとき $U_y$ は空であるか、または、$y$ の上にある
ある $u \in U$ に対して
$\Spec(\kappa(u)) \times_{\Spec(\kappa(y))} U_y$ は
$u$ 上の $U \times_Y U \to U$ のファイバーと同型である。
これは $U \to Y$ のファイバーが有限離散集合であることを含意する
（$U \times_Y U \to U$ はアフィンスキームのエタール射なので、
Morphisms, Lemma \ref{morphisms-lemma-etale-over-field} を参照せよ）。
したがって Morphisms, Lemma
\ref{morphisms-lemma-quasi-finite-at-point-characterize} により
$U \to Y$ は準有限である。$U$ がアフィンであるような任意のエタール射
$U \to X$ に対してこれが成り立つので、$X \to Y$ は局所準有限である。
\end{proof}













\section{有限表示の射}
\label{spaces-morphisms-section-finite-presentation}

\noindent
スキームの射が「局所有限表示」であるという性質は、始域・終域について
エタール局所的である。Descent, Remark
\ref{descent-remark-list-local-source-target} を参照せよ。
また、これは基底変換で安定で、標的上 fpqc 局所的でもある。
Morphisms, Lemma \ref{morphisms-lemma-base-change-finite-presentation} および
Descent, Lemma
\ref{descent-lemma-descending-property-locally-finite-presentation} を参照せよ。
したがって上の Lemma \ref{spaces-morphisms-lemma-local-source-target} により、
代数空間の射が局所有限表示であるとは何かを次のように定義できる。
射が表現可能ならば、これは Section \ref{spaces-morphisms-section-representable} で
すでに定義した概念と一致する。

\begin{definition}
\label{spaces-morphisms-definition-locally-finite-presentation}
$S$ をスキームとする。
$f : X \to Y$ を $S$ 上の代数空間の射とする。
\begin{enumerate}
\item Lemma \ref{spaces-morphisms-lemma-local-source-target} の同値な条件が
$\mathcal{P} =$「局所有限表示」に対して成り立つとき、$f$ は
{\it 局所有限表示}であるという。
\item $x \in |X|$ とする。$x$ の開近傍 $X' \subset X$ であって、
$f|_{X'} : X' \to Y$ が局所有限表示となるものが存在するとき、$f$ は
{\it $x$ において有限表示}であるという\footnote{「$x$ において局所有限表示」
という言い方は不自然に思われるが、現在の用語にも誤解を招く点がある。
「$x$ において有限表示」は、$f|_{X'}$ が有限表示となる開近傍
$X' \subset X$ が存在することを\textbf{意味しない}。}。
\item 代数空間の射 $f : X \to Y$ は、局所有限表示、準コンパクトかつ
準分離であるとき {\it 有限表示}であるという。
\end{enumerate}
\end{definition}

\noindent
有限表示射は、単に局所有限表示な準コンパクト射なのでは\textbf{ない}ことに注意する。

\begin{lemma}
\label{spaces-morphisms-lemma-composition-finite-presentation}
有限表示射の合成は有限表示である。局所有限表示についても同様である。
\end{lemma}

\begin{proof}
Remark \ref{spaces-morphisms-remark-composition-P} および Morphisms, Lemma
\ref{morphisms-lemma-composition-finite-presentation} を参照せよ。
また、準コンパクト射と準分離射に対する結果
（Lemmas \ref{spaces-morphisms-lemma-composition-quasi-compact} と
\ref{spaces-morphisms-lemma-composition-separated}）を用いる。
\end{proof}

\begin{lemma}
\label{spaces-morphisms-lemma-base-change-finite-presentation}
有限表示射の基底変換は有限表示である。局所有限表示についても同様である。
\end{lemma}

\begin{proof}
Remark \ref{spaces-morphisms-remark-base-change-P} および Morphisms, Lemma
\ref{morphisms-lemma-base-change-finite-presentation} を参照せよ。
また、準コンパクト射と準分離射に対する結果
（Lemmas \ref{spaces-morphisms-lemma-base-change-quasi-compact} と
\ref{spaces-morphisms-lemma-base-change-separated}）を用いる。
\end{proof}

\begin{lemma}
\label{spaces-morphisms-lemma-finite-presentation-local}
$S$ をスキームとする。
$f : X \to Y$ を $S$ 上の代数空間の射とする。
次の条件は同値である。
\begin{enumerate}
\item $f$ は局所有限表示である。
\item すべての $x \in |X|$ に対して、射 $f$ は $x$ において有限表示である。
\item 任意のスキーム $Z$ と任意の射 $Z \to Y$ に対して、射
$Z \times_Y X \to Z$ は局所有限表示である。
\item 任意のアフィンスキーム $Z$ と任意の射 $Z \to Y$ に対して、射
$Z \times_Y X \to Z$ は局所有限表示である。
\item スキーム $V$ と全射エタール射 $V \to Y$ であって、
$V \times_Y X \to V$ が局所有限表示となるものが存在する。
\item スキーム $U$ と全射エタール射 $\varphi : U \to X$ であって、
合成 $f \circ \varphi$ が局所有限表示となるものが存在する。
\item 任意の可換図式
$$
\xymatrix{
U \ar[d] \ar[r] & V \ar[d] \\
X \ar[r] & Y
}
$$
であって、$U$, $V$ はスキーム、垂直矢印はエタールであるものに対して、
上の水平矢印は局所有限表示である。
\item 可換図式
$$
\xymatrix{
U \ar[d] \ar[r] & V \ar[d] \\
X \ar[r] & Y
}
$$
であって、$U$, $V$ はスキーム、垂直矢印はエタール、$U \to X$ は全射で、
上の水平矢印は局所有限表示であるものが存在する。
\item ザリスキ被覆 $Y = \bigcup_{i \in I} Y_i$ および
$f^{-1}(Y_i) = \bigcup X_{ij}$ であって、各射 $X_{ij} \to Y_i$ が
局所有限表示となるものが存在する。
\end{enumerate}
\end{lemma}

\begin{proof}
省略する。
\end{proof}

\begin{lemma}
\label{spaces-morphisms-lemma-finite-presentation-finite-type}
局所有限表示な射は局所有限型である。有限表示射は有限型である。
\end{lemma}

\begin{proof}
$f : X \to Y$ を局所有限表示な代数空間の射とする。
これは Lemma \ref{spaces-morphisms-lemma-local-source-target} のような図式で、$h$ が
局所有限表示、垂直矢印 $a$ が全射であるものが存在することを意味する。
Morphisms, Lemma
\ref{morphisms-lemma-finite-presentation-finite-type} により、$h$ は局所有限型である。
したがって定義により $X \to Y$ は局所有限型である。
$f$ が有限表示ならば準コンパクトであり、したがって $f$ は有限型である。
\end{proof}

\begin{lemma}
\label{spaces-morphisms-lemma-finite-presentation-noetherian}
$S$ をスキームとする。
$f : X \to Y$ を $S$ 上の代数空間の射とする。
$f$ が有限表示で $Y$ がネーターならば、$X$ はネーターである。
\end{lemma}

\begin{proof}
$f$ が有限表示で $Y$ がネーターであると仮定する。Lemmas
\ref{spaces-morphisms-lemma-finite-presentation-finite-type} と
\ref{spaces-morphisms-lemma-locally-finite-type-locally-noetherian} により、$X$ は局所ネーターである。
$f$ は準コンパクトで $Y$ は準コンパクトなので、$X$ は準コンパクトである。
$f$ は有限表示なので準分離であり
（Definition \ref{spaces-morphisms-definition-locally-finite-presentation} 参照）、
$Y$ はネーターなので準分離である
（Properties of Spaces, Definition
\ref{spaces-properties-definition-noetherian} 参照）。
したがって Lemma \ref{spaces-morphisms-lemma-separated-over-separated} により $X$ は準分離である。
これで Properties of Spaces, Definition
\ref{spaces-properties-definition-noetherian} の三条件をすべて確認し、結論を得る。
\end{proof}

\begin{lemma}
\label{spaces-morphisms-lemma-noetherian-finite-type-finite-presentation}
$S$ をスキームとする。
$f : X \to Y$ を $S$ 上の代数空間の射とする。
\begin{enumerate}
\item $Y$ が局所ネーターで $f$ が局所有限型ならば、$f$ は局所有限表示である。
\item $Y$ が局所ネーターで $f$ が有限型かつ準分離ならば、$f$ は有限表示である。
\end{enumerate}
\end{lemma}

\begin{proof}
$f : X \to Y$ が局所有限型で、$Y$ が局所ネーターであると仮定する。
これは Lemma \ref{spaces-morphisms-lemma-local-source-target} のような図式で、$h$ が
局所有限型、垂直矢印 $a$ が全射であるものが存在することを意味する。
Morphisms, Lemma
\ref{morphisms-lemma-noetherian-finite-type-finite-presentation} により、
$h$ は局所有限表示である。したがって定義により $X \to Y$ は
局所有限表示である。これで (1) が示された。
$f$ が有限型かつ準分離ならば、準コンパクトかつ準分離でもあり、
(2) は直ちに従う。
\end{proof}

\begin{lemma}
\label{spaces-morphisms-lemma-finite-presentation-quasi-compact-quasi-separated}
$S$ をスキームとする。$Y$ を $S$ 上の準コンパクトかつ準分離な
代数空間とする。$X$ が $Y$ 上有限表示ならば、$X$ は準コンパクトかつ準分離である。
\end{lemma}

\begin{proof}
省略する。
\end{proof}

\begin{lemma}
\label{spaces-morphisms-lemma-finite-presentation-permanence}
$S$ をスキームとする。$f : X \to Y$ と $Y \to Z$ を
$S$ 上の代数空間の射とする。$X$ が $Z$ 上局所有限表示で、
$Y$ が $Z$ 上局所有限型ならば、$f$ は局所有限表示である。
\end{lemma}

\begin{proof}
スキーム $W$ と全射エタール射 $W \to Z$ を取る。
次にスキーム $V$ と全射エタール射 $V \to W \times_Z Y$ を取る。
最後にスキーム $U$ と全射エタール射 $U \to V \times_Y X$ を取る。
定義により、$U$ は $W$ 上局所有限表示で、$V$ は $W$ 上局所有限型である。
Morphisms, Lemma \ref{morphisms-lemma-finite-presentation-permanence} により、
射 $U \to V$ は局所有限表示である。したがって $f$ は局所有限表示である。
\end{proof}

\begin{lemma}
\label{spaces-morphisms-lemma-diagonal-morphism-finite-type}
$S$ をスキームとする。$f : X \to Y$ を、対角
$\Delta : X \to X \times_Y X$ をもつ $S$ 上の代数空間の射とする。
$f$ が局所有限型ならば、$\Delta$ は局所有限表示である。
$f$ が準分離かつ局所有限型ならば、$\Delta$ は有限表示である。
\end{lemma}

\begin{proof}
$\Delta$ は第二射影 $X \times_Y X \to X$ を介して $X$ 上の射であることに注意する。
$f$ が局所有限型であると仮定する。$X$ は $X$ 上有限表示であり、
$X \times_Y X$ は $X$ 上有限型である
（Lemma \ref{spaces-morphisms-lemma-base-change-finite-type} による）。
したがって第一の主張は Lemma
\ref{spaces-morphisms-lemma-finite-presentation-permanence} から従う。
第二の主張は、第一の主張、定義、および対角射が分離であること
（Lemma \ref{spaces-morphisms-lemma-properties-diagonal}）から従う。
\end{proof}

\begin{lemma}
\label{spaces-morphisms-lemma-open-immersion-locally-finite-presentation}
代数空間の開埋入は局所有限表示である。
\end{lemma}

\begin{proof}
開埋入は定義により表現可能なので、一般原理 Spaces, Lemma
\ref{spaces-lemma-representable-transformations-property-implication} および
Morphisms, Lemma
\ref{morphisms-lemma-open-immersion-locally-finite-presentation} を用いることができる。
\end{proof}

\begin{lemma}
\label{spaces-morphisms-lemma-closed-immersion-finite-presentation}
閉埋入 $i : Z \to X$ が有限表示であるための必要十分条件は、付随する
準連接イデアル層
$\mathcal{I} = \Ker(\mathcal{O}_X \to i_*\mathcal{O}_Z)$ が
$\mathcal{O}_X$-加群として有限型であることである。
\end{lemma}

\begin{proof}
$U$ をスキームとし、$U \to X$ を全射エタール射とする。
Lemma \ref{spaces-morphisms-lemma-finite-presentation-local} により、
$i' : Z \times_X U \to U$ が有限表示であることと $i$ が有限表示であることとは同値である。
Properties of Spaces, Section
\ref{spaces-properties-section-properties-modules} により、$\mathcal{I}$ が
有限型であることと
$\mathcal{I}|_U = \Ker(\mathcal{O}_U \to i'_*\mathcal{O}_{Z \times_X U})$ が
有限型であることとは同値である。したがってスキームの場合の結果から従う。
Morphisms, Lemma
\ref{morphisms-lemma-closed-immersion-finite-presentation} を参照せよ。
\end{proof}








\section{構成可能集合}
\label{spaces-morphisms-section-constructible}

\noindent
この節は Properties of Spaces, Section
\ref{spaces-properties-section-constructible} の続きである。

\begin{lemma}
\label{spaces-morphisms-lemma-inverse-image-constructible}
$S$ をスキームとする。
$f : X \to Y$ を $S$ 上の代数空間の射とする。
$E \subset |Y|$ を部分集合とする。
$E$ が $Y$ でエタール局所構成可能ならば、$f^{-1}(E)$ は
$X$ でエタール局所構成可能である。
\end{lemma}

\begin{proof}
スキーム $V$ と全射エタール射 $\varphi : V \to Y$ を取る。
スキーム $U$ と全射エタール射 $U \to V \times_Y X$ を取る。
すると $U \to X$ は全射エタールであり、$U$ における $f^{-1}(E)$ の
逆像は、$U \to V$ による $\varphi^{-1}(E)$ の逆像である。
したがって、$U \to V$ に対するスキームの場合
（Morphisms, Lemma \ref{morphisms-lemma-inverse-image-constructible}）と、
定義（Properties of Spaces, Definition
\ref{spaces-properties-definition-locally-constructible}）から補題が従う。
\end{proof}

\begin{theorem}[Chevalley の定理]
\label{spaces-morphisms-theorem-chevalley}
$S$ をスキームとする。
$f : X \to Y$ を $S$ 上の代数空間の射とする。
$f$ が準コンパクトかつ局所有限表示であると仮定する。
このとき $|X|$ の任意のエタール局所構成可能部分集合の像は、
$|Y|$ のエタール局所構成可能部分集合である。
\end{theorem}

\begin{proof}
$E \subset |X|$ をエタール局所構成可能とする。
$V$ がアフィンであるようなエタール射 $V \to Y$ を取る。
Properties of Spaces, Definition
\ref{spaces-properties-definition-locally-constructible} により、
$V$ における $f(E)$ の逆像が構成可能であることを示せば十分である。
$f$ は準コンパクトなので、$V \times_Y X$ は準コンパクトな代数空間である。
アフィンスキーム $U$ と全射エタール射 $U \to V \times_Y X$ を取る
（Properties of Spaces, Lemma
\ref{spaces-properties-lemma-quasi-compact-affine-cover}）。
Properties of Spaces, Lemma
\ref{spaces-properties-lemma-points-cartesian} により、$V$ における
$f(E)$ の逆像は、$U \to V$ による $U$ における $E$ の逆像の像である。
したがってスキームの場合から従う。Morphisms, Lemma
\ref{morphisms-lemma-chevalley} を参照せよ。
\end{proof}







\section{平坦射}
\label{spaces-morphisms-section-flat}

\noindent
スキームの射が「平坦」であるという性質は、始域・終域について
エタール局所的である。Descent, Remark
\ref{descent-remark-list-local-source-target} を参照せよ。
また、これは基底変換で安定で、標的上 fpqc 局所的でもある。
Morphisms, Lemma \ref{morphisms-lemma-base-change-flat} および
Descent, Lemma \ref{descent-lemma-descending-property-flat} を参照せよ。
したがって上の Lemma \ref{spaces-morphisms-lemma-local-source-target} により、
代数空間の平坦射の概念を次のように定義できる。
射が表現可能ならば、これは Section \ref{spaces-morphisms-section-representable} で
すでに定義した概念と一致する。

\begin{definition}
\label{spaces-morphisms-definition-flat}
$S$ をスキームとする。
$f : X \to Y$ を $S$ 上の代数空間の射とする。
\begin{enumerate}
\item Lemma \ref{spaces-morphisms-lemma-local-source-target} の同値な条件が
$\mathcal{P} =$「平坦」に対して成り立つとき、$f$ は {\it 平坦}であるという。
\item $x \in |X|$ とする。Lemma
\ref{spaces-morphisms-lemma-local-source-target-at-point} の同値な条件が
$\mathcal{Q} =$「誘導される局所環の写像が平坦」に対して成り立つとき、
$f$ は {\it $x$ において平坦}であるという。
\end{enumerate}
第二の部分は Descent, Lemma \ref{descent-lemma-flat-at-point} により意味をもつ。
\end{definition}

\noindent
簡単な健全性確認を行う。

\begin{lemma}
\label{spaces-morphisms-lemma-flat-is-flat-at-all-points}
$S$ をスキームとする。$f : X \to Y$ を $S$ 上の代数空間の射とする。
$f$ が平坦であるための必要十分条件は、$|X|$ のすべての点で $f$ が
平坦であることである。
\end{lemma}

\begin{proof}
$U$ と $V$ はスキーム、垂直矢印はエタール、$a$ は全射である可換図式
$$
\xymatrix{
U \ar[d]_a \ar[r]_h & V \ar[d]^b \\
X \ar[r]^f & Y
}
$$
を取る。定義により、$f$ が平坦であることと $h$ が平坦であることとは同値である
（Definition \ref{spaces-morphisms-definition-P}）。定義により、$f$ が
$x \in |X|$ において平坦であることと、$x$ に写るある（同値に任意の）
$u \in U$ において $h$ が平坦であることとは同値である
（Definition \ref{spaces-morphisms-definition-P-at-point}）。
したがって補題は、スキームの射が平坦であるための必要十分条件が
始域のすべての点で平坦であることから従う
（Morphisms, Definition \ref{morphisms-definition-flat}）。
\end{proof}

\begin{lemma}
\label{spaces-morphisms-lemma-composition-flat}
平坦射の合成は平坦である。
\end{lemma}

\begin{proof}
Remark \ref{spaces-morphisms-remark-composition-P} および Morphisms, Lemma
\ref{morphisms-lemma-composition-flat} を参照せよ。
\end{proof}

\begin{lemma}
\label{spaces-morphisms-lemma-base-change-flat}
平坦射の基底変換は平坦である。
\end{lemma}

\begin{proof}
Remark \ref{spaces-morphisms-remark-base-change-P} および Morphisms, Lemma
\ref{morphisms-lemma-base-change-flat} を参照せよ。
\end{proof}

\begin{lemma}
\label{spaces-morphisms-lemma-flat-local}
$S$ をスキームとする。
$f : X \to Y$ を $S$ 上の代数空間の射とする。
次の条件は同値である。
\begin{enumerate}
\item $f$ は平坦である。
\item すべての $x \in |X|$ に対して、射 $f$ は $x$ において平坦である。
\item 任意のスキーム $Z$ と任意の射 $Z \to Y$ に対して、射
$Z \times_Y X \to Z$ は平坦である。
\item 任意のアフィンスキーム $Z$ と任意の射 $Z \to Y$ に対して、射
$Z \times_Y X \to Z$ は平坦である。
\item スキーム $V$ と全射エタール射 $V \to Y$ であって、
$V \times_Y X \to V$ が平坦となるものが存在する。
\item スキーム $U$ と全射エタール射 $\varphi : U \to X$ であって、
合成 $f \circ \varphi$ が平坦となるものが存在する。
\item 任意の可換図式
$$
\xymatrix{
U \ar[d] \ar[r] & V \ar[d] \\
X \ar[r] & Y
}
$$
であって、$U$, $V$ はスキーム、垂直矢印はエタールであるものに対して、
上の水平矢印は平坦である。
\item 可換図式
$$
\xymatrix{
U \ar[d] \ar[r] & V \ar[d] \\
X \ar[r] & Y
}
$$
であって、$U$, $V$ はスキーム、垂直矢印はエタール、$U \to X$ は全射で、
上の水平矢印は平坦であるものが存在する。
\item ザリスキ被覆 $Y = \bigcup Y_i$ および
$f^{-1}(Y_i) = \bigcup X_{ij}$ であって、各射 $X_{ij} \to Y_i$ が
平坦となるものが存在する。
\end{enumerate}
\end{lemma}

\begin{proof}
省略する。
\end{proof}

\begin{lemma}
\label{spaces-morphisms-lemma-fppf-open}
局所有限表示な平坦射は普遍的開である。
\end{lemma}

\begin{proof}
$f : X \to Y$ を $S$ 上の代数空間の局所有限表示な平坦射とする。
$U$ と $V$ はスキーム、垂直矢印は全射かつエタールである図式
$$
\xymatrix{
U \ar[r]_\alpha \ar[d] & V \ar[d] \\
X \ar[r] & Y
}
$$
を取る。Spaces, Lemma
\ref{spaces-lemma-lift-morphism-presentations} を参照せよ。
Lemmas \ref{spaces-morphisms-lemma-flat-local} と
\ref{spaces-morphisms-lemma-finite-presentation-local} により、射 $\alpha$ は平坦かつ
局所有限表示である。したがって Morphisms, Lemma
\ref{morphisms-lemma-fppf-open} により、$\alpha$ は普遍的開である。
ゆえに Lemma \ref{spaces-morphisms-lemma-universally-open-local} により、
$X \to Y$ は普遍的開である。
\end{proof}

\begin{lemma}
\label{spaces-morphisms-lemma-fpqc-quotient-topology}
$S$ をスキームとする。
$f : X \to Y$ を $S$ 上の代数空間の平坦、準コンパクトかつ全射な射とする。
部分集合 $T \subset |Y|$ が開（それぞれ閉）であるための必要十分条件は、
$f^{-1}(|T|)$ が $|X|$ で開（それぞれ閉）であることである。
言い換えると、$f$ は submersive であり、実際 universally submersive である。
\end{lemma}

\begin{proof}
アフィンスキーム $V_i$ とエタール射 $V_i \to Y$ であって、
$V = \coprod V_i \to Y$ が全射となるものを取る。Properties of Spaces,
Lemma \ref{spaces-properties-lemma-cover-by-union-affines} を参照せよ。
各 $i$ に対して、代数空間 $V_i \times_Y X$ は準コンパクトである。
したがってアフィンスキーム $U_i$ と全射エタール射
$U_i \to V_i \times_Y X$ を取れる。Properties of Spaces, Lemma
\ref{spaces-properties-lemma-quasi-compact-affine-cover} を参照せよ。
すると合成 $U_i \to V_i \times_Y X \to V_i$ はアフィン間の全射平坦射である。
もちろん $U = \coprod U_i \to X$ は全射かつエタールであり、
$U \to V \times_Y X$ は全射である。さらに射 $U \to V$ は射
$U_i \to V_i$ の非交和である。したがって $U \to V$ は全射、準コンパクトかつ平坦である。
図式
$$
\xymatrix{
U \ar[r] \ar[d] & X \ar[d] \\
V \ar[r] & Y
}
$$
を考える。$|Y|$ の位相の定義により、集合 $T$ が閉（それぞれ開）であることと、
$g^{-1}(T) \subset |V|$ が閉（それぞれ開）であることとは同値である。
$f^{-1}(T)$ とその $|U|$ における逆像についても同じである。
$U \to V$ は準コンパクト、全射かつ平坦なので、Morphisms, Lemma
\ref{morphisms-lemma-fpqc-quotient-topology} により結論を得る。
\end{proof}

\begin{lemma}
\label{spaces-morphisms-lemma-flat-at-point-etale-local-rings}
$S$ をスキームとする。
$f : X \to Y$ を $S$ 上の代数空間の射とする。
$\overline{x}$ を点 $x \in |X|$ の上にある $X$ の幾何点とする。
$\overline{y} = f \circ \overline{x}$ とおく。次の条件は同値である。
\begin{enumerate}
\item $f$ は $x$ において平坦である。
\item エタール局所環の写像
$\mathcal{O}_{Y, \overline{y}} \to \mathcal{O}_{X, \overline{x}}$ は平坦である。
\end{enumerate}
\end{lemma}

\begin{proof}
$U$ と $V$ はスキーム、$a, b$ はエタールで、$u \in U$ が $x$ に写る
可換図式
$$
\xymatrix{
U \ar[d]_a \ar[r]_h & V \ar[d]^b \\
X \ar[r]^f & Y
}
$$
を取る。Properties of Spaces, Lemma
\ref{spaces-properties-lemma-geometric-lift-to-usual} により、
$u$ の上にある幾何点 $\overline{u} : \Spec(k) \to U$ であって、
$\overline{x} = a \circ \overline{u}$ となるものを取れる。
$\overline{v} = h \circ \overline{u}$ とおき、その像を $v \in V$ とする。
Properties of Spaces, Lemma
\ref{spaces-properties-lemma-describe-etale-local-ring} により
$$
\mathcal{O}_{X, \overline{x}} = \mathcal{O}_{U, u}^{sh}
\quad\text{かつ}\quad
\mathcal{O}_{Y, \overline{y}} = \mathcal{O}_{V, v}^{sh}
$$
である。平坦な水平矢印をもつ局所環の可換図式
$$
\xymatrix{
\mathcal{O}_{U, u} \ar[r] &
\mathcal{O}_{X, \overline{x}} \\
\mathcal{O}_{V, v} \ar[u] \ar[r] &
\mathcal{O}_{Y, \overline{y}} \ar[u]
}
$$
を得る。左の垂直矢印が平坦であることと右の垂直矢印が平坦であることが
同値であることを示さなければならない。Algebra, Lemma
\ref{algebra-lemma-flatness-descends-more-general} により、
$\mathcal{O}_{U, u}$ が $\mathcal{O}_{V, v}$ 上平坦であることと、
$\mathcal{O}_{X, \overline{x}}$ が $\mathcal{O}_{V, v}$ 上平坦であることとは同値である。
したがって More on Flatness, Lemma
\ref{flat-lemma-flat-up-down-henselization} から結論が従う。
\end{proof}

\begin{lemma}
\label{spaces-morphisms-lemma-flat-morphism-sites}
$S$ をスキームとする。
$f : X \to Y$ を $S$ 上の代数空間の射とする。
$f$ が平坦であるための必要十分条件は、$f$ に付随するサイトの射
$
(f_{small}, f^\sharp) :
(X_\etale, \mathcal{O}_X)
\to
(Y_\etale, \mathcal{O}_Y)
$
が平坦であることである。
\end{lemma}

\begin{proof}
$(f_{small}, f^\sharp)$ の平坦性は、$f_{small}^{-1}\mathcal{O}_Y$-加群としての
$\mathcal{O}_X$ の平坦性によって定義される。これは茎で確認できる。
Modules on Sites, Lemma \ref{sites-modules-lemma-check-flat-stalks} および
Properties of Spaces, Theorem
\ref{spaces-properties-theorem-exactness-stalks} を参照せよ。
しかし $f$ の平坦性を茎で確認できることはすでに見た。
Lemma \ref{spaces-morphisms-lemma-flat-at-point-etale-local-rings} を参照せよ。
\end{proof}

\begin{lemma}
\label{spaces-morphisms-lemma-flat-pullback-support}
$S$ をスキームとする。$f : Y \to X$ を $S$ 上の代数空間の射とする。
$\mathcal{F}$ を、スキーム論的台 $Z \subset X$ をもつ有限型準連接
$\mathcal{O}_X$-加群とする。$f$ が平坦ならば、$f^{-1}(Z)$ は
$f^*\mathcal{F}$ のスキーム論的台である。
\end{lemma}

\begin{proof}
Lemma \ref{spaces-morphisms-lemma-scheme-theoretic-support} のスキーム論的台の特徴づけと、
Lemma \ref{spaces-morphisms-lemma-flat-local} のエタール被覆による平坦射の特徴づけを用いて、
スキームの場合に帰着する。これは Morphisms, Lemma
\ref{morphisms-lemma-flat-pullback-support} である。
\end{proof}

\begin{lemma}
\label{spaces-morphisms-lemma-flat-morphism-scheme-theoretically-dense-open}
$S$ をスキームとする。
$f : X \to Y$ を $S$ 上の代数空間の平坦射とする。
$V \to Y$ を準コンパクトな開埋入とする。$V$ が $Y$ で
スキーム論的に稠密ならば、$f^{-1}V$ は $X$ でスキーム論的に稠密である。
\end{lemma}

\begin{proof}
Lemma \ref{spaces-morphisms-lemma-scheme-theoretically-dense} のスキーム論的に稠密な
開集合の特徴づけと、Lemma \ref{spaces-morphisms-lemma-flat-local} のエタール被覆による
平坦射の特徴づけを用いて、スキームの場合に帰着する。これは Morphisms, Lemma
\ref{morphisms-lemma-flat-morphism-scheme-theoretically-dense-open} である。
\end{proof}

\begin{lemma}
\label{spaces-morphisms-lemma-flat-base-change-scheme-theoretic-image}
$S$ をスキームとする。$f : X \to Y$ を $S$ 上の代数空間の平坦射とする。
$g : V \to Y$ を代数空間の準コンパクト射とする。
$Z \subset Y$ を $g$ のスキーム論的像とし、$Z' \subset X$ を基底変換
$V \times_Y X \to X$ のスキーム論的像とする。このとき $Z' = f^{-1}Z$ である。
\end{lemma}

\begin{proof}
$Y'$ がアフィンスキームの非交和となるような全射エタール射
$Y' \to Y$ を取る（Properties of Spaces, Lemma
\ref{spaces-properties-lemma-cover-by-union-affines}）。
$X'$ がアフィンスキームの非交和となるような全射エタール射
$X' \to X \times_Y Y'$ を取る。Lemma \ref{spaces-morphisms-lemma-flat-local} により、
射 $X' \to Y'$ は平坦である。$V' = V \times_Y Y'$ とおく。
Lemma \ref{spaces-morphisms-lemma-quasi-compact-scheme-theoretic-image} により、
$Y'$ における $Z$ の逆像は $V' \to Y'$ のスキーム論的像であり、
$X'$ における $Z'$ の逆像は
$V' \times_{Y'} X' \to X'$ のスキーム論的像である。
$X' \to X$ は全射エタールなので、射 $X' \to Y'$ と $V' \to Y'$ の場合に
結果を示せば十分である。したがって $X$ と $Y$ はアフィンスキームであると
仮定してよい。この場合 $V$ は準コンパクトな代数空間である。
アフィンスキーム $W$ と全射エタール射 $W \to V$ を取る
（Properties of Spaces, Lemma
\ref{spaces-properties-lemma-quasi-compact-affine-cover}）。
$V \to Y$ のスキーム論的像が $W \to Y$ のスキーム論的像と一致し、
同様に $V \times_Y X \to Y$ と $W \times_Y X \to X$ についても一致することは明らかである。
したがってスキームの場合に帰着する。これは Morphisms, Lemma
\ref{morphisms-lemma-flat-base-change-scheme-theoretic-image} である。
\end{proof}





\section{平坦加群}
\label{spaces-morphisms-section-flat-modules}

\noindent
この節では、加群が点において平坦であるとは何かを定義する。
そのため、代数空間の小エタールサイト $X_\etale$ 上の層の茎の概念を用いる。
Properties of Spaces, Definition
\ref{spaces-properties-definition-stalk} を参照せよ。

\begin{lemma}
\label{spaces-morphisms-lemma-flat-at-point}
$S$ をスキームとする。$f : X \to Y$ を $S$ 上の代数空間の射とする。
$\mathcal{F}$ を $X$ 上の準連接層とする。$x \in |X|$ とする。
次の条件は同値である。
\begin{enumerate}
\item ある可換図式
$$
\xymatrix{
U \ar[d]_a \ar[r]_h & V \ar[d]^b \\
X \ar[r]^f & Y
}
$$
であって、$U$ と $V$ はスキーム、$a, b$ はエタール、$u \in U$ は
$x$ に写り、加群 $a^*\mathcal{F}$ が $u$ で $V$ 上平坦となるものが存在する。
\item $x$ の上にある任意の幾何点 $\overline{x}$ と
$\overline{y} = f \circ \overline{x}$ に対して、茎
$\mathcal{F}_{\overline{x}}$ はエタール局所環
$\mathcal{O}_{Y, \overline{y}}$ 上平坦である。
\end{enumerate}
\end{lemma}

\begin{proof}
この証明を通じて、$x$ の上の幾何点
$\overline{x} : \Spec(k) \to X$ を固定し、その $Y$ における像を
$\overline{y} = f \circ \overline{x}$ と表す。
(1) のような図式が与えられると、Properties of Spaces, Lemma
\ref{spaces-properties-lemma-geometric-lift-to-usual} により、
$u$ の上にある幾何点 $\overline{u} : \Spec(k) \to U$ であって、
$\overline{x} = a \circ \overline{u}$ となるものを取れる。
$\overline{v} = h \circ \overline{u}$ とおき、その像を $v \in V$ とする。
Properties of Spaces, Lemma
\ref{spaces-properties-lemma-describe-etale-local-ring} により
$$
\mathcal{O}_{X, \overline{x}} = \mathcal{O}_{U, u}^{sh}
\quad\text{かつ}\quad
\mathcal{O}_{Y, \overline{y}} = \mathcal{O}_{V, v}^{sh}
$$
である。局所環の可換図式
$$
\xymatrix{
\mathcal{O}_{U, u} \ar[r] &
\mathcal{O}_{X, \overline{x}} \\
\mathcal{O}_{V, v} \ar[u] \ar[r] &
\mathcal{O}_{Y, \overline{y}} \ar[u]
}
$$
を得る。最後に Properties of Spaces, Lemma
\ref{spaces-properties-lemma-stalk-quasi-coherent} により
$$
\mathcal{F}_{\overline{x}} =
(a^*\mathcal{F})_u \otimes_{\mathcal{O}_{U, u}}
\mathcal{O}_{X, \overline{x}}
$$
である。したがって Algebra, Lemma
\ref{algebra-lemma-flatness-descends-more-general} により、
$(a^*\mathcal{F})_u$ が $\mathcal{O}_{V, v}$ 上平坦であることと、
$\mathcal{F}_{\overline{x}}$ が $\mathcal{O}_{V, v}$ 上平坦であることとは同値である。
ゆえに More on Flatness, Lemma
\ref{flat-lemma-flat-up-down-henselization} から結論が従う。
\end{proof}

\begin{definition}
\label{spaces-morphisms-definition-flat-module}
$S$ をスキームとする。
$f : X \to Y$ を $S$ 上の代数空間の射とする。
$\mathcal{F}$ を $X$ 上の準連接層とする。
\begin{enumerate}
\item $x \in |X|$ とする。Lemma \ref{spaces-morphisms-lemma-flat-at-point} の同値な条件が
成り立つとき、$\mathcal{F}$ は {\it $x$ において $Y$ 上平坦}であるという。
\item $\mathcal{F}$ がすべての $x \in |X|$ において $Y$ 上平坦であるとき、
$\mathcal{F}$ は {\it $Y$ 上平坦}であるという。
\end{enumerate}
\end{definition}

\noindent
この定義に伴う必須の基底変換補題を述べる。この補題は、準連接層が
平坦である軌跡の形成が平坦基底変換と可換することを含意する。

\begin{lemma}
\label{spaces-morphisms-lemma-base-change-module-flat}
$S$ をスキームとし、
$$
\xymatrix{
X' \ar[d]_{f'} \ar[r]_{g'} & X \ar[d]^f \\
Y' \ar[r]^g & Y
}
$$
を $S$ 上の代数空間のカルテジアン図式とする。
$x' \in |X'|$ とし、その像を $x \in |X|$ とする。
$\mathcal{F}$ を $X$ 上の準連接層とし、
$\mathcal{F}' = (g')^*\mathcal{F}$ と表す。
\begin{enumerate}
\item $\mathcal{F}$ が $x$ において $Y$ 上平坦ならば、
$\mathcal{F}'$ は $x'$ において $Y'$ 上平坦である。
\item $g$ が $f'(x')$ において平坦で、$\mathcal{F}'$ が
$x'$ において $Y'$ 上平坦ならば、$\mathcal{F}$ は $x$ において
$Y$ 上平坦である。
\end{enumerate}
特に、$\mathcal{F}$ が $Y$ 上平坦ならば、$\mathcal{F}'$ は $Y'$ 上平坦である。
\end{lemma}

\begin{proof}
スキーム $V$ と全射エタール射 $V \to Y$ を取る。
スキーム $U$ と全射エタール射 $U \to V \times_Y X$ を取る。
スキーム $V'$ と全射エタール射 $V' \to V \times_Y Y'$ を取る。
すると $U' = V' \times_V U$ は、全射エタール射
$U' = V' \times_V U \to Y' \times_Y X = X'$ を備えたスキームである。
$x' \in |X'|$ に写る $u' \in U'$ を取る。このとき
$x'$ における $\mathcal{F}'$ の $Y'$ 上の平坦性を、$u'$ における
$\mathcal{F}'|_{U'}$ の $V'$ 上の平坦性によって確認できる。
したがって補題は More on Morphisms, Lemma
\ref{more-morphisms-lemma-flat-locus-base-change} から従う。
\end{proof}

\noindent
次の補題は、加群を用いて平坦射の「合成」を論じる。
また、平坦性が一種の上から下への降下を満たすことも示す。

\begin{lemma}
\label{spaces-morphisms-lemma-composition-module-flat}
$S$ をスキームとする。$X \to Y \to Z$ を $S$ 上の代数空間の射とする。
$\mathcal{F}$ を $X$ 上の準連接層とする。
$x \in |X|$ とし、その像を $y \in |Y|$ とする。
\begin{enumerate}
\item $\mathcal{F}$ が $x$ において $Y$ 上平坦で、$Y$ が $y$ において
$Z$ 上平坦ならば、$\mathcal{F}$ は $x$ において $Z$ 上平坦である。
\item $x : \Spec(K) \to X$ を $x$ の代表とする。次の条件が成り立つとする。
\begin{enumerate}
\item $\mathcal{F}$ は $x$ において $Y$ 上平坦である。
\item $x^*\mathcal{F} \not = 0$ である。
\item $\mathcal{F}$ は $x$ において $Z$ 上平坦である。
\end{enumerate}
このとき $Y$ は $y$ において $Z$ 上平坦である。
\item $\overline{x}$ を $x$ の上にある $X$ の幾何点とし、その $Y$ における像を
$\overline{y}$ とする。$\mathcal{F}_{\overline{x}}$ が忠実平坦な
$\mathcal{O}_{Y, \overline{y}}$-加群で、$\mathcal{F}$ が $x$ において
$Z$ 上平坦ならば、$Y$ は $y$ において $Z$ 上平坦である。
\end{enumerate}
\end{lemma}

\begin{proof}
(3) のような $\overline{x}$ と $\overline{y}$ を取り、$Z$ に誘導される
幾何点を $\overline{z}$ と表す。Lemmas \ref{spaces-morphisms-lemma-flat-at-point} と
\ref{spaces-morphisms-lemma-flat-at-point-etale-local-rings} の平坦性の特徴づけにより、
補題は局所環の写像
$\mathcal{O}_{Z, \overline{z}} \to \mathcal{O}_{Y, \overline{y}}$ と加群
$\mathcal{F}_{\overline{x}}$ に関する純代数的な問題に帰着する。
(1) は Algebra, Lemma \ref{algebra-lemma-composition-flat} から従う。
条件 (2)(b) は
$\mathcal{F}_{\overline{x}}/
\mathfrak m_{\overline{y}} \mathcal{F}_{\overline{x}}$ が非零であることを保証する。
したがって (2)(a) $+$ (2)(b) は、$\mathcal{F}_{\overline{x}}$ が忠実平坦な
$\mathcal{O}_{Y, \overline{y}}$-加群であることを含意する。Algebra, Lemma
\ref{algebra-lemma-ff} を参照せよ。ゆえに (2) は (3) の特殊な場合である。
最後に、(3) は Algebra, Lemma
\ref{algebra-lemma-flat-permanence} から従う。
\end{proof}

\noindent
基底変換が「上側」で起こることもある。正確な主張を述べる。

\begin{lemma}
\label{spaces-morphisms-lemma-flat-permanence}
$S$ をスキームとする。$f : X \to Y$, $g : Y \to Z$ を
$S$ 上の代数空間の射とする。$\mathcal{G}$ を $Y$ 上の準連接層とする。
$x \in |X|$ とし、その像を $y \in |Y|$ とする。
$f$ が $x$ において平坦ならば、
$$
\mathcal{G}\text{ は次の上で平坦：}Z\text{、点：}y
\Leftrightarrow
f^*\mathcal{G}\text{ は次の上で平坦：}Z\text{、点：}x.
$$
特に、$f$ が全射かつ平坦ならば、$\mathcal{G}$ が $Z$ 上平坦であることと、
$f^*\mathcal{G}$ が $Z$ 上平坦であることとは同値である。
\end{lemma}

\begin{proof}
$X$ の幾何点 $\overline{x}$ を取り、$Y$ における像を $\overline{y}$、
$Z$ における像を $\overline{z}$ と表す。
Lemmas \ref{spaces-morphisms-lemma-flat-at-point} と
\ref{spaces-morphisms-lemma-flat-at-point-etale-local-rings} の平坦性の特徴づけ、および
Properties of Spaces, Lemma
\ref{spaces-properties-lemma-stalk-pullback-quasi-coherent} による
$\overline{x}$ における $f^*\mathcal{G}$ の茎の記述を用いると、補題は
局所環の写像
$\mathcal{O}_{Z, \overline{z}} \to \mathcal{O}_{Y, \overline{y}}
\to \mathcal{O}_{X, \overline{x}}$ と加群 $\mathcal{G}_{\overline{y}}$ に関する
純代数的な問題に帰着する。この代数的主張は Algebra, Lemma
\ref{algebra-lemma-flatness-descends-more-general} である。
\end{proof}

\begin{lemma}
\label{spaces-morphisms-lemma-pf-flat-module-open}
$S$ をスキームとする。
$f : X \to Y$ を $S$ 上の代数空間の射とする。
$\mathcal{F}$ を準連接 $\mathcal{O}_X$-加群とする。
$f$ が局所有限表示、$\mathcal{F}$ が有限型、
$X = \text{Supp}(\mathcal{F})$ であり、$\mathcal{F}$ が $Y$ 上平坦であると仮定する。
このとき $f$ は普遍的開である。
\end{lemma}

\begin{proof}
$V$ がスキームであるような全射エタール射 $\varphi : V \to Y$ を取る。
$U$ がスキームであるような全射エタール射 $U \to V \times_Y X$ を取る。
すると、$U \to V$ と準連接 $\mathcal{O}_V$-加群
$\varphi^*\mathcal{F}$ に対して補題を示せば十分である。
したがって補題はスキームの場合から従う。Morphisms, Lemma
\ref{morphisms-lemma-pf-flat-module-open} を参照せよ。
\end{proof}






\section{一般平坦性}
\label{spaces-morphisms-section-generic-flatness}

\noindent
この節は Morphisms, Section
\ref{morphisms-section-generic-flatness} に対応するものである。

\begin{proposition}
\label{spaces-morphisms-proposition-generic-flatness-reduced}
$S$ をスキームとする。
$f : X \to Y$ を $S$ 上の代数空間の射とする。
$\mathcal{F}$ を準連接 $\mathcal{O}_X$-加群層とする。次を仮定する。
\begin{enumerate}
\item $Y$ は被約である。
\item $f$ は有限型である。
\item $\mathcal{F}$ は有限型 $\mathcal{O}_X$-加群である。
\end{enumerate}
このとき、開稠密部分空間 $W \subset Y$ であって、$f$ の基底変換
$X_W \to W$ が平坦、局所有限表示かつ準コンパクトであり、
$\mathcal{F}|_{X_W}$ が $W$ 上平坦かつ $\mathcal{O}_{X_W}$ 上有限表示となるものが存在する。
\end{proposition}

\begin{proof}
$V$ をスキームとし、$V \to Y$ を全射エタール射とする。
$X_V = V \times_Y X$ とおき、$\mathcal{F}_V$ を $X_V$ への
$\mathcal{F}$ の制限とする。射 $X_V \to V$ と層 $\mathcal{F}_V$ に対して
結果が成り立つと仮定する。このとき開部分スキーム $V' \subset V$ であって、
$X_{V'} \to V'$ が平坦かつ有限表示であり、$\mathcal{F}_{V'}$ が
$V'$ 上平坦な有限表示 $\mathcal{O}_{X_{V'}}$-加群となるものが存在する。
$W \subset Y$ をエタール射 $V' \to Y$ の像とする。Properties of Spaces,
Lemma \ref{spaces-properties-lemma-etale-image-open} を参照せよ。
すると $V' \to W$ は全射エタール射なので、Lemmas
\ref{spaces-morphisms-lemma-finite-presentation-local}、\ref{spaces-morphisms-lemma-flat-local}、および
\ref{spaces-morphisms-lemma-quasi-compact-local} により、$X_W \to W$ は平坦、
局所有限表示かつ準コンパクトである。Properties of Spaces, Section
\ref{spaces-properties-section-properties-modules} の議論により、
$\mathcal{F}_W$ は有限表示 $\mathcal{O}_{X_W}$-加群であり、Lemma
\ref{spaces-morphisms-lemma-base-change-module-flat} により $\mathcal{F}_W$ は $W$ 上平坦である。
この議論により、命題は $Y$ がスキームの場合に帰着する。

\medskip\noindent
$Y$ がアフィンスキームの場合に命題を証明できると仮定する。
$f : X \to Y$ を、$Y$ がスキームであるような $S$ 上の代数空間の有限型射とし、
$\mathcal{F}$ を有限型準連接 $\mathcal{O}_X$-加群とする。
アフィン開被覆 $Y = \bigcup V_j$ を取る。仮定により、稠密開集合
$W_j \subset V_j$ であって、$X_{W_j} \to W_j$ が平坦、局所有限表示かつ
準コンパクトであり、$\mathcal{F}|_{X_{W_j}}$ が $W_j$ 上平坦かつ
有限表示 $\mathcal{O}_{X_{W_j}}$-加群となるものを取れる。
この状況では単に $W = \bigcup W_j$ と取れば結論を得る。
したがって命題は $Y$ がアフィンスキームである場合に帰着する。

\medskip\noindent
$Y$ を $S$ 上のアフィンスキーム、$f : X \to Y$ を $S$ 上の
代数空間の有限型射、$\mathcal{F}$ を有限型準連接 $\mathcal{O}_X$-加群とする。
$f$ は有限型なので準コンパクトであり、したがって $X$ は準コンパクトである。
よって、アフィンスキーム $U$ と全射エタール射 $U \to X$ を取れる。
Properties of Spaces, Lemma
\ref{spaces-properties-lemma-quasi-compact-affine-cover} を参照せよ。
$U \to Y$ は有限型であることに注意する（この場合、これが $f$ が
有限型であることの意味である）。したがって Morphisms, Proposition
\ref{morphisms-proposition-generic-flatness-reduced} を適用すると、
稠密開集合 $W \subset Y$ であって、$U_W \to W$ が平坦かつ有限表示であり、
$\mathcal{F}|_{U_W}$ が $W$ 上平坦な有限表示 $\mathcal{O}_{U_W}$-加群となるものが存在する。
定義により、これは $f$ の基底変換 $X_W \to W$ が平坦、局所有限表示かつ
準コンパクトであり、$\mathcal{F}|_{X_W}$ が $W$ 上平坦かつ
$\mathcal{O}_{X_W}$ 上有限表示であることを意味する。
\end{proof}

\noindent
上の補題の結果を $X_W \to W$ が有限表示であるという要請まで強めることはできない。
$\mathbf{A}^1_{\mathbf{Q}}/\mathbf{Z} \to \Spec(\mathbf{Q})$ が反例を与える。
問題は、対角射 $\Delta_{X/Y}$ が準コンパクトでない可能性、すなわち
$f$ が準分離でない可能性にある。明らかに、これが唯一の問題でもある。

\begin{proposition}
\label{spaces-morphisms-proposition-generic-flatness-reduced-quasi-separated}
$S$ をスキームとする。
$f : X \to Y$ を $S$ 上の代数空間の射とする。
$\mathcal{F}$ を準連接 $\mathcal{O}_X$-加群層とする。次を仮定する。
\begin{enumerate}
\item $Y$ は被約である。
\item $f$ は準分離である。
\item $f$ は有限型である。
\item $\mathcal{F}$ は有限型 $\mathcal{O}_X$-加群である。
\end{enumerate}
このとき、開稠密部分空間 $W \subset Y$ であって、$f$ の基底変換
$X_W \to W$ が平坦かつ有限表示であり、$\mathcal{F}|_{X_W}$ が
$W$ 上平坦かつ $\mathcal{O}_{X_W}$ 上有限表示となるものが存在する。
\end{proposition}

\begin{proof}
Proposition \ref{spaces-morphisms-proposition-generic-flatness-reduced} と、
「有限表示」$=$「局所有限表示」$+$「準コンパクト」$+$「準分離」
という事実から直ちに従う。
\end{proof}

\section{相対次元}
\label{spaces-morphisms-section-relative-dimension}

\noindent
この節では、代数空間の射の一点における相対次元と、
それに密接に関連するいくつかの性質を定義する。

\begin{definition}
\label{spaces-morphisms-definition-dimension-fibre}
$S$ をスキームとする。
$f : X \to Y$ を $S$ 上の代数空間の射とする。
$x \in |X|$ とする。
$d, r \in \{0, 1, 2, \ldots, \infty\}$ とする。
\begin{enumerate}
\item Lemma \ref{spaces-morphisms-lemma-local-source-target-at-point} の同値な条件が、
Descent, Lemma \ref{descent-lemma-dimension-local-ring-fibre} で述べた性質
$\mathcal{P}_d$ に対して成り立つとき、
{\it $x$ における $f$ のファイバーの局所環の次元}は $d$ であるという。
\item Lemma \ref{spaces-morphisms-lemma-local-source-target-at-point} の同値な条件が、
Descent, Lemma \ref{descent-lemma-transcendence-degree-at-point} で述べた性質
$\mathcal{P}_r$ に対して成り立つとき、
{\it $x/f(x)$ の超越次数}は $r$ であるという。
\item Lemma \ref{spaces-morphisms-lemma-local-source-target-at-point} の同値な条件が、
Descent, Lemma \ref{descent-lemma-dimension-at-point} で述べた性質
$\mathcal{P}_d$ に対して成り立つとき、
{\it $f$ は $x$ において相対次元 $d$ をもつ}という。
\end{enumerate}
\end{definition}

\noindent
この定義の意味を明記しておこう。すなわち、
Lemma \ref{spaces-morphisms-lemma-local-source-target-at-point} のような図式
$$
\xymatrix{
U \ar[d]_a \ar[r]_h & V \ar[d]^b \\
X \ar[r]^f & Y
}
\quad\quad
\xymatrix{
u \ar[d] \ar[r] & v \ar[d] \\
x \ar[r] & y
}
$$
を選ぶ。このとき
$$
\begin{matrix}
\text{相対次元：}f\text{、点 }x & = &
\dim_u (U_v) \\
\text{ファイバーの局所環の次元：}f\text{、点 }x & = &
\dim(\mathcal{O}_{U_v, u})\\
\text{超越次数：}x/f(x) & = &
\text{trdeg}_{\kappa(v)}(\kappa(u))
\end{matrix}
$$
を得る。$Y = \Spec(k)$ が体のスペクトルである場合、
$x$ における $X/Y$ の相対次元は $\dim_x(X)$ と同じであり、
$x/f(x)$ の超越次数は $k$ 上の超越次数であり、
$x$ における $f$ のファイバーの局所環の次元は
$x$ における局所環の次元にほかならない。
したがって、この場合には相対的な概念は絶対的な概念になる。

\begin{definition}
\label{spaces-morphisms-definition-relative-dimension}
$S$ をスキームとする。
$f : X \to Y$ を $S$ 上の代数空間の射とする。
$d \in \{0, 1, 2, \ldots\}$ とする。
\begin{enumerate}
\item すべての $x \in |X|$ において $f$ の相対次元が $\leq d$ であるとき、
$f$ は {\it 相対次元 $\leq d$} をもつという。
\item すべての $x \in |X|$ において $f$ の相対次元が $d$ であるとき、
$f$ は {\it 相対次元 $d$} をもつという。
\end{enumerate}
\end{definition}

\noindent
相対次元が $d$ に{\it 等しい}とは、大まかには、
空でないすべてのファイバーが次元 $d$ の等次元空間であることを意味する。

\begin{lemma}
\label{spaces-morphisms-lemma-compare-tr-deg}
$S$ をスキームとする。$X \to Y \to Z$ を $S$ 上の代数空間の射とする。
$x \in |X|$ とし、その像を $y \in |Y|$, $z \in |Z|$ とする。
$X \to Y$ は局所準有限であり、$Y \to Z$ は局所有限型であると仮定する。
このとき $x/z$ の超越次数は $y/z$ の超越次数に等しい。
\end{lemma}

\begin{proof}
可換図式
$$
\xymatrix{
U \ar[d] \ar[r] & V \ar[d] \ar[r] & W \ar[d] \\
X \ar[r] & Y \ar[r] & Z
}
\quad\quad
\xymatrix{
u \ar[d] \ar[r] & v \ar[d] \ar[r] & w \ar[d] \\
x \ar[r] & y \ar[r] & z
}
$$
で、$U, V, W$ がスキームであり、垂直矢印がエタールであるものを選べる。
定義により射 $U \to V$ は局所準有限であるから、
$\kappa(v) \subset \kappa(u)$ は有限である
（Morphisms, Lemma
\ref{morphisms-lemma-residue-field-quasi-finite}）。
したがって結果は明らかである。
\end{proof}

\begin{lemma}
\label{spaces-morphisms-lemma-jacobson-finite-type-points}
$S$ をスキームとする。$f : X \to Y$ を $S$ 上の代数空間の射とする。
$f$ が局所有限型であり、$Y$ が Jacobson であり
（Properties of Spaces, Remark
\ref{spaces-properties-remark-list-properties-local-etale-topology}）、
$x \in |X|$ が $X$ の有限型点であるならば、
$x/f(x)$ の超越次数は $0$ である。
\end{lemma}

\begin{proof}
$V$ をスキームとし、$V \to Y$ を全射エタール射として選ぶ。
$U$ をスキームとし、$U \to X \times_Y V$ を全射エタール射として選ぶ。
Lemma \ref{spaces-morphisms-lemma-finite-type-points-surjective-morphism} により、
$x$ に写る有限型点 $u \in U$ を見つけられる。
$U$ を縮小することにより、$u \in U$ は閉点であるとしてよい
（Morphisms, Lemma
\ref{morphisms-lemma-identify-finite-type-points}）。
$v \in V$ を $u$ の像とする。Morphisms, Lemma
\ref{morphisms-lemma-jacobson-finite-type-points} により、
拡大 $\kappa(u)/\kappa(v)$ は有限である。これで証明が完了する。
\end{proof}

\begin{lemma}
\label{spaces-morphisms-lemma-rel-dimension-dimension}
$S$ をスキームとする。$f : X \to Y$ を $S$ 上の局所 Noether
代数空間の射で、平坦、局所有限型、かつ相対次元 $d$ をもつものとする。
$x$ を $|X|$ の任意の点とし、その像 $y$ を $|Y|$ の点とする。このとき
$\dim_x(X) = \dim_y(Y) + d$ が成り立つ。
\end{lemma}

\begin{proof}
一点における代数空間の次元の定義
（Properties of Spaces, Definition
\ref{spaces-properties-definition-dimension-at-point}）と、
相対次元 $d$ をもつことの定義により、これはスキームに対する対応する主張
（Morphisms, Lemma
\ref{morphisms-lemma-rel-dimension-dimension}）に帰着する。
\end{proof}

\section{射とファイバーの次元}
\label{spaces-morphisms-section-dimension-fibres}

\noindent
この節は Morphisms, Section
\ref{morphisms-section-dimension-fibres} の類似である。
一点における代数空間の局所環をもたないため、
この節の定式化はやや不格好になる。

\begin{lemma}
\label{spaces-morphisms-lemma-dimension-fibre-at-a-point}
$S$ をスキームとする。
$f : X \to Y$ を $S$ 上の代数空間の射とする。
$x \in |X|$ とする。
$f$ は局所有限型であると仮定する。このとき
$$
\begin{matrix}
\text{相対次元：}f\text{、点 }x \\
= \\
\text{ファイバーの局所環の次元：}f\text{、点 }x \\
+ \\
\text{超越次数：}x/f(x)
\end{matrix}
$$
が成り立つ。記法は Definition
\ref{spaces-morphisms-definition-dimension-fibre} のとおりである。
\end{lemma}

\begin{proof}
Morphisms, Lemma
\ref{morphisms-lemma-dimension-fibre-at-a-point} を、
Lemma \ref{spaces-morphisms-lemma-local-source-target-at-point} のような
$h : U \to V$ と $u \in U$ に適用すれば直ちに従う。
\end{proof}

\begin{lemma}
\label{spaces-morphisms-lemma-dimension-fibre-at-a-point-additive}
$S$ をスキームとする。
$f : X \to Y$ と $g : Y \to Z$ を $S$ 上の代数空間の射とする。
$x \in |X|$ とし、$y = f(x)$ とおく。
$f$ と $g$ は局所有限型であると仮定する。このとき
\begin{enumerate}
\item
$$
\begin{matrix}
\text{相対次元：}g \circ f\text{、点 }x \\
\leq \\
\text{相対次元：}f\text{、点 }x \\
+ \\
\text{相対次元：}g\text{、点 }y
\end{matrix}
$$
\item $g(f(x)) = g(y)$ の類に属する体のスペクトルからのある射
$\Spec(k) \to Z$ に対して、射 $X_k \to Y_k$ が $x$ において平坦ならば、
(1) で等号が成り立つ。たとえば $f$ が $x$ において平坦ならば成り立つ。
\item
$$
\begin{matrix}
\text{超越次数：}x/g(f(x)) \\
= \\
\text{超越次数：}x/f(x) \\
+ \\
\text{超越次数：}f(x)/g(f(x))
\end{matrix}
$$
\end{enumerate}
\end{lemma}

\begin{proof}
垂直矢印がエタールかつ全射であり、$U, V, W$ がスキームである図式
$$
\xymatrix{
U \ar[d] \ar[r] & V \ar[d] \ar[r] & W \ar[d] \\
X \ar[r] & Y \ar[r] & Z
}
$$
を選ぶ（Spaces, Lemma
\ref{spaces-lemma-lift-morphism-presentations}）。
$x$ に写る $u \in U$ を選び、$u$ の $V, W$ における像をそれぞれ
$v, w$ とする。上段と点 $u, v, w$ に Morphisms, Lemma
\ref{morphisms-lemma-dimension-fibre-at-a-point-additive} を適用する。
詳細は省略する。
\end{proof}

\begin{lemma}
\label{spaces-morphisms-lemma-dimension-fibre-after-base-change}
$S$ をスキームとする。
$$
\xymatrix{
X' \ar[r]_{g'} \ar[d]_{f'} & X \ar[d]^f \\
Y' \ar[r]^g & Y
}
$$
を $S$ 上の代数空間のファイバー積図式とする。
$x' \in |X'|$ とし、$x = g'(x')$ とおく。
$f$ は局所有限型であると仮定する。このとき
\begin{enumerate}
\item
$$
\begin{matrix}
\text{相対次元：}f\text{、点 }x \\
= \\
\text{相対次元：}f'\text{、点 }x'
\end{matrix}
$$
\item
$$
\begin{matrix}
\text{ファイバーの局所環の次元：}f'\text{、点 }x' \\
- \\
\text{ファイバーの局所環の次元：}f\text{、点 }x \\
= \\
\text{超越次数：}x/f(x) \\
- \\
\text{超越次数：}x'/f'(x')
\end{matrix}
$$
が成り立ち、その共通値は $\geq 0$ である。
\item 同じ $y \in |Y|$ に写る $x$ と $y' \in |Y'|$ が与えられたとき、
(2) の整数が $0$ となる $x'$ を選ぶことができる。
\end{enumerate}
\end{lemma}

\begin{proof}
$V$ をスキームとし、$V \to Y$ を全射エタール射として選ぶ。
$U$ をスキームとし、$U \to V \times_Y X$ を全射エタール射として選ぶ。
$V'$ をスキームとし、$V' \to V \times_Y Y'$ を全射エタール射として選ぶ。
$U' = V' \times_V U$ とおく。このとき誘導射 $U' \to X'$ も
全射かつエタールである（議論は省略する）。
$x'$ に写る $u' \in U'$ を選ぶ。この時点で、$U', U, V', V$ と点 $u'$
からなるスキームの図式に Morphisms, Lemma
\ref{morphisms-lemma-dimension-fibre-after-base-change} を適用すれば
(1) と (2) が従う。
(3) を示すため、まず $y$ に写る $v \in V$ を選ぶ。
次に Properties of Spaces, Lemma
\ref{spaces-properties-lemma-points-cartesian} により、
$y'$ と $v$ に写る $v' \in V'$、および $x$ と $v$ に写る
$u \in U$ を選べる。最後に Morphisms, Lemma
\ref{morphisms-lemma-dimension-fibre-after-base-change} により、
$v'$ と $u$ に写り、その整数がゼロとなる $u' \in U'$ を選べる。
この $u'$ の像 $x' \in |X'|$ を取ればよい。
\end{proof}

\begin{lemma}
\label{spaces-morphisms-lemma-openness-bounded-dimension-fibres}
$S$ をスキームとする。
$f : X \to Y$ を $S$ 上の代数空間の射とする。
$n \geq 0$ とする。$f$ は局所有限型であると仮定する。集合
$$
W_n = \{x \in |X|
\text{ であって、\(x\) における }f\text{ の相対次元が } x \leq n\}
$$
は $|X|$ で開である。
\end{lemma}

\begin{proof}
図式
$$
\xymatrix{
U \ar[r]_h \ar[d]_a & V \ar[d] \\
X \ar[r] & Y
}
$$
で、$U$ と $V$ がスキームであり、垂直矢印が全射かつエタールであるものを選ぶ
（Spaces, Lemma
\ref{spaces-lemma-lift-morphism-presentations}）。
Morphisms, Lemma
\ref{morphisms-lemma-openness-bounded-dimension-fibres} により、
$h$ の相対次元が $\leq n$ となる点の集合 $U_n$ は $U$ で開である。
一点における代数空間の射の相対次元の定義から
$U_n = a^{-1}(W_n)$ である。
したがって、$|X|$ の位相の定義により補題が従う。
\end{proof}

\begin{lemma}
\label{spaces-morphisms-lemma-openness-bounded-dimension-fibres-finite-presentation}
$S$ をスキームとする。
$f : X \to Y$ を $S$ 上の代数空間の射とする。
$n \geq 0$ とする。$f$ は局所有限表示であると仮定する。
Lemma \ref{spaces-morphisms-lemma-openness-bounded-dimension-fibres} の開集合
$$
W_n = \{x \in |X|
\text{ であって、\(x\) における }f\text{ の相対次元が } x \leq n\}
$$
は $|X|$ 内でレトロコンパクトである
（Topology, Definition
\ref{topology-definition-quasi-compact}）。
\end{lemma}

\begin{proof}
図式
$$
\xymatrix{
U \ar[r]_h \ar[d]_a & V \ar[d] \\
X \ar[r] & Y
}
$$
で、$U$ と $V$ がスキームであり、垂直矢印が全射かつエタールであるものを選ぶ
（Spaces, Lemma
\ref{spaces-lemma-lift-morphism-presentations}）。
Lemma \ref{spaces-morphisms-lemma-openness-bounded-dimension-fibres} の証明で、
$a^{-1}(W_n) = U_n$ は射 $h$ に対する対応する集合であることを見た。
Morphisms, Lemma
\ref{morphisms-lemma-openness-bounded-dimension-fibres-finite-presentation}
により、$U_n$ は $U$ 内でレトロコンパクトである。
Properties of Spaces, Lemma
\ref{spaces-properties-lemma-space-locally-quasi-compact} とその証明も参照すれば、
$|X|$ の位相の定義から補題が従う。
\end{proof}

\begin{lemma}
\label{spaces-morphisms-lemma-locally-quasi-finite-rel-dimension-0}
$S$ をスキームとする。
$f : X \to Y$ を $S$ 上の代数空間の射とする。
$f$ は局所有限型であると仮定する。
このとき、$f$ が局所準有限であることと、
すべての $x \in |X|$ において $f$ の相対次元が $0$ であることは同値である。
\end{lemma}

\begin{proof}
図式
$$
\xymatrix{
U \ar[r]_h \ar[d]_a & V \ar[d] \\
X \ar[r] & Y
}
$$
で、$U$ と $V$ がスキームであり、垂直矢印が全射かつエタールであるものを選ぶ
（Spaces, Lemma
\ref{spaces-lemma-lift-morphism-presentations}）。
定義から、$h$ が局所準有限であることと $f$ が局所準有限であることは同値であり、
また、すべての $x \in |X|$ において $f$ が相対次元 $0$ をもつことと、
すべての $u \in U$ において $h$ が相対次元 $0$ をもつことは同値である。
したがって、$h$ に対する結果、すなわち Morphisms, Lemma
\ref{morphisms-lemma-locally-quasi-finite-rel-dimension-0} から従う。
\end{proof}

\begin{lemma}
\label{spaces-morphisms-lemma-locally-finite-type-quasi-finite-part}
$S$ をスキームとする。
$f : X \to Y$ を $S$ 上の代数空間の射とする。
$f$ は局所有限型であると仮定する。
このとき標準的な開部分空間 $X' \subset X$ であって、
$f|_{X'} : X' \to Y$ が局所準有限であり、
任意の $x \in |X|$, $x \not \in |X'|$ における $f$ の相対次元が
$\geq 1$ となるものが存在する。$X'$ の形成は任意の基底変換と可換である。
\end{lemma}

\begin{proof}
Lemmas \ref{spaces-morphisms-lemma-openness-bounded-dimension-fibres},
\ref{spaces-morphisms-lemma-locally-quasi-finite-rel-dimension-0}, および
\ref{spaces-morphisms-lemma-dimension-fibre-after-base-change} を組み合わせる。
\end{proof}

\begin{lemma}
\label{spaces-morphisms-lemma-quasi-finite-at-point}
$S$ をスキームとする。カルテジアン図式
$$
\xymatrix{
X \ar[d] & F \ar[l]^p \ar[d] \\
Y & \Spec(k) \ar[l]
}
$$
を考える。ここで $X \to Y$ は $S$ 上の代数空間の局所有限型射であり、
$k$ は $S$ 上の体である。
$z \in |F|$ が $\dim_z(F) = 0$ を満たすとする。
このとき、$X$ を $p(z)$ を含む開部分空間で置き換えた後、射
$$
X \longrightarrow Y
$$
は局所準有限である。
\end{lemma}

\begin{proof}
Lemma \ref{spaces-morphisms-lemma-locally-finite-type-quasi-finite-part} で得られる、
$f$ が局所準有限となる開部分空間を $X' \subset X$ とする。
$X'$ の形成は任意の基底変換と可換であるから、
$z \in X' \times_Y \Spec(k)$ である。したがって補題は明らかである。
\end{proof}

\section{次元公式}
\label{spaces-morphisms-section-dimension-formula}

\noindent
次元公式（Morphisms, Lemma
\ref{morphisms-lemma-dimension-formula}）の類似を定式化するのはやや難しい。
整代数空間（これは後で定義する）と普遍鎖状代数空間を
定義しなければならないからである。しかし、次の形は直接的である。

\begin{lemma}
\label{spaces-morphisms-lemma-dimension-formula-general}
$S$ をスキームとする。$f : X \to Y$ を $S$ 上の代数空間の射とする。
$Y$ は局所 Noether であり、$f$ は局所有限型であると仮定する。
$x \in |X|$ とし、その像を $y \in |Y|$ とする。このとき
\begin{align*}
& \text{局所環の次元：}X\text{、点 }x \leq \\
& \text{局所環の次元：}Y\text{、点 }y + E - \\
& \text{超越次数：}x/y
\end{align*}
が成り立つ。ここで $E$ は、$\xi \in |X|$ が $x$ に特殊化する点で、
$X$ の局所環の次元が $0$ となるものを動くときの、
$\xi/f(\xi)$ の超越次数の最大値である。
\end{lemma}

\begin{proof}
アフィンスキーム $V$、エタール射 $V \to Y$、および $y$ に写る点
$v \in V$ を選ぶ。アフィンスキーム $U$、エタール射
$U \to X \times_Y V$、および $V$ では $v$ に、$X$ では $x$ に写る点
$u \in U$ を選ぶ。Definition
\ref{spaces-morphisms-definition-dimension-fibre} と Properties of Spaces, Definition
\ref{spaces-properties-definition-dimension-local-ring} を展開すると、
示すべきことは
$$
\dim(\mathcal{O}_{U, u}) \leq
\dim(\mathcal{O}_{V, v}) + E - \text{trdeg}_{\kappa(v)}(\kappa(u))
$$
である。$\xi_U \in U$ を、$u$ を含む $U$ の既約成分の一般点とする。
すると $\xi_U$ は、量 $E$ の定義に用いた点のリストに属する
$\xi \in |X|$ に写る。実際、$E$ の定義に用いるすべての $\xi$ は
このようにして現れる（細部は省略する）。
とくに、そのような $\xi$ は有限個しかないので最大値を取れる
（すなわち、これは上限ではなく実際に最大値である）。
$\xi$ の $f(\xi)$ 上の超越次数は
$\text{trdeg}_{\kappa(\xi_V)}(\kappa(\xi_U))$ である。
ここで $\xi_V \in V$ は $\xi_U$ の像である。
したがって Morphisms, Lemma
\ref{morphisms-lemma-dimension-formula-general} から補題が従う。
\end{proof}

\begin{lemma}
\label{spaces-morphisms-lemma-alteration-dimension-general}
$S$ をスキームとする。
$f : X \to Y$ を $S$ 上の代数空間の射とする。
$Y$ は局所 Noether であり、$f$ は局所有限型であると仮定する。このとき
$$
\dim(X) \leq \dim(Y) + E
$$
が成り立つ。ここで $E$ は、$\xi$ が $X$ の局所環の次元が $0$
となる点を動くときの、$\xi/f(\xi)$ の超越次数の上限である。
\end{lemma}

\begin{proof}
Lemma \ref{spaces-morphisms-lemma-dimension-formula-general} と Properties of Spaces, Lemma
\ref{spaces-properties-lemma-dimension} の直ちに得られる帰結である。
\end{proof}

\section{シントミック射}
\label{spaces-morphisms-section-syntomic}

\noindent
スキームの射の「シントミック」という性質は始域と終域の上で
エタール局所的である（Descent, Remark
\ref{descent-remark-list-local-source-target}）。
また、基底変換で保たれ、終域上 fpqc 局所的でもある
（Morphisms, Lemma \ref{morphisms-lemma-base-change-syntomic} および
Descent, Lemma
\ref{descent-lemma-descending-property-syntomic}）。
したがって上の Lemma \ref{spaces-morphisms-lemma-local-source-target} により、
代数空間のシントミック射を次のように定義できる。
この定義は、射が表現可能であるとき、Section
\ref{spaces-morphisms-section-representable} ですでに定義した概念と一致する。

\begin{definition}
\label{spaces-morphisms-definition-syntomic}
$S$ をスキームとする。
$f : X \to Y$ を $S$ 上の代数空間の射とする。
\begin{enumerate}
\item Lemma \ref{spaces-morphisms-lemma-local-source-target} の同値な条件が
$\mathcal{P} =$「シントミック」に対して成り立つとき、
$f$ は {\it シントミック}であるという。
\item $x \in |X|$ とする。$x$ の開近傍 $X' \subset X$ であって、
$f|_{X'} : X' \to Y$ がシントミックとなるものが存在するとき、
$f$ は {\it $x$ においてシントミック}であるという。
\end{enumerate}
\end{definition}

\begin{lemma}
\label{spaces-morphisms-lemma-composition-syntomic}
シントミック射の合成はシントミックである。
\end{lemma}

\begin{proof}
Remark \ref{spaces-morphisms-remark-composition-P} および Morphisms, Lemma
\ref{morphisms-lemma-composition-syntomic} を参照。
\end{proof}

\begin{lemma}
\label{spaces-morphisms-lemma-base-change-syntomic}
シントミック射の基底変換はシントミックである。
\end{lemma}

\begin{proof}
Remark \ref{spaces-morphisms-remark-base-change-P} および Morphisms, Lemma
\ref{morphisms-lemma-base-change-syntomic} を参照。
\end{proof}

\begin{lemma}
\label{spaces-morphisms-lemma-syntomic-local}
$S$ をスキームとする。
$f : X \to Y$ を $S$ 上の代数空間の射とする。
次の条件は同値である。
\begin{enumerate}
\item $f$ はシントミックである。
\item すべての $x \in |X|$ に対して、射 $f$ は $x$ において
シントミックである。
\item 任意のスキーム $Z$ と任意の射 $Z \to Y$ に対して、
射 $Z \times_Y X \to Z$ はシントミックである。
\item 任意のアフィンスキーム $Z$ と任意の射 $Z \to Y$ に対して、
射 $Z \times_Y X \to Z$ はシントミックである。
\item スキーム $V$ と全射エタール射 $V \to Y$ であって、
$V \times_Y X \to V$ がシントミック射となるものが存在する。
\item スキーム $U$ と全射エタール射 $\varphi : U \to X$ であって、
合成 $f \circ \varphi$ がシントミックとなるものが存在する。
\item $U$, $V$ がスキームで垂直矢印がエタールである任意の可換図式
$$
\xymatrix{
U \ar[d] \ar[r] & V \ar[d] \\
X \ar[r] & Y
}
$$
に対して、上の水平矢印はシントミックである。
\item $U$, $V$ がスキーム、垂直矢印がエタール、$U \to X$ が全射で、
上の水平矢印がシントミックであるような可換図式
$$
\xymatrix{
U \ar[d] \ar[r] & V \ar[d] \\
X \ar[r] & Y
}
$$
が存在する。
\item Zariski 被覆 $Y = \bigcup_{i \in I} Y_i$ および
$f^{-1}(Y_i) = \bigcup X_{ij}$ であって、
各射 $X_{ij} \to Y_i$ がシントミックとなるものが存在する。
\end{enumerate}
\end{lemma}

\begin{proof}
省略する。
\end{proof}

\begin{lemma}
\label{spaces-morphisms-lemma-syntomic-locally-finite-presentation}
シントミック射は局所有限表示である。
\end{lemma}

\begin{proof}
スキームの場合（Morphisms, Lemma
\ref{morphisms-lemma-syntomic-locally-finite-presentation}）から直ちに従う。
\end{proof}

\begin{lemma}
\label{spaces-morphisms-lemma-syntomic-flat}
シントミック射は平坦である。
\end{lemma}

\begin{proof}
スキームの場合（Morphisms, Lemma
\ref{morphisms-lemma-syntomic-flat}）から直ちに従う。
\end{proof}

\begin{lemma}
\label{spaces-morphisms-lemma-syntomic-open}
シントミック射は普遍開である。
\end{lemma}

\begin{proof}
Lemmas \ref{spaces-morphisms-lemma-syntomic-locally-finite-presentation},
\ref{spaces-morphisms-lemma-syntomic-flat}, および
\ref{spaces-morphisms-lemma-fppf-open} を組み合わせる。
\end{proof}

\section{滑らかな射}
\label{spaces-morphisms-section-smooth}

\noindent
スキームの射の「滑らか」という性質は始域と終域の上でエタール局所的である
（Descent, Remark
\ref{descent-remark-list-local-source-target}）。
また、基底変換で保たれ、終域上 fpqc 局所的でもある
（Morphisms, Lemma \ref{morphisms-lemma-base-change-smooth} および
Descent, Lemma
\ref{descent-lemma-descending-property-smooth}）。
したがって上の Lemma \ref{spaces-morphisms-lemma-local-source-target} により、
代数空間の滑らかな射を次のように定義できる。
この定義は、射が表現可能であるとき、Section
\ref{spaces-morphisms-section-representable} ですでに定義した概念と一致する。

\begin{definition}
\label{spaces-morphisms-definition-smooth}
$S$ をスキームとする。
$f : X \to Y$ を $S$ 上の代数空間の射とする。
\begin{enumerate}
\item Lemma \ref{spaces-morphisms-lemma-local-source-target} の同値な条件が
$\mathcal{P} =$「滑らか」に対して成り立つとき、
$f$ は {\it 滑らか}であるという。
\item $x \in |X|$ とする。$x$ の開近傍 $X' \subset X$ であって、
$f|_{X'} : X' \to Y$ が滑らかとなるものが存在するとき、
$f$ は {\it $x$ において滑らか}であるという。
\end{enumerate}
\end{definition}

\begin{lemma}
\label{spaces-morphisms-lemma-composition-smooth}
滑らかな射の合成は滑らかである。
\end{lemma}

\begin{proof}
Remark \ref{spaces-morphisms-remark-composition-P} および Morphisms, Lemma
\ref{morphisms-lemma-composition-smooth} を参照。
\end{proof}

\begin{lemma}
\label{spaces-morphisms-lemma-base-change-smooth}
滑らかな射の基底変換は滑らかである。
\end{lemma}

\begin{proof}
Remark \ref{spaces-morphisms-remark-base-change-P} および Morphisms, Lemma
\ref{morphisms-lemma-base-change-smooth} を参照。
\end{proof}

\begin{lemma}
\label{spaces-morphisms-lemma-smooth-local}
$S$ をスキームとする。
$f : X \to Y$ を $S$ 上の代数空間の射とする。
次の条件は同値である。
\begin{enumerate}
\item $f$ は滑らかである。
\item すべての $x \in |X|$ に対して、射 $f$ は $x$ において滑らかである。
\item 任意のスキーム $Z$ と任意の射 $Z \to Y$ に対して、
射 $Z \times_Y X \to Z$ は滑らかである。
\item 任意のアフィンスキーム $Z$ と任意の射 $Z \to Y$ に対して、
射 $Z \times_Y X \to Z$ は滑らかである。
\item スキーム $V$ と全射エタール射 $V \to Y$ であって、
$V \times_Y X \to V$ が滑らかな射となるものが存在する。
\item スキーム $U$ と全射エタール射 $\varphi : U \to X$ であって、
合成 $f \circ \varphi$ が滑らかとなるものが存在する。
\item $U$, $V$ がスキームで垂直矢印がエタールである任意の可換図式
$$
\xymatrix{
U \ar[d] \ar[r] & V \ar[d] \\
X \ar[r] & Y
}
$$
に対して、上の水平矢印は滑らかである。
\item $U$, $V$ がスキーム、垂直矢印がエタール、$U \to X$ が全射で、
上の水平矢印が滑らかであるような可換図式
$$
\xymatrix{
U \ar[d] \ar[r] & V \ar[d] \\
X \ar[r] & Y
}
$$
が存在する。
\item Zariski 被覆 $Y = \bigcup_{i \in I} Y_i$ および
$f^{-1}(Y_i) = \bigcup X_{ij}$ であって、
各射 $X_{ij} \to Y_i$ が滑らかとなるものが存在する。
\end{enumerate}
\end{lemma}

\begin{proof}
省略する。
\end{proof}

\begin{lemma}
\label{spaces-morphisms-lemma-smooth-locally-finite-presentation}
代数空間の滑らかな射は局所有限表示である。
\end{lemma}

\begin{proof}
$X \to Y$ を代数空間の滑らかな射とする。
定義により、Lemma \ref{spaces-morphisms-lemma-local-source-target} のような図式で、
$h$ が滑らかで、垂直矢印 $a$ が全射となるものが存在する。
Morphisms, Lemma
\ref{morphisms-lemma-smooth-locally-finite-presentation} により、
$h$ は局所有限表示である。
したがって定義により $X \to Y$ は局所有限表示である。
\end{proof}

\begin{lemma}
\label{spaces-morphisms-lemma-smooth-locally-finite-type}
代数空間の滑らかな射は局所有限型である。
\end{lemma}

\begin{proof}
Lemmas \ref{spaces-morphisms-lemma-smooth-locally-finite-presentation} および
\ref{spaces-morphisms-lemma-finite-presentation-finite-type} を組み合わせる。
\end{proof}

\begin{lemma}
\label{spaces-morphisms-lemma-smooth-flat}
代数空間の滑らかな射は平坦である。
\end{lemma}

\begin{proof}
$X \to Y$ を代数空間の滑らかな射とする。
定義により、Lemma \ref{spaces-morphisms-lemma-local-source-target} のような図式で、
$h$ が滑らかで、垂直矢印 $a$ が全射となるものが存在する。
Morphisms, Lemma
\ref{morphisms-lemma-smooth-locally-finite-presentation} により、
$h$ は平坦である。したがって定義により $X \to Y$ は平坦である。
\end{proof}

\begin{lemma}
\label{spaces-morphisms-lemma-smooth-syntomic}
代数空間の滑らかな射はシントミックである。
\end{lemma}

\begin{proof}
$X \to Y$ を代数空間の滑らかな射とする。
定義により、Lemma \ref{spaces-morphisms-lemma-local-source-target} のような図式で、
$h$ が滑らかで、垂直矢印 $a$ が全射となるものが存在する。
Morphisms, Lemma \ref{morphisms-lemma-smooth-syntomic} により、
$h$ はシントミックである。
したがって定義により $X \to Y$ はシントミックである。
\end{proof}

\begin{lemma}
\label{spaces-morphisms-lemma-where-smooth}
$S$ をスキームとする。$f : X \to Y$ を $S$ 上の代数空間の射とする。
$f|_U : U \to Y$ が滑らかとなる最大の開部分空間 $U \subset X$ が存在する。
さらに、この開部分空間の形成は、次の射による基底変換と可換である。
\begin{enumerate}
\item 平坦かつ局所有限表示の射。
\item $f$ が局所有限表示である場合の平坦射。
\end{enumerate}
\end{lemma}

\begin{proof}
$U$ の存在は、滑らかであるという性質が始域上 Zariski 局所的
（さらにはエタール局所的）であることから従う
（Lemma \ref{spaces-morphisms-lemma-smooth-local}）。
さらに、この補題により、性質 (1) と (2) をスキームの射の場合へ移せる。
スキームの場合は Morphisms, Lemma
\ref{morphisms-lemma-set-points-where-fibres-smooth} である。
いくつかの細部を省略する。
\end{proof}

\begin{lemma}
\label{spaces-morphisms-lemma-smoothness-dimension-spaces}
$X$ と $Y$ をスキーム $S$ 上の局所 Noether 代数空間とし、
$f : X \to Y$ を滑らかな射とする。
すべての点 $x \in |X|$ とその像 $y \in |Y|$ に対して
$$
\dim_x(X) = \dim_y(Y) + \dim_x(X_y)
$$
が成り立つ。ここで $\dim_x(X_y)$ は Definition
\ref{spaces-morphisms-definition-dimension-fibre} の意味での、$x$ における $f$ の相対次元である。
\end{lemma}

\begin{proof}
一点における代数空間の次元の定義
（Properties of Spaces, Definition
\ref{spaces-properties-definition-dimension-at-point}）により、
これはスキームに対する対応する主張
（Morphisms, Lemma
\ref{morphisms-lemma-smoothness-dimension}）に帰着する。
\end{proof}

\section{非分岐射}
\label{spaces-morphisms-section-unramified}

\noindent
スキームの射の「非分岐」（それぞれ「G-非分岐」）という性質は、
始域と終域の上でエタール局所的である
（Descent, Remark
\ref{descent-remark-list-local-source-target}）。
また、基底変換で保たれ、終域上 fpqc 局所的でもある
（Morphisms, Lemma
\ref{morphisms-lemma-base-change-unramified} および Descent, Lemma
\ref{descent-lemma-descending-property-unramified}）。
したがって上の Lemma \ref{spaces-morphisms-lemma-local-source-target} により、
代数空間の非分岐射（それぞれ G-非分岐射）を次のように定義できる。
この定義は、射が表現可能であるとき、Section
\ref{spaces-morphisms-section-representable} ですでに定義した概念と一致する。

\begin{definition}
\label{spaces-morphisms-definition-unramified}
$S$ をスキームとする。
$f : X \to Y$ を $S$ 上の代数空間の射とする。
\begin{enumerate}
\item Lemma \ref{spaces-morphisms-lemma-local-source-target} の同値な条件が
$\mathcal{P} = \text{非分岐}$ に対して成り立つとき、
$f$ は {\it 非分岐}であるという。
\item $x \in |X|$ とする。$x$ の開近傍 $X' \subset X$ であって、
$f|_{X'} : X' \to Y$ が非分岐となるものが存在するとき、
$f$ は {\it $x$ において非分岐}であるという。
\item Lemma \ref{spaces-morphisms-lemma-local-source-target} の同値な条件が
$\mathcal{P} = \text{G-非分岐}$ に対して成り立つとき、
$f$ は {\it G-非分岐}であるという。
\item $x \in |X|$ とする。$x$ の開近傍 $X' \subset X$ であって、
$f|_{X'} : X' \to Y$ が G-非分岐となるものが存在するとき、
$f$ は {\it $x$ において G-非分岐}であるという。
\end{enumerate}
\end{definition}

\noindent
次の補題により、これ以降は非分岐射についてのみ理論を展開する。
G-非分岐射を用いたいときは、単に「局所有限表示な非分岐射」と呼ぶ。

\begin{lemma}
\label{spaces-morphisms-lemma-unramified-G-unramified}
$S$ をスキームとする。
$f : X \to Y$ を $S$ 上の代数空間の射とする。
このとき、$f$ が G-非分岐であることと、
$f$ が非分岐かつ局所有限表示であることは同値である。
\end{lemma}

\begin{proof}
Lemma \ref{spaces-morphisms-lemma-local-source-target} のような任意の図式を考える。
すると、主張は射 $h$ が G-非分岐であることと、
それが非分岐かつ局所有限表示であることが同値だというだけである。
これは Morphisms, Definition
\ref{morphisms-definition-unramified} から明らかである。
\end{proof}

\begin{lemma}
\label{spaces-morphisms-lemma-composition-unramified}
非分岐射の合成は非分岐である。
\end{lemma}

\begin{proof}
Remark \ref{spaces-morphisms-remark-composition-P} および Morphisms, Lemma
\ref{morphisms-lemma-composition-unramified} を参照。
\end{proof}

\begin{lemma}
\label{spaces-morphisms-lemma-base-change-unramified}
非分岐射の基底変換は非分岐である。
\end{lemma}

\begin{proof}
Remark \ref{spaces-morphisms-remark-base-change-P} および Morphisms, Lemma
\ref{morphisms-lemma-base-change-unramified} を参照。
\end{proof}

\begin{lemma}
\label{spaces-morphisms-lemma-unramified-local}
$S$ をスキームとする。
$f : X \to Y$ を $S$ 上の代数空間の射とする。
次の条件は同値である。
\begin{enumerate}
\item $f$ は非分岐である。
\item すべての $x \in |X|$ に対して、射 $f$ は $x$ において非分岐である。
\item 任意のスキーム $Z$ と任意の射 $Z \to Y$ に対して、
射 $Z \times_Y X \to Z$ は非分岐である。
\item 任意のアフィンスキーム $Z$ と任意の射 $Z \to Y$ に対して、
射 $Z \times_Y X \to Z$ は非分岐である。
\item スキーム $V$ と全射エタール射 $V \to Y$ であって、
$V \times_Y X \to V$ が非分岐射となるものが存在する。
\item スキーム $U$ と全射エタール射 $\varphi : U \to X$ であって、
合成 $f \circ \varphi$ が非分岐となるものが存在する。
\item $U$, $V$ がスキームで垂直矢印がエタールである任意の可換図式
$$
\xymatrix{
U \ar[d] \ar[r] & V \ar[d] \\
X \ar[r] & Y
}
$$
に対して、上の水平矢印は非分岐である。
\item $U$, $V$ がスキーム、垂直矢印がエタール、$U \to X$ が全射で、
上の水平矢印が非分岐であるような可換図式
$$
\xymatrix{
U \ar[d] \ar[r] & V \ar[d] \\
X \ar[r] & Y
}
$$
が存在する。
\item Zariski 被覆 $Y = \bigcup_{i \in I} Y_i$ および
$f^{-1}(Y_i) = \bigcup X_{ij}$ であって、
各射 $X_{ij} \to Y_i$ が非分岐となるものが存在する。
\end{enumerate}
\end{lemma}

\begin{proof}
省略する。
\end{proof}

\begin{lemma}
\label{spaces-morphisms-lemma-unramified-locally-finite-type}
代数空間の非分岐射は局所有限型である。
\end{lemma}

\begin{proof}
Lemma \ref{spaces-morphisms-lemma-local-source-target} のような図式を介して、
これは Morphisms, Lemma
\ref{morphisms-lemma-unramified-locally-finite-type} に移される。
\end{proof}

\begin{lemma}
\label{spaces-morphisms-lemma-unramified-quasi-finite}
$f$ が $x$ において非分岐ならば、$f$ は $x$ において準有限である。
とくに、非分岐射は局所準有限である。
\end{lemma}

\begin{proof}
Lemma \ref{spaces-morphisms-lemma-local-source-target} のような図式を介して、
これは Morphisms, Lemma
\ref{morphisms-lemma-unramified-quasi-finite} に移される。
\end{proof}

\begin{lemma}
\label{spaces-morphisms-lemma-immersion-unramified}
代数空間の埋入は非分岐である。
\end{lemma}

\begin{proof}
$i : X \to Y$ を代数空間の埋入とする。
$V$ をスキームとし、$V \to Y$ を全射エタール射として選ぶ。
すると $V \times_Y X \to V$ はスキームの埋入であり、したがって非分岐である
（Morphisms, Lemmas
\ref{morphisms-lemma-open-immersion-unramified} および
\ref{morphisms-lemma-closed-immersion-unramified}）。
ゆえに定義により $i$ は非分岐である。
\end{proof}

\begin{lemma}
\label{spaces-morphisms-lemma-diagonal-unramified-morphism}
$S$ をスキームとする。
$f : X \to Y$ を $S$ 上の代数空間の射とする。
\begin{enumerate}
\item $f$ が非分岐ならば、対角射
$\Delta_{X/Y} : X \to X \times_Y X$ は開埋入である。
\item $f$ が局所有限型で、$\Delta_{X/Y}$ が開埋入ならば、
$f$ は非分岐である。
\end{enumerate}
\end{lemma}

\begin{proof}
いずれの場合にも、$\Delta_{X/Y}$ は表現可能な単射である
（Lemma \ref{spaces-morphisms-lemma-properties-diagonal}）。
$V$ をスキームとし、$V \to Y$ を全射エタール射として選ぶ。
$U$ をスキームとし、$U \to X \times_Y V$ を全射エタール射として選ぶ。
可換図式
$$
\xymatrix{
U \ar[d] \ar[rr]_-{\Delta_{U/V}} & &
U \times_V U \ar[d] \ar[r] &
V \ar[d]^{\Delta_{V/Y}} \\
X \ar[rr]^-{\Delta_{X/Y}} & &
X \times_Y X \ar[r] &
V \times_Y V
}
$$
を考える。右の正方形はカルテジアンである。
左の垂直矢印は全射エタールである。
右の垂直矢印は、$Y$ 上エタールなスキーム間の射としてエタールである
（Properties of Spaces, Lemma
\ref{spaces-properties-lemma-etale-permanence}）。
したがって中央の垂直矢印もエタールである
（ただし全射であるとは限らない）。

\medskip\noindent
$f$ が非分岐であると仮定する。
すると $U \to V$ は非分岐なので、Morphisms, Lemma
\ref{morphisms-lemma-diagonal-unramified-morphism} により
$\Delta_{U/V}$ は開埋入である。
上の図式の左の正方形を見ると、Properties of Spaces, Lemma
\ref{spaces-properties-lemma-etale-local} により
$\Delta_{X/Y}$ はエタール射であるとわかる。
したがって $\Delta_{X/Y}$ は表現可能なエタール単射であり、
\'Etale Morphisms, Theorem
\ref{etale-theorem-etale-radicial-open} により開埋入である。
（スキームの言葉から関手の言葉への移行については Spaces, Lemma
\ref{spaces-lemma-representable-transformations-property-implication}
も参照。）

\medskip\noindent
$f$ が局所有限型で、$\Delta_{X/Y}$ が開埋入であると仮定する。
このとき $U \to V$ も局所有限型である
（代数空間の局所有限型射の定義による）。
上の表示図式を見ると、$\Delta_{U/V}$ は
$X \times_Y X$ 上エタールなスキーム間の射としてエタールである
（Properties of Spaces, Lemma
\ref{spaces-properties-lemma-etale-permanence}）。
しかし $\Delta_{U/V}$ はスキーム間の射の対角射なので、
いずれにしても埋入である
（Schemes, Lemma
\ref{schemes-lemma-diagonal-immersion}）。
したがって開埋入であり、Morphisms, Lemma
\ref{morphisms-lemma-diagonal-unramified-morphism} により
$U \to V$ は非分岐である。これは定義により $f$ が非分岐であることを意味する。
\end{proof}

\begin{lemma}
\label{spaces-morphisms-lemma-where-unramified}
$S$ をスキームとする。$S$ 上の代数空間の可換図式
$$
\xymatrix{
X \ar[rr]_f \ar[rd]_p & & Y \ar[ld]^q \\
& Z
}
$$
を考える。$X \to Z$ は局所有限型であると仮定する。
このとき、$|U(f)| \subset |X|$ が $f$ の非分岐点の集合となるような
開部分空間 $U(f) \subset X$ が存在する。
さらに、代数空間の任意の射 $Z' \to Z$ に対して、
$f' : X' \to Y'$ が $Z' \to Z$ による $f$ の基底変換ならば、
$U(f')$ は射影 $X' \to X$ による $U(f)$ の逆像である。
\end{lemma}

\begin{proof}
この補題は Morphisms, Lemma
\ref{morphisms-lemma-set-points-where-fibres-unramified} の類似であり、
実際それから導く。Definition \ref{spaces-morphisms-definition-unramified} により、集合
$\{x \in |X| : f \text{ が非分岐となる点 }x\}$ は $X$ で開である。
したがって最後の主張だけを示せばよい。
Lemma \ref{spaces-morphisms-lemma-permanence-finite-type} により、射
$X \to Y$ は局所有限型である。
Lemma \ref{spaces-morphisms-lemma-base-change-finite-type} により、射
$X' \to Y'$ は局所有限型である。

\medskip\noindent
$W$ をスキームとし、$W \to Z$ を全射エタール射として選ぶ。
$V$ をスキームとし、$V \to W \times_Z Y$ を全射エタール射として選ぶ。
$U$ をスキームとし、$U \to V \times_Y X$ を全射エタール射として選ぶ。
最後に $W'$ をスキームとし、
$W' \to W \times_Z Z'$ を全射エタール射として選ぶ。
$V' = W' \times_W V$, $U' = W' \times_W U$ とおく。
すると全射エタール射 $V' \to Y'$ と $U' \to X'$ を得る。
代数空間のエタール射は、付随する位相空間の開写像を誘導する
（Properties of Spaces, Lemma
\ref{spaces-properties-lemma-etale-open}）ことを、以下では断りなく用いる。
これと Lemma \ref{spaces-morphisms-lemma-unramified-local} から、
$U(f)$ は、射 $U \to V$ が非分岐となる $U$ の点の集合 $T$ の
$|X|$ における像である。同様に $U(f')$ は、射 $U' \to V'$ が
非分岐となる $U'$ の点の集合 $T'$ の $|X'|$ における像である。
構成により図式
$$
\xymatrix{
U' \ar[r] \ar[d] & U \ar[d] \\
V' \ar[r] & V
}
$$
は（スキームの圏で）カルテジアンである。
したがって、先に挙げた Morphisms, Lemma
\ref{morphisms-lemma-set-points-where-fibres-unramified} から、
$T'$ は $T$ の逆像であることが従う。
$|U'| \to |X'|$ は全射なので、これから補題が従う。
\end{proof}

\begin{lemma}
\label{spaces-morphisms-lemma-permanence-unramified}
$S$ をスキームとする。
$X \to Y \to Z$ を $S$ 上の代数空間の射とする。
$X \to Z$ が非分岐ならば、$X \to Y$ も非分岐である。
\end{lemma}

\begin{proof}
垂直矢印がエタールかつ全射である可換図式
$$
\xymatrix{
U \ar[d] \ar[r] & V \ar[d] \ar[r] & W \ar[d] \\
X \ar[r] & Y \ar[r] & Z
}
$$
を選ぶ（Spaces, Lemma
\ref{spaces-lemma-lift-morphism-presentations}）。
上段に Morphisms, Lemma
\ref{morphisms-lemma-unramified-permanence} を適用する。
\end{proof}

\section{エタール射}
\label{spaces-morphisms-section-etale}

\noindent
代数空間のエタール射という概念は Properties of Spaces, Definition
\ref{spaces-properties-definition-etale} で定義した。
一点において射がエタールであるとは次の意味である。

\begin{definition}
\label{spaces-morphisms-definition-etale}
$S$ をスキームとする。
$f : X \to Y$ を $S$ 上の代数空間の射とする。
$x \in |X|$ とする。$x$ の開近傍 $X' \subset X$ であって、
$f|_{X'} : X' \to Y$ がエタールとなるものが存在するとき、
$f$ は {\it $x$ においてエタール}であるという。
\end{definition}

\begin{lemma}
\label{spaces-morphisms-lemma-etale-local}
$S$ をスキームとする。
$f : X \to Y$ を $S$ 上の代数空間の射とする。
次の条件は同値である。
\begin{enumerate}
\item $f$ はエタールである。
\item すべての $x \in |X|$ に対して、射 $f$ は $x$ においてエタールである。
\item 任意のスキーム $Z$ と任意の射 $Z \to Y$ に対して、
射 $Z \times_Y X \to Z$ はエタールである。
\item 任意のアフィンスキーム $Z$ と任意の射 $Z \to Y$ に対して、
射 $Z \times_Y X \to Z$ はエタールである。
\item スキーム $V$ と全射エタール射 $V \to Y$ であって、
$V \times_Y X \to V$ がエタール射となるものが存在する。
\item スキーム $U$ と全射エタール射 $\varphi : U \to X$ であって、
合成 $f \circ \varphi$ がエタールとなるものが存在する。
\item $U$, $V$ がスキームで垂直矢印がエタールである任意の可換図式
$$
\xymatrix{
U \ar[d] \ar[r] & V \ar[d] \\
X \ar[r] & Y
}
$$
に対して、上の水平矢印はエタールである。
\item $U$, $V$ がスキーム、垂直矢印がエタール、$U \to X$ が全射で、
上の水平矢印がエタールであるような可換図式
$$
\xymatrix{
U \ar[d] \ar[r] & V \ar[d] \\
X \ar[r] & Y
}
$$
が存在する。
\item Zariski 被覆 $Y = \bigcup Y_i$ および
$f^{-1}(Y_i) = \bigcup X_{ij}$ であって、
各射 $X_{ij} \to Y_i$ がエタールとなるものが存在する。
\end{enumerate}
\end{lemma}

\begin{proof}
Properties of Spaces, Lemmas
\ref{spaces-properties-lemma-etale-local},
\ref{spaces-properties-lemma-base-change-etale}, および
\ref{spaces-properties-lemma-composition-etale} を組み合わせる。
いくつかの細部を省略する。
\end{proof}

\begin{lemma}
\label{spaces-morphisms-lemma-composition-etale}
代数空間の二つのエタール射の合成はエタールである。
\end{lemma}

\begin{proof}
Properties of Spaces, Lemma
\ref{spaces-properties-lemma-composition-etale} の写しである。
\end{proof}

\begin{lemma}
\label{spaces-morphisms-lemma-base-change-etale}
代数空間のエタール射を代数空間の任意の射で基底変換したものは
エタールである。
\end{lemma}

\begin{proof}
Properties of Spaces, Lemma
\ref{spaces-properties-lemma-base-change-etale} の写しである。
\end{proof}

\begin{lemma}
\label{spaces-morphisms-lemma-etale-locally-quasi-finite}
代数空間のエタール射は局所準有限である。
\end{lemma}

\begin{proof}
$X \to Y$ を代数空間のエタール射とする
（Properties of Spaces, Definition
\ref{spaces-properties-definition-etale}）。
Properties of Spaces, Lemma
\ref{spaces-properties-lemma-etale-local} により、これは
Lemma \ref{spaces-morphisms-lemma-local-source-target} のような図式で、
$h$ がエタールで、垂直矢印 $a$ が全射となるものが存在することを意味する。
Morphisms, Lemma
\ref{morphisms-lemma-etale-locally-quasi-finite} により、
$h$ は局所準有限である。
したがって定義により $X \to Y$ は局所準有限である。
\end{proof}

\begin{lemma}
\label{spaces-morphisms-lemma-etale-smooth}
代数空間のエタール射は滑らかである。
\end{lemma}

\begin{proof}
証明は Lemma \ref{spaces-morphisms-lemma-etale-locally-quasi-finite} の証明と同一である。
スキームのエタール射が滑らかであるという事実
（スキームのエタール射の定義による）を用いる。
\end{proof}

\begin{lemma}
\label{spaces-morphisms-lemma-etale-flat}
代数空間のエタール射は平坦である。
\end{lemma}

\begin{proof}
証明は Lemma \ref{spaces-morphisms-lemma-etale-locally-quasi-finite} の証明と同一であり、
Morphisms, Lemma \ref{morphisms-lemma-etale-flat} を用いる。
\end{proof}

\begin{lemma}
\label{spaces-morphisms-lemma-etale-locally-finite-presentation}
\begin{slogan}
エタールならば局所有限表示。
\end{slogan}
代数空間のエタール射は局所有限表示である。
\end{lemma}

\begin{proof}
証明は Lemma \ref{spaces-morphisms-lemma-etale-locally-quasi-finite} の証明と同一であり、
Morphisms, Lemma
\ref{morphisms-lemma-etale-locally-finite-presentation} を用いる。
\end{proof}

\begin{lemma}
\label{spaces-morphisms-lemma-etale-locally-finite-type}
代数空間のエタール射は局所有限型である。
\end{lemma}

\begin{proof}
エタール射は局所有限表示であり、局所有限表示の射は局所有限型である
（Lemmas \ref{spaces-morphisms-lemma-etale-locally-finite-presentation} および
\ref{spaces-morphisms-lemma-finite-presentation-finite-type}）。
\end{proof}

\begin{lemma}
\label{spaces-morphisms-lemma-etale-unramified}
代数空間のエタール射は非分岐である。
\end{lemma}

\begin{proof}
証明は Lemma \ref{spaces-morphisms-lemma-etale-locally-quasi-finite} の証明と同一であり、
Morphisms, Lemma
\ref{morphisms-lemma-etale-smooth-unramified} を用いる。
\end{proof}

\begin{lemma}
\label{spaces-morphisms-lemma-etale-permanence}
$S$ をスキームとする。$X, Y$ を代数空間 $Z$ 上エタールな代数空間とする。
$Z$ 上の任意の射 $X \to Y$ はエタールである。
\end{lemma}

\begin{proof}
Properties of Spaces, Lemma
\ref{spaces-properties-lemma-etale-permanence} の写しである。
\end{proof}

\begin{lemma}
\label{spaces-morphisms-lemma-unramified-flat-lfp-etale}
代数空間の局所有限表示、平坦、かつ非分岐な射はエタールである。
\end{lemma}

\begin{proof}
$X \to Y$ を代数空間の局所有限表示、平坦、かつ非分岐な射とする。
Properties of Spaces, Lemma
\ref{spaces-properties-lemma-etale-local} により、これは
Lemma \ref{spaces-morphisms-lemma-local-source-target} のような図式で、
$h$ が局所有限表示、平坦、かつ非分岐であり、
垂直矢印 $a$ が全射となるものが存在することを意味する。
Morphisms, Lemma
\ref{morphisms-lemma-flat-unramified-etale} により $h$ はエタールである。
したがって定義により $X \to Y$ はエタールである。
\end{proof}

\section{固有射}
\label{spaces-morphisms-section-proper}

\noindent
固有射という概念は代数幾何学で重要な役割を果たす。
代数空間の固有射を次のように定義する。

\begin{definition}
\label{spaces-morphisms-definition-proper}
$S$ をスキームとする。
$f : X \to Y$ を $S$ 上の代数空間の射とする。
$f$ が分離的、有限型、かつ普遍閉であるとき、
$f$ は {\it 固有}であるという。
\end{definition}

\begin{lemma}
\label{spaces-morphisms-lemma-proper-local}
$S$ をスキームとする。$f : X \to Y$ を $S$ 上の代数空間の射とする。
次の条件は同値である。
\begin{enumerate}
\item $f$ は固有である。
\item 任意のスキーム $Z$ と任意の射 $Z \to Y$ に対して、
射影 $Z \times_Y X \to Z$ は固有である。
\item 任意のアフィンスキーム $Z$ と任意の射 $Z \to Y$ に対して、
射影 $Z \times_Y X \to Z$ は固有である。
\item スキーム $V$ と全射エタール射 $V \to Y$ であって、
$V \times_Y X \to V$ が固有となるものが存在する。
\item Zariski 被覆 $Y = \bigcup Y_i$ であって、
各射 $f^{-1}(Y_i) \to Y_i$ が固有となるものが存在する。
\end{enumerate}
\end{lemma}

\begin{proof}
Lemmas \ref{spaces-morphisms-lemma-separated-local},
\ref{spaces-morphisms-lemma-finite-type-local},
\ref{spaces-morphisms-lemma-quasi-compact-local}, および
\ref{spaces-morphisms-lemma-universally-closed-local} を組み合わせる。
\end{proof}

\begin{lemma}
\label{spaces-morphisms-lemma-base-change-proper}
固有射の基底変換は固有である。
\end{lemma}

\begin{proof}
Lemmas \ref{spaces-morphisms-lemma-base-change-separated},
\ref{spaces-morphisms-lemma-base-change-finite-type}, および
\ref{spaces-morphisms-lemma-base-change-universally-closed} を参照。
\end{proof}

\begin{lemma}
\label{spaces-morphisms-lemma-composition-proper}
固有射の合成は固有である。
\end{lemma}

\begin{proof}
Lemmas \ref{spaces-morphisms-lemma-composition-separated},
\ref{spaces-morphisms-lemma-composition-finite-type}, および
\ref{spaces-morphisms-lemma-composition-universally-closed} を参照。
\end{proof}

\begin{lemma}
\label{spaces-morphisms-lemma-closed-immersion-proper}
代数空間の閉埋入は代数空間の固有射である。
\end{lemma}

\begin{proof}
閉埋入は定義により表現可能なので、Spaces, Lemma
\ref{spaces-lemma-representable-transformations-property-implication} と、
スキームの射に対する対応する結果
（Morphisms, Lemma
\ref{morphisms-lemma-closed-immersion-proper}）から従う。
\end{proof}

\begin{lemma}
\label{spaces-morphisms-lemma-universally-closed-permanence}
$S$ をスキームとする。$S$ 上の代数空間の可換図式
$$
\xymatrix{
X \ar[rr] \ar[rd] & &
Y \ar[ld] \\
& B &
}
$$
を考える。
\begin{enumerate}
\item $X \to B$ が普遍閉で、$Y \to B$ が分離的ならば、
射 $X \to Y$ は普遍閉である。
とくに、$|X|$ の $|Y|$ における像は閉である。
\item $X \to B$ が固有で、$Y \to B$ が分離的ならば、
射 $X \to Y$ は固有である。
\end{enumerate}
\end{lemma}

\begin{proof}
$X \to B$ は普遍閉で、$Y \to B$ は分離的であると仮定する。
射を $X \to X \times_B Y \to Y$ と分解する。
第一の射は閉埋入なので（Lemma
\ref{spaces-morphisms-lemma-semi-diagonal}）、普遍閉である。
射影 $X \times_B Y \to Y$ は普遍閉射の基底変換なので普遍閉である
（Lemma \ref{spaces-morphisms-lemma-base-change-universally-closed}）。
したがって $X \to Y$ は普遍閉射の合成として普遍閉である
（Lemma \ref{spaces-morphisms-lemma-composition-universally-closed}）。
これで (1) が証明された。(2) を導くには、(1) と Lemmas
\ref{spaces-morphisms-lemma-compose-after-separated},
\ref{spaces-morphisms-lemma-quasi-compact-permanence}, および
\ref{spaces-morphisms-lemma-permanence-finite-type} を組み合わせる。
\end{proof}

\begin{lemma}
\label{spaces-morphisms-lemma-image-proper-is-proper}
$S$ をスキームとする。$B$ を $S$ 上の代数空間とする。
$f : X \to Y$ を $B$ 上の代数空間の射とする。
$X$ が $B$ 上普遍閉で、$f$ が全射ならば、$Y$ は $B$ 上普遍閉である。
とくに、さらに $Y$ が $B$ 上分離的かつ有限型ならば、
$Y$ は $B$ 上固有である。
\end{lemma}

\begin{proof}
$X$ は普遍閉で、$f$ は全射であると仮定する。
構造射を $p : X \to B$, $q : Y \to B$ と書く。
$B' \to B$ を $S$ 上の代数空間の射とする。
基底変換 $f' : X_{B'} \to Y_{B'}$ は全射であり
（Lemma \ref{spaces-morphisms-lemma-base-change-surjective}）、
基底変換 $p' : X_{B'} \to B'$ は閉写像である。
$T \subset Y_{B'}$ が閉ならば、
$(f')^{-1}(T) \subset X_{B'}$ は閉であるから、
$p'((f')^{-1}(T)) = q'(T)$ も閉である。
したがって $q'$ は閉写像である。
\end{proof}

\begin{lemma}
\label{spaces-morphisms-lemma-scheme-theoretic-image-is-proper}
$S$ をスキームとする。
$$
\xymatrix{
X \ar[rr]_h \ar[rd]_f & & Y \ar[ld]^g \\
& B
}
$$
を $S$ 上の代数空間の射の可換図式とする。次を仮定する。
\begin{enumerate}
\item $X \to B$ は固有射である。
\item $Y \to B$ は分離的かつ局所有限型である。
\end{enumerate}
このとき $h$ のスキーム論的像 $Z \subset Y$ は $B$ 上固有であり、
$X \to Z$ は全射である。
\end{lemma}

\begin{proof}
$h$ のスキーム論的像は Section
\ref{spaces-morphisms-section-scheme-theoretic-image} で構成した。
$h$ は準コンパクトである
（Lemma \ref{spaces-morphisms-lemma-quasi-compact-quasi-separated-permanence}）から、
$|h|(|X|) \subset |Z|$ は稠密である
（Lemma \ref{spaces-morphisms-lemma-quasi-compact-scheme-theoretic-image}）。
一方、$|h|(|X|)$ は $|Y|$ で閉である
（Lemma \ref{spaces-morphisms-lemma-universally-closed-permanence}）から、
$X \to Z$ は全射である。
したがって $Z \to B$ は固有である
（Lemma \ref{spaces-morphisms-lemma-image-proper-is-proper}）。
\end{proof}

\begin{lemma}
\label{spaces-morphisms-lemma-separated-diagonal-proper}
$S$ をスキームとする。
$f : X \to Y$ を $S$ 上の代数空間の射とする。
次の条件は同値である。
\begin{enumerate}
\item $f$ は分離的である。
\item $\Delta_{X/Y} : X \to X \times_Y X$ は普遍閉である。
\item $\Delta_{X/Y} : X \to X \times_Y X$ は固有である。
\end{enumerate}
\end{lemma}

\begin{proof}
(1) $\Rightarrow$ (3) は Lemma
\ref{spaces-morphisms-lemma-closed-immersion-proper} から従う。
証明の残りでは Spaces, Lemma
\ref{spaces-lemma-representable-transformations-property-implication}
を断りなく用いる。
$\Delta_{X/Y}$ は表現可能な単射で、局所有限型であることを思い出そう
（Lemma \ref{spaces-morphisms-lemma-properties-diagonal}）。
スキームの射について固有 $\Rightarrow$ 普遍閉なので、(3) から (2) が従う。
$\Delta_{X/Y}$ が普遍閉ならば、\'Etale Morphisms, Lemma
\ref{etale-lemma-characterize-closed-immersion} により閉埋入である。
したがって (2) $\Rightarrow$ (1) であり、証明が完了する。
\end{proof}

\section{付値判定法}
\label{spaces-morphisms-section-valuative}

\noindent
この節では、代数空間の射に対する付値判定法の基礎を導入する。
さらに進んだ結果への参照を次に挙げる。
\begin{enumerate}
\item 普遍閉性の付値判定法は Section
\ref{spaces-morphisms-section-valuative-criterion-universally-closed} にある。
\item 分離性の付値判定法は Section
\ref{spaces-morphisms-section-valuative-separatedness} にある。
\item 固有性の付値判定法は Section
\ref{spaces-morphisms-section-valuative-criterion-properness} にある。
\item 追加の逆向きの主張は Decent Spaces, Section
\ref{decent-spaces-section-valuative-criterion-universally-closed}
および Decent Spaces, Lemma
\ref{decent-spaces-lemma-re-characterize-universally-closed} にある。
\item Noether の場合、離散付値環について判定法を確認すれば十分である。
これは Cohomology of Spaces, Section
\ref{spaces-cohomology-section-Noetherian-valuative-criterion} および
Limits of Spaces, Section
\ref{spaces-limits-section-Noetherian-valuative-criterion} で示される。
\item Noether の場合の付値判定法の精密化は Limits of Spaces, Section
\ref{spaces-limits-section-refined-valuative-criteria} にある。
\end{enumerate}
まず定義を形式的に述べ、次にスキームの射の場合との違いを論じる。

\begin{definition}
\label{spaces-morphisms-definition-valuative-criterion}
$S$ をスキームとする。
$f : X \to Y$ を $S$ 上の代数空間の射とする。
$A$ が分数体 $K$ をもつ付値環である任意の可換な実線図式
$$
\xymatrix{
\Spec(K) \ar[r] \ar[d] & X \ar[d] \\
\Spec(A) \ar[r] \ar@{-->}[ru] & Y
}
$$
が与えられたとき、点線矢印が高々一つ存在する
（存在は要求しない）ならば、
$f$ は {\it 付値判定法の一意性部分を満たす}という。
上のような任意の実線図式が与えられたとき、
体の拡大 $K'/K$、$A$ を支配する付値環 $A' \subset K'$、
および図式
$$
\xymatrix{
\Spec(K') \ar[r] \ar[d] & \Spec(K) \ar[r] & X \ar[d] \\
\Spec(A') \ar[r] \ar[rru] & \Spec(A) \ar[r] & Y
}
$$
を可換にする射 $\Spec(A') \to X$ が存在するならば、
$f$ は {\it 付値判定法の存在部分を満たす}という。
$f$ が存在部分と一意性部分の両方を満たすならば、
$f$ は {\it 付値判定法を満たす}という。
\end{definition}

\noindent
付値環の分数体を拡大する必要がありうるので、
付値判定法の存在部分の定式化は、代数空間の射に対して少し異なる。
実際には、この違いが役割を果たすことはほとんどない。
\begin{enumerate}
\item 付値判定法の一意性部分を確認するときは分数体の拡大を一切用いない。
したがって、これはスキームの場合とまったく同じである。
\item 一般には体の拡大を許す必要がある
（Example \ref{spaces-morphisms-example-finite-separable-needed}）。
\item 代数空間の射については、付値判定法の存在部分で
有限分離拡大 $K'/K$ を取れば常に十分である
（Lemma \ref{spaces-morphisms-lemma-finite-separable-enough}）。
\item $f : X \to Y$ が代数空間の分離射ならば、
付値判定法の存在部分を確認するとき、常に $K = K'$ と取れる
（Lemma \ref{spaces-morphisms-lemma-usual-enough}）。
\item 準コンパクトかつ準分離な射 $f : X \to Y$ に対しては、
「$f$ は分離的かつ普遍閉」と「$f$ は通常の付値判定法を満たす」が同値である
（Lemma
\ref{spaces-morphisms-lemma-characterize-separated-and-universally-closed}）。
固有性の付値判定法は通常のものである
（Lemma \ref{spaces-morphisms-lemma-characterize-proper}）。
\end{enumerate}
理論の第一段階として、代数空間の射が表現可能である場合には、
この判定法がスキームの射に対する定式化と同一であることを示す。

\begin{lemma}
\label{spaces-morphisms-lemma-valuative-criterion-representable}
$S$ をスキームとする。
$f : X \to Y$ を $S$ 上の代数空間の射とし、$f$ は表現可能であると仮定する。
次の条件は同値である。
\begin{enumerate}
\item $f$ は Definition \ref{spaces-morphisms-definition-valuative-criterion} の意味で
付値判定法の存在部分を満たす。
\item $A$ が分数体 $K$ をもつ付値環である任意の可換な実線図式
$$
\xymatrix{
\Spec(K) \ar[r] \ar[d] & X \ar[d] \\
\Spec(A) \ar[r] \ar@{-->}[ru] & Y
}
$$
に対して点線矢印が存在する。すなわち、$f$ は Schemes, Definition
\ref{schemes-definition-valuative-criterion} の意味で
付値判定法の存在部分を満たす。
\end{enumerate}
\end{lemma}

\begin{proof}
Definition \ref{spaces-morphisms-definition-valuative-criterion} のような可換図式
$$
\xymatrix{
\Spec(K') \ar[r] \ar[d] & \Spec(K) \ar[r] & X \ar[d] \\
\Spec(A') \ar[r] \ar[rru]^\varphi & \Spec(A) \ar[r] & Y
}
$$
が与えられたとき、同じ図式に入る射 $\Spec(A) \to X$ を見つけられることを
示せば十分である。
$X_A = \Spec(A) \times_Y Y$ とおく。
$f$ は表現可能なので、$X_A$ はスキームである。
射 $\varphi$ は射 $\varphi' : \Spec(A') \to X_A$ を与える。
$x \in X_A$ を $\varphi' : \Spec(A') \to X_A$ の閉点の像とする。
すると環の可換図式
$$
\xymatrix{
K' & K \ar[l] & \mathcal{O}_{X_A, x} \ar[l] \ar[lld] \\
A' \ar[u] & A \ar[l] & A \ar[l] \ar[u]
}
$$
を得る。$A$ は付値環であり、$A'$ は $A$ を支配するので、
$K \cap A' = A$ である。
したがって環写像 $\mathcal{O}_{X_A, x} \to K$ の像は $A$ に含まれる。
ゆえに望みどおり射 $\Spec(A) \to X_A$ を得る
（Schemes, Section
\ref{schemes-section-points}）。
\end{proof}

\begin{lemma}
\label{spaces-morphisms-lemma-finite-separable-enough}
$S$ をスキームとする。
$f : X \to Y$ を $S$ 上の代数空間の射とする。
次の条件は同値である。
\begin{enumerate}
\item $f$ は Definition \ref{spaces-morphisms-definition-valuative-criterion} の意味で
付値判定法の存在部分を満たす。
\item Definition \ref{spaces-morphisms-definition-valuative-criterion} の付値判定法の
存在部分で、拡大 $K'/K$ が有限分離であることを要求したものを $f$ が満たす。
\end{enumerate}
\end{lemma}

\begin{proof}
(1) が (2) を含意することを示さなければならない。
$K \subset K'$ が任意である Definition
\ref{spaces-morphisms-definition-valuative-criterion} のような図式
$$
\xymatrix{
\Spec(K') \ar[r] \ar[d] & \Spec(K) \ar[r] & X \ar[d] \\
\Spec(A') \ar[r] \ar[rru] & \Spec(A) \ar[r] & Y
}
$$
が与えられたとする。
$U$ をスキームとし、$U \to X$ を全射エタール射として選ぶ。このとき
$$
\Spec(A') \times_X U \longrightarrow \Spec(A')
$$
は全射エタールである。
$p$ を $\Spec(A') \times_X U$ の点で、
$\Spec(A')$ の閉点に写るものとする。
$p$ の一般化 $p' \leadsto p$ で、
$\Spec(A')$ の一般点に写るものを選ぶ。
そのような一般化が存在するのは、スキームの平坦射に沿って一般化が持ち上がる
からである（Morphisms, Lemma
\ref{morphisms-lemma-generalizations-lift-flat}）。
すると $p'$ はスキーム $\Spec(K') \times_X U$ の点に対応する。ここで
$$
\Spec(K') \times_X U
=
\Spec(K') \times_{\Spec(K)} (\Spec(K) \times_X U)
$$
である。したがって $p'$ は点 $q' \in \Spec(K) \times_X U$ に写り、
その剰余体は $K$ の有限分離拡大である。
最後に、$p' \leadsto p$ はスキーム $U$ 上の特殊化
$u' \leadsto u$ に写る。これらの記法により、環の図式
$$
\xymatrix{
\kappa(p') & & \kappa(q') \ar[ll] & \kappa(u') \ar[l] \\
& \mathcal{O}_{\Spec(A') \times_X U, p} \ar[lu] & &
\mathcal{O}_{U, u} \ar[ll] \ar[u] \\
K' \ar[uu] & A' \ar[l] \ar[u] & A \ar[l] \ar'[u][uu]
}
$$
を得る。これは、$A$ と $\mathcal{O}_{U, u}$ の像で生成される環
$B \subset \kappa(q')$ が $\kappa(p')$ の部分環へ写り、
その像が $\mathcal{O}_{\Spec(A') \times_X U, p} \to \kappa(p')$
の像 $B'$ に含まれることを意味する。
$B'$ は局所環であることに注意する。
$\mathfrak m \subset B$ を極大イデアルとする。
構成により、$A \cap \mathfrak m$
（それぞれ $\mathcal{O}_{U, u} \cap \mathfrak m$、
それぞれ $A' \cap \mathfrak m$）は、$A$
（それぞれ $\mathcal{O}_{U, u}$、それぞれ $A'$）の極大イデアルである。
$\mathfrak q = B \cap \mathfrak m$ とおく。
これは $A \cap \mathfrak q$ が $A$ の極大イデアルとなる素イデアルである。
したがって $B_{\mathfrak q} \subset \kappa(q')$ は $A$ を支配する局所環である。
Algebra, Lemma \ref{algebra-lemma-dominate} により、
分数体が $\kappa(q')$ であり $B_{\mathfrak q}$ を支配する付値環
$A_1 \subset \kappa(q')$ を見つけられる。
（局所的な）環写像 $\mathcal{O}_{U, u} \to A_1$ は射
$\Spec(A_1) \to U \to X$ を与え、図式
$$
\xymatrix{
\Spec(\kappa(q')) \ar[r] \ar[d] & \Spec(K) \ar[r] & X \ar[d] \\
\Spec(A_1) \ar[r] \ar[rru] & \Spec(A) \ar[r] & Y
}
$$
を可換にする。$A_1$ の分数体は $\kappa(q')$ であり、
構成により $\kappa(q')/K$ は有限分離なので、補題が証明された。
\end{proof}

\begin{lemma}
\label{spaces-morphisms-lemma-push-down-solution}
$S$ をスキームとする。$f : X \to Y$ を $S$ 上の代数空間の分離射とする。
$K \subset K'$ が任意である Definition
\ref{spaces-morphisms-definition-valuative-criterion} のような図式
$$
\xymatrix{
\Spec(K') \ar[r] \ar[d] & \Spec(K) \ar[r] & X \ar[d] \\
\Spec(A') \ar[r] \ar[rru] & \Spec(A) \ar[r] \ar@{-->}[ru] & Y
}
$$
が与えられたとする。
このとき点線矢印が存在して図式を可換にする。
\end{lemma}

\begin{proof}
図式に入る射 $\Spec(A) \to X$ を見つけられることを示さなければならない。

\medskip\noindent
$X$ の基底変換 $X_A = \Spec(A) \times_Y X$ を考える。
すると $X_A \to \Spec(A)$ は代数空間の分離射である
（Lemma \ref{spaces-morphisms-lemma-base-change-separated}）。
上の図式のすべての射を基底変換すると
$$
\xymatrix{
\Spec(K') \ar[r] \ar[d] & \Spec(K) \ar[r] & X_A \ar[d] \\
\Spec(A') \ar[r] \ar[rru] & \Spec(A) \ar@{=}[r] & \Spec(A)
}
$$
を得る。したがって $X$ を $X_A$ で置き換え、$Y = \Spec(A)$ と仮定し、
上のような図式があるとしてよい。
さらに $X$ を、$|\Spec(A')| \to |X|$ の像を含む
準コンパクトな開部分空間で置き換える。

\medskip\noindent
Lemma \ref{spaces-morphisms-lemma-quasi-compact-permanence} により、射
$\Spec(A') \to X$ は準コンパクトである。
$Z \subset X$ を $\Spec(A') \to X$ のスキーム論的像とする。
このとき $Z$ は $A$ 上の被約
（Lemma \ref{spaces-morphisms-lemma-scheme-theoretic-image-reduced}）、
準コンパクト（$X$ の閉部分空間として）、分離的
（$X$ の閉部分空間として）な代数空間である。
スキームの平坦射 $\Spec(K) \to \Spec(A)$ による射
$\Spec(A') \to X$ の基底変換
$$
\Spec(K') = \Spec(A') \times_{\Spec(A)} \Spec(K) \to
X \times_{\Spec(A)} \Spec(K) = X_K
$$
を考える。Lemma
\ref{spaces-morphisms-lemma-flat-base-change-scheme-theoretic-image} により、
この射のスキーム論的像は $Z$ の基底変換 $Z_K$ である。
一方、仮定（すなわち上の可換図式）により、この射は射
$\Spec(K) \to Z_K$ を経由し、この射は構造射
$Z_K \to \Spec(K)$ の切断である。
$Z_K$ は分離的なので、この切断は閉埋入である
（Lemma \ref{spaces-morphisms-lemma-section-immersion}）。
したがって $Z_K = \Spec(K)$ である。

\medskip\noindent
$V$ をアフィンスキームとし、$V \to Z$ を全射エタール射として選ぶ
（Properties of Spaces, Lemma
\ref{spaces-properties-lemma-quasi-compact-affine-cover}）。
$V = \Spec(B)$ と書く。
$Z$ は分離的なので、$V \times_Z \Spec(A') = \Spec(C)$ はアフィンである。
Lemma \ref{spaces-morphisms-lemma-quasi-compact-scheme-theoretic-image} により、
$V$ は $V \times_Z \Spec(A') \to V$ のスキーム論的像であるから、
$B \to C$ は単射である。
一方、$A' \to C$ は $V \to Z$ の基底変換に対応するのでエタールである。
$A'$ は捩れのない $A$-加群であるから、$A' \to C$ の平坦性により
$C$ は捩れのない $A$-加群であり、したがって
$B$ は捩れのない $A$-加群である。
$A$-加群として捩れがないことは平坦であることと同値である
（More on Algebra, Lemma
\ref{more-algebra-lemma-valuation-ring-torsion-free-flat}）。
次に
$$
V \times_Z V = \Spec(B')
$$
と書く。$V \to Z$ はエタールなので、二つの環写像 $B \to B'$ は
エタールである。標準全射 $B \otimes_A B \to B'$ は、
$A$ 上 $K$ とテンソルを取ると同型になる。これは
$Z_K = \Spec(K)$ だからである。
しかし上の考察により $B \otimes_A B$ は捩れのない $A$-加群である。
したがって $B' = B \otimes_A B$ である。
環写像 $A \to B$ を忠実平坦な環写像 $A \to B$ で基底変換したものは
エタールであることが従う
（$X \to \Spec(A)$ が全射なので
$\Spec(B) \to \Spec(A)$ は全射であることに注意する）。
したがって $A \to B$ はエタールである
（Descent, Lemma
\ref{descent-lemma-descending-property-etale}）。
言い換えれば、$V \to X$ はエタールである。
$V \times_Z V = V \times_{\Spec(A)} V$ なので、
代数空間として $Z = \Spec(A)$ である
（たとえば Spaces, Lemma
\ref{spaces-lemma-space-presentation}）。
これで証明が完了する。
\end{proof}

\begin{lemma}
\label{spaces-morphisms-lemma-usual-enough}
$S$ をスキームとする。$f : X \to Y$ を $S$ 上の代数空間の分離射とする。
次の条件は同値である。
\begin{enumerate}
\item $f$ は Definition \ref{spaces-morphisms-definition-valuative-criterion} の意味で
付値判定法の存在部分を満たす。
\item $A$ が分数体 $K$ をもつ付値環である任意の可換な実線図式
$$
\xymatrix{
\Spec(K) \ar[r] \ar[d] & X \ar[d] \\
\Spec(A) \ar[r] \ar@{-->}[ru] & Y
}
$$
に対して点線矢印が存在する。すなわち、$f$ は Schemes, Definition
\ref{schemes-definition-valuative-criterion} の意味で
付値判定法の存在部分を満たす。
\end{enumerate}
\end{lemma}

\begin{proof}
(1) が (2) を含意することを示さなければならない。
(2) のような可換図式
$$
\xymatrix{
\Spec(K) \ar[r] \ar[d] & X \ar[d] \\
\Spec(A) \ar[r] & Y
}
$$
が与えられたとする。(1) により、$K \subset K'$ が任意である
Definition \ref{spaces-morphisms-definition-valuative-criterion} のような可換図式
$$
\xymatrix{
\Spec(K') \ar[r] \ar[d] & \Spec(K) \ar[r] & X \ar[d] \\
\Spec(A') \ar[r] \ar[rru] & \Spec(A) \ar[r] & Y
}
$$
が存在する。Lemma \ref{spaces-morphisms-lemma-push-down-solution} により、
図式に入る射 $\Spec(A) \to X$ を見つけられる。すなわち (2) が成り立つ。
\end{proof}

\begin{example}
\label{spaces-morphisms-example-finite-separable-needed}
Spaces, Example
\ref{spaces-example-non-representable-descent} で構成した代数空間 $X$ を考える。
これは原点を二重化したアフィン直線の Galois ねじれであった。
この Galois ねじれは体の次数二の Galois 拡大 $k'/k$ に関するものである。
したがって、準コンパクトな射
$$
\pi : X \longrightarrow S = \mathbf{A}^1_k
$$
を備える。$\pi$ は普遍閉であると主張する。
実際、$\Spec(k') \to \Spec(k)$ による基底変換後、射 $\pi$ は射
$$
\text{原点を二重化したアフィン直線}
\longrightarrow
\text{アフィン直線}
$$
と同一視され、これは普遍閉である（細部を省略する）。
射 $\Spec(k') \to \Spec(k)$ は普遍閉かつ全射なので、
図式追跡により $\pi$ は普遍閉である。
一方、図式
$$
\xymatrix{
\Spec(k((x))) \ar[r] \ar[d] & X \ar[d]^\pi \\
\Spec(k[[x]]) \ar[r] \ar@{..>}[ru] & \mathbf{A}^1_k
}
$$
を考える。$0 \in \mathbf{A}^1_k$ 上の $X$ の唯一の点は単射
$\Spec(k') \to X$ に対応するから、点線矢印が存在しえないことは明らかである。
これは一般に有限分離体拡大が必要であることを示す。
\end{example}

\begin{lemma}
\label{spaces-morphisms-lemma-base-change-valuative-criteria}
付値判定法の存在部分（それぞれ一意性部分）を満たす代数空間の射を、
代数空間の任意の射で基底変換したものは、
付値判定法の存在部分（それぞれ一意性部分）を満たす。
\end{lemma}

\begin{proof}
$f : X \to Y$ をスキーム $S$ 上の代数空間の射とする。
$Z \to Y$ を $S$ 上の代数空間の任意の射とする。
次の形の可換な実線図式を考える。
$$
\xymatrix{
\Spec(K) \ar[r] \ar[d] & Z \times_Y X \ar[r] \ar[d] & X \ar[d] \\
\Spec(A) \ar[r] \ar@{-->}[ru] \ar@{-->}[rru] & Z \ar[r] & Y
}
$$
図式を可換にする北西向きの点線矢印の集合は、
図式を可換にする西北西向きの点線矢印の集合と一対一に対応する。
これで「一意性」の場合の補題が証明された。
存在部分について、$f$ は付値判定法の存在部分を満たすと仮定する。
上のような可換な実線図式が与えられたとき、仮定により、
体の拡大 $K'/K$、$A$ を支配する付値環 $A' \subset K'$、
および次の可換図式に入る射 $\Spec(A') \to X$ が存在する。
$$
\xymatrix{
\Spec(K') \ar[r] \ar[d] &
\Spec(K) \ar[r] & Z \times_Y X \ar[r] & X \ar[d] \\
\Spec(A') \ar[r] \ar[rrru] & \Spec(A) \ar[r] & Z \ar[r] & Y
}
$$
上の考察により、斜めの矢印は望みどおり射
$\Spec(A') \to Z \times_Y X$ に対応する。
\end{proof}

\begin{lemma}
\label{spaces-morphisms-lemma-composition-valuative-criteria}
付値判定法（の存在部分、それぞれ一意性部分）を満たす
代数空間の二つの射の合成は、
付値判定法（の存在部分、それぞれ一意性部分）を満たす。
\end{lemma}

\begin{proof}
$f : X \to Y$, $g : Y \to Z$ をスキーム $S$ 上の代数空間の射とする。
次の形の可換な実線図式を考える。
$$
\xymatrix{
\Spec(K) \ar[dd] \ar[r] & X \ar[d]^f \\
& Y \ar[d]^g \\
\Spec(A) \ar[r] \ar@{-->}[ru] \ar@{-->}[ruu] & Z
}
$$
$g$ について一意性部分が成り立つならば、
図式を可換にする北西向きの点線矢印は高々一つである。
さらに $f$ について一意性部分が成り立つならば、
図式を可換にする北北西向きの点線矢印は高々一つである。
存在の場合の証明は、次の図式を考えれば得られる。
$$
\xymatrix{
\Spec(K'') \ar[r] \ar[dd] &
\Spec(K') \ar[r] &
\Spec(K) \ar[r] &
X \ar[d]^f \\
& & & Y \ar[d]^g \\
\Spec(A'') \ar[r] \ar[rrruu] &
\Spec(A') \ar[r] \ar[rru] &
\Spec(A) \ar[r] &
Z
}
$$
すなわち、$g$ に対する存在部分は拡大 $K'$、付値環 $A'$、
および矢印 $\Spec(A') \to Y$ を与える。
次いで $f$ に対する存在部分が拡大 $K''$、付値環 $A''$、
および矢印 $\Spec(A'') \to X$ を与える。
\end{proof}

\section{普遍閉性の付値判定法}
\label{spaces-morphisms-section-valuative-criterion-universally-closed}

\noindent
付値判定法の存在部分は、準コンパクト射に対して普遍閉性を含意する
（Lemma
\ref{spaces-morphisms-lemma-quasi-compact-existence-universally-closed}）。
スキームの場合、これは「必要十分」な主張であるが、
代数空間の射については正しくない。
Example \ref{spaces-morphisms-example-strange-universally-closed} は
$\mathbf{A}^1_k/\mathbf{Z} \to \Spec(k)$ が普遍閉であることを示すが、
付値判定法の存在部分が成り立たないことは容易にわかる。
この話題は Decent Spaces, Section
\ref{decent-spaces-section-valuative-criterion-universally-closed}
で再考し、射の始域が decent space ならば逆が成り立つことを示す
（相対版については Decent Spaces, Lemma
\ref{decent-spaces-lemma-re-characterize-universally-closed} も参照）。

\begin{lemma}
\label{spaces-morphisms-lemma-quasi-compact-existence-universally-closed}
$S$ をスキームとする。
$f : X \to Y$ を $S$ 上の代数空間の射とする。次を仮定する。
\begin{enumerate}
\item $f$ は準コンパクトである。
\item $f$ は付値判定法の存在部分を満たす。
\end{enumerate}
このとき $f$ は普遍閉である。
\end{lemma}

\begin{proof}
Lemmas \ref{spaces-morphisms-lemma-base-change-quasi-compact} および
\ref{spaces-morphisms-lemma-base-change-valuative-criteria} により、
性質 (1) と (2) は任意の基底変換で保たれる。
Lemma \ref{spaces-morphisms-lemma-universally-closed-local} により、
$T$ が $S$ 上のアフィンスキームで $Y$ へ写るとき、
$|T \times_Y X| \to |T|$ が閉写像であることだけを示せばよい。
したがって、次を証明すれば十分である。
$Y$ がアフィンスキームで、$f : X \to Y$ が準コンパクトかつ
付値判定法の存在部分を満たすならば、
$f : |X| \to |Y|$ は閉写像である。
この状況で $X$ は準コンパクトな代数空間である。
Properties of Spaces, Lemma
\ref{spaces-properties-lemma-quasi-compact-affine-cover} により、
アフィンスキーム $U$ と全射エタール射 $\varphi : U \to X$ が存在する。
$T \subset |X|$ を閉集合とする。
逆像 $\varphi^{-1}(T) \subset U$ は閉であるから、
あるアフィン閉部分スキーム $Z \subset U$ の点の集合である。
したがって Algebra, Lemma
\ref{algebra-lemma-image-stable-specialization-closed} により、
$f(T) = f(\varphi(|Z|)) \subset |Y|$ が特殊化で閉じていれば、それは閉である。

\medskip\noindent
$y' \leadsto y$ を $Y$ の特殊化で $y' \in f(T)$ となるものとする。
$f$ により $y'$ に写る点 $x' \in T \subset |X|$ を選ぶ。
ある体 $K$ に対する射 $\Spec(K) \to X$ で $x'$ を表せる。
したがって図式
$$
\xymatrix{
\Spec(K) \ar[r]_-{x'} \ar[d] & X \ar[d]^f \\
\Spec(\mathcal{O}_{Y, y}) \ar[r] & Y,
}
$$
を得る。左の垂直写像の存在については Schemes, Section
\ref{schemes-section-points} を参照。
環写像 $\mathcal{O}_{Y, y} \to K$ の像を支配する付値環
$A \subset K$ を選ぶ
（$y' \not = y$ なので像は体ではなく局所環であり、これは可能である。
Algebra, Lemma \ref{algebra-lemma-dominate} を参照）。
仮定により、体の拡大 $K'/K$、$A$ を支配する付値環
$A' \subset K'$、および可換図式に入る射
$\Spec(A') \to X$ が存在する。
$A'$ は $A$ を支配し、$A$ は $\mathcal{O}_{Y, y}$ を支配するので、
$\Spec(A')$ の閉点は点 $x \in X$ に写る。
この点は $f(x) = y$ を満たし、$x'$ の特殊化である。
$T$ は閉なので $x \in T$、したがって望みどおり $y \in f(T)$ である。
\end{proof}

\noindent
次の補題は Decent Spaces, Lemma
\ref{decent-spaces-lemma-re-characterize-universally-closed} で一般化される。

\begin{lemma}
\label{spaces-morphisms-lemma-characterize-universally-closed-quasi-separated}
$S$ をスキームとする。
$f : X \to Y$ を $S$ 上の代数空間の射とする。
\begin{enumerate}
\item $f$ が準分離かつ普遍閉ならば、
$f$ は付値判定法の存在部分を満たす。
\item $f$ が準コンパクトかつ準分離ならば、
$f$ が普遍閉であることと、付値判定法の存在部分が成り立つことは同値である。
\end{enumerate}
\end{lemma}

\begin{proof}
(1) が成り立つならば、Lemma
\ref{spaces-morphisms-lemma-quasi-compact-existence-universally-closed} と組み合わせて
(2) を得る。
$f$ は準分離かつ普遍閉であると仮定する。
Definition \ref{spaces-morphisms-definition-valuative-criterion} のような図式
$$
\xymatrix{
\Spec(K) \ar[r] \ar[d] & X \ar[d] \\
\Spec(A) \ar[r] & Y
}
$$
が与えられたとする。
形式的な議論により、望む図式
$$
\xymatrix{
\Spec(K') \ar[r] \ar[d] & \Spec(K) \ar[r] & X \ar[d] \\
\Spec(A') \ar[r] \ar[rru] & \Spec(A) \ar[r] & Y
}
$$
の存在は、射 $X_A \to \Spec(A)$ の場合の存在から従う。
準分離かつ普遍閉であることは基底変換で保たれるので、
次の段落の結果から補題が従う。

\medskip\noindent
$A$ が分数体 $K$ をもつ付値環であるような実線図式
$$
\xymatrix{
\Spec(K) \ar[r]_-x \ar[d] & X \ar[d]^f \\
\Spec(A) \ar@{=}[r] \ar@{..>}[ru] & \Spec(A)
}
$$
を考える。Lemma \ref{spaces-morphisms-lemma-quasi-compact-permanence} と
$f$ が準分離であることから、射 $x$ は準コンパクトである。
$f$ は普遍閉なので、とくに
$|f|(\overline{\{x\}})$ は $\Spec(A)$ で閉である。
この像は $\Spec(A)$ の一般点を含むから、
$x$ の閉包の中に、$\Spec(A)$ の閉点に写る点 $x' \in |X|$ が存在する。
Lemma \ref{spaces-morphisms-lemma-reach-points-scheme-theoretic-image} により、
可換図式
$$
\xymatrix{
\Spec(K') \ar[r] \ar[d] & \Spec(K) \ar[d] \\
\Spec(A') \ar[r] & X
}
$$
であって、$\Spec(A')$ の閉点が $x' \in |X|$ に写るものを見つけられる。
したがって $\Spec(A') \to \Spec(A)$ は閉点を閉点へ写す。
すなわち $A'$ は $A$ を支配する。これで証明が完了する。
\end{proof}

\begin{lemma}
\label{spaces-morphisms-lemma-characterize-universally-closed-separated}
$S$ をスキームとする。$f : X \to Y$ を $S$ 上の代数空間の射とする。
$f$ は準コンパクトかつ分離的であると仮定する。
次の条件は同値である。
\begin{enumerate}
\item $f$ は普遍閉である。
\item Definition \ref{spaces-morphisms-definition-valuative-criterion} の意味で
付値判定法の存在部分が成り立つ。
\item $A$ が分数体 $K$ をもつ付値環である任意の可換な実線図式
$$
\xymatrix{
\Spec(K) \ar[r] \ar[d] & X \ar[d] \\
\Spec(A) \ar[r] \ar@{-->}[ru] & Y
}
$$
に対して点線矢印が存在する。すなわち、$f$ は Schemes, Definition
\ref{schemes-definition-valuative-criterion} の意味で
付値判定法の存在部分を満たす。
\end{enumerate}
\end{lemma}

\begin{proof}
$f$ は分離的なので、Lemma \ref{spaces-morphisms-lemma-usual-enough} により
(2) と (3) は同値である。
(3) と (1) の同値性は Lemma
\ref{spaces-morphisms-lemma-characterize-universally-closed-quasi-separated} から従う。
\end{proof}

\begin{lemma}
\label{spaces-morphisms-lemma-lift-valuation-ring-through-flat-morphism}
$S$ をスキームとする。
$f : X \to Y$ を $S$ 上の代数空間の平坦射とする。
$A$ が付値環であるような射 $\Spec(A) \to Y$ を取る。
$\Spec(A)$ の閉点が、$|X| \to |Y|$ の像に属する $|Y|$ の点へ写るならば、
可換図式
$$
\xymatrix{
\Spec(A') \ar[r] \ar[d] & X \ar[d] \\
\Spec(A) \ar[r] & Y
}
$$
で、$A \to A'$ が付値環の拡大となるものが存在する
（More on Algebra, Definition
\ref{more-algebra-definition-extension-valuation-rings}）。
\end{lemma}

\begin{proof}
基底変換 $X_A \to \Spec(A)$ は平坦であり
（Lemma \ref{spaces-morphisms-lemma-base-change-flat}）、
$\Spec(A)$ の閉点は $|X_A| \to |\Spec(A)|$ の像に属する
（Properties of Spaces, Lemma
\ref{spaces-properties-lemma-points-cartesian}）。
したがって $Y = \Spec(A)$ としてよい。
$U$ をスキームとし、$U \to X$ を全射エタール射とする。
$u \in U$ を $\Spec(A)$ の閉点に写る点とする。
平坦な局所環写像 $A \to B = \mathcal{O}_{U, u}$ を考える。
Algebra, Lemma \ref{algebra-lemma-ff-rings} により、
素イデアル $\mathfrak q \subset B$ であって、
$\mathfrak q$ が $(0) \subset A$ の上にあるものが存在する。
Algebra, Lemma \ref{algebra-lemma-dominate} により、
$B/\mathfrak q$ を支配する付値環
$A' \subset \kappa(\mathfrak q)$ を見つけられる。
誘導射 $\Spec(A') \to U \to X$ が、補題の問題の解である。
\end{proof}

\begin{lemma}
\label{spaces-morphisms-lemma-refined-valuative-criterion-universally-closed}
$S$ をスキームとする。
$f : X \to Y$ と $h : U \to X$ を $S$ 上の代数空間の射とする。次を仮定する。
\begin{enumerate}
\item $f$ と $h$ は準コンパクトである。
\item $|h|(|U|)$ は $|X|$ で稠密である。
\end{enumerate}
$A$ が分数体 $K$ をもつ付値環である任意の可換な実線図式
$$
\xymatrix{
\Spec(K) \ar[r] \ar[d] & U \ar[r] & X \ar[d] \\
\Spec(A) \ar[rr] \ar@{-->}[rru] & & Y
}
$$
に対して、次を仮定する。
\begin{enumerate}
\item[(3)] 図式を可換にする点線矢印は高々一つである。
\item[(4)] 体の拡大 $K'/K$、$A$ を支配する付値環
$A' \subset K'$、および図式
$$
\xymatrix{
\Spec(K') \ar[r] \ar[d] & \Spec(K) \ar[r] & U \ar[r] & X \ar[d] \\
\Spec(A') \ar[r] \ar[rrru] & \Spec(A) \ar[rr] & & Y
}
$$
を可換にする射 $\Spec(A') \to X$ が存在する。
\end{enumerate}
このとき $f$ は普遍閉である。さらに
\begin{enumerate}
\item[(5)] $f$ は準分離である
\end{enumerate}
ならば、$f$ は分離的かつ普遍閉である。
\end{lemma}

\begin{proof}
(1), (2), (3), (4) を仮定する。
$f$ に対する付値判定法の存在部分を確認する。
Lemma \ref{spaces-morphisms-lemma-quasi-compact-existence-universally-closed} により、
これから $f$ が普遍閉であることが従う。
そのため、可換図式
\begin{equation}
\label{spaces-morphisms-equation-start-with}
\vcenter{
\xymatrix{
\Spec(K) \ar[r] \ar[d] & X \ar[d] \\
\Spec(A) \ar[r] & Y
}
}
\end{equation}
を考える。ここで $A$ は付値環で、$K$ は $A$ の分数体である。
付値環と体は被約なので、Properties of Spaces, Lemma
\ref{spaces-properties-lemma-map-into-reduction} により、
$U$, $X$, $S$ をそれぞれの被約化で置き換えてよい。
この場合、$h(U)$ が稠密であるという仮定は、
$h : U \to X$ のスキーム論的像が $X$ であることを意味する
（Lemma \ref{spaces-morphisms-lemma-scheme-theoretic-image-reduced}）。

\medskip\noindent
$Y$ がアフィンの場合への帰着。
$\Spec(A)$ の閉点が
$\Im(|\Spec(R)| \to |Y|)$ の元へ写るようなエタール射
$\Spec(R) \to Y$ を選ぶ。
Lemma \ref{spaces-morphisms-lemma-lift-valuation-ring-through-flat-morphism} により、
付値環の局所環写像 $A \to A'$ と、可換図式
$$
\xymatrix{
\Spec(A') \ar[r] \ar[d] & \Spec(R) \ar[d] \\
\Spec(A) \ar[r] & Y
}
$$
に入る射 $\Spec(A') \to \Spec(R)$ を見つけられる。
Definition \ref{spaces-morphisms-definition-valuative-criterion} では付値環の拡大を許すので、
$A$ を $A'$、$Y$ を $\Spec(R)$、
$X$ を $X \times_Y \Spec(R)$、
$U$ を $U \times_Y \Spec(R)$ で置き換えてよいことは明らかである。

\medskip\noindent
以後、$Y = \Spec(R)$ はアフィンスキームであると仮定する。
$\Spec(K) \to X$ が $|\Spec(B)| \to |X|$ の像に属するような、
アフィンスキームからのエタール射 $\Spec(B) \to X$ を取る。
$K$ を拡大 $K' \supset K$ で、$A$ を $A$ を支配する付値環
$A' \subset K'$ で置き換えてよいから
（その存在は Algebra, Lemma
\ref{algebra-lemma-dominate} による）、
射 $\Spec(K) \to X$ は $\Spec(B)$ を経由すると仮定してよい
（$|X|$ の定義による）。
言い換えれば、$K$ を $B$-代数とみなしてよい。
$B$ 上の多項式環 $P$ と $B$-代数の全射 $P \to K$ を選ぶ。
すると $\Spec(P) \to X$ は合成
$\Spec(P) \to \Spec(B) \to X$ として平坦である。
したがって Lemma
\ref{spaces-morphisms-lemma-flat-base-change-scheme-theoretic-image} により、
射 $U \times_X \Spec(P) \to \Spec(P)$ のスキーム論的像は
$\Spec(P)$ である。
Lemma \ref{spaces-morphisms-lemma-reach-points-scheme-theoretic-image} により、
可換図式
$$
\xymatrix{
\Spec(K') \ar[r] \ar[d] & U \times_X \Spec(P) \ar[d] \\
\Spec(A') \ar[r] & \Spec(P)
}
$$
であって、$A'$ は付値環、$K'$ は $A'$ の分数体であり、
$\Spec(A')$ の閉点が $\Spec(K) \subset \Spec(P)$ に写るものを
見つけられる。言い換えれば、$B$-代数写像
$\varphi : K \to A'/\mathfrak m_{A'}$ がある。
$\varphi(A)$ を支配し、分数体が
$K'' = A'/\mathfrak m_{A'}$ である付値環
$A'' \subset A'/\mathfrak m_{A'}$ を選ぶ
（Algebra, Lemma
\ref{algebra-lemma-dominate}）。次のようにおく。
$$
C = \{\lambda \in A' \mid \lambda \bmod \mathfrak m_{A'} \in A''\}.
$$
これは Algebra, Lemma
\ref{algebra-lemma-stack-valuation-rings} により付値環である。
$C$ は分数体 $K'$ をもつ $R$-代数なので、補題の主張のような
可換な実線図式
$$
\xymatrix{
\Spec(K'_1) \ar@{-->}[r] \ar@{-->}[d] &
\Spec(K') \ar[r] \ar[d] & U \ar[r] & X \ar[d] \\
\Spec(C_1) \ar@{-->}[r] \ar@{-->}[rrru] & \Spec(C) \ar[rr] & & Y
}
$$
を得る。したがって仮定 (4) は $C \to C_1$ と、
図式を可換にする点線矢印を与える。
$A_1' = (C_1)_\mathfrak p$ を、$\mathfrak m_{A'} \subset C$ の上にある
素イデアル $\mathfrak p \subset C_1$ における $C_1$ の局所化とする。
More on Algebra, Lemma
\ref{more-algebra-lemma-valuation-ring-torsion-free-flat} により
$C \to C_1$ は平坦なので、そのような素イデアル $\mathfrak p$ が存在する
（Algebra, Lemmas
\ref{algebra-lemma-local-flat-ff} および
\ref{algebra-lemma-ff-rings}）。
$A'$ は $\mathfrak m_{A'}$ における $C$ の局所化であり、
$A'_1$ は付値環であることに注意する
（Algebra, Lemma
\ref{algebra-lemma-make-valuation-rings}）。
言い換えれば、$A' \to A'_1$ は付値環の局所環写像である。
仮定 (3) は図式
$$
\xymatrix{
\Spec(A'_1) \ar[r] \ar[d] & \Spec(C_1) \ar[r] & X \\
\Spec(A') \ar[r] & \Spec(P) \ar[r] & \Spec(B) \ar[u]
}
$$
が可換であることを含意する。
したがって、射 $\Spec(C_1) \to X$ の
$\Spec(C_1/\mathfrak p)$ への制限は、
$\Spec(C_1/\mathfrak p)$ の一般点上で合成
$$
\Spec(\kappa(\mathfrak p)) \to
\Spec(A'/\mathfrak m_{A'}) = \Spec(K'') \to
\Spec(K) \to X
$$
に制限される。さらに $C_1/\mathfrak p$ は、
$A$ を支配する $A''$ を支配する付値環である
（Algebra, Lemma
\ref{algebra-lemma-make-valuation-rings}）。
したがって射 $\Spec(C_1/\mathfrak p) \to X$ は、
望みどおり図式 (\ref{spaces-morphisms-equation-start-with}) に対する
付値判定法の存在部分を証示する。

\medskip\noindent
次に (5) も成り立つ、すなわち射
$\Delta : X \to X \times_S X$ が準コンパクトであると仮定する。
この場合、$h$ と $\Delta$ に対して仮定 (1) -- (4) が成り立つ。
したがって証明の第一部から $\Delta$ は普遍閉である。
Lemma \ref{spaces-morphisms-lemma-separated-diagonal-proper} により、
$f$ は分離的である。
\end{proof}

\section{分離性の付値判定法}
\label{spaces-morphisms-section-valuative-separatedness}

\noindent
まず逆向きの主張を証明し、次に判定法を述べる。

\begin{lemma}
\label{spaces-morphisms-lemma-separated-implies-valuative}
$S$ をスキームとする。
$f : X \to Y$ を $S$ 上の代数空間の射とする。
$f$ が分離的ならば、$f$ は付値判定法の一意性部分を満たす。
\end{lemma}

\begin{proof}
Definition \ref{spaces-morphisms-definition-valuative-criterion} のような図式が
与えられたとする。
図式に入る相異なる二つの射 $a, b : \Spec(A) \to X$ があると仮定する。
$Z \subset \Spec(A)$ を $a$ と $b$ の等化子とする。このとき
$Z = \Spec(A) \times_{(a, b), X \times_Y X, \Delta} X$ である。
$f$ が分離的ならば $\Delta$ は閉埋入なので、
これは $\Spec(A)$ の閉部分スキームである。
仮定により $\Spec(A)$ の一般点を含む。
$A$ は整域なので、これは $Z = \Spec(A)$ を含意する。
したがって望みどおり $a = b$ である。
\end{proof}

\begin{lemma}[分離性の付値判定法]
\label{spaces-morphisms-lemma-valuative-criterion-separatedness}
$S$ をスキームとする。
$f : X \to Y$ を $S$ 上の代数空間の射とする。次を仮定する。
\begin{enumerate}
\item 射 $f$ は準分離である。
\item 射 $f$ は付値判定法の一意性部分を満たす。
\end{enumerate}
このとき $f$ は分離的である。
\end{lemma}

\begin{proof}
仮定 (1) は $\Delta_{X/Y}$ が準コンパクトであることを意味する。
射 $\Delta_{X/Y} : X \to X \times_Y X$ は付値判定法の存在部分を
満たすと主張する。可換な実線図式
$$
\xymatrix{
\Spec(K) \ar[r] \ar[d] & X \ar[d] \\
\Spec(A) \ar[r] \ar@{-->}[ru] & X \times_Y X
}
$$
が与えられたとする。右下の矢印は、$Y$ 上の射の組
$a, b : \Spec(A) \to X$ に対応する。
仮定 (2) により $a = b$ である。
したがって $a$ を点線矢印として用いればよい。
ゆえに Lemma
\ref{spaces-morphisms-lemma-quasi-compact-existence-universally-closed} を適用でき、
$\Delta_{X/Y}$ は普遍閉であることがわかる。
$\Delta_{X/Y}$ は常に局所有限型かつ分離的なので、
More on Morphisms, Lemma
\ref{more-morphisms-lemma-characterize-finite} から
$\Delta_{X/Y}$ は有限射であると結論する
（Spaces, Lemma
\ref{spaces-lemma-representable-transformations-property-implication}
の一般原理も用いる）。
この時点で $\Delta_{X/Y}$ は表現可能な有限単射なので、
Morphisms, Lemma
\ref{morphisms-lemma-finite-monomorphism-closed} により閉埋入である。
\end{proof}

\begin{lemma}
\label{spaces-morphisms-lemma-characterize-separated-and-universally-closed}
$S$ をスキームとする。$f : X \to Y$ を $S$ 上の代数空間の射とする。
$f$ は準コンパクトかつ準分離であると仮定する。
次の条件は同値である。
\begin{enumerate}
\item $f$ は分離的かつ普遍閉である。
\item Definition \ref{spaces-morphisms-definition-valuative-criterion} の意味で
付値判定法が成り立つ。
\item $A$ が分数体 $K$ をもつ付値環である任意の可換な実線図式
$$
\xymatrix{
\Spec(K) \ar[r] \ar[d] & X \ar[d] \\
\Spec(A) \ar[r] \ar@{-->}[ru] & Y
}
$$
に対して一意な点線矢印が存在する。
すなわち、$f$ は Schemes, Definition
\ref{schemes-definition-valuative-criterion} の意味で付値判定法を満たす。
\end{enumerate}
\end{lemma}

\begin{proof}
$f$ は準分離なので、付値判定法の一意性部分が成り立てば
$f$ は分離的である
（Lemma \ref{spaces-morphisms-lemma-valuative-criterion-separatedness}）。
逆に $f$ が分離的ならば、付値判定法の一意性部分を満たす
（Lemma \ref{spaces-morphisms-lemma-separated-implies-valuative}）。
以上から、三つの各場合において、射 $f$ は分離的であり、
付値判定法の一意性部分を満たすことがわかる。
この場合、補題は Lemma
\ref{spaces-morphisms-lemma-characterize-universally-closed-separated} の形式的な帰結である。
\end{proof}

\section{固有性の付値判定法}
\label{spaces-morphisms-section-valuative-criterion-properness}

\noindent
主張を述べる。

\begin{lemma}[固有性の付値判定法]
\label{spaces-morphisms-lemma-characterize-proper}
$S$ をスキームとする。$f : X \to Y$ を $S$ 上の代数空間の射とする。
$f$ は有限型かつ準分離であると仮定する。
次の条件は同値である。
\begin{enumerate}
\item $f$ は固有である。
\item Definition \ref{spaces-morphisms-definition-valuative-criterion} の意味で
付値判定法が成り立つ。
\item $A$ が分数体 $K$ をもつ付値環である任意の可換な実線図式
$$
\xymatrix{
\Spec(K) \ar[r] \ar[d] & X \ar[d] \\
\Spec(A) \ar[r] \ar@{-->}[ru] & Y
}
$$
に対して一意な点線矢印が存在する。
すなわち、$f$ は Schemes, Definition
\ref{schemes-definition-valuative-criterion} の意味で付値判定法を満たす。
\end{enumerate}
\end{lemma}

\begin{proof}
Lemma \ref{spaces-morphisms-lemma-characterize-separated-and-universally-closed} と
定義から形式的に従う。
\end{proof}

\section{整射と有限射}
\label{spaces-morphisms-section-integral}

\noindent
代数空間の表現可能な射が整（それぞれ有限）であることの意味は、
Section \ref{spaces-morphisms-section-representable} ですでに定義した。

\begin{lemma}
\label{spaces-morphisms-lemma-integral-representable}
$S$ をスキームとする。
$f : X \to Y$ を $S$ 上の代数空間の表現可能な射とする。
$f$ が（Section \ref{spaces-morphisms-section-representable} の意味で）
整（それぞれ有限）であるための必要十分条件は、
任意のアフィンスキーム $Z$ と射 $Z \to Y$ に対して、
スキーム $X \times_Y Z$ がアフィンで、$Z$ 上整
（それぞれ有限）となることである。
\end{lemma}

\begin{proof}
スキームの整射（それぞれ有限射）の定義
（Morphisms, Definition
\ref{morphisms-definition-integral}）から直接従う。
\end{proof}

\noindent
これにより、次の定義を置ける。

\begin{definition}
\label{spaces-morphisms-definition-integral}
$S$ をスキームとする。
$f : X \to Y$ を $S$ 上の代数空間の射とする。
\begin{enumerate}
\item 任意のアフィンスキーム $Z$ と射 $Z \to Y$ に対して、
代数空間 $X \times_Y Z$ が $Z$ 上整なアフィンスキームで表現されるとき、
$f$ は {\it 整}であるという。
\item 任意のアフィンスキーム $Z$ と射 $Z \to Y$ に対して、
代数空間 $X \times_Y Z$ が $Z$ 上有限なアフィンスキームで表現されるとき、
$f$ は {\it 有限}であるという。
\end{enumerate}
\end{definition}

\begin{lemma}
\label{spaces-morphisms-lemma-integral-local}
$S$ をスキームとする。
$f : X \to Y$ を $S$ 上の代数空間の射とする。
次の条件は同値である。
\begin{enumerate}
\item $f$ は表現可能かつ整（それぞれ有限）である。
\item $f$ は整（それぞれ有限）である。
\item スキーム $V$ と全射エタール射 $V \to Y$ であって、
$V \times_Y X \to V$ が整（それぞれ有限）となるものが存在する。
\item Zariski 被覆 $Y = \bigcup Y_i$ であって、
各射 $f^{-1}(Y_i) \to Y_i$ が整（それぞれ有限）となるものが存在する。
\end{enumerate}
\end{lemma}

\begin{proof}
(1) が (2) を含意することは明らかである。
$V$ を $Y$ 上エタールなアフィンの非交和とすれば、
(2) が (3) を含意することも明らかである
（Properties of Spaces, Lemma
\ref{spaces-properties-lemma-cover-by-union-affines}）。
$V \to Y$ は (3) のようであると仮定する。
すると $V$ の任意のアフィン開集合 $W$ に対して、
$W \times_Y X$ は $V \times_Y X$ のアフィン開集合である。
したがって Properties of Spaces, Lemma
\ref{spaces-properties-lemma-subscheme} により、
$V \times_Y X$ はスキームである。
さらに射 $V \times_Y X \to V$ はアフィンである。
ゆえに Spaces, Lemma
\ref{spaces-lemma-morphism-sheaves-with-P-effective-descent-etale}
を適用できる。整射（それぞれ有限射）の類は、必要な性質をすべて満たすからである
（Morphisms, Lemmas
\ref{morphisms-lemma-base-change-finite} および Descent, Lemmas
\ref{descent-lemma-descending-property-integral},
\ref{descent-lemma-descending-property-finite}, および
\ref{descent-lemma-affine} を参照）。
この補題を適用した結論は、$f$ が表現可能かつ整
（それぞれ有限）であること、すなわち (1) である。

\medskip\noindent
(1) と (4) の同値性は、整（それぞれ有限）であることが
終域上 Zariski 局所的であることから従う
（上の参照は、整または有限であることが実際には終域上 fpqc 局所的であることを
示している）。
\end{proof}

\begin{lemma}
\label{spaces-morphisms-lemma-composition-integral}
整射（それぞれ有限射）の合成は整（それぞれ有限）である。
\end{lemma}

\begin{proof}
省略する。
\end{proof}

\begin{lemma}
\label{spaces-morphisms-lemma-base-change-integral}
整射（それぞれ有限射）の基底変換は整（それぞれ有限）である。
\end{lemma}

\begin{proof}
省略する。
\end{proof}

\begin{lemma}
\label{spaces-morphisms-lemma-finite-integral}
代数空間の有限射は整である。
代数空間の局所有限型な整射は有限である。
\end{lemma}

\begin{proof}
どちらの場合にも射は表現可能であり、$Y$ へ写るアフィンスキームによる
基底変換後に条件を確認できる
（Lemmas \ref{spaces-morphisms-lemma-integral-local}）。
したがって、スキームの場合の同じ補題、すなわち Morphisms, Lemma
\ref{morphisms-lemma-finite-integral} から従う。
\end{proof}

\begin{lemma}
\label{spaces-morphisms-lemma-integral-universally-closed}
$S$ をスキームとする。
$f : X \to Y$ を $S$ 上の代数空間の射とする。
次の条件は同値である。
\begin{enumerate}
\item $f$ は整である。
\item $f$ はアフィンかつ普遍閉である。
\end{enumerate}
\end{lemma}

\begin{proof}
どちらの場合にも射は表現可能であり、$Y$ へ写るアフィンスキームによる
基底変換後に条件を確認できる
（Lemmas \ref{spaces-morphisms-lemma-integral-local},
\ref{spaces-morphisms-lemma-affine-local}, および
\ref{spaces-morphisms-lemma-universally-closed-local}）。
したがって Morphisms, Lemma
\ref{morphisms-lemma-integral-universally-closed} から従う。
\end{proof}

\begin{lemma}
\label{spaces-morphisms-lemma-finite-quasi-finite}
代数空間の有限射は準有限である。
\end{lemma}

\begin{proof}
$f : X \to Y$ を代数空間の射とする。
Definition \ref{spaces-morphisms-definition-integral} および Lemmas
\ref{spaces-morphisms-lemma-quasi-compact-local} と
\ref{spaces-morphisms-lemma-quasi-finite-local} により、
両性質は $Y$ 上のアフィンへの基底変換後に確認できる。
すなわち $Y$ はアフィンであるとしてよい。
$f$ が有限ならば $X$ はスキームである。
したがってスキームに対する対応する結果
（Morphisms, Lemma
\ref{morphisms-lemma-finite-quasi-finite}）から従う。
\end{proof}

\begin{lemma}
\label{spaces-morphisms-lemma-finite-proper}
$S$ をスキームとする。
$f : X \to Y$ を $S$ 上の代数空間の射とする。
次の条件は同値である。
\begin{enumerate}
\item $f$ は有限である。
\item $f$ はアフィンかつ固有である。
\end{enumerate}
\end{lemma}

\begin{proof}
どちらの場合にも射は表現可能であり、$Y$ へ写るアフィンスキームへの
基底変換後に条件を確認できる
（Lemmas \ref{spaces-morphisms-lemma-integral-local},
\ref{spaces-morphisms-lemma-affine-local}, および
\ref{spaces-morphisms-lemma-proper-local}）。
したがって Morphisms, Lemma
\ref{morphisms-lemma-finite-proper} から従う。
\end{proof}

\begin{lemma}
\label{spaces-morphisms-lemma-closed-immersion-finite}
閉埋入は有限である（したがって、もちろん整である）。
\end{lemma}

\begin{proof}
省略する。
\end{proof}

\begin{lemma}
\label{spaces-morphisms-lemma-finite-union-finite}
$S$ をスキームとする。
$X_i \to Y$, $i = 1, \ldots, n$ を $S$ 上の代数空間の有限射とする。
このとき $X_1 \amalg \ldots \amalg X_n \to Y$ も有限である。
\end{lemma}

\begin{proof}
エタール局所化により、スキームの場合
（Morphisms, Lemma
\ref{morphisms-lemma-finite-union-finite}）から従う。
\end{proof}

\begin{lemma}
\label{spaces-morphisms-lemma-finite-permanence}
$S$ をスキームとする。
$f : X \to Y$ と $g : Y \to Z$ を $S$ 上の代数空間の射とする。
\begin{enumerate}
\item $g \circ f$ が有限で $g$ が分離的ならば、$f$ は有限である。
\item $g \circ f$ が整で $g$ が分離的ならば、$f$ は整である。
\end{enumerate}
\end{lemma}

\begin{proof}
$g \circ f$ が有限（それぞれ整）で、$g$ が分離的であると仮定する。
Lemma \ref{spaces-morphisms-lemma-base-change-integral} により、基底変換
$X \times_Z Y \to Y$ は有限（それぞれ整）である。
$Y \to Z$ は分離的なので、射 $X \to X \times_Z Y$ は閉埋入である
（Lemma \ref{spaces-morphisms-lemma-section-immersion}）。
閉埋入は有限（それぞれ整）である
（Lemma \ref{spaces-morphisms-lemma-closed-immersion-finite}）。
有限射（それぞれ整射）の合成は有限（それぞれ整）である
（Lemma \ref{spaces-morphisms-lemma-composition-integral}）。
したがって主張が従う。
\end{proof}





\section{有限局所自由射}
\label{spaces-morphisms-section-finite-locally-free}

\noindent
代数空間の表現可能射が有限局所自由であることの意味は、
Section \ref{spaces-morphisms-section-representable} ですでに定義した。

\begin{lemma}
\label{spaces-morphisms-lemma-finite-locally-free-representable}
$S$ をスキームとする。$f : X \to Y$ を $S$ 上の代数空間の
表現可能射とする。このとき、$f$ が
（Section \ref{spaces-morphisms-section-representable} の意味で）有限局所自由であることと、
$f$ がアフィンであり、層 $f_*\mathcal{O}_X$ が有限局所自由な
$\mathcal{O}_Y$-加群であることとは同値である。
\end{lemma}

\begin{proof}
$f$ が（Section \ref{spaces-morphisms-section-representable} で定義された意味で）
有限局所自由であると仮定する。これは、始域がスキームである任意の射
$V \to Y$ に対し、基底変換
$f' : V \times_Y X \to V$ がスキームの有限局所自由射であることを意味する。
さらにこれは（スキームの有限局所自由射の定義により）、
$f'_*\mathcal{O}_{V \times_Y X}$ が有限局所自由な
$\mathcal{O}_V$-加群であることを意味する。$V \to Y$ は全射エタールとなるように
選んでよい。Properties of Spaces,
Lemma \ref{spaces-properties-lemma-pushforward-etale-base-change-modules}
により、$f_*\mathcal{O}_X$ の $V$ への制限は有限局所自由である。
したがって、$Y_{spaces, \etale}$ 上の層 $f_*\mathcal{O}_X$ に
Modules on Sites, Lemma \ref{sites-modules-lemma-local-final-object}
を適用すると、$f_*\mathcal{O}_X$ は有限局所自由であることが分かる。

\medskip\noindent
逆に、$f$ がアフィンであり、$f_*\mathcal{O}_X$ が有限局所自由な
$\mathcal{O}_Y$-加群であると仮定する。$V$ をスキームとし、
$V \to Y$ を全射エタール射とする。再び Properties of Spaces,
Lemma \ref{spaces-properties-lemma-pushforward-etale-base-change-modules}
により、$f'_*\mathcal{O}_{V \times_Y X}$ は有限局所自由である。
したがって $f' : V \times_Y X \to V$ は（アフィンでもあるので）
有限局所自由である。
Spaces,
Lemma \ref{spaces-lemma-morphism-sheaves-with-P-effective-descent-etale}
により、$f$ は有限局所自由であると結論できる（Morphisms,
Lemma \ref{morphisms-lemma-base-change-finite-locally-free},
Descent, Lemmas \ref{descent-lemma-descending-property-finite-locally-free}
および \ref{descent-lemma-affine} を用いる）。これで証明された。
\end{proof}

\noindent
これにより、次の定義を与える準備が整った。

\begin{definition}
\label{spaces-morphisms-definition-finite-locally-free}
$S$ をスキームとする。
$f : X \to Y$ を $S$ 上の代数空間の射とする。
$f$ がアフィンであり、$f_*\mathcal{O}_X$ が有限局所自由な
$\mathcal{O}_Y$-加群であるとき、$f$ は \emph{有限局所自由} であるという。
この場合、層 $f_*\mathcal{O}_X$ が階数 $d$ の有限局所自由層であるならば、
$f$ は \emph{階数} または \emph{次数} $d$ をもつという。
\end{definition}

\begin{lemma}
\label{spaces-morphisms-lemma-finite-locally-free-local}
$S$ をスキームとする。
$f : X \to Y$ を $S$ 上の代数空間の射とする。
次は同値である：
\begin{enumerate}
\item $f$ は表現可能かつ有限局所自由である。
\item $f$ は有限局所自由である。
\item スキーム $V$ と全射エタール射 $V \to Y$ が存在して、
$V \times_Y X \to V$ は有限局所自由である。
\item $Y = \bigcup Y_i$ というザリスキー被覆が存在して、各射
$f^{-1}(Y_i) \to Y_i$ は有限局所自由である。
\end{enumerate}
\end{lemma}

\begin{proof}
(1) が (2) を含意することは明らかであり、$V$ を $Y$ 上エタールな
アフィンの非交和として取れば、(2) は (3) を含意する；
Properties of Spaces,
Lemma \ref{spaces-properties-lemma-cover-by-union-affines} を参照せよ。
$V \to Y$ が (3) のとおりであると仮定する。このとき、$V$ の任意の
アフィン開集合 $W$ に対して、$W \times_Y X$ は
$V \times_Y X$ のアフィン開集合である。したがって Properties of Spaces,
Lemma \ref{spaces-properties-lemma-subscheme}
により、$V \times_Y X$ はスキームである。さらに射
$V \times_Y X \to V$ はアフィンである。したがって Spaces,
Lemma \ref{spaces-lemma-morphism-sheaves-with-P-effective-descent-etale}
を適用できる。実際、有限局所自由射の類は必要な性質をすべて満たす
（Morphisms,
Lemma \ref{morphisms-lemma-base-change-finite-locally-free},
Descent, Lemmas \ref{descent-lemma-descending-property-finite-locally-free}
および \ref{descent-lemma-affine} を参照）。
この補題を適用した結論は、$f$ が表現可能かつ有限局所自由、
すなわち (1) が成り立つということである。

\medskip\noindent
(1) と (4) の同値性は、有限局所自由であることが終域上ザリスキー局所的であることから従う
（上の参照は、実際には有限局所自由であることが終域上 fpqc 局所的であることを示している）。
\end{proof}

\begin{lemma}
\label{spaces-morphisms-lemma-composition-finite-locally-free}
有限局所自由射の合成は有限局所自由である。
\end{lemma}

\begin{proof}
省略する。
\end{proof}

\begin{lemma}
\label{spaces-morphisms-lemma-base-change-finite-locally-free}
有限局所自由射の基底変換は有限局所自由である。
\end{lemma}

\begin{proof}
省略する。
\end{proof}

\begin{lemma}
\label{spaces-morphisms-lemma-finite-flat}
$S$ をスキームとする。
$f : X \to Y$ を $S$ 上の代数空間の射とする。
次は同値である：
\begin{enumerate}
\item $f$ は有限局所自由である。
\item $f$ は有限、平坦、かつ有限表示局所的である。
\end{enumerate}
$Y$ が局所ネーター的ならば、これらはさらに
\begin{enumerate}
\item[(3)] $f$ は有限かつ平坦である。
\end{enumerate}
とも同値である。
\end{lemma}

\begin{proof}
三つの場合のいずれでも射は表現可能であり、全射エタール射
$V \to Y$ による基底変換の後で性質を確認できる；
Lemmas \ref{spaces-morphisms-lemma-integral-local},
\ref{spaces-morphisms-lemma-finite-locally-free-local},
\ref{spaces-morphisms-lemma-flat-local}, および
\ref{spaces-morphisms-lemma-finite-presentation-local} を参照せよ。
$Y$ が局所ネーター的ならば、$V$ も局所ネーター的である。
したがって、スキームの場合の対応する結果、すなわち
Morphisms, Lemma \ref{morphisms-lemma-finite-flat} から従う。
\end{proof}





\section{有理写像}
\label{spaces-morphisms-section-rational-maps}

\noindent
本節は Morphisms, Section \ref{morphisms-section-rational-maps}
に対応する。位相空間の稠密開集合どうしの共通部分が稠密開集合であることを、
以下では断りなく用いる。

\begin{definition}
\label{spaces-morphisms-definition-rational-map}
$S$ をスキームとする。$X$, $Y$ を $S$ 上の代数空間とする。
\begin{enumerate}
\item $X$ の稠密開部分空間 $U$, $V$ 上で定義された $S$ 上の代数空間の射
$f : U \to Y$, $g : V \to Y$ を考える。ある稠密開部分空間
$W \subset U \cap V$ 上で $f|_W = g|_W$ となるとき、$f$ と $g$ は
\emph{同値}であるという。
\item \emph{$X$ から $Y$ への有理写像}とは、(1) で定義した同値関係による
同値類である。
\item $S$ 上の代数空間の射 $X \to B$ および $Y \to B$ が与えられているとする。
$X$ から $Y$ への有理写像が \emph{$X$ から $Y$ への $B$-有理写像}であるとは、
同値類の代表 $f : U \to Y$ であって $B$ 上の射であるものが存在することをいう。
\end{enumerate}
\end{definition}

\noindent
この定義の (1) における二つの射 $f$, $g$ については、同値であるという代わりに、
同じ有理写像を定めるともいう。今後考える多くの場合に、代数空間
$X$ と $Y$ は、スキームである稠密開部分空間 $X'$ と $Y'$ を含む。
このとき $X$ から $Y$ への有理写像は、Morphisms,
Definition \ref{spaces-morphisms-definition-rational-map} の意味での
$X'$ から $Y'$ への $S$-有理写像と同じである。
したがって、スキームについて展開された理論をすべて利用できる。

\begin{definition}
\label{spaces-morphisms-definition-rational-function}
$S$ をスキームとする。$X$ を $S$ 上の代数空間とする。
\emph{$X$ 上の有理関数}とは、$X$ から $\mathbf{A}^1_S$ への有理写像である。
\end{definition}

\noindent
Morphisms, Definition \ref{morphisms-definition-rational-function}
に続く議論を見ると、$X$ がスキームである場合には、これはそこで定義された概念と
同じであることが分かる。

\medskip\noindent
$S$ 上の任意のスキーム $T$ に対し、標準的な同一視
$$
\Mor_S(T, \mathbf{A}^1_S) = \Mor(T, \mathbf{A}^1_\mathbf{Z}) =
\Gamma(T, \mathcal{O}_T)
$$
があることを思い出そう；Schemes,
Example \ref{schemes-example-global-sections} を参照せよ。
したがって $\mathbf{A}^1_S$ は $S$ 上のスキームの圏における環対象である。
言い換えると、加法と乗法は射
$$
+ : \mathbf{A}^1_S \times_S \mathbf{A}^1_S \to \mathbf{A}^1_S
\quad\text{and}\quad
* : \mathbf{A}^1_S \times_S \mathbf{A}^1_S \to \mathbf{A}^1_S
$$
を定め、環における加法と乗法の公理（いつものとおり $1$ をもつ可換環）を満たす。
したがって、$\mathbf{A}^1_S$ への有理写像の集合も自然な環構造をもつ。

\begin{definition}
\label{spaces-morphisms-definition-ring-of-rational-functions}
$S$ をスキームとする。
$X$ を $S$ 上の代数空間とする。
\emph{$X$ 上の有理関数環}とは、元が有理関数であり、加法と乗法を上のように定めた
環 $R(X)$ のことである。
\end{definition}

\noindent
整代数空間の関数体は後で定義する；Spaces over Fields,
Section \ref{spaces-over-fields-section-integral-spaces} を参照せよ。

\begin{definition}
\label{spaces-morphisms-definition-domain-of-definition}
$S$ をスキームとする。$\varphi$ を $S$ 上の二つの代数空間 $X$ と $Y$ の間の
有理写像とする。$x \in |U|$ となる $\varphi$ の代表 $(U, f)$ が存在するとき、
$\varphi$ は \emph{点 $x \in |X|$ で定義されている}という。
$\varphi$ の \emph{定義域}とは、$\varphi$ が定義されているすべての点の集合である。
\end{definition}

\noindent
定義域は Properties of Spaces,
Lemma \ref{spaces-properties-lemma-open-subspaces} により、$X$ の開部分空間とみなす。
この定義のもとでは、一般に $\varphi$ が定義域全体で定義された代表をもつとは限らない。

\begin{lemma}
\label{spaces-morphisms-lemma-rational-map-from-reduced-to-separated}
$S$ をスキームとする。$X$ と $Y$ を $S$ 上の代数空間とする。
$X$ は被約であり、$Y$ は $S$ 上分離的であると仮定する。
$\varphi$ を $X$ から $Y$ への有理写像とし、その定義域を $U \subset X$ とする。
このとき、$\varphi$ を表す代数空間の射 $f : U \to Y$ が一意に存在する。
\end{lemma}

\begin{proof}
$(V, g)$ と $(V', g')$ を $\varphi$ の代表とする。このとき $g, g'$ は
ある稠密開部分空間 $W \subset V \cap V'$ 上で一致する。
一方、$g|_{V \cap V'}$ と $g'|_{V \cap V'}$ の等化子 $E$ は
$V \cap V'$ の閉部分空間である。実際、それは
$g|_{V \cap V'}$ と $g'|_{V \cap V'}$ が定める射
$V \cap V' \to Y \times_S Y$ による
$\Delta : Y \to Y \times_S Y$ の基底変換である。
いま $W \subset E$ から $|E| = |V \cap V'|$ が従う。
$V \cap V'$ は被約なので、スキーム論的に $E = V \cap V'$、すなわち
$g|_{V \cap V'} = g'|_{V \cap V'}$ である；Properties of Spaces,
Lemma \ref{spaces-properties-lemma-map-into-reduction} を参照せよ。
したがって、$\coprod V \to U$ は fppf 層の全射であり、
$\coprod_{V, V'} V \cap V' = (\coprod V) \times_U (\coprod V)$ であるから、
$\varphi$ の代表 $g : V \to Y$ たちを貼り合わせて射
$f : U \to Y$ を得ることができる。
\end{proof}

\noindent
一般には有理写像を合成することには意味がない。第一の有理写像の代表の像と、
第二の有理写像の定義域との共通部分が空でありうるからである。
しかし、空間が既約であると仮定して支配的有理写像を考えるならば、
有理写像を合成できる。

\begin{definition}
\label{spaces-morphisms-definition-dominant-rational}
$S$ をスキームとする。$X$ と $Y$ を $S$ 上の代数空間とする。
$|X|$ と $|Y|$ が既約であると仮定する。$X$ から $Y$ への有理写像が
\emph{支配的}であるとは、その任意の代表 $f : U \to Y$ が
Definition \ref{spaces-morphisms-definition-dominant} の意味で支配的射であることをいう。
\end{definition}

\noindent
既約代数空間 $X$ と $Y$ の間の支配的有理写像 $\varphi$ と、
$Y$ から $Z$ への任意の有理写像 $\psi$ とを合成できる。
実際、$|U| \subset |X|$ が稠密開である代表 $f : U \to Y$ と、
$|V| \subset |Y|$ が稠密開である代表 $g : V \to Z$ を選ぶ。
すると $W = |f|^{-1}(V) \subset |X|$ は開かつ空でない（$|f|$ の像は稠密なので、
空でない開集合 $V$ と交わらなければならない）ため、$|X|$ が既約であることから稠密である。
$\psi \circ \varphi$ を $g \circ f|_W : W \to Z$ の同値類として定義する。
これが良定義であることの確認は省略する。

\medskip\noindent
このようにして、対象が $S$ 上の既約代数空間で、射が支配的有理写像である圏を得る。

\begin{definition}
\label{spaces-morphisms-definition-birational-spaces}
$S$ をスキームとする。$X$ と $Y$ を $S$ 上の代数空間とし、
$|X|$ と $|Y|$ は既約であるとする。$X$ と $Y$ が、$S$ 上の既約代数空間と
支配的有理写像の圏で同型であるとき、$X$ と $Y$ は \emph{双有理}であるという。
\end{definition}

\noindent
$X$ と $Y$ が双有理な既約代数空間ならば、任意の代数空間 $Z$ に対して、
$X$ から $Z$ への有理写像の集合と $Y$ から $Z$ への有理写像の集合との間に
全単射がある（$Z$ に関して関手的である）。「一般の」既約代数空間に対しては、
これは可能な定義の一つにすぎない。別の定義として、$X$ と $Y$ の有理関数環が
同型であることを要求することもできる；これら二つの概念は場合によっては同値である
（将来の参照をここに挿入する）。

\begin{lemma}
\label{spaces-morphisms-lemma-birational}
$S$ をスキームとする。$X$ と $Y$ を $S$ 上の代数空間とし、
$|X|$ と $|Y|$ は既約であるとする。このとき、$X$ と $Y$ が双有理であることと、
$S$ 上の代数空間として同型な、空でない開部分空間
$U \subset X$ と $V \subset Y$ が存在することとは同値である。
\end{lemma}

\begin{proof}
$X$ と $Y$ が双有理であると仮定する。$f : U \to Y$ と $g : V \to X$ が、
$X$ から $Y$ への支配的有理写像と $Y$ から $X$ への支配的有理写像で、
互いに逆なものを定めるとする。$U$ を縮小することにより、
$f : U \to Y$ は $V$ を経由すると仮定してよい。
$g \circ f$ は支配的有理写像として恒等写像なので、合成
$U \to V \to X$ は $U$ のある稠密開部分上で恒等写像である。
したがって $U$ をより小さい開集合で置き換えることで、
$U \to V \to X$ は $U$ の $X$ への包含写像であると仮定してよい。
対称性により、開部分空間 $V' \subset V$ が存在して、
$g|_{V'} : V' \to X$ は $U \subset X$ を経由し、
$V' \to U \to Y$ は恒等写像となる。
$|U| \to |V|$ による $|V'|$ の逆像は $|U|$ の開集合であり、したがって
ある開部分空間 $U' \subset U$ に対する $|U'|$ に等しい；
Properties of Spaces, Lemma \ref{spaces-properties-lemma-open-subspaces} を参照せよ。
このとき $U' \subset U \to V$ は $U' \to V'$ として分解する。
同様に $V' \to U$ は $V' \to U'$ として分解する。
$U' \to V'$ と $V' \to U'$ が $S$ 上の代数空間の互いに逆な射であることを
読者は確認でき、証明が完了する。
\end{proof}










\section{代数空間の相対正規化}
\label{spaces-morphisms-section-normalization-X-in-Y}

\noindent
本節は Morphisms, Section \ref{morphisms-section-normalization-X-in-Y}
に対応する。

\begin{lemma}
\label{spaces-morphisms-lemma-integral-closure}
$S$ をスキームとする。$X$ を $S$ 上の代数空間とする。
$\mathcal{A}$ を準連接な $\mathcal{O}_X$-代数の層とする。
準連接な $\mathcal{O}_X$-代数の層 $\mathcal{A}' \subset \mathcal{A}$ であって、
$X_\etale$ の任意のアフィン対象 $U$ に対し、環
$\mathcal{A}'(U) \subset \mathcal{A}(U)$ が $\mathcal{A}(U)$ における
$\mathcal{O}_X(U)$ の整閉包となるものが存在する。
\end{lemma}

\begin{proof}
$U$ を $X_\etale$ の対象とする。このとき $U$ はスキームである。
ザリスキーサイトへの制限を $\mathcal{A}|_U$ と書く。
$\mathcal{A}|_U$ は準連接な $\mathcal{O}_U$-代数の層なので、
Morphisms, Lemma \ref{morphisms-lemma-integral-closure}
を適用して準連接部分代数 $\mathcal{A}'_U \subset \mathcal{A}|_U$ を得る。
この $\mathcal{A}'_U$ は任意のアフィン開集合 $W \subset U$ 上で、
補題の主張に記された値をもつ。
$f : U' \to U$ が $X_\etale$ における射ならば、
$\mathcal{A}|_{U'} = f^*(\mathcal{A}|_U)$ である。ここで $f^*$ はザリスキー位相における
射 $f$ による引き戻しを表す；これは $\mathcal{A}$ が準連接だからである
（Properties of Spaces, Section
\ref{spaces-properties-section-quasi-coherent} の導入部と、そこで参照されている
スキーム上の降下の章の議論を参照せよ）。$f$ はエタールなので、
More on Morphisms, Lemma
\ref{more-morphisms-lemma-integral-closure-smooth-pullback}
により標準同型 $f^*(\mathcal{A}'_U) = \mathcal{A}'_{U'}$ を得る。
これから直ちに、$X_\etale$ 上の $\mathcal{O}_X$-代数の部分前層
$\mathcal{A}' \subset \mathcal{A}$ であって、ザリスキー位相に関する層であり、
アフィン対象上で正しい値をもつものが得られる。
さらに、各 $\mathcal{A}'_U$ がスキーム $U$ 上準連接であり、
エタール射 $f : U' \to U$ に対して
$\mathcal{A}'_{U'} = f^*(\mathcal{A}'_U)$ であることから、
$\mathcal{A}'$ は $X_\etale$ 上でも準連接である
（これは局所的性質であり、上の参照は $U_\etale$ 上の準連接加群を
まさにこの仕方で記述している）。
\end{proof}

\begin{definition}
\label{spaces-morphisms-definition-integral-closure}
$S$ をスキームとする。$X$ を $S$ 上の代数空間とする。
$\mathcal{A}$ を準連接な $\mathcal{O}_X$-代数の層とする。
\emph{$\mathcal{A}$ における $\mathcal{O}_X$ の整閉包}とは、上の
Lemma \ref{spaces-morphisms-lemma-integral-closure} で構成した準連接な
$\mathcal{O}_X$-部分代数 $\mathcal{A}' \subset \mathcal{A}$ のことである。
\end{definition}

\noindent
特に、代数空間の準コンパクトかつ準分離的な射
$f : Y \to X$ に対して $\mathcal{A} = f_*\mathcal{O}_Y$ である場合に、
これを適用する（Lemma \ref{spaces-morphisms-lemma-pushforward} を参照せよ）。
この準連接 $\mathcal{O}_X$-代数の相対スペクトルを取れば
（Lemma \ref{spaces-morphisms-lemma-affine-equivalence-algebras}）、$Y$ における $X$ の正規化を得る。

\begin{definition}
\label{spaces-morphisms-definition-normalization-X-in-Y}
$S$ をスキームとする。$f : Y \to X$ を $S$ 上の代数空間の
準コンパクトかつ準分離的な射とする。$\mathcal{O}'$ を
$f_*\mathcal{O}_Y$ における $\mathcal{O}_X$ の整閉包とする。
\emph{$Y$ における $X$ の正規化}とは、$S$ 上の代数空間の射
$$
\nu : X' = \underline{\Spec}_X(\mathcal{O}') \to X
$$
である。これには、もとの射 $f$ の自然な分解
$$
Y \xrightarrow{f'} X' \xrightarrow{\nu} X
$$
が備わっている。
\end{definition}

\noindent
この分解を得るには、Remark
\ref{spaces-morphisms-remark-factorization-quasi-compact-quasi-separated}
と $\underline{\Spec}$ 構成の関手性を用いる。

\begin{lemma}
\label{spaces-morphisms-lemma-properties-normalization}
$S$ をスキームとする。$f : Y \to X$ を $S$ 上の代数空間の
準コンパクトかつ準分離的な射とする。$Y \to X' \to X$ を
$Y$ における $X$ の正規化とする。
\begin{enumerate}
\item $W \to X$ が $S$ 上の代数空間のエタール射ならば、
$W \times_X X'$ は $W \times_X Y$ における $W$ の正規化である。
\item $Y$ と $X$ が表現可能ならば、$X'$ も表現可能であり、
Morphisms, Section \ref{morphisms-section-normalization} で構成された
スキーム $Y$ におけるスキーム $X$ の正規化と標準的に同型である。
\end{enumerate}
\end{lemma}

\begin{proof}
$Y$ における $X$ の正規化の形成がエタール基底変換と可換することは
構成から直ちに分かる。すなわち (1) が成り立つ。
一方、$X$ と $Y$ がスキームであるとき、アフィン開集合 $U \subset X$ に対して
$f_*\mathcal{O}_Y(U) = \mathcal{O}_Y(f^{-1}(U))$ であるから、
$\nu^{-1}(U)$ はスキームの場合の対応する構成で得られるものと
まったく同じ環のスペクトルである。
\end{proof}

\noindent
この構成の特徴づけを与えよう。

\begin{lemma}
\label{spaces-morphisms-lemma-characterize-normalization}
$S$ をスキームとする。$f : Y \to X$ を $S$ 上の代数空間の
準コンパクトかつ準分離的な射とする。
$Y$ における $X$ の正規化を $\nu : X' \to X$ とするとき、分解
$f = \nu \circ f'$ は次の二つの性質によって特徴づけられる：
\begin{enumerate}
\item 射 $\nu$ は整である。
\item $\pi : Z \to X$ が整である任意の分解 $f = \pi \circ g$ に対して、
可換図式
$$
\xymatrix{
Y \ar[d]_{f'} \ar[r]_g & Z \ar[d]^\pi \\
X' \ar[ru]^h \ar[r]^\nu & X
}
$$
を満たす一意な射 $h : X' \to Z$ が存在する。
\end{enumerate}
さらに、(2) の射 $h : X' \to Z$ は $Y$ における $Z$ の正規化である。
\end{lemma}

\begin{proof}
Definition \ref{spaces-morphisms-definition-normalization-X-in-Y} のとおり、
$\mathcal{O}' \subset f_*\mathcal{O}_Y$ を $\mathcal{O}_X$ の整閉包とする。
射 $\nu$ は構成により整であり、これで (1) が示された。
(2) のように、$\pi : Z \to X$ が整である分解
$f = \pi \circ g$ が与えられたと仮定する。
Definition \ref{spaces-morphisms-definition-integral} により $\pi$ はアフィンなので、
$Z$ は準連接な $\mathcal{O}_X$-代数の層 $\mathcal{B}$ の相対スペクトルである。
射 $g : X \to Z$ は $\mathcal{O}_X$-代数の準同型
$\chi : \mathcal{B} \to f_*\mathcal{O}_Y$ に対応する。
$X$ 上エタールな任意のアフィン $U$ に対し、$\mathcal{B}(U)$ は
$\mathcal{O}_X(U)$ 上整なので（Definition \ref{spaces-morphisms-definition-integral}）、
Lemma \ref{spaces-morphisms-lemma-integral-closure} から
$\chi(\mathcal{B}) \subset \mathcal{O}'$ が分かる。
相対スペクトルの関手性 Lemma \ref{spaces-morphisms-lemma-affine-equivalence-algebras}
により、これは一意な射 $h : X' \to Z$ を与える。
図式が可換であることの確認は省略する。

\medskip\noindent
(1) と (2) が分解 $f = \nu \circ f'$ を特徴づけることは明らかである。
実際、それらはこの分解をある圏の始対象として特徴づけている。
(2) の射 $h$ は Lemma \ref{spaces-morphisms-lemma-finite-permanence} により整である。
$\pi' : Z' \to Z$ が整である分解 $g = \pi' \circ g'$ が与えられると、
分解 $f = (\pi \circ \pi') \circ g'$ と射 $h' : X' \to Z'$ が得られる。
一意性から $\pi' \circ h' = h$ である。
したがって特徴づけ (1), (2) を射 $h : X' \to Z$ に適用でき、
これが補題の最後の主張を与える。
\end{proof}

\begin{lemma}
\label{spaces-morphisms-lemma-normalization-in-reduced}
$S$ をスキームとする。$f : Y \to X$ を $S$ 上の代数空間の
準コンパクトかつ準分離的な射とする。$X' \to X$ を $Y$ における
$X$ の正規化とする。$Y$ が被約ならば、$X'$ も被約である。
\end{lemma}

\begin{proof}
被約環の部分環は被約であるという事実から従う。
いくつかの詳細は省略する。
\end{proof}

\begin{lemma}
\label{spaces-morphisms-lemma-normalization-generic}
$S$ をスキームとする。$f : Y \to X$ をスキームの準コンパクトかつ
準分離的な射とする。$X' \to X$ を $Y$ における $X$ の正規化とする。
$x' \in |X'|$ が余次元 $0$ の点ならば（Properties of Spaces,
Definition \ref{spaces-properties-definition-dimension-local-ring}）、
$x'$ は余次元 $0$ のある $y \in |Y|$ の像である。
\end{lemma}

\begin{proof}
Lemma \ref{spaces-morphisms-lemma-properties-normalization} と定義により、
$X = \Spec(A)$ はアフィンであると仮定してよい。このとき
$X' = \Spec(A')$ であり、$A'$ は $\Gamma(Y, \mathcal{O}_Y)$ における
$A$ の整閉包で、$x'$ は $A'$ の極小素イデアルに対応する。
$V = \Spec(B)$ がアフィンであるような全射エタール射 $V \to Y$ を選ぶ。
このとき $A' \to B$ は単射なので、$A'$ の各極小素イデアルは
$B$ のある極小素イデアルの像である；Algebra,
Lemma \ref{algebra-lemma-injective-minimal-primes-in-image} を参照せよ。
補題が従う。
\end{proof}

\begin{lemma}
\label{spaces-morphisms-lemma-normalization-in-disjoint-union}
$S$ をスキームとする。
$f : Y \to X$ を $S$ 上の代数空間の準コンパクトかつ準分離的な射とする。
$Y = Y_1 \amalg Y_2$ が二つの代数空間の非交和であると仮定する。
$f_i = f|_{Y_i}$ と書く。$X_i'$ を $Y_i$ における $X$ の正規化とする。
このとき $X_1' \amalg X_2'$ は $Y$ における $X$ の正規化である。
\end{lemma}

\begin{proof}
省略する。
\end{proof}

\begin{lemma}
\label{spaces-morphisms-lemma-normalization-in-universally-closed}
$S$ をスキームとする。
$f : X \to Y$ を $S$ 上の代数空間の準コンパクト、準分離的、かつ普遍閉な射とする。
このとき $f_*\mathcal{O}_X$ は $\mathcal{O}_Y$ 上整である。
言い換えると、$X$ における $Y$ の正規化は、Remark
\ref{spaces-morphisms-remark-factorization-quasi-compact-quasi-separated} の分解
$$
X \longrightarrow \underline{\Spec}_Y(f_*\mathcal{O}_X)
\longrightarrow Y
$$
に等しい。
\end{lemma}

\begin{proof}
問題は $Y$ 上エタール局所的なので、$Y = \Spec(R)$ がアフィンである場合に帰着してよい。
$h \in \Gamma(X, \mathcal{O}_X)$ とする。$h$ が $R$ 上のモニック方程式を満たすことを
示さなければならない。$h$ を次の可換図式における射と考える：
$$
\xymatrix{
X \ar[rr]_h \ar[rd]_f & & \mathbf{A}^1_Y \ar[ld] \\
& Y &
}
$$
$Z \subset \mathbf{A}^1_Y$ を $h$ のスキーム論的像とする；
Definition \ref{spaces-morphisms-definition-scheme-theoretic-image} を参照せよ。
$f$ が準コンパクトで $\mathbf{A}^1_Y \to Y$ が分離的なので、射 $h$ は
準コンパクトである；Lemma \ref{spaces-morphisms-lemma-quasi-compact-permanence} を参照せよ。
Lemma \ref{spaces-morphisms-lemma-quasi-compact-scheme-theoretic-image} により、射
$X \to Z$ の台位相空間上の像は稠密である。
Lemma \ref{spaces-morphisms-lemma-universally-closed-permanence} により、射 $X \to Z$ は閉である。
したがって（集合論的に）$h(X) = Z$ である。よって
Lemma \ref{spaces-morphisms-lemma-image-proper-is-proper} を用いて、$Z \to Y$ は普遍閉
（実際には固有）であると結論できる。$Z \subset \mathbf{A}^1_Y$ なので、
$Z \to Y$ はアフィンかつ固有であり、したがって
Lemma \ref{spaces-morphisms-lemma-integral-universally-closed} により整である。
$\mathbf{A}^1_Y = \Spec(R[T])$ と書けば、$Z$ のイデアル
$I \subset R[T]$ はモニック多項式 $P(T) \in R[T]$ を含む。
したがって $P(h) = 0$ であり、証明された。
\end{proof}

\begin{lemma}
\label{spaces-morphisms-lemma-normalization-in-integral}
$S$ をスキームとする。$f : Y \to X$ を $S$ 上の代数空間の整射とする。
このとき $Y$ における $X$ の整閉包は $Y$ に等しい。
\end{lemma}

\begin{proof}
Lemma \ref{spaces-morphisms-lemma-integral-universally-closed} により、これは
Lemma \ref{spaces-morphisms-lemma-normalization-in-universally-closed} の特別な場合である。
\end{proof}

\begin{lemma}
\label{spaces-morphisms-lemma-nagata-normalization-finite}
$S$ をスキームとする。$f : X \to Y$ を $S$ 上の代数空間の射とする。
次を仮定する：
\begin{enumerate}
\item $Y$ は永田である。
\item $f$ は有限型かつ準分離的である。
\item $X$ は被約である。
\end{enumerate}
このとき、$X$ における $Y$ の正規化 $\nu : Y' \to Y$ は有限である。
\end{lemma}

\begin{proof}
Lemma \ref{spaces-morphisms-lemma-properties-normalization} により、問題は $Y$ 上エタール局所的である。
したがって $Y = \Spec(R)$ はアフィンであると仮定してよい。
すると $R$ はネーター永田環であり、$\Gamma(X, \mathcal{O}_X)$ における
$R$ の整閉包が $R$ 上有限であることを示せばよい。
$f$ は準コンパクトなので、$X$ も準コンパクトである。
アフィンスキーム $U$ と全射エタール射 $U \to X$ を選ぶ
（Properties of Spaces, Lemma
\ref{spaces-properties-lemma-quasi-compact-affine-cover}）。
このとき $\Gamma(X, \mathcal{O}_X) \subset \Gamma(U, \mathcal{O}_X)$ である。
$R$ はネーター的なので、$\Gamma(U, \mathcal{O}_U)$ における
$R$ の整閉包が $R$ 上有限であることを示せば十分である。
$U \to Y$ は有限型なので、これは Morphisms,
Lemma \ref{morphisms-lemma-nagata-normalization-finite} から従う。
\end{proof}






\section{正規化}
\label{spaces-morphisms-section-normalization}

\noindent
本節は Morphisms, Section \ref{morphisms-section-normalization}
に対応する。

\begin{lemma}
\label{spaces-morphisms-lemma-prepare-normalization}
$S$ をスキームとする。$X$ を $S$ 上の代数空間とする。
次は同値である：
\begin{enumerate}
\item スキーム $U$ からの全射エタール射 $U \to X$ であって、
$U$ の各準コンパクト開集合が有限個の既約成分をもつものが存在する。
\item 任意のスキーム $U$ と任意のエタール射 $U \to X$ に対して、
$U$ の各準コンパクト開集合は有限個の既約成分をもつ。
\item $X$ 上エタールな任意の準コンパクト代数空間 $Y$ に対して、
$Y$ の余次元 $0$ の点の集合は有限である（Properties of Spaces,
Definition \ref{spaces-properties-definition-dimension-local-ring}）。
\item $X$ 上エタールな任意の準コンパクト代数空間 $Y$ に対して、
空間 $|Y|$ は有限個の既約成分をもつ。
\end{enumerate}
$X$ が表現可能ならば、これは $X$ の各準コンパクト開集合が有限個の
既約成分をもつことを意味する。
\end{lemma}

\begin{proof}
(1) と (2) の同値性および最後の主張は、Descent,
Lemma \ref{descent-lemma-locally-finite-nr-irred-local-fppf} と
Properties of Spaces, Lemma \ref{spaces-properties-lemma-type-property}
から従う。スキームである $Y$ だけを考えれば、(4) が (1) と (2) を
含意することは明らかである。同様に、スキームの場合、余次元 $0$ の点は
その既約成分の生成点なので、(3) は (1) と (2) を含意する；例えば
Properties of Spaces,
Lemma \ref{spaces-properties-lemma-codimension-0-points} を参照せよ。

\medskip\noindent
逆に (2) を仮定し、$Y \to X$ を $Y$ が準コンパクトであるような
代数空間のエタール射とする。このとき、アフィンスキーム $V$ と
全射エタール射 $V \to Y$ を選べる
（Properties of Spaces, Lemma
\ref{spaces-properties-lemma-quasi-compact-affine-cover}）。
(2) により $V$ は有限個の既約成分をもち、$|V| \to |Y|$ は全射連続なので、
Topology, Lemma
\ref{topology-lemma-surjective-continuous-irreducible-components}
により $|Y|$ は有限個の既約成分をもつ。したがって (4) が成り立つ。
同様に、Properties of Spaces,
Lemma \ref{spaces-properties-lemma-codimension-0-points} により、
$V$ の既約成分の生成点の像は $Y$ の余次元 $0$ の点であり、
それらは有限個である。すなわち (3) が成り立つ。
\end{proof}

\begin{lemma}
\label{spaces-morphisms-lemma-locally-noetherian-can-be-normalized}
$S$ をスキームとする。$X$ を $S$ 上の局所ネーター代数空間とする。
このとき $X$ は Lemma \ref{spaces-morphisms-lemma-prepare-normalization} の同値な条件を満たす。
\end{lemma}

\begin{proof}
$U \to X$ がエタールで $U$ がスキームならば、$U$ は局所ネータースキームである；
Properties of Spaces, Section
\ref{spaces-properties-section-types-properties} を参照せよ。
局所ネータースキームの既約成分の集合は局所有限である
（Divisors, Lemma \ref{divisors-lemma-components-locally-finite}）。
したがって、$X$ は補題の条件 (2) を満たす。
\end{proof}

\begin{lemma}
\label{spaces-morphisms-lemma-flat-image-point-codimension-0}
$S$ をスキームとする。$f : X \to Y$ を $S$ 上の代数空間の平坦射とする。
このとき $x \in |X|$ に対して、$x$ が余次元 $0$ ならば、
$X \Rightarrow f(x)$ という含意のとおり、その像は $Y$ において余次元 $0$ である。
\end{lemma}

\begin{proof}
Properties of Spaces,
Lemma \ref{spaces-properties-lemma-codimension-0-points} とエタール局所化により、
これはスキームの射と既約成分の生成点の場合に翻訳される。
この場合、一般化はスキームの平坦射に沿って持ち上がることから結果が従う；
Morphisms, Lemma \ref{morphisms-lemma-generalizations-lift-flat} を参照せよ。
\end{proof}

\begin{lemma}
\label{spaces-morphisms-lemma-flat-locally-finite-type-points-codimension-0}
$S$ をスキームとする。$f : X \to Y$ を $S$ 上の代数空間の射とする。
$f$ は平坦かつ局所有限型であり、$Y$ は
Lemma \ref{spaces-morphisms-lemma-prepare-normalization} の同値な条件を満たすと仮定する。
このとき $X$ も Lemma \ref{spaces-morphisms-lemma-prepare-normalization} の同値な条件を満たし、
$x \in |X|$ に対して、$x$ が余次元 $0$ ならば、
$X \Rightarrow f(x)$ という含意のとおり、その像は $Y$ において余次元 $0$ である。
\end{lemma}

\begin{proof}
最後の主張は Lemma \ref{spaces-morphisms-lemma-flat-image-point-codimension-0} から従う。
$V$ がスキームであるような全射エタール射 $V \to Y$ を選ぶ。
$U$ がスキームであるような全射エタール射 $U \to X \times_Y V$ を選ぶ。
$U$ の各準コンパクト開集合が有限個の既約成分をもつことを示せば十分である。
以下、Properties of Spaces,
Lemma \ref{spaces-properties-lemma-codimension-0-points} の結果を断りなく用いる。
すでに示したことにより、$U$ の余次元 $0$ の点は $U$ の余次元 $0$ の点の上にあり、
これらは仮定により局所有限である。したがって、余次元 $0$ の $v \in V$ に対して、
スキーム論的ファイバー $U_v = U \times_V v$ の余次元 $0$ の点が局所有限であることを
示せば十分である。これは $U_v$ が $\kappa(v)$ 上局所有限型のスキームであり、
したがって局所ネーター的であることから成り立つ。例えば
Lemma \ref{spaces-morphisms-lemma-locally-noetherian-can-be-normalized} を適用できる。
\end{proof}

\begin{lemma}
\label{spaces-morphisms-lemma-normalization}
$S$ をスキームとする。Lemma \ref{spaces-morphisms-lemma-prepare-normalization} の同値な条件を
満たす任意の $S$ 上の代数空間 $X$ に対して、次の性質をもつ代数空間の射
$$
\nu_X : X^\nu \longrightarrow X
$$
が存在する：
\begin{enumerate}
\item $X$ が Lemma \ref{spaces-morphisms-lemma-prepare-normalization} の同値な条件を満たすならば、
$X^\nu$ は正規であり、$\nu_X$ は整である。
\item $X$ が各準コンパクト開集合に有限個の既約成分をもつスキームならば、
$\nu_X : X^\nu \to X$ は Morphisms,
Section \ref{morphisms-section-normalization} で構成された $X$ の正規化である。
\item $f : X \to Y$ が $S$ 上の代数空間の射で、両者が
Lemma \ref{spaces-morphisms-lemma-prepare-normalization} の同値な条件を満たし、かつ
$X$ の各余次元 $0$ の点が $f$ によって $Y$ の余次元 $0$ の点へ写されるならば、
$\nu_Y \circ f^\nu = f \circ \nu_X$ を満たす $S$ 上の代数空間の射
$f^\nu : X^\nu \to Y^\nu$ が一意に存在する。
\item $f : X \to Y$ が代数空間のエタール射または滑らかな射であり、$Y$ が
Lemma \ref{spaces-morphisms-lemma-prepare-normalization} の同値な条件を満たすならば、
(3) の仮定が成り立ち、射 $f^\nu$ は同型
$X^\nu  \to X \times_Y Y^\nu$ を誘導する。
\end{enumerate}
\end{lemma}

\begin{proof}
圏 $\mathcal{C}$ を次のように考える。その対象は $S$ 上のスキーム $U$ であって、
$U$ の各準コンパクト開集合が有限個の既約成分をもつものであり、その射は、
$U$ の各既約成分の生成点を $V$ のある既約成分の生成点に写す
$S$ 上のスキームの射 $g : U \to V$ である。
すでに次を示している：
\begin{enumerate}
\item[(a)] $U \in \Ob(\mathcal{C})$ に対し、Morphisms,
Definition \ref{morphisms-definition-normalization} の正規化射
$\nu_U : U^\nu \to U$ がある。
\item[(b)] $U \in \Ob(\mathcal{C})$ に対して、射 $\nu_U$ は整であり、
$U^\nu$ は正規スキームである；Morphisms,
Lemma \ref{morphisms-lemma-normalization-normal} を参照せよ。
\item[(c)] 任意の $g : U \to V \in \text{Arrows}(\mathcal{C})$ に対して、
$\nu_V \circ g^\nu = g \circ \nu_U$ を満たす一意な射
$g^\nu : U^\nu \to V^\nu$ が存在する；Morphisms,
Lemma \ref{morphisms-lemma-normalization-normal} の (4) を合成
$X^\nu \to X \to Y$ に適用せよ。
\item[(d)] $V \in \Ob(\mathcal{C})$ で $g : U \to V$ がエタールまたは滑らかならば、
$U \in \Ob(\mathcal{C})$ かつ $g \in \text{Arrows}(\mathcal{C})$ であり、
射 $g^\nu$ は同型 $U^\nu \to U \times_V V^\nu$ を誘導する；
Lemma \ref{spaces-morphisms-lemma-flat-locally-finite-type-points-codimension-0} および
More on Morphisms, Lemma
\ref{more-morphisms-lemma-normalization-and-smooth} を参照せよ。
\end{enumerate}
我々の課題は、この構成を対応する $S$ 上の代数空間 $X$ の圏へ拡張することである。

\medskip\noindent
$X$ を Lemma \ref{spaces-morphisms-lemma-prepare-normalization} の同値な条件を満たす
$S$ 上の代数空間とする。$U$ がスキームであるような全射エタール射
$U \to X$ を取る。射影を $s, t : R \to U$ として
$R = U \times_X U$ と置き、さらに $X = U/R$ となるように
$j = (t, s) : R \to U \times_S U$ と置く；Spaces,
Lemma \ref{spaces-lemma-space-presentation} を参照せよ。
$X$ に関する仮定により $U$ と $R$ は $\mathcal{C}$ の対象であり、
射 $s$ と $t$ は $S$ 上のスキームのエタール射であることに注意する。
(a) により正規化射 $\nu_U : U^\nu \to U$ と
$\nu_R : R^\nu \to R$ があり、(d) により射
$s^\nu : R^\nu \to U^\nu$, $t^\nu : R^\nu \to U^\nu$ があり、これらは
同型 $R^\nu \to R \times_{s, U} U^\nu$ および
$R^\nu \to U^\nu \times_{U, t} R$ を定める。
したがって $s^\nu$ と $t^\nu$ はエタールである
（エタール射の基底変換と同型だからである）。誘導される射
$j^\nu = (t^\nu, s^\nu) : R^\nu \to U^\nu \times_S U^\nu$ は単射である。
実際、それは次の合成に等しい：
\begin{align*}
R^\nu
& \to
(U^\nu \times_{U, t} R) \times_R (R \times_{s, U} U^\nu) \\
& =
U^\nu \times_{U, t} R \times_{s, U} U^\nu \\
& \xrightarrow{j}
U^\nu \times_U (U \times_S U) \times_U U^\nu \\
& =
U^\nu \times_S U^\nu
\end{align*}
第一の矢印は $\nu_R$ の対角射である。
（これは、$R^\nu$ が $U^\nu$ への $R$ の制限の部分スキームであることを示す。）
性質 (d) を用いたファイバー積の形式的計算により、
$R^\nu \times_{s^\nu, U^\nu, t^\nu} R^\nu$ は
$R \times_{s, U, t} R$ の正規化であることが分かる。
したがってエタール射 $c : R \times_{s, U, t} R \to R$ は、(d) により
$c^\nu$ へ一意に延長される。射 $c^\nu$ は射影
$\text{pr}_{13} : U^\nu \times_S U^\nu \times_S U^\nu \to U^\nu \times_S U^\nu$
と両立する。同様に、因子を入れ替える射
$U^\nu \times_S U^\nu \to U^\nu \times_S U^\nu$ と両立する射
$i^\nu : R^\nu \to R^\nu$ があり、対角射
$U^\nu \to U^\nu \times_S U^\nu$ と両立する射
$e^\nu : U^\nu \to R^\nu$ がある。以上から、
$j^\nu : R^\nu \to U^\nu \times_S U^\nu$ はエタール同値関係である。
ここで $X^\nu = U^\nu/R^\nu$ と置く
（Spaces, Theorem \ref{spaces-theorem-presentation}）。
すると Groupoids, Lemma
\ref{groupoids-lemma-criterion-fibre-product} により
$U^\nu = X^\nu \times_X U$ である。

\medskip\noindent
前段落で示したことは次のとおりである。
Lemma \ref{spaces-morphisms-lemma-prepare-normalization} の同値な条件を満たす
$S$ 上の任意の代数空間 $X$ に対して、$U$ がスキームであるような
全射エタール射 $g : U \to X$ を選ぶと、代数空間のカルテジアン図式
$$
\xymatrix{
X^\nu \ar[d]_{\nu_X} & U^\nu \ar[l]^{g^\nu} \ar[d]^{\nu_U} \\
X & U \ar[l]_g
}
$$
を得る。これから直ちに $X^\nu$ が正規代数空間であり、$\nu_X$ が整射であることが
従う。これで補題の (1) が得られた。

\medskip\noindent
射 $\nu_X : X^\nu \to X$ が一意な同型を除いて $g$ の選択に依存しないことを
後で示すが、いま $X$ がスキームならば $\text{id} : X \to X$ を選ぶ。
これにより補題の (2) が成り立つことは明らかである。

\medskip\noindent
まだ (3) と (4) を証明する必要がある。$g : U \to X$、
$\nu_X : X^\nu \to X$、および $g^\nu : U^\nu \to X^\nu$ を上のとおりとする。
$Z$ を正規スキームとし、$h : Z \to U$ と $a : Z \to X^\nu$ を
$g \circ h = \nu_X \circ a$ を満たす $S$ 上の射とする。また $Z$ の各既約成分は、
$h$ を介して $U$ のある既約成分を支配すると仮定する。
Morphisms, Lemma \ref{morphisms-lemma-normalization-normal} の (4) により、
$h = \nu_U \circ h^\nu$ を満たす一意な射 $h^\nu : Z \to U^\nu$ を得る。
図示すれば：
$$
\xymatrix{
X^\nu \ar[d]_{\nu_X} &
U^\nu \ar[l]^{g^\nu} \ar[d]^{\nu_U} &
Z \ar[l]^{h^\nu} \ar@/_1em/[ll]_a \ar[dl]^h
\\
X & U \ar[l]_g
}
$$
$a = g^\nu \circ h^\nu$ であることに注意する。実際、四隅が
$X^\nu$, $X$, $U^\nu$, $U$ である正方形はカルテジアンなので、
これは $h^\nu$ が（$h$ を与えれば）一意であることから直ちに従う。
言い換えると、上のような $h : Z \to U$ が与えられれば（$a$ は与えなくても）、
$\nu_X \circ a = g \circ h$ を満たす射 $a : Z \to X^\nu$ が一意に存在する。

\medskip\noindent
$f : X \to Y$ を補題の主張の (3) のとおりとする。
$U$ と $V$ がスキームであり、$f$ が射 $g : U \to V$ に持ち上がるような
全射エタール射 $U \to X$ と $V \to Y$ を選んだと仮定する。
このとき $g \in \text{Arrows}(\mathcal{C})$ であり、$\nu_U$ と $\nu_V$ と両立する
一意な射 $g^\nu : U^\nu \to V^\nu$ を得る。しかしこのとき、二つの射
$$
R^\nu = U^\nu \times_{X^\nu} U^\nu
\to U^\nu \to V^\nu \to Y^\nu
$$
は、前段落の注意を（$X$ の代わりに $Y$ に適用して）用いると同じでなければならない。
$X^\nu$ は $U^\nu$ を $R^\nu$ で割った商として構成されているので、
(3) に述べた（唯一の）射 $f^\nu : X^\nu \to Y^\nu$ が得られる。

\medskip\noindent
$X^\nu$ の構成が全射エタール射 $g : U \to X$ の選択に依存しないことを見るため、
前段落の構成を $\text{id} : X \to X$ と、$X$ のエタール被覆間の射
$U' \to U$ に適用する。$X$ の任意の二つのエタール被覆は、両方を支配する
第三の被覆をもつので、これで十分である。$U' \to X$ または $U \to X$ を用いて
構成された二つの正規化間の射は、$U'$ への基底変換後に同型となり、したがって
もともと同型であったことを読者は確認できる。詳細は省略する。

\medskip\noindent
同様に証明される (4) の証明は省略する；上の (d) を用いるのがヒントである。
\end{proof}

\noindent
これにより次の定義に至る。

\begin{definition}
\label{spaces-morphisms-definition-normalization}
$S$ をスキームとする。$X$ を Lemma \ref{spaces-morphisms-lemma-prepare-normalization} の
同値な条件を満たす $S$ 上の代数空間とする。
$X$ の \emph{正規化}を、Lemma \ref{spaces-morphisms-lemma-normalization} で構成した射
$$
\nu_X : X^\nu \longrightarrow X
$$
として定義する。
\end{definition}

\noindent
この定義は局所ネーター代数空間に適用できる；
Lemma \ref{spaces-morphisms-lemma-locally-noetherian-can-be-normalized} を参照せよ。
通常、正規化は被約代数空間についてのみ定義される。
上の定義では、$X$ の正規化は $X$ の簡約 $X_{red}$ の正規化と同じである。

\begin{lemma}
\label{spaces-morphisms-lemma-normalization-reduced}
$S$ をスキームとする。$X$ を Lemma \ref{spaces-morphisms-lemma-prepare-normalization} の
同値な条件を満たす $S$ 上の代数空間とする。
正規化射 $\nu$ は簡約 $X_{red}$ を経由し、$X^\nu \to X_{red}$ は
$X_{red}$ の正規化である。
\end{lemma}

\begin{proof}
これは $X$ 上エタール局所的に確認できるので、スキームの場合、すなわち
Morphisms, Lemma \ref{morphisms-lemma-normalization-reduced} に帰着する。
いくつかの詳細は省略する。
\end{proof}

\begin{lemma}
\label{spaces-morphisms-lemma-normalization-normal}
$S$ をスキームとする。$X$ を Lemma \ref{spaces-morphisms-lemma-prepare-normalization} の
同値な条件を満たす $S$ 上の代数空間とする。
\begin{enumerate}
\item 正規化 $X^\nu$ は正規である。
\item 射 $\nu : X^\nu \to X$ は整かつ全射である。
\item 写像 $|\nu| : |X^\nu| \to |X|$ は余次元 $0$ の点の集合の間の全単射を誘導する
（Properties of Spaces,
Definition \ref{spaces-properties-definition-dimension-local-ring}）。
\item $Z \to X$ を射とする。$Z$ は正規代数空間であり、$z \in |Z|$ に対して、
$z$ が $Z \Rightarrow f(z)$ において余次元 $0$ ならば、$X$ において余次元 $0$
であると仮定する。このとき一意な分解 $Z \to X^\nu \to X$ が存在する。
\end{enumerate}
\end{lemma}

\begin{proof}
性質 (1), (2), (3) は、スキームに対する対応する結果
（Morphisms, Lemma \ref{morphisms-lemma-normalization-normal}）と、
スキームの点が既約成分の生成点であることと、その局所環の次元が零であることとの
同値性（Properties, Lemma \ref{properties-lemma-generic-point}）を合わせれば従う。

\medskip\noindent
$Z \to X$ を (4) のような射とする。$U$ をスキームとし、
$U \to X$ を全射エタール射とする。スキーム $V$ と全射エタール射
$V \to U \times_X Z$ を選ぶ。余次元 $0$ の点に関する条件により、
$V \to U$ は $V$ の既約成分の生成点を $U$ の既約成分の生成点へ写す。
したがって Morphisms,
Lemma \ref{morphisms-lemma-normalization-normal} により一意な分解
$V \to U^\nu \to U$ を得る。一意性により、二つの写像
$$
V \times_{U \times_X Z} V \to V \to U^\nu
$$
は一致する。実際、これらは写像
$V \times_{U \times_X Z} V \to V \to U$ の一意な分解である。
代数空間 $U \times_X Z$ は商 $V/V \times_{U \times_X Z} V$ に等しいので
（Spaces, Section \ref{spaces-section-presentations} を参照せよ）、
標準的な射 $U \times_X Z \to U^\nu$ を得る。図示すれば
$$
\xymatrix{
U \times_X Z \ar[r] \ar[d] & U^\nu \ar[r] \ar[d] & U \ar[d] \\
Z \ar@/_/[rr] \ar@{..>}[r] & X^\nu \ar[r] & X
}
$$
点線の矢印を得るため、射 $U \times_X Z \to U^\nu$ の構成がエタール射
$U \to X$ に関して関手的であることに注意する
（正確な定式化と証明は省略する）。したがって、射影を
$s, t : R \to U$ として $R = U \times_X U$ と置くと、
$s, t : R \to U$ および $s^\nu, t^\nu : R^\nu \to U^\nu$ と可換する射
$R \times_X Z \to R^\nu$ を得る。
Lemma \ref{spaces-morphisms-lemma-normalization} の証明で見たように
$X^\nu = U^\nu/R^\nu$ であることを思い出そう。$X = U/R$ なので、
簡単な層論的議論により $Z = (U \times_X Z)/(R \times_X Z)$ である。
したがって射 $U \times_X Z \to U^\nu$ と $R \times_X Z \to R^\nu$ は、
所望の射 $Z \to X^\nu$ を定める。
\end{proof}

\begin{lemma}
\label{spaces-morphisms-lemma-nagata-normalization}
$S$ をスキームとする。$X$ を $S$ 上の永田代数空間とする。
正規化 $\nu : X^\nu \to X$ は有限射である。
\end{lemma}

\begin{proof}
$X$ が永田であることは局所ネーター的であることを含むので、
Definition \ref{spaces-morphisms-definition-normalization} を適用できる。
Lemma \ref{spaces-morphisms-lemma-normalization} における $X^\nu$ の構成により、直ちに
スキームの場合、すなわち Morphisms,
Lemma \ref{morphisms-lemma-nagata-normalization} に帰着する。
\end{proof}













\section{分離的な局所準有限射}
\label{spaces-morphisms-section-schemehood}

\noindent
本節では、スキーム上局所準有限かつ分離的な代数空間は表現可能であることを証明する。
これから、分離的かつ局所準有限な射が表現可能であることが従う
（Lemma \ref{spaces-morphisms-lemma-locally-quasi-finite-separated-representable} を参照せよ）。
しかしまず、補題を一つ示す（これは Proposition
\ref{spaces-morphisms-proposition-locally-quasi-finite-separated-over-scheme}
によって不要になる）。

\begin{lemma}
\label{spaces-morphisms-lemma-neighbourhood-scheme}
$S$ をスキームとする。$S$ 上の代数空間の可換図式
$$
\xymatrix{
V' \ar[r] \ar[rd] & T' \times_T X \ar[r] \ar[d] & X \ar[d] \\
& T' \ar[r] & T
}
$$
を考える。次を仮定する：
\begin{enumerate}
\item $T' \to T$ はアフィンスキームのエタール射である。
\item $X \to T$ は分離的な局所準有限射である。
\item $V'$ は $T' \times_T X$ の開部分空間である。
\item $V' \to T'$ は準アフィンである。
\end{enumerate}
この状況で、$X$ における $V'$ の像 $U$ は、表現可能な
$X$ の準コンパクト開部分空間である。
\end{lemma}

\begin{proof}
まず自明な観察をいくつか行う。
Lemma \ref{spaces-morphisms-lemma-quasi-affine-local} により $V'$ は表現可能である。
また準コンパクトでもある（アフィンスキーム上の準アフィンスキームだからである；
Morphisms, Lemma \ref{morphisms-lemma-quasi-affine-separated} を参照せよ）。
$T' \times_T X \to X$ はエタールなので
（Properties of Spaces,
Lemma \ref{spaces-properties-lemma-base-change-etale}）、写像
$|T' \times_T X| \to |X|$ は開である；Properties of Spaces,
Lemma \ref{spaces-properties-lemma-etale-open} を参照せよ。
$|V'|$ の像に対応する開部分空間を $U \subset X$ とする；
Properties of Spaces, Lemma \ref{spaces-properties-lemma-open-subspaces}
を参照せよ。$|V'|$ は準コンパクトなので $|U|$ も準コンパクトであり、したがって
Properties of Spaces,
Lemma \ref{spaces-properties-lemma-quasi-compact-space} により
$U$ は準コンパクト代数空間である。

\medskip\noindent
Morphisms,
Lemma \ref{morphisms-lemma-locally-quasi-finite-qc-source-universally-bounded}
により、射 $T' \to T$ は普遍的に有界である。したがって、
$T' \to T$ のファイバーの次数を抑える整数 $n$ に関する帰納法を行える；
エタール射の場合のこの整数の記述については Morphisms,
Lemma \ref{morphisms-lemma-etale-universally-bounded} を参照せよ。
$n = 1$ ならば、$T' \to T$ は開埋入であり
（\'Etale Morphisms, Theorem \ref{etale-theorem-etale-radicial-open} を参照）、
結果は明らかである。$n > 1$ と仮定する。

\medskip\noindent
アフィンスキーム $T'' = T' \times_T T'$ を考える。
$T' \to T$ はエタールなので、（開かつ閉なアフィン部分スキームへの）分解
$T'' = \Delta(T') \amalg T^*$ がある。実際、
$\Delta = \Delta_{T'/T}$ は Morphisms,
Lemma \ref{morphisms-lemma-diagonal-unramified-morphism} により開であり、
$T' \to T$ はアフィン間の射として分離的なので閉である。
第二射影のファイバーの次数は、基底変換
$\text{pr}_1 : T' \times_T T' \to T'$ として $n$ で抑えられる；
Morphisms, Lemma \ref{morphisms-lemma-base-change-universally-bounded}
を参照せよ。一方、$\text{pr}_1|_{\Delta(T')} : \Delta(T') \to T'$ は同型で、
各ファイバーはちょうど一点をもつ。したがって Morphisms,
Lemma \ref{morphisms-lemma-etale-universally-bounded} を適用すると、制限
$\text{pr}_1|_{T^*} : T^* \to T'$ のファイバーの次数は $n - 1$ で抑えられる。
よって、図式
$$
\xymatrix{
p_0^{-1}(V') \cap X^* \ar[r] \ar[rd] &
X^* \ar[r]_{p_1|_{X^*}} \ar[d] &
X' \ar[d] \\
& T^* \ar[r]^{\text{pr}_1|_{T^*}} & T'
}
$$
に帰納法の仮定を適用すると、$p_1(p_0^{-1}(V') \cap X^*)$ は
準コンパクトスキームであることが分かる。ここで
$X'' = T'' \times_T X$, $X^* = T^* \times_T X$, $X' = T' \times_T X$ と置き、
$p_0, p_1 : X'' \to X'$ は $\text{pr}_0, \text{pr}_1$ の基底変換である。
補題の仮定の大部分は、基底変換により上の図式に対応する仮定を与える。
例えば $p_0^{-1}(V') = T'' \times_{T'} V'$ は、基底変換として
$T''$ 上準アフィンなスキームである。いくつかの確認は省略する。

\medskip\noindent
Properties of Spaces, Lemma \ref{spaces-properties-lemma-subscheme}
により、
$$
p_1(p_0^{-1}(V')) =
V' \cup p_1(p_0^{-1}(V') \cap X^*)
$$
は準コンパクトスキームである。さらに、$p_1(p_0^{-1}(V'))$ は、
証明の第一段落で論じた準コンパクト開部分空間 $U \subset X$ の逆像であることは明らかである。
言い換えると、$T' \times_T U$ はスキームである。
$T' \times_T U$ は準コンパクトであり、かつ $T'$ 上分離的かつ局所準有限であることに
注意する。実際、$T' \times_T X \to T'$ はもとの射 $X \to T$ の基底変換なので、
局所準有限かつ分離的である（Lemmas \ref{spaces-morphisms-lemma-base-change-separated} および
\ref{spaces-morphisms-lemma-base-change-quasi-finite} を参照せよ）。
したがって More on Morphisms,
Lemma \ref{more-morphisms-lemma-quasi-finite-separated-quasi-affine}
により $T' \times_T U \to T'$ は準アフィンである。

\medskip\noindent
Descent, Lemma \ref{descent-lemma-descent-data-sheaves}
により、これはエタール被覆 $\{T' \to W\}$ に関する
$T' \times_T U / T'$ 上の降下データを与える。ここで $W \subset T$ は
射 $T' \to T$ の像である。$U'$ は $T'$ 上準アフィンなので、Descent,
Lemma \ref{descent-lemma-quasi-affine} によりこのデータは有効である。
さらに Descent, Lemma \ref{descent-lemma-descent-data-sheaves} の最後の部分により、
これは所望どおり $U$ がスキームであることを含意する。
いくつかの細部は省略する。
\end{proof}

\begin{proposition}
\label{spaces-morphisms-proposition-locally-quasi-finite-separated-over-scheme}
$S$ をスキームとする。
$f : X \to T$ を $S$ 上の代数空間の射とする。
次を仮定する：
\begin{enumerate}
\item $T$ は表現可能である。
\item $f$ は局所準有限である。
\item $f$ は分離的である。
\end{enumerate}
このとき $X$ は表現可能である。
\end{proposition}

\begin{proof}
$T = \bigcup T_i$ をスキーム $T$ のアフィン開被覆とする。
開部分空間 $X_i = f^{-1}(T_i)$ が表現可能であることを示せれば、
Properties of Spaces, Lemma \ref{spaces-properties-lemma-subscheme}
により $X$ は表現可能である。$X_i = T_i \times_T X$ であり、
局所準有限性と分離性はいずれも基底変換で安定であることに注意する；
Lemmas \ref{spaces-morphisms-lemma-base-change-separated} および
\ref{spaces-morphisms-lemma-base-change-quasi-finite} を参照せよ。
したがって $T$ はアフィンスキームであると仮定してよい。

\medskip\noindent
Properties of Spaces,
Lemma \ref{spaces-properties-lemma-union-of-quasi-compact}
により、各 $X_i$ がアフィンスキームによる全射エタール被覆をもつような
ザリスキー被覆 $X = \bigcup X_i$ が存在する。再び Properties of Spaces,
Lemma \ref{spaces-properties-lemma-subscheme} により、各 $X_i$ について命題を
証明すれば十分である。したがって、アフィンスキーム $U$ と全射エタール射
$U \to X$ が存在すると仮定してよい。これで次段落の状況に帰着する。

\medskip\noindent
$$
U \longrightarrow X \longrightarrow T
$$
があり、$U$ と $T$ はアフィンスキーム、$U \to X$ は全射エタール、
$X \to T$ は分離的かつ局所準有限であると仮定する。
Lemmas \ref{spaces-morphisms-lemma-etale-locally-quasi-finite} および
\ref{spaces-morphisms-lemma-composition-quasi-finite} により、射 $U \to T$ は局所準有限である。
$U$ と $T$ はアフィンなので、これは準有限である。
$R = U \times_X U$ と置く。このとき $X = U/R$ である；Spaces,
Lemma \ref{spaces-lemma-space-presentation} を参照せよ。
$X \to T$ は分離的なので、射 $R \to U \times_T U$ は閉埋入である；
Lemma \ref{spaces-morphisms-lemma-fibre-product-after-map} を参照せよ。
特に $R$ もアフィンスキームである。$U \to X$ はエタールなので、射影
$t, s : R \to U$ もエタールである。特に $s$ と $t$ は準有限、平坦、かつ有限表示である
（Morphisms, Lemmas \ref{morphisms-lemma-etale-locally-quasi-finite},
\ref{morphisms-lemma-etale-flat} および
\ref{morphisms-lemma-etale-locally-finite-presentation} を参照せよ）。

\medskip\noindent
$(U, R, s, t, c)$ を $U$ 上のエタール同値関係 $R$ に付随する亜群とする。
$u \in U$ を点とし、その像を $p \in T$ と書く。スキーム $T$ 上の亜群
$(U, R, s, t, c)$ と上の点 $p$ および $u$ に対して、More on Groupoids,
Lemma \ref{more-groupoids-lemma-quasi-finite-over-base-j-proper}
を用いる。前段落の議論により、この補題の仮定 (1) -- (7) はすべて満たされる。
したがってエタール近傍 $(T', p') \to (T, p)$、非交和分解
$$
U_{T'} = U' \amalg W, \quad
R_{T'} = R' \amalg W'
$$
および $u' \in U'$ であって、上記の More on Groupoids,
Lemma \ref{more-groupoids-lemma-quasi-finite-over-base-j-proper}
の結論 (a), (b), (c), (d), (e), (f), (g), (h) を満たすものを得る。
（必要なら $T'$ を縮小して）$T'$ はアフィンであると仮定してよい。
結論 (h) により $R' = U' \times_{X_{T'}} U'$ であり、射影は $s'$ と $t'$ の
制限と同一視される。したがって、結論 (g) の
$(U', R', s'|_{R'}, t'|_{R'}, c'|_{R' \times_{t', U', s'} R'})$ は
エタール同値関係である。Spaces, Lemma \ref{spaces-lemma-finding-opens}
により、$U'/R'$ は $X_{T'}$ の開部分空間である。
結論 (d) によりスキーム $U'$, $R'$ はアフィンで、射
$s'|_{R'}, t'|_{R'}$ は有限エタールである。よって Groupoids,
Proposition \ref{groupoids-proposition-finite-flat-equivalence}
を適用でき、$U'/R'$ はアフィンスキームである。

\medskip\noindent
以上から、上のような点の各対 $(u, p)$ に対して、
$\kappa(p) = \kappa(p')$ となるエタール近傍 $(T', p') \to (T, p)$ と、
$u$ へ写る点 $u' \in U_{T'}$ であって、$|X_{T'}|$ における $u'$ の像 $x'$ が
アフィンスキームである開近傍 $V'$ を $X_{T'}$ 内にもつものを見つけられる。
Lemma \ref{spaces-morphisms-lemma-neighbourhood-scheme} を適用して、スキームである開部分空間
$W \subset X$ で、$x$（$|X|$ における $u$ の像）を含むものを得る。
これは各 $x$ に対して成り立つので、Properties of Spaces,
Lemma \ref{spaces-properties-lemma-subscheme} により $X$ はスキームである。
これで証明を終える。
\end{proof}




\section{応用}
\label{spaces-morphisms-section-applications}

\noindent
次の補題の別証明は、More on Morphisms of Spaces, Section
\ref{spaces-more-morphisms-section-structure-quasi-finite} で論じられている
代数空間の（非表現可能）射に対するザリスキーの主定理の帰結とみなすことで得られる。
実際、More on Morphisms of Spaces,
Lemma \ref{spaces-more-morphisms-lemma-quasi-finite-separated-quasi-affine}
により、代数空間の準有限かつ分離的な射は準アフィンであり、したがって表現可能である。

\begin{lemma}
\label{spaces-morphisms-lemma-locally-quasi-finite-separated-representable}
$S$ をスキームとする。$f : X \to Y$ を $S$ 上の代数空間の射とする。
$f$ が局所準有限かつ分離的ならば、$f$ は表現可能である。
\end{lemma}

\begin{proof}
Proposition \ref{spaces-morphisms-proposition-locally-quasi-finite-separated-over-scheme}
と、局所準有限かつ分離的であることが任意の基底変換で保たれるという事実から
直ちに従う；Lemmas \ref{spaces-morphisms-lemma-base-change-quasi-finite} および
\ref{spaces-morphisms-lemma-base-change-separated} を参照せよ。
\end{proof}

\begin{lemma}
\label{spaces-morphisms-lemma-etale-universally-injective-open}
\begin{slogan}
普遍単射エタール写像は開埋入である。
\end{slogan}
$S$ をスキームとする。$f : X \to Y$ を $S$ 上の代数空間の
エタールかつ普遍単射な射とする。このとき $f$ は開埋入である。
\end{lemma}

\begin{proof}
$T \to Y$ をスキームから $Y$ への射とする。
$X \times_Y T \to T$ が開埋入であることを示せれば、証明は終わる。
エタール性と普遍単射性はいずれも基底変換で安定な射の性質なので
（Lemmas \ref{spaces-morphisms-lemma-base-change-etale} および
\ref{spaces-morphisms-lemma-base-change-universally-injective} を参照）、$Y$ はスキームであると
仮定してよい。対角射 $\Delta_{X/Y} : X \to X \times_Y X$ はエタール、単射、かつ
Lemma \ref{spaces-morphisms-lemma-universally-injective} により全射であることに注意する。
したがって $\Delta_{X/Y}$ は同型である
（Spaces, Lemma
\ref{spaces-lemma-surjective-flat-locally-finite-presentation} を参照）。
特に $X$ は $Y$ 上分離的である。よって Proposition
\ref{spaces-morphisms-proposition-locally-quasi-finite-separated-over-scheme}
により $X$ もスキームである。最後に、スキームのエタール射に関する基本定理
\'Etale Morphisms, Theorem \ref{etale-theorem-etale-radicial-open}
により $X \to Y$ は開埋入である。
\end{proof}



\section{ザリスキーの主定理（表現可能な場合）}
\label{spaces-morphisms-section-Zariski}

\noindent
これは、正規化がエタール局所化と可換することを用いて証明できる版である。
より強力な版（非表現可能射に対するもの）を証明する前に、さらに道具を発展させる
必要がある。More on Morphisms of Spaces, Section
\ref{spaces-more-morphisms-section-structure-quasi-finite} を参照せよ。

\begin{lemma}
\label{spaces-morphisms-lemma-finite-type-separated}
$S$ をスキームとする。$f : X \to Y$ を $S$ 上の代数空間の
表現可能、有限型、かつ分離的な射とする。$Y'$ を $X$ における $Y$ の正規化とする。
図示すれば：
$$
\xymatrix{
X \ar[rd]_f \ar[rr]_{f'} & & Y' \ar[ld]^\nu \\
& Y &
}
$$
このとき、次を満たす開部分空間 $U' \subset Y'$ が存在する：
\begin{enumerate}
\item $(f')^{-1}(U') \to U'$ は同型である。
\item $(f')^{-1}(U') \subset X$ は $f$ が準有限である点の集合である。
\end{enumerate}
\end{lemma}

\begin{proof}
$W$ がスキームであるような全射エタール射 $W \to Y$ を取る。
このとき $W \times_Y X$ もスキームである。Lemma
\ref{spaces-morphisms-lemma-properties-normalization} により、代数空間 $W \times_Y Y'$ は
表現可能であり、スキーム $W \times_Y X$ におけるスキーム $W$ の正規化である。
図示すれば
$$
\xymatrix{
W \times_Y X \ar[rd]_{(1, f)} \ar[rr]_{(1, f')} & &
W \times_Y Y' \ar[ld]^{(1, \nu)} \\
& W &
}
$$
More on Morphisms, Lemma \ref{more-morphisms-lemma-finite-type-separated}
により、補題の結果は $W$ 上で成り立つ。開部分スキーム
$V' \subset W \times_Y Y'$ を次を満たすものとする：
\begin{enumerate}
\item $(1, f')^{-1}(V') \to V'$ は同型である。
\item $(1, f')^{-1}(V') \subset W \times_Y X$ は $(1, f)$ が準有限である点の集合である。
\end{enumerate}
Lemma \ref{spaces-morphisms-lemma-locally-finite-type-quasi-finite-part} により、
$f$ が準有限である点の最大開集合 $U \subset X$ があり、
$W \times_Y U = (1, f')^{-1}(V')$ である。
射 $f'|_U : U \to Y'$ は Lemma \ref{spaces-morphisms-lemma-closed-immersion-local}
により開埋入である。実際、その $W$ への基底変換は同型
$(1, f')^{-1}(V') \to V'$ に開埋入 $V' \to W \times_Y Y'$ を続けたものである。
$U' = \Im(U \to Y')$ と置けば証明が終わる
（$(f')^{-1}(U') = U$ の確認は省略する）。
\end{proof}

\noindent
次の補題では、局所準有限な分離的射が表現可能であることをすでに示したので、
表現可能性の仮定を外すことができる。

\begin{lemma}
\label{spaces-morphisms-lemma-quasi-finite-separated-quasi-affine}
$S$ をスキームとする。
$f : X \to Y$ を $S$ 上の代数空間の射とする。
$f$ は準有限かつ分離的であると仮定する。
$Y'$ を $X$ における $Y$ の正規化とする。図示すれば：
$$
\xymatrix{
X \ar[rd]_f \ar[rr]_{f'} & & Y' \ar[ld]^\nu \\
& Y &
}
$$
このとき $f'$ は準コンパクトな開埋入であり、$\nu$ は整である。
特に $f$ は準アフィンである。
\end{lemma}

\begin{proof}
Lemma \ref{spaces-morphisms-lemma-locally-quasi-finite-separated-representable}
により射 $f$ は表現可能である。したがって
Lemma \ref{spaces-morphisms-lemma-finite-type-separated} を適用できる。
よって開部分空間 $U' \subset Y'$ が存在して、
$(f')^{-1}(U') = X$（！）であり、$X \to U'$ は同型である。
言い換えると、$f'$ は開埋入である。$f$ は準コンパクトで、
$\nu : Y' \to Y$ は分離的なので、$f'$ は準コンパクトであることに注意する
（Lemma \ref{spaces-morphisms-lemma-quasi-compact-permanence}）。
したがって、任意のアフィンスキーム $Z$ と射 $Z \to Y$ に対し、
ファイバー積 $Z \times_Y X$ はアフィンスキーム $Z \times_Y Y'$ の
準コンパクト開部分スキームである。よって定義により $f$ は準アフィンである。
\end{proof}








\section{普遍同相射}
\label{spaces-morphisms-section-universal-homeomorphisms}

\noindent
スキームの普遍同相射の類は、合成および任意の基底変換について閉じており、
底上 fppf 局所的である。Morphisms, Lemmas
\ref{morphisms-lemma-composition-universal-homeomorphism} および
\ref{morphisms-lemma-base-change-universal-homeomorphism}、さらに Descent,
Lemma \ref{descent-lemma-descending-property-universal-homeomorphism} を参照せよ。
したがって、この概念に Section \ref{spaces-morphisms-section-representable} の議論を適用すれば、
代数空間の表現可能射が普遍同相射であることの意味が分かる。

\begin{lemma}
\label{spaces-morphisms-lemma-universal-homeomorphism-representable}
$S$ をスキームとする。
$f : X \to Y$ を $S$ 上の代数空間の表現可能射とする。
このとき $f$ が（Section \ref{spaces-morphisms-section-representable} の意味で）普遍同相射であることと、
代数空間の任意の射 $Z \to Y$ に対して、基底変換写像
$Z \times_Y X \to Z$ が同相写像 $|Z \times_Y X| \to |Z|$ を誘導することとは同値である。
\end{lemma}

\begin{proof}
代数空間の任意の射 $Z \to Y$ に対して、基底変換写像
$Z \times_Y X \to Z$ が同相写像 $|Z \times_Y X| \to |Z|$ を誘導するならば、
$Z$ がスキームである場合にも同じことが成り立つ。これは形式的に、$f$ が
Section \ref{spaces-morphisms-section-representable} の意味で普遍同相射であることを含意する。
逆に $f$ が Section \ref{spaces-morphisms-section-representable} の意味で普遍同相射ならば、
$X \to Y$ は整、普遍単射、かつ全射である
（Spaces, Lemma
\ref{spaces-lemma-representable-transformations-property-implication}
および Morphisms, Lemma
\ref{morphisms-lemma-universal-homeomorphism} による）。
したがって $f$ は Lemma \ref{spaces-morphisms-lemma-integral-universally-closed} により普遍閉であり、
また普遍単射かつ（普遍的に）全射である。すなわち $f$ は普遍同相射である。
\end{proof}

\begin{definition}
\label{spaces-morphisms-definition-universal-homeomorphism}
$S$ をスキームとする。
$S$ 上の代数空間の射 $f : X \to Y$ が \emph{普遍同相射}であるとは、
代数空間の任意の射 $Z \to Y$ に対して、基底変換
$Z \times_Y X \to Z$ が同相写像 $|Z \times_Y X| \to |Z|$ を誘導することをいう。
\end{definition}

\noindent
Lemma \ref{spaces-morphisms-lemma-universal-homeomorphism-representable} により、この定義は
代数空間の表現可能射に対する既存の定義と衝突しない。
代数空間の射については、普遍同相射が常に整であるとは限らない。

\begin{example}
\label{spaces-morphisms-example-universal-homeomorphism}
これは Remark \ref{spaces-morphisms-remark-universally-injective-not-separated} の続きである。
代数空間
$X = \mathbf{A}^1_k/\{x \sim -x \mid x \not = 0\}$
を考える。射
$$
\mathbf{A}^1_k \longrightarrow X \longrightarrow \mathbf{A}^1_k
$$
であって、第一の矢印が全射エタール、第二の矢印が普遍単射で、合成が写像
$x \mapsto x^2$ となるものがある。したがって合成は普遍閉である。
よって写像 $X \to \mathbf{A}^1_k$ は普遍同相射であるが、
$X \to \mathbf{A}^1_k$ は分離的ではない。
\end{example}

\noindent
$S$ をスキームとする。
$f : X \to Y$ を $S$ 上の代数空間の普遍同相射とする。
このとき $f$ は普遍閉であり、したがって準コンパクトである；
Lemma \ref{spaces-morphisms-lemma-universally-closed-quasi-compact} を参照せよ。
しかし $f$ は分離的とは限らず（上の例を参照）、準分離的ですらないことがある。
一例は、無限次元アフィン空間
$\mathbf{A}^\infty = \Spec(k[x_1, x_2, \ldots])$ を、非零座標のうち有限個の
符号を反転することで与えられる同値関係で割ったものである（詳細は省略する）。

\medskip\noindent
まず必須の補題を述べる。

\begin{lemma}
\label{spaces-morphisms-lemma-base-change-universal-homeomorphism}
代数空間の普遍同相射を代数空間の任意の射で基底変換したものは普遍同相射である。
\end{lemma}

\begin{proof}
定義から直ちに従う。
\end{proof}

\begin{lemma}
\label{spaces-morphisms-lemma-composition-universal-homeomorphism}
代数空間の二つの普遍同相射の合成は普遍同相射である。
\end{lemma}

\begin{proof}
省略する。
\end{proof}

\begin{lemma}
\label{spaces-morphisms-lemma-reduction-universal-homeomorphism}
$S$ をスキームとする。$X$ を $S$ 上の代数空間とする。
標準的な閉埋入 $X_{red} \to X$（Properties of Spaces, Definition
\ref{spaces-properties-definition-reduced-induced-space} を参照）は普遍同相射である。
\end{lemma}

\begin{proof}
省略する。
\end{proof}

\noindent
現在これを置くのにより良い場所がないので、次の結果をここに置く。

\begin{lemma}
\label{spaces-morphisms-lemma-integral-universally-injective-push-pull}
$S$ をスキームとする。$f : Y \to X$ を $S$ 上の代数空間の
普遍単射な整射とする。
\begin{enumerate}
\item 関手
$$
f_{small, *} : \Sh(Y_\etale) \longrightarrow \Sh(X_\etale)
$$
は充満忠実であり、その本質像は、$X_\etale$ 上の集合の層 $\mathcal{F}$ であって、
$|X| \setminus f(|Y|)$ への制限が $*$ と同型であるものからなる。
\item 関手
$$
f_{small, *} : \textit{Ab}(Y_\etale) \longrightarrow \textit{Ab}(X_\etale)
$$
は充満忠実であり、その本質像は、台が $f(|Y|)$ に含まれる
$Y_\etale$ 上のアーベル層からなる。
\end{enumerate}
いずれの場合にも $f_{small}^{-1}$ は関手 $f_{small, *}$ の左逆である。
\end{lemma}

\begin{proof}
$f$ は整なので普遍閉である
（Lemma \ref{spaces-morphisms-lemma-integral-universally-closed}）。
特に $f(|Y|)$ は $|X|$ の閉部分集合であり、主張には意味がある。
証明の残りは Lemma \ref{spaces-morphisms-lemma-closed-immersion-push-pull} の証明と同一であるが、
\'Etale Cohomology, Proposition
\ref{etale-cohomology-proposition-integral-universally-injective-pushforward}
を、\'Etale Cohomology, Proposition
\ref{etale-cohomology-proposition-closed-immersion-pushforward}
の代わりに用いる。
\end{proof}
